% ============================================================================
% Appendices (Phase B, 2026-11-02). Order: A label sets, B training settings (Sophon App. B order),
% C missing input features, D method checks, E complete results. Every number is a generated macro;
% "% CHECK:" marks a claim to recheck against the final macros.
% ============================================================================
% \appendix is called in main.tex before this file is input.

% ----------------------------------------------------------------------------
\section{Label sets}\label{sec:app-labels}
% ----------------------------------------------------------------------------

To give what each JetClass-II class becomes at every label set, table~\ref{tab:app-levels} lists the
levels of the class hierarchy and the groups of the \Kcoarse-class label set with the number of classes
each merges, and table~\ref{tab:app-vocabulary} gives the group of every class at every level of the
hierarchy, in each random partition and in each random partition with controlled flavor boundaries. The
label maps behind both tables, one row per class and one column per label set, are part of the code
release.

\begin{table}[p]
\centering
\small
\caption{The levels of the class hierarchy, finest first: the number of classes, what each level labels, and whether models were pretrained on it. Below, the groups of the 17-class label set, by the group number used in Table~\ref{tab:app-vocabulary}, with the number of JetClass-II classes each merges. $\ell$ is an electron or a muon; $\tau_\ell$ a leptonic $\tau$ decay. A four-prong group with a neutrino has three visible prongs.}
\label{tab:app-levels}
\begin{tabular}{r p{0.62\linewidth} l}
\toprule
classes & labels & pretrained \\
\midrule
188 & the JetClass-II classes & yes \\
162 & one QCD class; every resonant class kept & yes \\
64 & each decay topology by its numbers of $b$ and $c$ quarks & no \\
43 & each decay topology by whether it holds a $b$ quark, else a $c$, else neither & yes \\
30 & each decay topology by whether it holds a $b$ or $c$ quark & no \\
17 & the number of prongs and whether each is a quark or gluon (without distinguishing them), an electron or muon, a $\tau$ lepton or a neutrino, without flavor & yes \\
4 & two, three or four visible prongs (a neutrino is not counted), or QCD & no \\
2 & resonance decay or QCD & no \\
\midrule
\multicolumn{3}{@{}l}{\itshape the 17 classes of the 17-class label set} \\
0 & two-prong: two quarks or gluons & 10 classes \\
3 & two-prong: $\ell\ell$ & 2 classes \\
2 & two-prong: $\tau_h$, $\tau_\ell$ & 2 classes \\
1 & two-prong: $\tau_h$, $\tau_h$ & 1 class \\
9 & four-prong: four quarks or gluons & 27 classes \\
11 & four-prong: two quarks or gluons, $\ell\ell$ & 10 classes \\
12 & four-prong: two quarks or gluons, $\tau_h$, $\tau_\ell$ & 10 classes \\
10 & four-prong: two quarks or gluons, $\tau_h$, $\tau_h$ & 5 classes \\
4 & three-prong: three quarks or gluons & 29 classes \\
6 & three-prong: two quarks or gluons, $\ell$ & 10 classes \\
7 & three-prong: one quark or gluon, $\ell\ell$ & 10 classes \\
8 & three-prong: one quark or gluon, $\tau_h$, $\tau_\ell$ & 10 classes \\
5 & three-prong: one quark or gluon, $\tau_h$, $\tau_h$ & 5 classes \\
13 & four-prong: two quarks or gluons, $\ell$, $\nu$ & 12 classes \\
14 & four-prong: two quarks or gluons, $\tau_\ell$, $\nu$ & 12 classes \\
15 & four-prong: two quarks or gluons, $\tau_h$, $\nu$ & 6 classes \\
16 & QCD & 27 classes \\
\bottomrule
\end{tabular}
\end{table}

{\footnotesize
\setlength{\tabcolsep}{3pt}
\begin{longtable}{l rrrrrrrrrrrr}
\caption{Every JetClass-II class and its group at each coarser level of the class hierarchy (columns headed by the number of classes at that level: 64, 43, 30, 17), in each random partition, numbered as in Table~\ref{tab:pairs}, and in each random partition with controlled flavor boundaries: no flavor boundary, a $b$-quark boundary and a random flavor boundary (Section~\ref{sec:controls}). Group numbers are those of the released label maps; classes that share a number in a column share a class in that label set. At the 162-class level every resonant class keeps its own label and the QCD classes share one; at 4 classes the groups are decays with two, three or four visible prongs (a neutrino is not counted) and QCD. Rows are ordered by the 17-class group (Table~\ref{tab:app-levels}), with a gap between groups. $q$ is a $u$ or $d$ quark; $\tau_h$ a hadronic and $\tau_e$, $\tau_\mu$ leptonic $\tau$ decays.}\label{tab:app-vocabulary}\\
\toprule
 & \multicolumn{4}{c}{level} & \multicolumn{5}{c}{random partition} & \multicolumn{3}{c}{flavor boundary} \\
JetClass-II class & 64 & 43 & 30 & 17 & 1 & 2 & 3 & 4 & 5 & none & $b$ quark & random \\
\midrule
\endfirsthead
\toprule
 & \multicolumn{4}{c}{level} & \multicolumn{5}{c}{random partition} & \multicolumn{3}{c}{flavor boundary} \\
JetClass-II class & 64 & 43 & 30 & 17 & 1 & 2 & 3 & 4 & 5 & none & $b$ quark & random \\
\midrule
\endhead
\bottomrule
\endlastfoot
$X\to bb$ & 0 & 0 & 0 & 0 & 0 & 9 & 0 & 9 & 0 & 0 & 0 & 0 \\
$X\to bc$ & 1 & 0 & 0 & 0 & 0 & 0 & 0 & 9 & 0 & 0 & 0 & 0 \\
$X\to bq$ & 2 & 0 & 0 & 0 & 0 & 0 & 9 & 0 & 0 & 0 & 0 & 0 \\
$X\to cc$ & 3 & 1 & 0 & 0 & 0 & 0 & 0 & 0 & 0 & 0 & 0 & 0 \\
$X\to cs$ & 4 & 1 & 0 & 0 & 0 & 9 & 9 & 0 & 0 & 0 & 0 & 0 \\
$X\to cq$ & 4 & 1 & 0 & 0 & 9 & 0 & 9 & 0 & 0 & 0 & 0 & 0 \\
$X\to ss$ & 5 & 2 & 1 & 0 & 9 & 9 & 9 & 9 & 0 & 0 & 0 & 0 \\
$X\to qq$ & 5 & 2 & 1 & 0 & 0 & 0 & 9 & 0 & 0 & 0 & 0 & 0 \\
$X\to sq$ & 5 & 2 & 1 & 0 & 9 & 0 & 9 & 0 & 0 & 0 & 0 & 0 \\
$X\to gg$ & 5 & 2 & 1 & 0 & 3 & 3 & 9 & 13 & 9 & 15 & 15 & 15 \\
\addlinespace[2pt]
$X\to ee$ & 8 & 5 & 4 & 3 & 11 & 9 & 13 & 11 & 0 & 0 & 0 & 0 \\
$X\to \mu\mu$ & 8 & 5 & 4 & 3 & 11 & 9 & 9 & 11 & 9 & 11 & 9 & 9 \\
\addlinespace[2pt]
$X\to \tau_h\tau_e$ & 7 & 4 & 3 & 2 & 15 & 3 & 12 & 12 & 1 & 4 & 4 & 4 \\
$X\to \tau_h\tau_\mu$ & 7 & 4 & 3 & 2 & 4 & 4 & 10 & 10 & 12 & 0 & 0 & 0 \\
\addlinespace[2pt]
$X\to \tau_h\tau_h$ & 6 & 3 & 2 & 1 & 4 & 10 & 12 & 3 & 4 & 7 & 7 & 7 \\
\addlinespace[2pt]
$X\to YY\to bbbb$ & 31 & 21 & 15 & 9 & 0 & 9 & 11 & 0 & 11 & 9 & 11 & 9 \\
$X\to YY\to bbcc$ & 32 & 21 & 15 & 9 & 9 & 0 & 11 & 0 & 9 & 9 & 11 & 9 \\
$X\to YY\to bbss$ & 33 & 21 & 15 & 9 & 9 & 0 & 0 & 0 & 9 & 9 & 11 & 11 \\
$X\to YY\to bbqq$ & 33 & 21 & 15 & 9 & 9 & 0 & 11 & 0 & 9 & 9 & 11 & 11 \\
$X\to YY\to bbgg$ & 33 & 21 & 15 & 9 & 15 & 15 & 3 & 3 & 13 & 13 & 13 & 13 \\
$X\to YY\to qqqb$ & 34 & 21 & 15 & 9 & 0 & 0 & 0 & 0 & 9 & 9 & 11 & 11 \\
$X\to YY\to qqbs$ & 34 & 21 & 15 & 9 & 0 & 9 & 11 & 9 & 11 & 9 & 11 & 9 \\
$X\to YY\to bbcq$ & 35 & 21 & 15 & 9 & 9 & 0 & 11 & 9 & 11 & 9 & 11 & 11 \\
$X\to YY\to ccbs$ & 36 & 21 & 15 & 9 & 0 & 9 & 0 & 0 & 9 & 9 & 11 & 9 \\
$X\to YY\to ccbq$ & 36 & 21 & 15 & 9 & 0 & 9 & 11 & 0 & 11 & 9 & 11 & 11 \\
$X\to YY\to qqbc$ & 37 & 21 & 15 & 9 & 0 & 9 & 0 & 9 & 9 & 9 & 11 & 11 \\
$X\to YY\to bcsq$ & 37 & 21 & 15 & 9 & 0 & 9 & 0 & 0 & 11 & 9 & 11 & 9 \\
$X\to YY\to cccc$ & 38 & 22 & 15 & 9 & 9 & 9 & 0 & 9 & 9 & 9 & 9 & 9 \\
$X\to YY\to ccss$ & 39 & 22 & 15 & 9 & 0 & 0 & 0 & 0 & 11 & 9 & 9 & 11 \\
$X\to YY\to ccqq$ & 39 & 22 & 15 & 9 & 9 & 9 & 11 & 9 & 11 & 9 & 9 & 11 \\
$X\to YY\to ccgg$ & 39 & 22 & 15 & 9 & 10 & 12 & 4 & 13 & 4 & 13 & 13 & 13 \\
$X\to YY\to ccsq$ & 39 & 22 & 15 & 9 & 9 & 0 & 0 & 9 & 9 & 9 & 9 & 9 \\
$X\to YY\to qqqc$ & 40 & 22 & 15 & 9 & 0 & 9 & 0 & 0 & 11 & 9 & 9 & 11 \\
$X\to YY\to sscq$ & 40 & 22 & 15 & 9 & 9 & 0 & 0 & 9 & 9 & 9 & 9 & 11 \\
$X\to YY\to qqcs$ & 40 & 22 & 15 & 9 & 0 & 9 & 11 & 9 & 11 & 9 & 9 & 9 \\
$X\to YY\to ssss$ & 41 & 23 & 16 & 9 & 9 & 0 & 11 & 9 & 11 & 9 & 9 & 11 \\
$X\to YY\to ssqq$ & 41 & 23 & 16 & 9 & 9 & 0 & 0 & 9 & 9 & 9 & 9 & 9 \\
$X\to YY\to ssgg$ & 41 & 23 & 16 & 9 & 3 & 15 & 3 & 7 & 13 & 13 & 13 & 13 \\
$X\to YY\to qqqq$ & 41 & 23 & 16 & 9 & 0 & 9 & 11 & 9 & 9 & 9 & 9 & 9 \\
$X\to YY\to qqgg$ & 41 & 23 & 16 & 9 & 12 & 10 & 13 & 12 & 12 & 13 & 13 & 13 \\
$X\to YY\to gggg$ & 41 & 23 & 16 & 9 & 4 & 14 & 12 & 12 & 9 & 4 & 4 & 4 \\
$X\to YY\to qqqs$ & 41 & 23 & 16 & 9 & 9 & 0 & 11 & 0 & 11 & 9 & 9 & 9 \\
\addlinespace[2pt]
$X\to YY\to bbee$ & 45 & 27 & 19 & 11 & 9 & 12 & 0 & 4 & 11 & 3 & 3 & 3 \\
$X\to YY\to bb\mu\mu$ & 45 & 27 & 19 & 11 & 11 & 11 & 0 & 11 & 9 & 11 & 9 & 9 \\
$X\to YY\to ccee$ & 46 & 28 & 19 & 11 & 11 & 11 & 0 & 11 & 9 & 3 & 3 & 3 \\
$X\to YY\to cc\mu\mu$ & 46 & 28 & 19 & 11 & 11 & 11 & 11 & 4 & 9 & 11 & 9 & 9 \\
$X\to YY\to ssee$ & 47 & 29 & 20 & 11 & 9 & 13 & 0 & 9 & 9 & 3 & 3 & 3 \\
$X\to YY\to ss\mu\mu$ & 47 & 29 & 20 & 11 & 11 & 13 & 0 & 11 & 13 & 11 & 9 & 9 \\
$X\to YY\to qqee$ & 47 & 29 & 20 & 11 & 11 & 11 & 11 & 11 & 11 & 3 & 3 & 3 \\
$X\to YY\to qq\mu\mu$ & 47 & 29 & 20 & 11 & 11 & 11 & 0 & 9 & 9 & 11 & 9 & 9 \\
$X\to YY\to ggee$ & 47 & 29 & 20 & 11 & 9 & 11 & 9 & 9 & 4 & 4 & 4 & 4 \\
$X\to YY\to gg\mu\mu$ & 47 & 29 & 20 & 11 & 10 & 11 & 3 & 11 & 3 & 9 & 9 & 9 \\
\addlinespace[2pt]
$X\to YY\to bb\tau_h\tau_e$ & 48 & 30 & 21 & 12 & 14 & 13 & 3 & 10 & 12 & 13 & 13 & 13 \\
$X\to YY\to bb\tau_h\tau_\mu$ & 48 & 30 & 21 & 12 & 4 & 9 & 3 & 15 & 7 & 4 & 4 & 4 \\
$X\to YY\to cc\tau_h\tau_e$ & 49 & 31 & 21 & 12 & 15 & 2 & 10 & 1 & 9 & 13 & 13 & 13 \\
$X\to YY\to cc\tau_h\tau_\mu$ & 49 & 31 & 21 & 12 & 12 & 3 & 4 & 10 & 0 & 4 & 4 & 4 \\
$X\to YY\to ss\tau_h\tau_e$ & 50 & 32 & 22 & 12 & 15 & 10 & 12 & 6 & 12 & 13 & 13 & 13 \\
$X\to YY\to ss\tau_h\tau_\mu$ & 50 & 32 & 22 & 12 & 12 & 3 & 12 & 6 & 14 & 4 & 4 & 4 \\
$X\to YY\to qq\tau_h\tau_e$ & 50 & 32 & 22 & 12 & 2 & 4 & 4 & 12 & 1 & 13 & 13 & 13 \\
$X\to YY\to qq\tau_h\tau_\mu$ & 50 & 32 & 22 & 12 & 4 & 0 & 12 & 12 & 4 & 4 & 4 & 4 \\
$X\to YY\to gg\tau_h\tau_e$ & 50 & 32 & 22 & 12 & 12 & 4 & 10 & 1 & 14 & 1 & 1 & 1 \\
$X\to YY\to gg\tau_h\tau_\mu$ & 50 & 32 & 22 & 12 & 2 & 6 & 2 & 3 & 12 & 14 & 14 & 14 \\
\addlinespace[2pt]
$X\to YY\to bb\tau_h\tau_h$ & 42 & 24 & 17 & 10 & 4 & 7 & 4 & 14 & 15 & 11 & 9 & 9 \\
$X\to YY\to cc\tau_h\tau_h$ & 43 & 25 & 17 & 10 & 12 & 13 & 4 & 3 & 15 & 11 & 9 & 9 \\
$X\to YY\to ss\tau_h\tau_h$ & 44 & 26 & 18 & 10 & 13 & 4 & 15 & 2 & 15 & 11 & 9 & 9 \\
$X\to YY\to qq\tau_h\tau_h$ & 44 & 26 & 18 & 10 & 13 & 11 & 4 & 9 & 10 & 11 & 9 & 9 \\
$X\to YY\to gg\tau_h\tau_h$ & 44 & 26 & 18 & 10 & 1 & 11 & 7 & 13 & 10 & 1 & 1 & 1 \\
\addlinespace[2pt]
$X\to YY\to bbb$ & 9 & 6 & 5 & 4 & 7 & 1 & 3 & 15 & 10 & 10 & 10 & 10 \\
$X\to YY\to bbc$ & 10 & 6 & 5 & 4 & 4 & 1 & 13 & 10 & 1 & 10 & 10 & 10 \\
$X\to YY\to bbs$ & 11 & 6 & 5 & 4 & 3 & 4 & 10 & 15 & 6 & 10 & 10 & 10 \\
$X\to YY\to bbq$ & 11 & 6 & 5 & 4 & 4 & 15 & 10 & 6 & 1 & 10 & 10 & 10 \\
$X\to YY\to bbg$ & 11 & 6 & 5 & 4 & 6 & 0 & 1 & 14 & 13 & 14 & 14 & 14 \\
$X\to YY\to ccb$ & 12 & 6 & 5 & 4 & 6 & 12 & 14 & 3 & 2 & 10 & 10 & 10 \\
$X\to YY\to ssb$ & 13 & 6 & 5 & 4 & 10 & 4 & 13 & 10 & 6 & 10 & 10 & 10 \\
$X\to YY\to qqb$ & 13 & 6 & 5 & 4 & 9 & 12 & 12 & 10 & 3 & 10 & 10 & 10 \\
$X\to YY\to ggb$ & 13 & 6 & 5 & 4 & 12 & 7 & 7 & 1 & 13 & 6 & 6 & 6 \\
$X\to YY\to bsq$ & 13 & 6 & 5 & 4 & 3 & 10 & 12 & 15 & 7 & 10 & 10 & 10 \\
$X\to YY\to bcs$ & 14 & 6 & 5 & 4 & 1 & 13 & 15 & 3 & 10 & 10 & 10 & 10 \\
$X\to YY\to bcq$ & 14 & 6 & 5 & 4 & 9 & 12 & 3 & 3 & 10 & 10 & 10 & 10 \\
$X\to YY\to ccc$ & 15 & 7 & 5 & 4 & 3 & 1 & 10 & 3 & 6 & 10 & 10 & 10 \\
$X\to YY\to ccs$ & 16 & 7 & 5 & 4 & 7 & 9 & 13 & 12 & 7 & 10 & 10 & 10 \\
$X\to YY\to ccq$ & 16 & 7 & 5 & 4 & 7 & 4 & 10 & 10 & 7 & 10 & 10 & 10 \\
$X\to YY\to ccg$ & 16 & 7 & 5 & 4 & 12 & 6 & 14 & 12 & 15 & 14 & 14 & 14 \\
$X\to YY\to ssc$ & 17 & 7 & 5 & 4 & 10 & 13 & 4 & 12 & 10 & 10 & 10 & 10 \\
$X\to YY\to qqc$ & 17 & 7 & 5 & 4 & 3 & 1 & 4 & 13 & 15 & 10 & 10 & 10 \\
$X\to YY\to ggc$ & 17 & 7 & 5 & 4 & 6 & 2 & 1 & 14 & 0 & 6 & 6 & 6 \\
$X\to YY\to csq$ & 17 & 7 & 5 & 4 & 7 & 1 & 15 & 15 & 6 & 10 & 10 & 10 \\
$X\to YY\to sss$ & 18 & 8 & 6 & 4 & 9 & 12 & 9 & 10 & 3 & 10 & 10 & 10 \\
$X\to YY\to ssq$ & 18 & 8 & 6 & 4 & 3 & 1 & 6 & 10 & 15 & 10 & 10 & 10 \\
$X\to YY\to ssg$ & 18 & 8 & 6 & 4 & 6 & 10 & 14 & 12 & 14 & 14 & 14 & 14 \\
$X\to YY\to qqs$ & 18 & 8 & 6 & 4 & 4 & 15 & 10 & 10 & 6 & 10 & 10 & 10 \\
$X\to YY\to qqq$ & 18 & 8 & 6 & 4 & 4 & 4 & 13 & 3 & 7 & 10 & 10 & 10 \\
$X\to YY\to qqg$ & 18 & 8 & 6 & 4 & 12 & 7 & 7 & 13 & 0 & 14 & 14 & 14 \\
$X\to YY\to ggs$ & 18 & 8 & 6 & 4 & 12 & 6 & 14 & 1 & 13 & 6 & 6 & 6 \\
$X\to YY\to ggq$ & 18 & 8 & 6 & 4 & 1 & 2 & 2 & 14 & 4 & 6 & 6 & 6 \\
$X\to YY\to ggg$ & 18 & 8 & 6 & 4 & 14 & 14 & 6 & 9 & 2 & 2 & 2 & 2 \\
\addlinespace[2pt]
$X\to YY\to bbe$ & 22 & 12 & 9 & 6 & 7 & 12 & 10 & 10 & 10 & 14 & 14 & 14 \\
$X\to YY\to bb\mu$ & 22 & 12 & 9 & 6 & 12 & 12 & 12 & 10 & 4 & 0 & 0 & 0 \\
$X\to YY\to cce$ & 23 & 13 & 9 & 6 & 14 & 4 & 1 & 3 & 13 & 14 & 14 & 14 \\
$X\to YY\to cc\mu$ & 23 & 13 & 9 & 6 & 10 & 0 & 14 & 10 & 13 & 0 & 0 & 0 \\
$X\to YY\to sse$ & 24 & 14 & 10 & 6 & 12 & 4 & 14 & 15 & 10 & 14 & 14 & 14 \\
$X\to YY\to ss\mu$ & 24 & 14 & 10 & 6 & 4 & 10 & 6 & 12 & 0 & 0 & 0 & 0 \\
$X\to YY\to qqe$ & 24 & 14 & 10 & 6 & 7 & 6 & 1 & 15 & 15 & 14 & 14 & 14 \\
$X\to YY\to qq\mu$ & 24 & 14 & 10 & 6 & 6 & 6 & 2 & 10 & 14 & 0 & 0 & 0 \\
$X\to YY\to gge$ & 24 & 14 & 10 & 6 & 3 & 7 & 7 & 12 & 15 & 2 & 2 & 2 \\
$X\to YY\to gg\mu$ & 24 & 14 & 10 & 6 & 6 & 2 & 14 & 1 & 13 & 2 & 2 & 2 \\
\addlinespace[2pt]
$X\to YY\to bee$ & 25 & 15 & 11 & 7 & 14 & 10 & 1 & 10 & 5 & 2 & 2 & 2 \\
$X\to YY\to b\mu\mu$ & 25 & 15 & 11 & 7 & 4 & 4 & 6 & 10 & 2 & 4 & 4 & 4 \\
$X\to YY\to cee$ & 26 & 16 & 11 & 7 & 10 & 10 & 6 & 10 & 4 & 2 & 2 & 2 \\
$X\to YY\to c\mu\mu$ & 26 & 16 & 11 & 7 & 5 & 10 & 1 & 10 & 15 & 4 & 4 & 4 \\
$X\to YY\to see$ & 27 & 17 & 12 & 7 & 14 & 14 & 6 & 9 & 15 & 2 & 2 & 2 \\
$X\to YY\to qee$ & 27 & 17 & 12 & 7 & 10 & 10 & 4 & 10 & 2 & 2 & 2 & 2 \\
$X\to YY\to gee$ & 27 & 17 & 12 & 7 & 10 & 14 & 5 & 9 & 2 & 5 & 5 & 5 \\
$X\to YY\to s\mu\mu$ & 27 & 17 & 12 & 7 & 10 & 14 & 5 & 9 & 15 & 4 & 4 & 4 \\
$X\to YY\to q\mu\mu$ & 27 & 17 & 12 & 7 & 5 & 10 & 6 & 4 & 15 & 4 & 4 & 4 \\
$X\to YY\to g\mu\mu$ & 27 & 17 & 12 & 7 & 14 & 5 & 1 & 9 & 15 & 5 & 5 & 5 \\
\addlinespace[2pt]
$X\to YY\to b\tau_h\tau_e$ & 28 & 18 & 13 & 8 & 8 & 5 & 1 & 5 & 2 & 8 & 8 & 8 \\
$X\to YY\to b\tau_h\tau_\mu$ & 28 & 18 & 13 & 8 & 10 & 5 & 8 & 8 & 5 & 8 & 8 & 8 \\
$X\to YY\to c\tau_h\tau_e$ & 29 & 19 & 13 & 8 & 10 & 5 & 1 & 5 & 2 & 8 & 8 & 8 \\
$X\to YY\to c\tau_h\tau_\mu$ & 29 & 19 & 13 & 8 & 14 & 14 & 1 & 5 & 2 & 8 & 8 & 8 \\
$X\to YY\to s\tau_h\tau_e$ & 30 & 20 & 14 & 8 & 10 & 14 & 1 & 5 & 5 & 8 & 8 & 8 \\
$X\to YY\to q\tau_h\tau_e$ & 30 & 20 & 14 & 8 & 10 & 8 & 1 & 5 & 5 & 8 & 8 & 8 \\
$X\to YY\to g\tau_h\tau_e$ & 30 & 20 & 14 & 8 & 14 & 8 & 1 & 8 & 8 & 8 & 8 & 8 \\
$X\to YY\to s\tau_h\tau_\mu$ & 30 & 20 & 14 & 8 & 14 & 14 & 1 & 5 & 5 & 8 & 8 & 8 \\
$X\to YY\to q\tau_h\tau_\mu$ & 30 & 20 & 14 & 8 & 14 & 5 & 1 & 5 & 8 & 8 & 8 & 8 \\
$X\to YY\to g\tau_h\tau_\mu$ & 30 & 20 & 14 & 8 & 8 & 14 & 8 & 5 & 2 & 8 & 8 & 8 \\
\addlinespace[2pt]
$X\to YY\to b\tau_h\tau_h$ & 19 & 9 & 7 & 5 & 8 & 14 & 8 & 8 & 8 & 4 & 4 & 4 \\
$X\to YY\to c\tau_h\tau_h$ & 20 & 10 & 7 & 5 & 14 & 8 & 8 & 5 & 8 & 4 & 4 & 4 \\
$X\to YY\to s\tau_h\tau_h$ & 21 & 11 & 8 & 5 & 8 & 8 & 8 & 8 & 8 & 4 & 4 & 4 \\
$X\to YY\to q\tau_h\tau_h$ & 21 & 11 & 8 & 5 & 10 & 5 & 1 & 8 & 5 & 4 & 4 & 4 \\
$X\to YY\to g\tau_h\tau_h$ & 21 & 11 & 8 & 5 & 8 & 8 & 1 & 5 & 2 & 4 & 4 & 4 \\
\addlinespace[2pt]
$X\to YY\to bce\nu$ & 51 & 33 & 23 & 13 & 9 & 0 & 10 & 11 & 3 & 11 & 11 & 11 \\
$X\to YY\to bc\mu\nu$ & 51 & 33 & 23 & 13 & 9 & 13 & 0 & 4 & 13 & 9 & 9 & 9 \\
$X\to YY\to bqe\nu$ & 52 & 33 & 23 & 13 & 0 & 4 & 13 & 11 & 12 & 11 & 11 & 11 \\
$X\to YY\to bq\mu\nu$ & 52 & 33 & 23 & 13 & 15 & 13 & 10 & 11 & 3 & 9 & 9 & 9 \\
$X\to YY\to cse\nu$ & 53 & 34 & 23 & 13 & 13 & 15 & 0 & 0 & 3 & 11 & 11 & 11 \\
$X\to YY\to cqe\nu$ & 53 & 34 & 23 & 13 & 15 & 3 & 15 & 15 & 10 & 11 & 11 & 11 \\
$X\to YY\to cs\mu\nu$ & 53 & 34 & 23 & 13 & 9 & 9 & 15 & 3 & 3 & 9 & 9 & 9 \\
$X\to YY\to cq\mu\nu$ & 53 & 34 & 23 & 13 & 15 & 4 & 10 & 3 & 13 & 9 & 9 & 9 \\
$X\to YY\to sqe\nu$ & 54 & 35 & 24 & 13 & 13 & 0 & 15 & 15 & 3 & 11 & 11 & 11 \\
$X\to YY\to qqe\nu$ & 54 & 35 & 24 & 13 & 13 & 4 & 15 & 4 & 3 & 11 & 11 & 11 \\
$X\to YY\to sq\mu\nu$ & 54 & 35 & 24 & 13 & 13 & 13 & 0 & 3 & 3 & 9 & 9 & 9 \\
$X\to YY\to qq\mu\nu$ & 54 & 35 & 24 & 13 & 9 & 11 & 3 & 4 & 13 & 9 & 9 & 9 \\
\addlinespace[2pt]
$X\to YY\to bc\tau_e\nu$ & 55 & 36 & 25 & 14 & 7 & 0 & 7 & 14 & 14 & 6 & 6 & 6 \\
$X\to YY\to bc\tau_\mu\nu$ & 55 & 36 & 25 & 14 & 6 & 12 & 12 & 3 & 10 & 9 & 9 & 9 \\
$X\to YY\to bq\tau_e\nu$ & 56 & 36 & 25 & 14 & 14 & 10 & 6 & 1 & 4 & 6 & 6 & 6 \\
$X\to YY\to bq\tau_\mu\nu$ & 56 & 36 & 25 & 14 & 6 & 7 & 2 & 10 & 4 & 9 & 9 & 9 \\
$X\to YY\to cs\tau_e\nu$ & 57 & 37 & 25 & 14 & 4 & 6 & 14 & 12 & 14 & 6 & 6 & 6 \\
$X\to YY\to cq\tau_e\nu$ & 57 & 37 & 25 & 14 & 3 & 4 & 14 & 15 & 15 & 6 & 6 & 6 \\
$X\to YY\to cs\tau_\mu\nu$ & 57 & 37 & 25 & 14 & 7 & 12 & 2 & 13 & 0 & 9 & 9 & 9 \\
$X\to YY\to cq\tau_\mu\nu$ & 57 & 37 & 25 & 14 & 6 & 6 & 1 & 10 & 15 & 9 & 9 & 9 \\
$X\to YY\to sq\tau_e\nu$ & 58 & 38 & 26 & 14 & 1 & 4 & 14 & 12 & 10 & 6 & 6 & 6 \\
$X\to YY\to qq\tau_e\nu$ & 58 & 38 & 26 & 14 & 12 & 2 & 14 & 10 & 15 & 6 & 6 & 6 \\
$X\to YY\to sq\tau_\mu\nu$ & 58 & 38 & 26 & 14 & 10 & 10 & 10 & 10 & 13 & 9 & 9 & 9 \\
$X\to YY\to qq\tau_\mu\nu$ & 58 & 38 & 26 & 14 & 12 & 0 & 1 & 15 & 13 & 9 & 9 & 9 \\
\addlinespace[2pt]
$X\to YY\to bc\tau_h\nu$ & 59 & 39 & 27 & 15 & 0 & 11 & 3 & 9 & 9 & 12 & 12 & 12 \\
$X\to YY\to bq\tau_h\nu$ & 60 & 39 & 27 & 15 & 13 & 11 & 15 & 15 & 9 & 12 & 12 & 12 \\
$X\to YY\to cs\tau_h\nu$ & 61 & 40 & 27 & 15 & 13 & 11 & 13 & 9 & 3 & 12 & 12 & 12 \\
$X\to YY\to cq\tau_h\nu$ & 61 & 40 & 27 & 15 & 13 & 3 & 0 & 0 & 9 & 12 & 12 & 12 \\
$X\to YY\to sq\tau_h\nu$ & 62 & 41 & 28 & 15 & 13 & 11 & 9 & 0 & 10 & 12 & 12 & 12 \\
$X\to YY\to qq\tau_h\nu$ & 62 & 41 & 28 & 15 & 0 & 15 & 0 & 15 & 12 & 12 & 12 & 12 \\
\addlinespace[2pt]
QCD, $bbccss$ & 63 & 42 & 29 & 16 & 16 & 16 & 16 & 16 & 16 & 16 & 16 & 16 \\
QCD, $bbccs$ & 63 & 42 & 29 & 16 & 16 & 16 & 16 & 16 & 16 & 16 & 16 & 16 \\
QCD, $bbcc$ & 63 & 42 & 29 & 16 & 16 & 16 & 16 & 16 & 16 & 16 & 16 & 16 \\
QCD, $bbcss$ & 63 & 42 & 29 & 16 & 16 & 16 & 16 & 16 & 16 & 16 & 16 & 16 \\
QCD, $bbcs$ & 63 & 42 & 29 & 16 & 16 & 16 & 16 & 16 & 16 & 16 & 16 & 16 \\
QCD, $bbc$ & 63 & 42 & 29 & 16 & 16 & 16 & 16 & 16 & 16 & 16 & 16 & 16 \\
QCD, $bbss$ & 63 & 42 & 29 & 16 & 16 & 16 & 16 & 16 & 16 & 16 & 16 & 16 \\
QCD, $bbs$ & 63 & 42 & 29 & 16 & 16 & 16 & 16 & 16 & 16 & 16 & 16 & 16 \\
QCD, $bb$ & 63 & 42 & 29 & 16 & 16 & 16 & 16 & 16 & 16 & 16 & 16 & 16 \\
QCD, $bccss$ & 63 & 42 & 29 & 16 & 16 & 16 & 16 & 16 & 16 & 16 & 16 & 16 \\
QCD, $bccs$ & 63 & 42 & 29 & 16 & 16 & 16 & 16 & 16 & 16 & 16 & 16 & 16 \\
QCD, $bcc$ & 63 & 42 & 29 & 16 & 16 & 16 & 16 & 16 & 16 & 16 & 16 & 16 \\
QCD, $bcss$ & 63 & 42 & 29 & 16 & 16 & 16 & 16 & 16 & 16 & 16 & 16 & 16 \\
QCD, $bcs$ & 63 & 42 & 29 & 16 & 16 & 16 & 16 & 16 & 16 & 16 & 16 & 16 \\
QCD, $bc$ & 63 & 42 & 29 & 16 & 16 & 16 & 16 & 16 & 16 & 16 & 16 & 16 \\
QCD, $bss$ & 63 & 42 & 29 & 16 & 16 & 16 & 16 & 16 & 16 & 16 & 16 & 16 \\
QCD, $bs$ & 63 & 42 & 29 & 16 & 16 & 16 & 16 & 16 & 16 & 16 & 16 & 16 \\
QCD, $b$ & 63 & 42 & 29 & 16 & 16 & 16 & 16 & 16 & 16 & 16 & 16 & 16 \\
QCD, $ccss$ & 63 & 42 & 29 & 16 & 16 & 16 & 16 & 16 & 16 & 16 & 16 & 16 \\
QCD, $ccs$ & 63 & 42 & 29 & 16 & 16 & 16 & 16 & 16 & 16 & 16 & 16 & 16 \\
QCD, $cc$ & 63 & 42 & 29 & 16 & 16 & 16 & 16 & 16 & 16 & 16 & 16 & 16 \\
QCD, $css$ & 63 & 42 & 29 & 16 & 16 & 16 & 16 & 16 & 16 & 16 & 16 & 16 \\
QCD, $cs$ & 63 & 42 & 29 & 16 & 16 & 16 & 16 & 16 & 16 & 16 & 16 & 16 \\
QCD, $c$ & 63 & 42 & 29 & 16 & 16 & 16 & 16 & 16 & 16 & 16 & 16 & 16 \\
QCD, $ss$ & 63 & 42 & 29 & 16 & 16 & 16 & 16 & 16 & 16 & 16 & 16 & 16 \\
QCD, $s$ & 63 & 42 & 29 & 16 & 16 & 16 & 16 & 16 & 16 & 16 & 16 & 16 \\
QCD, light & 63 & 42 & 29 & 16 & 16 & 16 & 16 & 16 & 16 & 16 & 16 & 16 \\
\end{longtable}
}


Three rules fixed before any pretraining select the random partitions. A draw assigns the resonant classes at random to as many groups as the
\Kcoarse-class level has resonant groups, plus one group for all QCD jets, and matches each group to a
resonant group of the \Kcoarse-class level whose share of the reweighted training jets it reproduces
within \RandVtwoShareTol, so that every group carries the weight in the loss of the group it replaces.
Each of \RandVtwoTrials{} draws is then refined by a search that lowers the number of class pairs it
merges that the \Kcoarse-class level also merges, and the draw with the fewest such pairs becomes a
candidate. Of \RandVtwoPoolDraws{} candidates built from consecutive seeds, \RandVtwoPoolPass{} keep
every realized share within the tolerance after the refinement. The \RandVtwoNPartitions{} partitions
are a random sample of these that meets a balance rule: each of \RandVtwoNPairs{} class pairs, among
them the pairs of the linear probes on class pairs, is merged in \RandVtwoMergedMin{} or
\RandVtwoMergedMax{} partitions and kept apart in the others. The balance rule makes every one of these
pairs available both merged and kept apart, which the comparison of merged against separated classes
needs. Because the refinement pushes the partitions away from the grouping by decay topology, a group can
mix decays with different numbers of prongs and the partitions keep many flavor boundaries;
table~\ref{tab:randprong} lists the prong-count boundaries each one keeps. Each partition is pretrained
once, the first and third with the initial weights and training examples of seed one and the other three
with those of seed two, and is compared with the \Kfull-class model of its seed.
% CHECK: RandVtwoRuns = one; partitions p1, p3 at seed 1 and p2, p4, p5 at seed 2 (configs/arms/v2_paper_runs.txt:23-27)

The three random partitions with controlled flavor boundaries share one grouping and differ only in where
the four-quark decays \RandFlavSplitOrbit{} go. In the partition with no flavor boundary, all decays that
differ only in quark flavor share a group. The partition with a $b$-quark boundary starts from it and
moves the \RandFlavCutClasses{} of the \RandFlavSplitOrbitClasses{} four-quark decays that contain a $b$
quark into another group, while the \RandFlavOrbitClassesMoved{} decays of the forms
\RandFlavOrbitsMovedList{} move in the opposite direction, so that every group keeps its share exactly.
Its only flavor boundary lies between the four-quark decays with and without a $b$ quark, and it keeps
\RandFlavFoneBcSplit{} of the \RandFlavBcPairs{} pairs of four-quark decays that differ by exchanging
$b$ and $c$ quarks apart. The partition with a random flavor boundary makes the same compensating move but
moves \RandFlavCutClasses{} four-quark decays drawn at random, \RandFlavFonerCutWithB{} of them with a $b$
quark, under rules fixed before the draw: about as many of the moved decays contain a $b$ quark as do not;
no pair that differs by exchanging $b$ and $c$ quarks is kept apart; the pairs of the balance rule are
merged or kept apart as in the partition with no flavor boundary; and every group's share of the training
examples stays within \RandFlavFonerShareTol{} of that in the partition with a $b$-quark boundary. It
therefore keeps \RandFlavFonerBcSplit{} of the $b$--$c$ exchange pairs apart, and the two partitions differ
in whether the moved decays are chosen by flavor, not in how many move. Each of the three is pretrained
with \RandFlavRuns{} seeds, one and two.

% ----------------------------------------------------------------------------
\section{Training settings}\label{sec:app-settings}
% ----------------------------------------------------------------------------

The settings follow the order of Sophon's appendix~B~\cite{sophon}: architecture, inputs, reweighting,
optimization and training length, then the checkpoint rule. Each subsection after the first gives only
what differs from pretraining, so that every model can be rebuilt from the data and the code release.

\subsection{Pretraining at every label set}\label{sec:app-pretrain}

Every model is the ParT network~\cite{part} in the configuration of Sophon's released
code~\cite{sophoncode}, as implemented in weaver-core \DesignWeaverVersion{}, the training framework of
ParT and Sophon. It has \DesignParticleBlocks{} particle-attention blocks with a pairwise interaction
bias, then \DesignClassBlocks{} class-attention blocks, with \DesignEmbedDim-dimensional particle
representations and \DesignHeads{} attention heads. The head has one hidden layer of \DesignHeadHidden{}
units with a ReLU activation and dropout of \MlstFtHeadDropout{} after it, followed by an output layer as
wide as the label set.
% CHECK: MlstFtHeadDropout (fine-tuning spec) equals the pretraining head dropout, fc_params (512, 0.1); interface request DesignHeadDropout
Where Sophon's paper and its released code differ, in the number of particle-attention blocks
(section~\ref{sec:data}) and in the paper's GELU activation throughout~\cite[App.~B.1]{sophon}, we follow
the code.

The initial weights follow PyTorch's default initialization, except for the class token, which is drawn
from a narrow truncated normal distribution. The weaver-core version used by Sophon's code instead draws
the linear layers from a Xavier-uniform distribution with zero biases and rescales the residual
projections by depth; we did not reproduce that initialization (section~\ref{sec:data}).

Each particle carries \FillNFeatures{} input features: \FillTopNPresent{} computed from its
four-momentum and that of the jet, its electric charge, five binary flags for its type (charged hadron,
neutral hadron, photon, electron and muon) and four impact-parameter features, the transverse and
longitudinal impact parameters with their uncertainties. The pairwise features of the attention bias are
computed from the four-momenta, which are rescaled so that every jet has the same transverse momentum, as
in Sophon~\cite[Sec.~II]{sophon}. The features first pass a BatchNorm layer whose running mean and
variance are accumulated in pretraining; the substitutions of appendix~\ref{sec:app-inputs} are defined
with respect to this normalization.

The training jets are drawn with Sophon's reweighting in $(p_{\mathrm{T}}, \msd)$ within each
of its reweighting categories of the \Kfull{} classes~\cite[Sec.~II]{sophon}, with the same weights for
every label set.

Training proceeds in cycles of \DesignVtwoNWindows{} epochs. In each cycle the rows of every training
file are divided at random into \DesignVtwoNWindows{} disjoint parts, and each epoch reads one part, or
\DesignVtwoDataFraction{} of the rows, of nearly every file; a row is therefore read in one epoch per
cycle, and within that epoch the reweighting can draw it several times. The order of the training
examples and the seeds of the dropout masks are functions of the seed and the epoch number alone, which is
what makes models with the same seed see the same training examples in the same order whatever their
labels.

Every model is trained for \DesignEpochs{} epochs of \DesignSamplesPerEpoch{} jets, \DesignExamplesSeen{}
training examples in all, in batches of \DesignBatchSize{} with mixed precision. Sophon's paper gives
\LitSophonSamplesPerEpochMillion{} million jets per epoch~\cite[App.~B.1]{sophon}; the exact number is
that of its code~\cite{sophoncode}. The optimizer is Ranger: Lookahead~\cite{lookahead} with
$k = \LitSophonLookaheadK$ and $\alpha = \LitSophonLookaheadAlpha$ over RAdam~\cite{radam} with
$\beta_1 = \LitSophonRadamBetaOne$, $\beta_2 = \LitSophonRadamBetaTwo$ and $\epsilon$ of
\LitSophonRadamEps, the values of Sophon~\cite[App.~B.1]{sophon}. There is no weight decay, as in
Sophon's code; Sophon applies weight decay only to its ParT reference~\cite[App.~B.3]{sophon}. The
learning rate is \DesignStartLr{} for the first \DesignLrFlatPercent{} of the epochs. From epoch
\DesignLrDecayFirstEpoch{}, counting from zero, it falls by a constant factor at every epoch, reaching
\DesignEndLr{}, \DesignLrEndPercent{} of its initial value, in the last epoch.

The validation accuracy is computed after every epoch on every selected jet of \DesignVtwoNValFiles{}
validation files that no run trains on. The checkpoint evaluated throughout this paper is the first
maximum of the validation accuracy among epochs \DesignWavgFirstEpoch--\DesignWavgLastEpoch, where the
learning rate is lowest, so that every model is read at the same stage of its schedule; Sophon keeps the
checkpoint with the highest validation accuracy~\cite[App.~B.1]{sophon}. For the
check of appendix~\ref{sec:app-bn}, a second copy of each selected checkpoint has its BatchNorm
statistics recomputed on \DesignBnJets{} training jets, and the saved initial weights of the randomly
initialized network are treated the same way.

\subsection{Auxiliary mass regression and kinematic inputs only}\label{sec:app-variants}

The models with the auxiliary mass regression differ from the others by one output, which regresses the
target $\theta$ of equation~(\ref{eq:masstarget}) with a log-cosh loss, and one loss term. A jet without a
generator-level match has a stored generator mass of exactly zero and is left out of the regression term,
though not out of the classification. The total loss is $L = \Lcls + \lambda\,\Lreg$ of
equation~(\ref{eq:masstarget}) with $\lambda = \MlstMassLambda$ at both label sets. Every other setting is that of appendix~\ref{sec:app-pretrain}, and each of these
models shares its initial weights, except the output layer, and its training examples with the model of
the same label set and seed trained without the regression.

The \Kfull-class model with kinematic inputs only keeps the \FillTopNPresent{} kinematic features of
each particle and the pairwise features, and drops the charge, the type flags and the impact-parameter
features. It is trained once, on the training examples of seed one in their order and with the settings of
appendix~\ref{sec:app-pretrain}. Its first layer is narrower than that of the other models, so its
initial weights are its own.
% CHECK: the kinematic-only run reads the seed-1 example order (mtx2-l188kin-s1, RUNS.csv:1200)

\subsection{Linear and MLP probes}\label{sec:app-probes}

Each linear probe, equation~(\ref{eq:probe}), is fitted as a multinomial logistic regression on the
standardized frozen features by full-batch L-BFGS in double precision. A fit counts as converged when the largest component of its
gradient falls below \LpConvergedGrad, and its inverse regularization strength is chosen on validation
jets from a logarithmic grid between \LpCMin{} and \LpVfourCMax.
The frozen features are stored in half precision.
Each MLP probe has one hidden layer of \LpMlpHidden{} units and is trained on the same frozen features at
the largest training set of each task; its result is the average over \LpMlpSeedsWord{} seeds.

On JetClass-II the linear probes are scored on a class-balanced test set of \JcBalNTest{} jets from the
test files: \JcBalTestPerClass{} jets of each class, or every selected jet of the \JcBalNShort{} classes
that have fewer, the rarest with \JcBalRarestTest. The training sets hold equal numbers of jets per class
from files that no pretraining run trained or validated on, and they are nested from \FtSizeEThree{} to
\FtSizeESix{} jets, each contained in the next. At \FtSizeESix{} jets, \JcBalShortTrainESix{} classes
have fewer held-out jets than their share and contribute all they have. The regularization is chosen on
up to \JcBalValPerClass{} validation jets per class. The AUC averaged over classes takes each of the
\Kfull{} classes against the rest on every test jet, and table~\ref{tab:probes} also gives it over the
classes with a full test sample and over the resonant and the QCD classes.

The linear probes on class pairs are binary. They are trained on up to \JcBalPairTrainPerSide{} jets
per side and tested on up to \JcBalPairTestPerSide, each side at its natural composition, and choose
their regularization on \LpPairNValMin{} to \LpPairNValMax{} jets held out from the training jets. The pair
$X\to ee$ against $X\to\mu\mu$ is the exception, with the sets of section~\ref{sec:evaluation}.

On JetClass the linear probes are trained on up to \LpNTrainMaxJci{} jets, choose their regularization on
\LpNValJci{} validation jets and are scored on \LpNTestJci{} test jets, taken from the first test files of
each class with equal numbers per class; the AUC averaged over pairs of classes uses \LpNTestAucJci{} of
them. On top tagging the largest training set holds \LpNTrainMaxTop{} jets and the test set
\LpNTestTop; on quark/gluon discrimination they hold \LpNTrainMaxQg{} Pythia jets and \LpNTestQg. The Herwig
jets are only scored, by the linear probe trained on Pythia jets. On both benchmarks the missing input
features are substituted as appendix~\ref{sec:app-inputs} describes.
% CHECK: LpNTestTop (404,000 in fits_v3) against the landscape test split of 400,000 (inventory_facts 0.3)
% CHECK: Herwig linear probes reuse the Pythia-trained probe (as fine-tuning does)

The full-set linear probe on top is not limited by the number of training jets: for the \Kfull-class
model, its rejection at 30\% signal efficiency with \FtSizeEFive{} training jets and with the full set
agree within the spread over pretraining runs (table~\ref{tab:probesn}).
% CHECK: pending (no macro) top R30 of the 188-class linear probe at 1e5 vs full set within one SD over runs, chosen substitution

\subsection{Fine-tuning and training from scratch}\label{sec:app-ft}

Each pretrained model is fine-tuned once per task and training-set size, from its selected checkpoint,
whose hash every fine-tuning run records. The whole head, hidden layer included, is initialized afresh,
with dropout \MlstFtHeadDropout{} after the hidden layer. Every label set therefore starts fine-tuning
from its class token with the same head, whereas a pretrained hidden layer was trained to serve its own
label set's outputs. The optimizer is Ranger, as in pretraining, with mixed precision, and the missing
input features of top and quark/gluon are set to zero. On top tagging and quark/gluon discrimination all
runs of one size use the same training jets, so the \MlstFtRefSeedsWord{} networks trained from scratch
differ only in their training seed. On JetClass every pretrained model is fine-tuned on the first training
subset of each size, and the network trained from scratch with seed $k$ trains on the $k$-th, so its spread
includes that of the training subset.
% CHECK: on JetClass the from-scratch seed k reads subset k at every size (ref_v100/jc1/leg2/scratch-v2/N*/s{1,2,3}; ROC_BRIEF 6 states it for 10^6)

On the full training sets of top and quark/gluon, and on \FtSizeESix{} JetClass jets, fine-tuning uses
ParT's settings~\cite[Sec.~5.2]{part}. The pretrained weights train at a constant learning rate of
\MlstFtLr{} and the head at \MlstFtHeadMult{} times that rate, with weight decay \MlstFtWeightDecay{} and
batches of \MlstFtBatch{} jets. Training lasts \MlstFtBenchEpochs{} epochs on top and quark/gluon and
\MlstFtEpochsESix{} on JetClass, where the head's learning rate falls to \MlstFtHeadDecayTo{} of its
value over the last \MlstFtHeadDecayShare{} of the epochs. The epoch with the highest validation accuracy
is kept, the earliest in a tie. The networks trained from scratch on these sets follow ParT's settings
for training from scratch~\cite[Secs.~5.1 and~5.2]{part}: a learning rate of \MlstFtRefLr{} for every
weight, constant and then lowered to \MlstFtHeadDecayTo{} of its value over the last
\MlstFtHeadDecayShare{} of the epochs; weight decay \MlstFtRefWeightDecayBench{} on top and quark/gluon
and none on JetClass; and the epochs and batches of fine-tuning. The kinematic-only model is fine-tuned
on top with the same settings and has no missing input.

Below the full training sets, and on JetClass-II at every size, pretrained networks and networks
trained from scratch follow one rule. They train in batches of \FtBudgetBatch{} jets for at most
\FtBudgetSteps{} steps and are evaluated on a fixed validation set every \FtBudgetEvalSmall{} steps below
\FtBudgetEvalSmallBelow{} training jets and every \FtBudgetEval{} steps at and above that size. The
validation set holds \FtValJets{} jets, or on JetClass-II the \FtValJetsBal{} jets of the class-balanced
validation set. Training stops after \FtBudgetPatience{} evaluations without improvement, and the
checkpoint with the best validation metric is kept. The pretrained weights train at a constant learning
rate of \FtLrPre{} and the head at \FtHeadMult{} times that rate, at most \FtHeadCap. The networks trained
from scratch train every weight at \FtLrScratch{} with the training framework's default schedule, which is
constant and lowered only near the end of the maximum training length. The weight decay is
\FtBudgetWeightDecay{} throughout. Top and quark/gluon are fine-tuned on \FtSizeEThree, \FtSizeEFour{}
and \FtSizeEFive{} jets below their full sets; JetClass on \FtSizeEFour{} and \FtSizeEFive{}; and
class-balanced JetClass-II on \FtSizeEFour, \FtSizeEFive{} and \FtSizeESix{}. On JetClass-II the
fine-tuned models are the levels of the class hierarchy and the models with the auxiliary mass
regression.

The two learning rates of these smaller sets were chosen on validation (table~\ref{tab:lrsweep}). The
\Kfull-class model of seed one was fine-tuned with its pretrained weights at \FtSweepPre{}, and the
network trained from scratch at \FtSweepScratch{}, each on \FtSweepN{} jets of top and of quark/gluon.
The rate with the largest best validation accuracy summed over the two tasks was kept; the test jets play
no part. In both grids the chosen rate is the largest one swept, so the sweep does not bracket a
maximum, and the summed accuracies of the three rates lie close together.
% CHECK: FtLrPre and FtLrScratch are the largest values of FtSweepPre and FtSweepScratch (true in budget/lrsweep_metrics.json, 2026-11-02: sums 1.742/1.741/1.739 and 1.724/1.722/1.720)

\providecommand{\TableContinued}[2]{\label{#2}\ifnum\getpagerefnumber{#2}>\getpagerefnumber{#1}\relax{\footnotesize\textit{Table~\ref{#1} (continued)}}\par\smallskip\fi}
\begin{center}
\small
\ifdefined\tablabw\else\newlength{\tablabw}\newlength{\tablabtmp}\newlength{\tablabcol}\newlength{\tablabnat}\newlength{\tablabatom}\fi
\setlength{\tablabw}{\linewidth}\setlength{\tablabnat}{0pt}\setlength{\tablabatom}{0pt}
\setlength{\tablabcol}{1\linewidth}\addtolength{\tablabcol}{-2\tabcolsep}
\setlength{\tablabnat}{0pt}\settowidth{\tablabtmp}{learning rate}\ifdim\tablabtmp>\tablabnat\setlength{\tablabnat}{\tablabtmp}\fi\settowidth{\tablabtmp}{\ensuremath{3\times10^{-3}}}\ifdim\tablabtmp>\tablabnat\setlength{\tablabnat}{\tablabtmp}\fi\settowidth{\tablabtmp}{\ensuremath{3\times10^{-4}}}\ifdim\tablabtmp>\tablabnat\setlength{\tablabnat}{\tablabtmp}\fi\addtolength{\tablabcol}{-\tablabnat}\addtolength{\tablabcol}{-2\tabcolsep}
\setlength{\tablabnat}{0pt}\settowidth{\tablabtmp}{0.9188$^{\ast}$}\ifdim\tablabtmp>\tablabnat\setlength{\tablabnat}{\tablabtmp}\fi\settowidth{\tablabtmp}{0.9152}\ifdim\tablabtmp>\tablabnat\setlength{\tablabnat}{\tablabtmp}\fi\settowidth{\tablabtmp}{0.9177}\ifdim\tablabtmp>\tablabnat\setlength{\tablabnat}{\tablabtmp}\fi\addtolength{\tablabcol}{-\tablabnat}\addtolength{\tablabcol}{-2\tabcolsep}
\ifdim\tablabcol<\tablabw\setlength{\tablabw}{\tablabcol}\fi
\setlength{\tablabcol}{1\linewidth}\addtolength{\tablabcol}{-2\tabcolsep}
\setlength{\tablabnat}{0pt}\settowidth{\tablabtmp}{learning rate}\ifdim\tablabtmp>\tablabnat\setlength{\tablabnat}{\tablabtmp}\fi\settowidth{\tablabtmp}{\ensuremath{3\times10^{-3}}}\ifdim\tablabtmp>\tablabnat\setlength{\tablabnat}{\tablabtmp}\fi\settowidth{\tablabtmp}{\ensuremath{3\times10^{-4}}}\ifdim\tablabtmp>\tablabnat\setlength{\tablabnat}{\tablabtmp}\fi\addtolength{\tablabcol}{-\tablabnat}\addtolength{\tablabcol}{-2\tabcolsep}
\setlength{\tablabnat}{0pt}\settowidth{\tablabtmp}{quark/gluon}\ifdim\tablabtmp>\tablabnat\setlength{\tablabnat}{\tablabtmp}\fi\settowidth{\tablabtmp}{0.8059$^{\ast}$}\ifdim\tablabtmp>\tablabnat\setlength{\tablabnat}{\tablabtmp}\fi\settowidth{\tablabtmp}{0.8029}\ifdim\tablabtmp>\tablabnat\setlength{\tablabnat}{\tablabtmp}\fi\addtolength{\tablabcol}{-\tablabnat}\addtolength{\tablabcol}{-2\tabcolsep}
\setlength{\tablabnat}{0pt}\settowidth{\tablabtmp}{1.7423$^{\ast}$}\ifdim\tablabtmp>\tablabnat\setlength{\tablabnat}{\tablabtmp}\fi\settowidth{\tablabtmp}{1.7204}\ifdim\tablabtmp>\tablabnat\setlength{\tablabnat}{\tablabtmp}\fi\settowidth{\tablabtmp}{1.7217}\ifdim\tablabtmp>\tablabnat\setlength{\tablabnat}{\tablabtmp}\fi\addtolength{\tablabcol}{-\tablabnat}\addtolength{\tablabcol}{-2\tabcolsep}
\ifdim\tablabcol<\tablabw\setlength{\tablabw}{\tablabcol}\fi
\setlength{\tablabnat}{0pt}\settowidth{\tablabtmp}{\mbox{pretrained on \Kfull{} classes,} \mbox{fine-tuned,} \mbox{run 1}}\ifdim\tablabtmp>\tablabnat\setlength{\tablabnat}{\tablabtmp}\fi\settowidth{\tablabtmp}{\mbox{no pretraining,} \mbox{trained from scratch}}\ifdim\tablabtmp>\tablabnat\setlength{\tablabnat}{\tablabtmp}\fi
\settowidth{\tablabtmp}{pretrained on \Kfull{} classes,}\ifdim\tablabtmp>\tablabatom\setlength{\tablabatom}{\tablabtmp}\fi\settowidth{\tablabtmp}{trained from scratch}\ifdim\tablabtmp>\tablabatom\setlength{\tablabatom}{\tablabtmp}\fi\settowidth{\tablabtmp}{no pretraining,}\ifdim\tablabtmp>\tablabatom\setlength{\tablabatom}{\tablabtmp}\fi
\ifdim\tablabw>7cm\setlength{\tablabw}{7cm}\fi\ifdim\tablabnat<\tablabw\setlength{\tablabw}{\tablabnat}\fi\ifdim\tablabatom>\tablabw\setlength{\tablabw}{\tablabatom}\fi
\begin{minipage}{\linewidth}\centering
\captionof{table}{Learning-rate choice for fine-tuning Sophon's network, pretrained on JetClass-II with \Kfull{} classes (run 1) or not pretrained, on $10^{4}$ top and quark/gluon jets: the best validation accuracy of each rate on top and quark/gluon and their sum. Pretrained: the rate of the pretrained layers, with the new output head at 50 times that rate, at most $5\times10^{-3}$. $^{\ast}$: the rate chosen, on the sum for the pretrained weights and on each task's accuracy for the networks trained from scratch; the test jets play no part.}
\label{tab:lrsweep}
\par\medskip
\begin{adjustbox}{max width=\linewidth}
\begin{tabular}{>{\raggedright\arraybackslash}p{\tablabw}lc}
\toprule
 & learning rate & top \\
\midrule
\mbox{pretrained on \Kfull{} classes,} \mbox{fine-tuned,} \mbox{run 1} & \ensuremath{10^{-4}} & 0.9204 \\
\mbox{pretrained on \Kfull{} classes,} \mbox{fine-tuned,} \mbox{run 1} & \ensuremath{3\times10^{-4}} & 0.9192 \\
\mbox{pretrained on \Kfull{} classes,} \mbox{fine-tuned,} \mbox{run 1} & \ensuremath{10^{-3}} & 0.9212 \\
\mbox{no pretraining,} \mbox{trained from scratch} & \ensuremath{3\times10^{-4}} & 0.9152 \\
\mbox{no pretraining,} \mbox{trained from scratch} & \ensuremath{10^{-3}} & 0.9188$^{\ast}$ \\
\mbox{no pretraining,} \mbox{trained from scratch} & \ensuremath{3\times10^{-3}} & 0.9177 \\
\bottomrule
\end{tabular}
\end{adjustbox}
\end{minipage}
\par\bigskip
\begin{minipage}{\linewidth}\centering
\TableContinued{tab:lrsweep}{tab:lrsweep:2}
\begin{adjustbox}{max width=\linewidth, max totalheight=0.92\textheight}
\begin{tabular}{>{\raggedright\arraybackslash}p{\tablabw}lcc}
\toprule
 & learning rate & quark/gluon & sum \\
\midrule
\mbox{pretrained on \Kfull{} classes,} \mbox{fine-tuned,} \mbox{run 1} & \ensuremath{10^{-4}} & 0.8204 & 1.7408 \\
\mbox{pretrained on \Kfull{} classes,} \mbox{fine-tuned,} \mbox{run 1} & \ensuremath{3\times10^{-4}} & 0.8202 & 1.7394 \\
\mbox{pretrained on \Kfull{} classes,} \mbox{fine-tuned,} \mbox{run 1} & \ensuremath{10^{-3}} & 0.8211 & 1.7423$^{\ast}$ \\
\mbox{no pretraining,} \mbox{trained from scratch} & \ensuremath{3\times10^{-4}} & 0.8052 & 1.7204 \\
\mbox{no pretraining,} \mbox{trained from scratch} & \ensuremath{10^{-3}} & 0.8029 & 1.7217 \\
\mbox{no pretraining,} \mbox{trained from scratch} & \ensuremath{3\times10^{-3}} & 0.8059$^{\ast}$ & 1.7236 \\
\bottomrule
\end{tabular}
\end{adjustbox}
\end{minipage}
\end{center}


Fine-tuning keeps the missing inputs at zero for every model. Replacing zero by the training mean in the
fine-tuning of the \Kfull-class model of seed one on the full top training set moves its rejection at
30\% signal efficiency by less than the spread over pretraining runs, so the benchmark comparison does
not depend on that choice.
% CHECK: pending (no macro) fine-tuning zero-versus-mean check on top (188-class seed 1, full set): |R30 change| < seed SD of MlstFtTopRthirtyFull

\subsection{Transfer to $X\to b\bar b$ and $X\to bs$}\label{sec:app-xbx}
% CHECK: pending XbbNTrainPerClass, XbsNTrainPerClass, XbbRepetition, XbsRepetition (samples due 2026-11-03; drop with F9 if not staged by 2026-11-12)

For $X\to b\bar b$ and $X\to bs$ the transfer head follows Sophon's~\cite[App.~B.5]{sophon}: a
hidden layer of \DesignHeadHidden{} units with a ReLU activation and two outputs on the frozen class
token, with the hidden layer initialized from that of the pretrained head. It is trained for one pass
over \LitSophonTransferSamplesMillion{} million samples in batches of \LitSophonTransferBatch{} at a
constant learning rate of \LitSophonTransferLr. Sophon states neither an optimizer nor a
regularization; we use RAdam with the settings of pretraining, dropout after the hidden layer and no
weight decay, and train several seeds per pretraining run, which the tables average.

The signal and QCD training jets are drawn in matched bins of $(p_{\mathrm{T}}, \msd)$,
\XbbNTrainPerClass{} per class for $X\to b\bar b$ and \XbsNTrainPerClass{} for $X\to bs$. They are kept
in a fixed shuffled order, repeated \XbbRepetition{} and \XbsRepetition{} times to reach Sophon's number
of samples. The training QCD jets come from JetClass-II files that no pretraining run read and the evaluation QCD jets
from its test files, so no pretraining jet enters either. Whether the signal mass is fixed or variable is
recorded for every fit; Sophon trains its $X\to bs$ transfer head on a variable mass and evaluates it at
\LitSophonXMass~GeV~\cite[Sec.~III]{sophon}.

Sophon's direct discriminant can be formed only by the \Kfull- and \Kres-class models, in which
$X\to b\bar b$ is a class of its own. The closest-output-node control scores the pretrained $X\to bq$
output against QCD; where a label set merges $X\to bq$ with other decays, it scores the output of that
group. The network trained from scratch for $X\to bs$ is the Sophon network on the same inputs and jets,
trained for Sophon's number of epochs for its reference, \LitSophonScratchEpochs, with weight
decay~\cite[App.~B.3]{sophon}, its epoch chosen on our validation jets, with three seeds. No network is
trained from scratch for $X\to b\bar b$.

\subsection{Compute}\label{sec:app-compute}

To report what each result cost and on which hardware, table~\ref{tab:compute} gives every pretraining
run's GPU types with the hours spent on each, and its total GPU-hours. The hours sum the training and
validation passes of the run's epochs as each epoch recorded its device. A run could resume on another
GPU type after an interruption, so several runs list more than one. The
pretraining runs used \MlstComputeTotalGpuHours{} GPU-hours on \MlstComputeGpuTypes{} GPUs; the
\ComputeMainGpu{} carries the most hours, between \ComputeMainGpuShareMin{} and \ComputeMainGpuShareMax{}
of each label set's. The fine-tuning records hold the device but not the time, so the table gives the GPU
types of each run's fine-tuning without hours.

\begin{table*}[t]
\centering
\small
\caption{GPU types and GPU-hours of every pretraining run, summed over the training and validation passes of its epochs as each epoch recorded its device, and the GPU types its fine-tuning ran on.}
\label{tab:compute}
\begin{adjustbox}{max width=\linewidth}
\begin{tabular}{lllr}
\toprule
run & GPU types (hours) & GPU-hours & fine-tuning GPU types \\
\midrule
\mbox{pretrained on \Kfull{} classes,} \mbox{run 1} & V100 (135) & 135 & L4, RTX 3090, TITAN RTX, V100 \\
\mbox{pretrained on \Kfull{} classes,} \mbox{run 2} & V100 (145) & 145 & L4, RTX 4090, TITAN RTX, V100 \\
\mbox{pretrained on \Kfull{} classes,} \mbox{run 3} & RTX 3090 (116) & 116 & L4, L40, RTX 5000 Ada Generation, V100 \\
\mbox{pretrained on \Kres{} classes,} \mbox{run 1} & V100 (75), RTX 3090 (54) & 129 & L4, L40, RTX 3090, RTX 5000 Ada Generation, TITAN RTX, V100 \\
\mbox{pretrained on \Kres{} classes,} \mbox{run 2} & V100 (86), RTX 3090 (40) & 126 & L4, L40, Quadro RTX 6000, RTX 5000 Ada Generation, V100 \\
\mbox{pretrained on \Kres{} classes,} \mbox{run 3} & V100 (123), L4 (39) & 162 & Quadro RTX 6000, RTX 3090, TITAN RTX, V100 \\
\mbox{pretrained on \Kmid{} classes,} \mbox{run 1} & RTX 3090 (113) & 113 & L4, RTX 3090, TITAN RTX, V100 \\
\mbox{pretrained on \Kmid{} classes,} \mbox{run 2} & V100 (101), RTX 3090 (24) & 125 & L4, L40, Quadro RTX 6000, RTX 3090, TITAN RTX, V100 \\
\mbox{pretrained on \Kmid{} classes,} \mbox{run 3} & V100 (107), RTX 3090 (20) & 127 & L4, L40, RTX 3090, V100 \\
\mbox{pretrained on \Kcoarse{} classes,} \mbox{run 1} & RTX 3090 (108), L4 (3) & 111 & L4, L40, RTX 5000 Ada Generation, V100 \\
\mbox{pretrained on \Kcoarse{} classes,} \mbox{run 2} & V100 (133) & 133 & L4, TITAN RTX, V100 \\
\mbox{pretrained on \Kcoarse{} classes,} \mbox{run 3} & V100 (101), L4 (72) & 174 & L4, RTX 3090, RTX 4090, TITAN RTX, V100 \\
\mbox{pretrained on \Kres{} classes + mass,} \mbox{run 1} & V100 (129), RTX 3090 (26) & 155 & L4, Quadro RTX 6000, RTX 3090, RTX 4090, TITAN RTX, V100 \\
\mbox{pretrained on \Kres{} classes + mass,} \mbox{run 2} & V100 (130), RTX 3090 (15) & 145 & L4, Quadro RTX 6000, RTX 3090, RTX 5000 Ada Generation, V100 \\
\mbox{pretrained on \Kcoarse{} classes + mass,} \mbox{run 1} & V100 (126), RTX 3090 (16) & 142 & L4, L40, RTX 3090, RTX 5000 Ada Generation, TITAN RTX, V100 \\
\mbox{pretrained on \Kcoarse{} classes + mass,} \mbox{run 2} & RTX 3090 (101) & 101 & L4, L40, Quadro RTX 6000, V100 \\
\mbox{pretrained on 17 random classes (partition 1),} \mbox{run 1} & RTX 3090 (78), V100 (38) & 116 & L4, L40, Quadro RTX 6000, RTX 3090, TITAN RTX, V100 \\
\mbox{pretrained on 17 random classes (partition 2),} \mbox{run 2} & RTX 3090 (110) & 110 & L4, TITAN RTX, V100 \\
\mbox{pretrained on 17 random classes (partition 3),} \mbox{run 1} & RTX 3090 (109) & 109 & L4, L40, RTX 3090, V100 \\
\mbox{pretrained on 17 random classes (partition 4),} \mbox{run 2} & RTX 3090 (107) & 107 & L4, L40, RTX 3090, V100 \\
\mbox{pretrained on 17 random classes (partition 5),} \mbox{run 2} & RTX 3090 (98) & 98 & L4, L40, Quadro RTX 6000, V100 \\
\mbox{pretrained on 17 classes (no flavor boundary),} \mbox{run 1} & RTX 3090 (99) & 99 & L4, TITAN RTX, V100 \\
\mbox{pretrained on 17 classes (no flavor boundary),} \mbox{run 2} & RTX 3090 (130) & 130 & L4, Quadro RTX 6000, RTX 4090, RTX 5000 Ada Generation, V100 \\
\mbox{pretrained on 17 classes ($b$-quark boundary),} \mbox{run 1} & V100 (98), RTX 3090 (24) & 122 & L4, Quadro RTX 6000, TITAN RTX, V100 \\
\mbox{pretrained on 17 classes ($b$-quark boundary),} \mbox{run 2} & RTX 3090 (137) & 137 & L4, L40, Quadro RTX 6000, RTX 4090, TITAN RTX, V100 \\
\mbox{pretrained on 17 classes (random flavor boundary),} \mbox{run 1} & RTX 3090 (125) & 125 & L4, RTX 3090, RTX 5000 Ada Generation, V100 \\
\mbox{pretrained on 17 classes (random flavor boundary),} \mbox{run 2} & V100 (135), RTX 3090 (7) & 142 & L4, L40, Quadro RTX 6000, V100 \\
\mbox{pretrained on \Kfull{} classes (kinematic inputs only)} & RTX 4090 (75) & 75 & L40, V100 \\
\bottomrule
\end{tabular}
\end{adjustbox}
\end{table*}


% ----------------------------------------------------------------------------
\section{Missing input features}\label{sec:app-inputs}
% ----------------------------------------------------------------------------

To show what each substitution puts in each missing slot, table~\ref{tab:inputs} lists the
\FillNFeatures{} per-particle inputs in the layout of ParT's table of input features~\cite[Table~2]{part}
and marks those that the top and the quark/gluon data provide. It also gives where a zero lies after the
input normalization of the \Kfull-class model of seed one, in standard deviations of the training
particles. The top data lack \FillTopNZero{} inputs: the charge, the five type flags and the four impact-parameter
features. The quark/gluon data provide the particle type and the charge, ParT's exp
inputs~\cite{part}, and lack the \FillQgNZero{} impact-parameter features.

The three substitutions differ in what they put in a missing slot. Zero is also the value every model
sees in fine-tuning. The training mean is the running mean of the model's own input normalization,
which the normalization maps to its bias: it is the centroid of the training particles, with fractional
flags that no particle has, and it differs from model to model. A random draw replaces the missing
features of a particle by those of a real JetClass-II particle. For top, the draw takes a species with the
JetClass-II species frequencies, a charge given the species and the impact-parameter features of a random
particle of that species; for quark/gluon, the flags and the charge are kept and only the impact-parameter
features are drawn, from particles of the same species. Each particle receives four draws whose features
are averaged. The draws come from the class-balanced JetClass-II training set and depend only on the
batch counter and the seed, so every model sees the same draws.

Each single draw is a state a real particle takes, but the average of four is not, and neither are zero
and the mean: in JetClass-II every particle has exactly one type flag set, and only neutral particles have
zero impact parameters. A zeroed flag alone lies at most \MlstInputBnZeroSigmaMax{} standard deviations
from the training particles after the normalization (table~\ref{tab:inputs}); the state that no
pretraining particle has is the joint one, with no flag set.
% CHECK: pending MlstInputBnZeroSigmaMax (0.96 from the 188-class seed-1 input BatchNorm, docs/lit/missing_inputs_2026-11-01.md; must stay at or below about 1)

The substitution for each task is chosen on validation jets and applied to every model, the randomly
initialized network included: \LpSubstitutionTop{} for top and \LpSubstitutionQg{} for quark/gluon.
Table~\ref{tab:fill} gives all three substitutions for every configuration and so completes
figure~\ref{fig:f5}, which shows the \Kfull-class model and the randomly initialized network.
% CHECK: pending LpSubstitutionTop, LpSubstitutionQg (fits_v4_fill)

\providecommand{\TableContinued}[2]{\label{#2}\ifnum\getpagerefnumber{#2}>\getpagerefnumber{#1}\relax{\footnotesize\textit{Table~\ref{#1} (continued)}}\par\smallskip\fi}
\begin{center}
\small
\ifdefined\tablabw\else\newlength{\tablabw}\newlength{\tablabtmp}\newlength{\tablabcol}\newlength{\tablabnat}\newlength{\tablabatom}\fi
\setlength{\tablabw}{\linewidth}\setlength{\tablabnat}{0pt}\setlength{\tablabatom}{0pt}
\setlength{\tablabcol}{1\linewidth}\addtolength{\tablabcol}{-2\tabcolsep}
\setlength{\tablabnat}{0pt}\settowidth{\tablabtmp}{---}\ifdim\tablabtmp>\tablabnat\setlength{\tablabnat}{\tablabtmp}\fi\settowidth{\tablabtmp}{top}\ifdim\tablabtmp>\tablabnat\setlength{\tablabnat}{\tablabtmp}\fi\settowidth{\tablabtmp}{$\bullet$}\ifdim\tablabtmp>\tablabnat\setlength{\tablabnat}{\tablabtmp}\fi\addtolength{\tablabcol}{-\tablabnat}\addtolength{\tablabcol}{-2\tabcolsep}
\setlength{\tablabnat}{0pt}\settowidth{\tablabtmp}{quark/gluon}\ifdim\tablabtmp>\tablabnat\setlength{\tablabnat}{\tablabtmp}\fi\settowidth{\tablabtmp}{---}\ifdim\tablabtmp>\tablabnat\setlength{\tablabnat}{\tablabtmp}\fi\settowidth{\tablabtmp}{$\bullet$}\ifdim\tablabtmp>\tablabnat\setlength{\tablabnat}{\tablabtmp}\fi\addtolength{\tablabcol}{-\tablabnat}\addtolength{\tablabcol}{-2\tabcolsep}
\setlength{\tablabnat}{0pt}\settowidth{\tablabtmp}{zero after normalization}\ifdim\tablabtmp>\tablabnat\setlength{\tablabnat}{\tablabtmp}\fi\settowidth{\tablabtmp}{$-0.08$}\ifdim\tablabtmp>\tablabnat\setlength{\tablabnat}{\tablabtmp}\fi\settowidth{\tablabtmp}{$-0.10$}\ifdim\tablabtmp>\tablabnat\setlength{\tablabnat}{\tablabtmp}\fi\addtolength{\tablabcol}{-\tablabnat}\addtolength{\tablabcol}{-2\tabcolsep}
\ifdim\tablabcol<\tablabw\setlength{\tablabw}{\tablabcol}\fi
\setlength{\tablabcol}{1\linewidth}\addtolength{\tablabcol}{-2\tabcolsep}
\setlength{\tablabnat}{0pt}\settowidth{\tablabtmp}{running mean of the input normalization}\ifdim\tablabtmp>\tablabnat\setlength{\tablabnat}{\tablabtmp}\fi\settowidth{\tablabtmp}{mean}\ifdim\tablabtmp>\tablabnat\setlength{\tablabnat}{\tablabtmp}\fi\addtolength{\tablabcol}{-\tablabnat}\addtolength{\tablabcol}{-2\tabcolsep}
\ifdim\tablabcol<\tablabw\setlength{\tablabw}{\tablabcol}\fi
\setlength{\tablabcol}{1\linewidth}\addtolength{\tablabcol}{-2\tabcolsep}
\setlength{\tablabnat}{0pt}\settowidth{\tablabtmp}{a species drawn from JetClass-II's frequencies (top)}\ifdim\tablabtmp>\tablabnat\setlength{\tablabnat}{\tablabtmp}\fi\settowidth{\tablabtmp}{a JetClass-II particle of the same species}\ifdim\tablabtmp>\tablabnat\setlength{\tablabnat}{\tablabtmp}\fi\settowidth{\tablabtmp}{drawn given the species (top)}\ifdim\tablabtmp>\tablabnat\setlength{\tablabnat}{\tablabtmp}\fi\addtolength{\tablabcol}{-\tablabnat}\addtolength{\tablabcol}{-2\tabcolsep}
\ifdim\tablabcol<\tablabw\setlength{\tablabw}{\tablabcol}\fi
\setlength{\tablabnat}{0pt}\settowidth{\tablabtmp}{$\log(p_{\mathrm{T}}/p_{\mathrm{T}}^{\mathrm{jet}})$}\ifdim\tablabtmp>\tablabnat\setlength{\tablabnat}{\tablabtmp}\fi\settowidth{\tablabtmp}{$\log p_{\mathrm{T}}$, rescaled}\ifdim\tablabtmp>\tablabnat\setlength{\tablabnat}{\tablabtmp}\fi\settowidth{\tablabtmp}{electric charge}\ifdim\tablabtmp>\tablabnat\setlength{\tablabnat}{\tablabtmp}\fi
\settowidth{\tablabtmp}{$\log(p_{\mathrm{T}}/p_{\mathrm{T}}^{\mathrm{jet}})$}\ifdim\tablabtmp>\tablabatom\setlength{\tablabatom}{\tablabtmp}\fi\settowidth{\tablabtmp}{$\log p_{\mathrm{T}}$, rescaled}\ifdim\tablabtmp>\tablabatom\setlength{\tablabatom}{\tablabtmp}\fi\settowidth{\tablabtmp}{electric charge}\ifdim\tablabtmp>\tablabatom\setlength{\tablabatom}{\tablabtmp}\fi
\ifdim\tablabw>7cm\setlength{\tablabw}{7cm}\fi\ifdim\tablabnat<\tablabw\setlength{\tablabw}{\tablabnat}\fi\ifdim\tablabatom>\tablabw\setlength{\tablabw}{\tablabatom}\fi
\begin{minipage}{\linewidth}\centering
\captionof{table}{The per-particle inputs: which the top and the quark/gluon data provide ($\bullet$), where a zero lies after the network's input normalization (in standard deviations of the training particles), and what the mean and random-draw substitutions put in a missing slot. A random draw is averaged over four draws.}
\label{tab:inputs}
\par\medskip
\begin{adjustbox}{max width=\linewidth}
\begin{tabular}{>{\raggedright\arraybackslash}p{\tablabw}ccc}
\toprule
input & top & quark/gluon & zero after normalization \\
\midrule
$\log p_{\mathrm{T}}$, rescaled & $\bullet$ & $\bullet$ & -- \\
$\log E$, rescaled & $\bullet$ & $\bullet$ & -- \\
$\log(p_{\mathrm{T}}/p_{\mathrm{T}}^{\mathrm{jet}})$ & $\bullet$ & $\bullet$ & -- \\
$\log(E/E^{\mathrm{jet}})$ & $\bullet$ & $\bullet$ & -- \\
$\Delta R$ & $\bullet$ & $\bullet$ & -- \\
electric charge & --- & $\bullet$ & $0.00$ \\
charged hadron & --- & $\bullet$ & $-0.96$ \\
neutral hadron & --- & $\bullet$ & $-0.35$ \\
photon & --- & $\bullet$ & $-0.81$ \\
electron & --- & $\bullet$ & $-0.10$ \\
muon & --- & $\bullet$ & $-0.08$ \\
$d_0$ & --- & --- & $0.00$ \\
$\sigma_{d_0}$ & --- & --- & $-0.50$ \\
$d_z$ & --- & --- & $0.00$ \\
$\sigma_{d_z}$ & --- & --- & $-0.44$ \\
$\Delta\eta$ & $\bullet$ & $\bullet$ & -- \\
$\Delta\phi$ & $\bullet$ & $\bullet$ & -- \\
\bottomrule
\end{tabular}
\end{adjustbox}
\end{minipage}
\par\bigskip
\begin{minipage}{\linewidth}\centering
\TableContinued{tab:inputs}{tab:inputs:2}
\begin{adjustbox}{max width=\linewidth, max totalheight=0.92\textheight}
\begin{tabular}{>{\raggedright\arraybackslash}p{\tablabw}c}
\toprule
input & mean \\
\midrule
electric charge & running mean of the input normalization \\
charged hadron & running mean of the input normalization \\
neutral hadron & running mean of the input normalization \\
photon & running mean of the input normalization \\
electron & running mean of the input normalization \\
muon & running mean of the input normalization \\
$d_0$ & running mean of the input normalization \\
$\sigma_{d_0}$ & running mean of the input normalization \\
$d_z$ & running mean of the input normalization \\
$\sigma_{d_z}$ & running mean of the input normalization \\
\bottomrule
\end{tabular}
\end{adjustbox}
\end{minipage}
\par\bigskip
\begin{minipage}{\linewidth}\centering
\TableContinued{tab:inputs}{tab:inputs:3}
\begin{adjustbox}{max width=\linewidth, max totalheight=0.92\textheight}
\begin{tabular}{>{\raggedright\arraybackslash}p{\tablabw}c}
\toprule
input & random draws \\
\midrule
electric charge & drawn given the species (top) \\
charged hadron & a species drawn from JetClass-II's frequencies (top) \\
neutral hadron & a species drawn from JetClass-II's frequencies (top) \\
photon & a species drawn from JetClass-II's frequencies (top) \\
electron & a species drawn from JetClass-II's frequencies (top) \\
muon & a species drawn from JetClass-II's frequencies (top) \\
$d_0$ & a JetClass-II particle of the same species \\
$\sigma_{d_0}$ & a JetClass-II particle of the same species \\
$d_z$ & a JetClass-II particle of the same species \\
$\sigma_{d_z}$ & a JetClass-II particle of the same species \\
\bottomrule
\end{tabular}
\end{adjustbox}
\end{minipage}
\end{center}

\begin{table*}[t]
\centering
\caption{$R_{30}$, the inverse background efficiency at 30\% signal efficiency, of the linear probes of table~\ref{tab:summary} on top and quark/gluon tagging with Pythia jets, with the missing inputs (table~\ref{tab:inputs}) at zero, the training mean (mean) or random draws of a real particle (draws). Table~\ref{tab:summary} uses \LpSubstitutionTop{} for top and \LpSubstitutionQg{} for quark/gluon, so those columns are left out. The random network here has its batch-normalization statistics recomputed on \DesignBnJets{} training jets, unlike that of table~\ref{tab:summary}. Values are the mean $\pm$ standard deviation over \MlstNRunsWord{} runs, \MlstMassNRunsWord{} per mass network and flavor-boundary partition, and over five random partitions. Bold marks values within three standard deviations (both added in quadrature, before rounding) of the best in their column. The random network is a reference, never bold.}
\label{tab:fill}
\begin{tabular*}{\linewidth}{@{\extracolsep{\fill}}lcccc@{}}
\toprule
 & \multicolumn{2}{c}{top, $R_{30}$} & \multicolumn{2}{c}{quark/gluon (Pythia), $R_{30}$} \\
\cmidrule(lr){2-3}\cmidrule(lr){4-5}
 & zero & mean & zero & draws \\
\midrule
\multicolumn{5}{@{}l}{\textit{linear probe on the frozen features}} \\
\Kfull{} classes & {\fontseries{b}\selectfont 59.6\,$\pm$\,11.6} & {\fontseries{b}\selectfont 73\,$\pm$\,15} & {\fontseries{b}\selectfont 65.6\,$\pm$\,7.3} & {\fontseries{b}\selectfont 68.1\,$\pm$\,1.8} \\
\Kres{} classes & {\fontseries{b}\selectfont 75.5\,$\pm$\,16.9} & {\fontseries{b}\selectfont 64\,$\pm$\,11} & {\fontseries{b}\selectfont 57.9\,$\pm$\,9.3} & {\fontseries{b}\selectfont 66.8\,$\pm$\,0.4} \\
\Kres{} classes + mass & {\fontseries{b}\selectfont 85.7\,$\pm$\,12.9} & {\fontseries{b}\selectfont 68\,$\pm$\,15} & {\fontseries{b}\selectfont 63.2\,$\pm$\,3.1} & {\fontseries{b}\selectfont 71.2\,$\pm$\,1.2} \\
\Kmid{} classes & {\fontseries{b}\selectfont 56.9\,$\pm$\,12.6} & {\fontseries{b}\selectfont 58\,$\pm$\,14} & {\fontseries{b}\selectfont 60.9\,$\pm$\,8.5} & {\fontseries{b}\selectfont 69.0\,$\pm$\,1.1} \\
\Kcoarse{} classes & {\fontseries{b}\selectfont 54.6\,$\pm$\,4.6} & {\fontseries{b}\selectfont 55\,$\pm$\,5} & {\fontseries{b}\selectfont 62.0\,$\pm$\,7.0} & {\fontseries{b}\selectfont 71.6\,$\pm$\,2.8} \\
\Kcoarse{} classes + mass & {\fontseries{b}\selectfont 85.9\,$\pm$\,24.6} & {\fontseries{b}\selectfont 68\,$\pm$\,18} & {\fontseries{b}\selectfont 65.8\,$\pm$\,0.3} & {\fontseries{b}\selectfont 73.3\,$\pm$\,2.8} \\
\multicolumn{5}{@{}l}{17-class partitions} \\
\quad no flavor boundary & {\fontseries{b}\selectfont 69.5\,$\pm$\,0.2} & {\fontseries{b}\selectfont 54\,$\pm$\,15} & {\fontseries{b}\selectfont 64.8\,$\pm$\,2.3} & {\fontseries{b}\selectfont 71.3\,$\pm$\,0.3} \\
\quad $b$-quark boundary & {\fontseries{b}\selectfont 77.2\,$\pm$\,5.0} & {\fontseries{b}\selectfont 57\,$\pm$\,16} & {\fontseries{b}\selectfont 57.0\,$\pm$\,3.4} & {\fontseries{b}\selectfont 70.7\,$\pm$\,0.3} \\
\quad random flavor boundary & {\fontseries{b}\selectfont 52.7\,$\pm$\,16.7} & {\fontseries{b}\selectfont 37\,$\pm$\,5} & {\fontseries{b}\selectfont 59.9\,$\pm$\,8.0} & {\fontseries{b}\selectfont 70.5\,$\pm$\,1.0} \\
\quad random partitions & {\fontseries{b}\selectfont 51.8\,$\pm$\,16.5} & {\fontseries{b}\selectfont 51\,$\pm$\,6} & {\fontseries{b}\selectfont 59.2\,$\pm$\,6.5} & {\fontseries{b}\selectfont 64.7\,$\pm$\,2.8} \\
random network & 57.6\,$\pm$\,12.7 & 60\,$\pm$\,2 & 54.2\,$\pm$\,2.4 & 52.4\,$\pm$\,0.4 \\
\bottomrule
\end{tabular*}
\end{table*}


% ----------------------------------------------------------------------------
\section{Method checks}\label{sec:app-checks}
% ----------------------------------------------------------------------------

\subsection{MLP probes and the hidden layer of the head}\label{sec:app-mlp}

MLP probes answer whether a linear null means that the information is absent or only that a linear
classifier cannot read it. Wherever the main text reports that a linear probe finds no difference between
label sets, table~\ref{tab:mlp} gives the MLP probe on the same frozen features at the largest training
set of JetClass-II, JetClass and the class pairs, and of top and quark/gluon with the missing inputs set
to zero.

The last rows of the same table address a second question: whether the hidden layer of the head, where
Sophon's transfer head starts, holds information that a linear probe cannot read from the class
token. For each \Kfull-class model the pretrained first layer of the head and its activation,
$\mathrm{ReLU}(W_1 x + b_1)$, are applied to the saved class token $x$. A linear probe on these
\DesignHeadHidden{} outputs is trained on JetClass-II with \FtSizeESix{} training jets and on the full
quark/gluon training set, and the table sets it beside the linear probe on the class token itself.

\providecommand{\TableContinued}[2]{\label{#2}\ifnum\getpagerefnumber{#2}>\getpagerefnumber{#1}\relax{\footnotesize\textit{Table~\ref{#1} (continued)}}\par\smallskip\fi}
\begin{center}
\small
\ifdefined\tablabw\else\newlength{\tablabw}\newlength{\tablabtmp}\newlength{\tablabcol}\newlength{\tablabnat}\newlength{\tablabatom}\fi
\setlength{\tablabw}{\linewidth}\setlength{\tablabnat}{0pt}\setlength{\tablabatom}{0pt}
\setlength{\tablabcol}{1\linewidth}\addtolength{\tablabcol}{-2\tabcolsep}
\setlength{\tablabnat}{0pt}\settowidth{\tablabtmp}{0.9806\,$\pm$\,$<$0.0001}\ifdim\tablabtmp>\tablabnat\setlength{\tablabnat}{\tablabtmp}\fi\settowidth{\tablabtmp}{0.8648\,$\pm$\,0.0014}\ifdim\tablabtmp>\tablabnat\setlength{\tablabnat}{\tablabtmp}\fi\settowidth{\tablabtmp}{0.8974\,$\pm$\,0.0018}\ifdim\tablabtmp>\tablabnat\setlength{\tablabnat}{\tablabtmp}\fi\addtolength{\tablabcol}{-\tablabnat}\addtolength{\tablabcol}{-2\tabcolsep}
\setlength{\tablabnat}{0pt}\settowidth{\tablabtmp}{1.019\,$\pm$\,0.007}\ifdim\tablabtmp>\tablabnat\setlength{\tablabnat}{\tablabtmp}\fi\settowidth{\tablabtmp}{1.158, 1.150}\ifdim\tablabtmp>\tablabnat\setlength{\tablabnat}{\tablabtmp}\fi\settowidth{\tablabtmp}{1.168, 1.159}\ifdim\tablabtmp>\tablabnat\setlength{\tablabnat}{\tablabtmp}\fi\addtolength{\tablabcol}{-\tablabnat}\addtolength{\tablabcol}{-2\tabcolsep}
\ifdim\tablabcol<\tablabw\setlength{\tablabw}{\tablabcol}\fi
\setlength{\tablabcol}{1\linewidth}\addtolength{\tablabcol}{-2\tabcolsep}
\setlength{\tablabnat}{0pt}\settowidth{\tablabtmp}{0.9834\,$\pm$\,$<$0.0001}\ifdim\tablabtmp>\tablabnat\setlength{\tablabnat}{\tablabtmp}\fi\settowidth{\tablabtmp}{0.8694\,$\pm$\,0.0017}\ifdim\tablabtmp>\tablabnat\setlength{\tablabnat}{\tablabtmp}\fi\settowidth{\tablabtmp}{0.9035\,$\pm$\,0.0020}\ifdim\tablabtmp>\tablabnat\setlength{\tablabnat}{\tablabtmp}\fi\addtolength{\tablabcol}{-\tablabnat}\addtolength{\tablabcol}{-2\tabcolsep}
\setlength{\tablabnat}{0pt}\settowidth{\tablabtmp}{0.8372\,$\pm$\,0.0015}\ifdim\tablabtmp>\tablabnat\setlength{\tablabnat}{\tablabtmp}\fi\settowidth{\tablabtmp}{0.8610\,$\pm$\,0.0008}\ifdim\tablabtmp>\tablabnat\setlength{\tablabnat}{\tablabtmp}\fi\settowidth{\tablabtmp}{0.9524\,$\pm$\,0.0004}\ifdim\tablabtmp>\tablabnat\setlength{\tablabnat}{\tablabtmp}\fi\addtolength{\tablabcol}{-\tablabnat}\addtolength{\tablabcol}{-2\tabcolsep}
\ifdim\tablabcol<\tablabw\setlength{\tablabw}{\tablabcol}\fi
\setlength{\tablabcol}{1\linewidth}\addtolength{\tablabcol}{-2\tabcolsep}
\setlength{\tablabnat}{0pt}\settowidth{\tablabtmp}{0.8859\,$\pm$\,0.0021}\ifdim\tablabtmp>\tablabnat\setlength{\tablabnat}{\tablabtmp}\fi\settowidth{\tablabtmp}{0.9146\,$\pm$\,0.0012}\ifdim\tablabtmp>\tablabnat\setlength{\tablabnat}{\tablabtmp}\fi\settowidth{\tablabtmp}{0.9764, 0.9763}\ifdim\tablabtmp>\tablabnat\setlength{\tablabnat}{\tablabtmp}\fi\addtolength{\tablabcol}{-\tablabnat}\addtolength{\tablabcol}{-2\tabcolsep}
\ifdim\tablabcol<\tablabw\setlength{\tablabw}{\tablabcol}\fi
\setlength{\tablabcol}{1\linewidth}\addtolength{\tablabcol}{-2\tabcolsep}
\setlength{\tablabnat}{0pt}\settowidth{\tablabtmp}{0.9667\,$\pm$\,0.0006}\ifdim\tablabtmp>\tablabnat\setlength{\tablabnat}{\tablabtmp}\fi\settowidth{\tablabtmp}{0.9668, 0.9682}\ifdim\tablabtmp>\tablabnat\setlength{\tablabnat}{\tablabtmp}\fi\settowidth{\tablabtmp}{0.9670, 0.9626}\ifdim\tablabtmp>\tablabnat\setlength{\tablabnat}{\tablabtmp}\fi\addtolength{\tablabcol}{-\tablabnat}\addtolength{\tablabcol}{-2\tabcolsep}
\setlength{\tablabnat}{0pt}\settowidth{\tablabtmp}{210, 220}\ifdim\tablabtmp>\tablabnat\setlength{\tablabnat}{\tablabtmp}\fi\settowidth{\tablabtmp}{210\,$\pm$\,$<$10}\ifdim\tablabtmp>\tablabnat\setlength{\tablabnat}{\tablabtmp}\fi\settowidth{\tablabtmp}{220, 160}\ifdim\tablabtmp>\tablabnat\setlength{\tablabnat}{\tablabtmp}\fi\addtolength{\tablabcol}{-\tablabnat}\addtolength{\tablabcol}{-2\tabcolsep}
\ifdim\tablabcol<\tablabw\setlength{\tablabw}{\tablabcol}\fi
\setlength{\tablabcol}{1\linewidth}\addtolength{\tablabcol}{-2\tabcolsep}
\setlength{\tablabnat}{0pt}\settowidth{\tablabtmp}{0.8831\,$\pm$\,0.0003}\ifdim\tablabtmp>\tablabnat\setlength{\tablabnat}{\tablabtmp}\fi\settowidth{\tablabtmp}{0.8844\,$\pm$\,0.0001}\ifdim\tablabtmp>\tablabnat\setlength{\tablabnat}{\tablabtmp}\fi\settowidth{\tablabtmp}{0.8870\,$\pm$\,0.0055}\ifdim\tablabtmp>\tablabnat\setlength{\tablabnat}{\tablabtmp}\fi\addtolength{\tablabcol}{-\tablabnat}\addtolength{\tablabcol}{-2\tabcolsep}
\setlength{\tablabnat}{0pt}\settowidth{\tablabtmp}{77.5, 83.1}\ifdim\tablabtmp>\tablabnat\setlength{\tablabnat}{\tablabtmp}\fi\settowidth{\tablabtmp}{80.6, 77.9}\ifdim\tablabtmp>\tablabnat\setlength{\tablabnat}{\tablabtmp}\fi\settowidth{\tablabtmp}{80.8, 78.6}\ifdim\tablabtmp>\tablabnat\setlength{\tablabnat}{\tablabtmp}\fi\addtolength{\tablabcol}{-\tablabnat}\addtolength{\tablabcol}{-2\tabcolsep}
\ifdim\tablabcol<\tablabw\setlength{\tablabw}{\tablabcol}\fi
\setlength{\tablabcol}{1\linewidth}\addtolength{\tablabcol}{-2\tabcolsep}
\setlength{\tablabnat}{0pt}\settowidth{\tablabtmp}{0.8060\,$\pm$\,0.0014}\ifdim\tablabtmp>\tablabnat\setlength{\tablabnat}{\tablabtmp}\fi\settowidth{\tablabtmp}{0.8062\,$\pm$\,0.0024}\ifdim\tablabtmp>\tablabnat\setlength{\tablabnat}{\tablabtmp}\fi\settowidth{\tablabtmp}{0.8069, 0.8067}\ifdim\tablabtmp>\tablabnat\setlength{\tablabnat}{\tablabtmp}\fi\addtolength{\tablabcol}{-\tablabnat}\addtolength{\tablabcol}{-2\tabcolsep}
\setlength{\tablabnat}{0pt}\settowidth{\tablabtmp}{35.3, 34.5}\ifdim\tablabtmp>\tablabnat\setlength{\tablabnat}{\tablabtmp}\fi\settowidth{\tablabtmp}{35.3, 35.0}\ifdim\tablabtmp>\tablabnat\setlength{\tablabnat}{\tablabtmp}\fi\settowidth{\tablabtmp}{35.9, 36.5}\ifdim\tablabtmp>\tablabnat\setlength{\tablabnat}{\tablabtmp}\fi\addtolength{\tablabcol}{-\tablabnat}\addtolength{\tablabcol}{-2\tabcolsep}
\ifdim\tablabcol<\tablabw\setlength{\tablabw}{\tablabcol}\fi
\setlength{\tablabcol}{1\linewidth}\addtolength{\tablabcol}{-2\tabcolsep}
\setlength{\tablabnat}{0pt}\settowidth{\tablabtmp}{0.6304\,$\pm$\,0.0038}\ifdim\tablabtmp>\tablabnat\setlength{\tablabnat}{\tablabtmp}\fi\settowidth{\tablabtmp}{0.6865\,$\pm$\,0.0030}\ifdim\tablabtmp>\tablabnat\setlength{\tablabnat}{\tablabtmp}\fi\settowidth{\tablabtmp}{0.9077, 0.9046}\ifdim\tablabtmp>\tablabnat\setlength{\tablabnat}{\tablabtmp}\fi\addtolength{\tablabcol}{-\tablabnat}\addtolength{\tablabcol}{-2\tabcolsep}
\setlength{\tablabnat}{0pt}\settowidth{\tablabtmp}{1.002\,$\pm$\,0.006}\ifdim\tablabtmp>\tablabnat\setlength{\tablabnat}{\tablabtmp}\fi\settowidth{\tablabtmp}{1.104\,$\pm$\,0.017}\ifdim\tablabtmp>\tablabnat\setlength{\tablabnat}{\tablabtmp}\fi\settowidth{\tablabtmp}{1.359, 1.335}\ifdim\tablabtmp>\tablabnat\setlength{\tablabnat}{\tablabtmp}\fi\addtolength{\tablabcol}{-\tablabnat}\addtolength{\tablabcol}{-2\tabcolsep}
\ifdim\tablabcol<\tablabw\setlength{\tablabw}{\tablabcol}\fi
\setlength{\tablabcol}{1\linewidth}\addtolength{\tablabcol}{-2\tabcolsep}
\setlength{\tablabnat}{0pt}\settowidth{\tablabtmp}{0.9975\,$\pm$\,$<$0.0001}\ifdim\tablabtmp>\tablabnat\setlength{\tablabnat}{\tablabtmp}\fi\settowidth{\tablabtmp}{0.9977\,$\pm$\,$<$0.0001}\ifdim\tablabtmp>\tablabnat\setlength{\tablabnat}{\tablabtmp}\fi\settowidth{\tablabtmp}{0.9978\,$\pm$\,$<$0.0001}\ifdim\tablabtmp>\tablabnat\setlength{\tablabnat}{\tablabtmp}\fi\addtolength{\tablabcol}{-\tablabnat}\addtolength{\tablabcol}{-2\tabcolsep}
\setlength{\tablabnat}{0pt}\settowidth{\tablabtmp}{1.023\,$\pm$\,0.021}\ifdim\tablabtmp>\tablabnat\setlength{\tablabnat}{\tablabtmp}\fi\settowidth{\tablabtmp}{1.125\,$\pm$\,0.011}\ifdim\tablabtmp>\tablabnat\setlength{\tablabnat}{\tablabtmp}\fi\settowidth{\tablabtmp}{1.744, 1.637}\ifdim\tablabtmp>\tablabnat\setlength{\tablabnat}{\tablabtmp}\fi\addtolength{\tablabcol}{-\tablabnat}\addtolength{\tablabcol}{-2\tabcolsep}
\ifdim\tablabcol<\tablabw\setlength{\tablabw}{\tablabcol}\fi
\setlength{\tablabcol}{1\linewidth}\addtolength{\tablabcol}{-2\tabcolsep}
\setlength{\tablabnat}{0pt}\settowidth{\tablabtmp}{0.9935\,$\pm$\,$<$0.0001}\ifdim\tablabtmp>\tablabnat\setlength{\tablabnat}{\tablabtmp}\fi\settowidth{\tablabtmp}{0.7472\,$\pm$\,0.0056}\ifdim\tablabtmp>\tablabnat\setlength{\tablabnat}{\tablabtmp}\fi\settowidth{\tablabtmp}{0.8373\,$\pm$\,0.0061}\ifdim\tablabtmp>\tablabnat\setlength{\tablabnat}{\tablabtmp}\fi\addtolength{\tablabcol}{-\tablabnat}\addtolength{\tablabcol}{-2\tabcolsep}
\setlength{\tablabnat}{0pt}\settowidth{\tablabtmp}{1.011\,$\pm$\,0.030}\ifdim\tablabtmp>\tablabnat\setlength{\tablabnat}{\tablabtmp}\fi\settowidth{\tablabtmp}{1.112\,$\pm$\,0.006}\ifdim\tablabtmp>\tablabnat\setlength{\tablabnat}{\tablabtmp}\fi\settowidth{\tablabtmp}{1.442, 1.372}\ifdim\tablabtmp>\tablabnat\setlength{\tablabnat}{\tablabtmp}\fi\addtolength{\tablabcol}{-\tablabnat}\addtolength{\tablabcol}{-2\tabcolsep}
\ifdim\tablabcol<\tablabw\setlength{\tablabw}{\tablabcol}\fi
\setlength{\tablabcol}{1\linewidth}\addtolength{\tablabcol}{-2\tabcolsep}
\setlength{\tablabnat}{0pt}\settowidth{\tablabtmp}{0.7228\,$\pm$\,0.0031}\ifdim\tablabtmp>\tablabnat\setlength{\tablabnat}{\tablabtmp}\fi\settowidth{\tablabtmp}{0.8222\,$\pm$\,0.0049}\ifdim\tablabtmp>\tablabnat\setlength{\tablabnat}{\tablabtmp}\fi\settowidth{\tablabtmp}{0.9680, 0.9647}\ifdim\tablabtmp>\tablabnat\setlength{\tablabnat}{\tablabtmp}\fi\addtolength{\tablabcol}{-\tablabnat}\addtolength{\tablabcol}{-2\tabcolsep}
\setlength{\tablabnat}{0pt}\settowidth{\tablabtmp}{1.085\,$\pm$\,0.030}\ifdim\tablabtmp>\tablabnat\setlength{\tablabnat}{\tablabtmp}\fi\settowidth{\tablabtmp}{1.121\,$\pm$\,0.017}\ifdim\tablabtmp>\tablabnat\setlength{\tablabnat}{\tablabtmp}\fi\settowidth{\tablabtmp}{1.432, 1.372}\ifdim\tablabtmp>\tablabnat\setlength{\tablabnat}{\tablabtmp}\fi\addtolength{\tablabcol}{-\tablabnat}\addtolength{\tablabcol}{-2\tabcolsep}
\ifdim\tablabcol<\tablabw\setlength{\tablabw}{\tablabcol}\fi
\setlength{\tablabcol}{1\linewidth}\addtolength{\tablabcol}{-2\tabcolsep}
\setlength{\tablabnat}{0pt}\settowidth{\tablabtmp}{0.9882\,$\pm$\,$<$0.0001}\ifdim\tablabtmp>\tablabnat\setlength{\tablabnat}{\tablabtmp}\fi\settowidth{\tablabtmp}{0.8185\,$\pm$\,0.0019}\ifdim\tablabtmp>\tablabnat\setlength{\tablabnat}{\tablabtmp}\fi\settowidth{\tablabtmp}{0.8879\,$\pm$\,0.0033}\ifdim\tablabtmp>\tablabnat\setlength{\tablabnat}{\tablabtmp}\fi\addtolength{\tablabcol}{-\tablabnat}\addtolength{\tablabcol}{-2\tabcolsep}
\setlength{\tablabnat}{0pt}\settowidth{\tablabtmp}{1.050\,$\pm$\,0.004}\ifdim\tablabtmp>\tablabnat\setlength{\tablabnat}{\tablabtmp}\fi\settowidth{\tablabtmp}{1.055\,$\pm$\,0.005}\ifdim\tablabtmp>\tablabnat\setlength{\tablabnat}{\tablabtmp}\fi\settowidth{\tablabtmp}{1.158, 1.140}\ifdim\tablabtmp>\tablabnat\setlength{\tablabnat}{\tablabtmp}\fi\addtolength{\tablabcol}{-\tablabnat}\addtolength{\tablabcol}{-2\tabcolsep}
\ifdim\tablabcol<\tablabw\setlength{\tablabw}{\tablabcol}\fi
\setlength{\tablabcol}{1\linewidth}\addtolength{\tablabcol}{-2\tabcolsep}
\setlength{\tablabnat}{0pt}\settowidth{\tablabtmp}{0.6823\,$\pm$\,0.0051}\ifdim\tablabtmp>\tablabnat\setlength{\tablabnat}{\tablabtmp}\fi\settowidth{\tablabtmp}{0.7143\,$\pm$\,0.0013}\ifdim\tablabtmp>\tablabnat\setlength{\tablabnat}{\tablabtmp}\fi\settowidth{\tablabtmp}{0.9145\,$\pm$\,0.0004}\ifdim\tablabtmp>\tablabnat\setlength{\tablabnat}{\tablabtmp}\fi\addtolength{\tablabcol}{-\tablabnat}\addtolength{\tablabcol}{-2\tabcolsep}
\setlength{\tablabnat}{0pt}\settowidth{\tablabtmp}{1.039\,$\pm$\,0.005}\ifdim\tablabtmp>\tablabnat\setlength{\tablabnat}{\tablabtmp}\fi\settowidth{\tablabtmp}{1.055\,$\pm$\,0.006}\ifdim\tablabtmp>\tablabnat\setlength{\tablabnat}{\tablabtmp}\fi\settowidth{\tablabtmp}{1.141, 1.120}\ifdim\tablabtmp>\tablabnat\setlength{\tablabnat}{\tablabtmp}\fi\addtolength{\tablabcol}{-\tablabnat}\addtolength{\tablabcol}{-2\tabcolsep}
\ifdim\tablabcol<\tablabw\setlength{\tablabw}{\tablabcol}\fi
\setlength{\tablabcol}{1\linewidth}\addtolength{\tablabcol}{-2\tabcolsep}
\setlength{\tablabnat}{0pt}\settowidth{\tablabtmp}{0.9674, 0.9679}\ifdim\tablabtmp>\tablabnat\setlength{\tablabnat}{\tablabtmp}\fi\settowidth{\tablabtmp}{0.9682\,$\pm$\,0.0003}\ifdim\tablabtmp>\tablabnat\setlength{\tablabnat}{\tablabtmp}\fi\settowidth{\tablabtmp}{0.9736\,$\pm$\,0.0005}\ifdim\tablabtmp>\tablabnat\setlength{\tablabnat}{\tablabtmp}\fi\addtolength{\tablabcol}{-\tablabnat}\addtolength{\tablabcol}{-2\tabcolsep}
\setlength{\tablabnat}{0pt}\settowidth{\tablabtmp}{1.033\,$\pm$\,0.016}\ifdim\tablabtmp>\tablabnat\setlength{\tablabnat}{\tablabtmp}\fi\settowidth{\tablabtmp}{1.061, 1.051}\ifdim\tablabtmp>\tablabnat\setlength{\tablabnat}{\tablabtmp}\fi\settowidth{\tablabtmp}{1.135\,$\pm$\,0.018}\ifdim\tablabtmp>\tablabnat\setlength{\tablabnat}{\tablabtmp}\fi\addtolength{\tablabcol}{-\tablabnat}\addtolength{\tablabcol}{-2\tabcolsep}
\ifdim\tablabcol<\tablabw\setlength{\tablabw}{\tablabcol}\fi
\setlength{\tablabcol}{1\linewidth}\addtolength{\tablabcol}{-2\tabcolsep}
\setlength{\tablabnat}{0pt}\settowidth{\tablabtmp}{$X\to bc$ vs $bq$, $cs$}\ifdim\tablabtmp>\tablabnat\setlength{\tablabnat}{\tablabtmp}\fi\settowidth{\tablabtmp}{1.003\,$\pm$\,0.004}\ifdim\tablabtmp>\tablabnat\setlength{\tablabnat}{\tablabtmp}\fi\settowidth{\tablabtmp}{1.057, 1.018}\ifdim\tablabtmp>\tablabnat\setlength{\tablabnat}{\tablabtmp}\fi\addtolength{\tablabcol}{-\tablabnat}\addtolength{\tablabcol}{-2\tabcolsep}
\setlength{\tablabnat}{0pt}\settowidth{\tablabtmp}{$X\to b\bar b$ vs $c\bar c$}\ifdim\tablabtmp>\tablabnat\setlength{\tablabnat}{\tablabtmp}\fi\settowidth{\tablabtmp}{1.00\,$\pm$\,0.09}\ifdim\tablabtmp>\tablabnat\setlength{\tablabnat}{\tablabtmp}\fi\settowidth{\tablabtmp}{1.11\,$\pm$\,0.01}\ifdim\tablabtmp>\tablabnat\setlength{\tablabnat}{\tablabtmp}\fi\addtolength{\tablabcol}{-\tablabnat}\addtolength{\tablabcol}{-2\tabcolsep}
\ifdim\tablabcol<\tablabw\setlength{\tablabw}{\tablabcol}\fi
\setlength{\tablabcol}{1\linewidth}\addtolength{\tablabcol}{-2\tabcolsep}
\setlength{\tablabnat}{0pt}\settowidth{\tablabtmp}{$X\to YY\to b\bar{b}qq$ vs $c\bar{c}qq$}\ifdim\tablabtmp>\tablabnat\setlength{\tablabnat}{\tablabtmp}\fi\settowidth{\tablabtmp}{0.97\,$\pm$\,0.08}\ifdim\tablabtmp>\tablabnat\setlength{\tablabnat}{\tablabtmp}\fi\settowidth{\tablabtmp}{1.02\,$\pm$\,0.06}\ifdim\tablabtmp>\tablabnat\setlength{\tablabnat}{\tablabtmp}\fi\addtolength{\tablabcol}{-\tablabnat}\addtolength{\tablabcol}{-2\tabcolsep}
\ifdim\tablabcol<\tablabw\setlength{\tablabw}{\tablabcol}\fi
\setlength{\tablabcol}{1\linewidth}\addtolength{\tablabcol}{-2\tabcolsep}
\setlength{\tablabnat}{0pt}\settowidth{\tablabtmp}{QCD $b\bar{b}$ vs $c\bar{c}$}\ifdim\tablabtmp>\tablabnat\setlength{\tablabnat}{\tablabtmp}\fi\settowidth{\tablabtmp}{1.08\,$\pm$\,0.09}\ifdim\tablabtmp>\tablabnat\setlength{\tablabnat}{\tablabtmp}\fi\settowidth{\tablabtmp}{1.13\,$\pm$\,0.05}\ifdim\tablabtmp>\tablabnat\setlength{\tablabnat}{\tablabtmp}\fi\addtolength{\tablabcol}{-\tablabnat}\addtolength{\tablabcol}{-2\tabcolsep}
\setlength{\tablabnat}{0pt}\settowidth{\tablabtmp}{QCD with $b$ vs light}\ifdim\tablabtmp>\tablabnat\setlength{\tablabnat}{\tablabtmp}\fi\settowidth{\tablabtmp}{1.03, 0.94}\ifdim\tablabtmp>\tablabnat\setlength{\tablabnat}{\tablabtmp}\fi\settowidth{\tablabtmp}{1.05\,$\pm$\,0.09}\ifdim\tablabtmp>\tablabnat\setlength{\tablabnat}{\tablabtmp}\fi\addtolength{\tablabcol}{-\tablabnat}\addtolength{\tablabcol}{-2\tabcolsep}
\ifdim\tablabcol<\tablabw\setlength{\tablabw}{\tablabcol}\fi
\setlength{\tablabcol}{1\linewidth}\addtolength{\tablabcol}{-2\tabcolsep}
\setlength{\tablabnat}{0pt}\settowidth{\tablabtmp}{QCD with $c$, no $b$, vs light}\ifdim\tablabtmp>\tablabnat\setlength{\tablabnat}{\tablabtmp}\fi\settowidth{\tablabtmp}{1.041\,$\pm$\,0.021}\ifdim\tablabtmp>\tablabnat\setlength{\tablabnat}{\tablabtmp}\fi\settowidth{\tablabtmp}{1.057\,$\pm$\,0.016}\ifdim\tablabtmp>\tablabnat\setlength{\tablabnat}{\tablabtmp}\fi\addtolength{\tablabcol}{-\tablabnat}\addtolength{\tablabcol}{-2\tabcolsep}
\ifdim\tablabcol<\tablabw\setlength{\tablabw}{\tablabcol}\fi
\setlength{\tablabcol}{1\linewidth}\addtolength{\tablabcol}{-2\tabcolsep}
\setlength{\tablabnat}{0pt}\settowidth{\tablabtmp}{0.9801\,$\pm$\,$<$0.0001}\ifdim\tablabtmp>\tablabnat\setlength{\tablabnat}{\tablabtmp}\fi\settowidth{\tablabtmp}{0.9802\,$\pm$\,$<$0.0001}\ifdim\tablabtmp>\tablabnat\setlength{\tablabnat}{\tablabtmp}\fi\settowidth{\tablabtmp}{JetClass-II AUC}\ifdim\tablabtmp>\tablabnat\setlength{\tablabnat}{\tablabtmp}\fi\addtolength{\tablabcol}{-\tablabnat}\addtolength{\tablabcol}{-2\tabcolsep}
\ifdim\tablabcol<\tablabw\setlength{\tablabw}{\tablabcol}\fi
\setlength{\tablabcol}{1\linewidth}\addtolength{\tablabcol}{-2\tabcolsep}
\setlength{\tablabnat}{0pt}\settowidth{\tablabtmp}{quark/gluon AUC}\ifdim\tablabtmp>\tablabnat\setlength{\tablabnat}{\tablabtmp}\fi\settowidth{\tablabtmp}{0.8780\,$\pm$\,0.0059}\ifdim\tablabtmp>\tablabnat\setlength{\tablabnat}{\tablabtmp}\fi\settowidth{\tablabtmp}{0.8802\,$\pm$\,0.0070}\ifdim\tablabtmp>\tablabnat\setlength{\tablabnat}{\tablabtmp}\fi\addtolength{\tablabcol}{-\tablabnat}\addtolength{\tablabcol}{-2\tabcolsep}
\setlength{\tablabnat}{0pt}\settowidth{\tablabtmp}{quark/gluon $R_{30}$}\ifdim\tablabtmp>\tablabnat\setlength{\tablabnat}{\tablabtmp}\fi\settowidth{\tablabtmp}{66\,$\pm$\,7}\ifdim\tablabtmp>\tablabnat\setlength{\tablabnat}{\tablabtmp}\fi\settowidth{\tablabtmp}{70\,$\pm$\,6}\ifdim\tablabtmp>\tablabnat\setlength{\tablabnat}{\tablabtmp}\fi\addtolength{\tablabcol}{-\tablabnat}\addtolength{\tablabcol}{-2\tabcolsep}
\ifdim\tablabcol<\tablabw\setlength{\tablabw}{\tablabcol}\fi
\setlength{\tablabnat}{0pt}\settowidth{\tablabtmp}{\mbox{pretrained on \Kfull{} classes,} \mbox{linear probe on the 512-unit hidden layer}}\ifdim\tablabtmp>\tablabnat\setlength{\tablabnat}{\tablabtmp}\fi\settowidth{\tablabtmp}{\mbox{pretrained on 17 random classes (partition 1),} \mbox{MLP probe,} \mbox{run 1}}\ifdim\tablabtmp>\tablabnat\setlength{\tablabnat}{\tablabtmp}\fi\settowidth{\tablabtmp}{\mbox{pretrained on 17 random classes (partition 2),} \mbox{MLP probe,} \mbox{run 2}}\ifdim\tablabtmp>\tablabnat\setlength{\tablabnat}{\tablabtmp}\fi
\settowidth{\tablabtmp}{pretrained on 17 classes (random flavor boundary),}\ifdim\tablabtmp>\tablabatom\setlength{\tablabatom}{\tablabtmp}\fi\settowidth{\tablabtmp}{pretrained on \Kfull{} classes (kinematic inputs only),}\ifdim\tablabtmp>\tablabatom\setlength{\tablabatom}{\tablabtmp}\fi\settowidth{\tablabtmp}{pretrained on 17 classes (no flavor boundary),}\ifdim\tablabtmp>\tablabatom\setlength{\tablabatom}{\tablabtmp}\fi
\ifdim\tablabw>7cm\setlength{\tablabw}{7cm}\fi\ifdim\tablabnat<\tablabw\setlength{\tablabw}{\tablabnat}\fi\ifdim\tablabatom>\tablabw\setlength{\tablabw}{\tablabatom}\fi
\begin{minipage}{\linewidth}\centering
\captionof{table}{Nonlinear probes on the frozen features of networks pretrained on JetClass-II with 188, 162, 43 and 17 classes and of the controls, trained and scored on each task (a multilayer perceptron with one hidden layer of \LpMlpHidden{} units, \LpMlpSeedsWord{} seeds averaged) at the largest training set: class-balanced JetClass-II AUC (all, resonant and QCD classes) with the ratio $\rho$ of $1-\mathrm{AUC}$, equation~(\ref{eq:ratio}), to the \Kfull-class run of the same seed, JetClass AUC, and AUC and $R_{30}$ on top, quark/gluon and Herwig jets (missing inputs set to zero); then the class pairs of table~\ref{tab:pairs}, class-balanced, except $X\to ee$ vs $\mu\mu$, which has no nonlinear probe. The models with the mass regression and the numbered random partitions have no nonlinear probe on the class-balanced test set: they show a dash in its columns and are left out of the pair blocks; the next block gives them on the natural-composition test set beside the levels of the hierarchy on that set, and no $\rho$ crosses the two test sets. Last, the linear probe on the 512-unit hidden layer of the \Kfull-class models' pretraining head beside the one on the class token, with the missing inputs also set to zero. Mean $\pm$ standard deviation over the \MlstNRunsWord{} pretraining runs of each level of the class hierarchy; configurations with \RandFlavRuns{} runs give both values. Every random partition has 17 classes; each numbered one (1--5) has a single run, numbered by the pretraining seed it shares with the \Kfull-class run of that seed. The random network without pretraining: three initializations. The \Kfull-class model with kinematic inputs only, a diagnostic of the missing inputs, has one run and no input missing, and is evaluated on top tagging only.}
\label{tab:mlp}
\par\medskip
\begin{adjustbox}{max width=\linewidth, max totalheight=0.55\textheight}
\begin{tabular}{>{\raggedright\arraybackslash}p{\tablabw}cc}
\toprule
 & \multicolumn{2}{c}{JetClass-II} \\
\cmidrule(lr){2-3}
 & AUC & $\rho$ \\
\midrule
\mbox{pretrained on \Kfull{} classes,} \mbox{MLP probe} & 0.9806\,$\pm$\,$<$0.0001 & 1.000 \\
\mbox{pretrained on \Kres{} classes,} \mbox{MLP probe} & 0.9802\,$\pm$\,0.0001 & 1.019\,$\pm$\,0.007 \\
\mbox{pretrained on \Kmid{} classes,} \mbox{MLP probe} & 0.9773\,$\pm$\,0.0001 & 1.170\,$\pm$\,0.006 \\
\mbox{pretrained on \Kcoarse{} classes,} \mbox{MLP probe} & 0.9731\,$\pm$\,0.0004 & 1.388\,$\pm$\,0.021 \\
\mbox{pretrained on 17 classes (no flavor boundary),} \mbox{MLP probe} & 0.9749, 0.9746 & 1.294, 1.310 \\
\mbox{pretrained on 17 classes ($b$-quark boundary),} \mbox{MLP probe} & 0.9774, 0.9775 & 1.168, 1.159 \\
\mbox{pretrained on 17 classes (random flavor boundary),} \mbox{MLP probe} & 0.9776, 0.9777 & 1.158, 1.150 \\
\mbox{random network,} \mbox{MLP probe} & 0.8648\,$\pm$\,0.0014 & -- \\
\mbox{random network, BatchNorm recomputed,} \mbox{MLP probe} & 0.8974\,$\pm$\,0.0018 & -- \\
\bottomrule
\end{tabular}
\end{adjustbox}
\end{minipage}
\par\bigskip
\begin{minipage}{\linewidth}\centering
\TableContinued{tab:mlp}{tab:mlp:2}
\begin{adjustbox}{max width=\linewidth, max totalheight=0.92\textheight}
\begin{tabular}{>{\raggedright\arraybackslash}p{\tablabw}cc}
\toprule
 & \multicolumn{2}{c}{JetClass-II} \\
\cmidrule(lr){2-3}
 & AUC, resonant & AUC, QCD \\
\midrule
\mbox{pretrained on \Kfull{} classes,} \mbox{MLP probe} & 0.9834\,$\pm$\,$<$0.0001 & 0.9643\,$\pm$\,0.0001 \\
\mbox{pretrained on \Kres{} classes,} \mbox{MLP probe} & 0.9833\,$\pm$\,0.0001 & 0.9618\,$\pm$\,0.0003 \\
\mbox{pretrained on \Kmid{} classes,} \mbox{MLP probe} & 0.9802\,$\pm$\,0.0002 & 0.9605\,$\pm$\,0.0001 \\
\mbox{pretrained on \Kcoarse{} classes,} \mbox{MLP probe} & 0.9766\,$\pm$\,0.0004 & 0.9524\,$\pm$\,0.0004 \\
\mbox{pretrained on 17 classes (no flavor boundary),} \mbox{MLP probe} & 0.9786, 0.9783 & 0.9528, 0.9522 \\
\mbox{pretrained on 17 classes ($b$-quark boundary),} \mbox{MLP probe} & 0.9806, 0.9807 & 0.9582, 0.9584 \\
\mbox{pretrained on 17 classes (random flavor boundary),} \mbox{MLP probe} & 0.9810, 0.9811 & 0.9571, 0.9576 \\
\mbox{random network,} \mbox{MLP probe} & 0.8694\,$\pm$\,0.0017 & 0.8372\,$\pm$\,0.0015 \\
\mbox{random network, BatchNorm recomputed,} \mbox{MLP probe} & 0.9035\,$\pm$\,0.0020 & 0.8610\,$\pm$\,0.0008 \\
\bottomrule
\end{tabular}
\end{adjustbox}
\end{minipage}
\par\bigskip
\begin{minipage}{\linewidth}\centering
\TableContinued{tab:mlp}{tab:mlp:3}
\begin{adjustbox}{max width=\linewidth, max totalheight=0.92\textheight}
\begin{tabular}{>{\raggedright\arraybackslash}p{\tablabw}c}
\toprule
 & JetClass \\
 & AUC \\
\midrule
\mbox{pretrained on \Kfull{} classes,} \mbox{MLP probe} & 0.9786\,$\pm$\,0.0004 \\
\mbox{pretrained on \Kres{} classes,} \mbox{MLP probe} & 0.9786\,$\pm$\,0.0002 \\
\mbox{pretrained on \Kmid{} classes,} \mbox{MLP probe} & 0.9782\,$\pm$\,0.0003 \\
\mbox{pretrained on \Kcoarse{} classes,} \mbox{MLP probe} & 0.9776\,$\pm$\,0.0005 \\
\mbox{pretrained on \Kres{} classes + mass,} \mbox{MLP probe} & 0.9794, 0.9799 \\
\mbox{pretrained on \Kcoarse{} classes + mass,} \mbox{MLP probe} & 0.9787, 0.9785 \\
\mbox{pretrained on 17 random classes (partition 1),} \mbox{MLP probe,} \mbox{run 1} & 0.9739 \\
\mbox{pretrained on 17 random classes (partition 2),} \mbox{MLP probe,} \mbox{run 2} & 0.9750 \\
\mbox{pretrained on 17 random classes (partition 3),} \mbox{MLP probe,} \mbox{run 1} & 0.9744 \\
\mbox{pretrained on 17 random classes (partition 4),} \mbox{MLP probe,} \mbox{run 2} & 0.9742 \\
\mbox{pretrained on 17 random classes (partition 5),} \mbox{MLP probe,} \mbox{run 2} & 0.9752 \\
\mbox{pretrained on 17 classes (no flavor boundary),} \mbox{MLP probe} & 0.9764, 0.9763 \\
\mbox{pretrained on 17 classes ($b$-quark boundary),} \mbox{MLP probe} & 0.9767, 0.9776 \\
\mbox{pretrained on 17 classes (random flavor boundary),} \mbox{MLP probe} & 0.9767, 0.9763 \\
\mbox{random network,} \mbox{MLP probe} & 0.8859\,$\pm$\,0.0021 \\
\mbox{random network, BatchNorm recomputed,} \mbox{MLP probe} & 0.9146\,$\pm$\,0.0012 \\
\bottomrule
\end{tabular}
\end{adjustbox}
\end{minipage}
\par\bigskip
\begin{minipage}{\linewidth}\centering
\TableContinued{tab:mlp}{tab:mlp:4}
\begin{adjustbox}{max width=\linewidth, max totalheight=0.92\textheight}
\begin{tabular}{>{\raggedright\arraybackslash}p{\tablabw}cc}
\toprule
 & \multicolumn{2}{c}{top} \\
\cmidrule(lr){2-3}
 & AUC & $R_{30}$ \\
\midrule
\mbox{pretrained on \Kfull{} classes,} \mbox{MLP probe} & 0.9713\,$\pm$\,0.0034 & 290\,$\pm$\,60 \\
\mbox{pretrained on \Kres{} classes,} \mbox{MLP probe} & 0.9711\,$\pm$\,0.0042 & 280\,$\pm$\,90 \\
\mbox{pretrained on \Kmid{} classes,} \mbox{MLP probe} & 0.9688\,$\pm$\,0.0029 & 250\,$\pm$\,60 \\
\mbox{pretrained on \Kcoarse{} classes,} \mbox{MLP probe} & 0.9688\,$\pm$\,0.0018 & 220\,$\pm$\,20 \\
\mbox{pretrained on \Kres{} classes + mass,} \mbox{MLP probe} & 0.9731, 0.9754 & 320, 380 \\
\mbox{pretrained on \Kcoarse{} classes + mass,} \mbox{MLP probe} & 0.9747, 0.9704 & 360, 260 \\
\mbox{pretrained on 17 random classes (partition 1),} \mbox{MLP probe,} \mbox{run 1} & 0.9635 & 160 \\
\mbox{pretrained on 17 random classes (partition 2),} \mbox{MLP probe,} \mbox{run 2} & 0.9676 & 210 \\
\mbox{pretrained on 17 random classes (partition 3),} \mbox{MLP probe,} \mbox{run 1} & 0.9667 & 190 \\
\mbox{pretrained on 17 random classes (partition 4),} \mbox{MLP probe,} \mbox{run 2} & 0.9587 & 120 \\
\mbox{pretrained on 17 random classes (partition 5),} \mbox{MLP probe,} \mbox{run 2} & 0.9701 & 250 \\
\mbox{pretrained on 17 classes (no flavor boundary),} \mbox{MLP probe} & 0.9668, 0.9682 & 210, 220 \\
\mbox{pretrained on 17 classes ($b$-quark boundary),} \mbox{MLP probe} & 0.9718, 0.9724 & 300, 310 \\
\mbox{pretrained on 17 classes (random flavor boundary),} \mbox{MLP probe} & 0.9670, 0.9626 & 220, 160 \\
\mbox{random network,} \mbox{MLP probe} & 0.9667\,$\pm$\,0.0006 & 210\,$\pm$\,$<$10 \\
\mbox{random network, BatchNorm recomputed,} \mbox{MLP probe} & 0.9712\,$\pm$\,0.0024 & 300\,$\pm$\,50 \\
\mbox{pretrained on \Kfull{} classes (kinematic inputs only),} \mbox{MLP probe} & 0.9811 & 720 \\
\bottomrule
\end{tabular}
\end{adjustbox}
\end{minipage}
\par\bigskip
\begin{minipage}{\linewidth}\centering
\TableContinued{tab:mlp}{tab:mlp:5}
\begin{adjustbox}{max width=\linewidth, max totalheight=0.92\textheight}
\begin{tabular}{>{\raggedright\arraybackslash}p{\tablabw}cc}
\toprule
 & \multicolumn{2}{c}{quark/gluon} \\
\cmidrule(lr){2-3}
 & AUC & $R_{30}$ \\
\midrule
\mbox{pretrained on \Kfull{} classes,} \mbox{MLP probe} & 0.8937\,$\pm$\,0.0043 & 82.6\,$\pm$\,3.5 \\
\mbox{pretrained on \Kres{} classes,} \mbox{MLP probe} & 0.8889\,$\pm$\,0.0048 & 77.0\,$\pm$\,2.2 \\
\mbox{pretrained on \Kmid{} classes,} \mbox{MLP probe} & 0.8904\,$\pm$\,0.0031 & 80.9\,$\pm$\,4.3 \\
\mbox{pretrained on \Kcoarse{} classes,} \mbox{MLP probe} & 0.8870\,$\pm$\,0.0055 & 78.9\,$\pm$\,1.7 \\
\mbox{pretrained on \Kres{} classes + mass,} \mbox{MLP probe} & 0.8932, 0.8914 & 81.8, 78.4 \\
\mbox{pretrained on \Kcoarse{} classes + mass,} \mbox{MLP probe} & 0.8933, 0.8927 & 80.9, 81.9 \\
\mbox{pretrained on 17 random classes (partition 1),} \mbox{MLP probe,} \mbox{run 1} & 0.8927 & 80.9 \\
\mbox{pretrained on 17 random classes (partition 2),} \mbox{MLP probe,} \mbox{run 2} & 0.8882 & 79.0 \\
\mbox{pretrained on 17 random classes (partition 3),} \mbox{MLP probe,} \mbox{run 1} & 0.8938 & 80.5 \\
\mbox{pretrained on 17 random classes (partition 4),} \mbox{MLP probe,} \mbox{run 2} & 0.8811 & 72.0 \\
\mbox{pretrained on 17 random classes (partition 5),} \mbox{MLP probe,} \mbox{run 2} & 0.8858 & 75.5 \\
\mbox{pretrained on 17 classes (no flavor boundary),} \mbox{MLP probe} & 0.8912, 0.8921 & 77.5, 83.1 \\
\mbox{pretrained on 17 classes ($b$-quark boundary),} \mbox{MLP probe} & 0.8887, 0.8911 & 80.8, 78.6 \\
\mbox{pretrained on 17 classes (random flavor boundary),} \mbox{MLP probe} & 0.8906, 0.8870 & 80.6, 77.9 \\
\mbox{random network,} \mbox{MLP probe} & 0.8844\,$\pm$\,0.0001 & 76.6\,$\pm$\,0.8 \\
\mbox{random network, BatchNorm recomputed,} \mbox{MLP probe} & 0.8831\,$\pm$\,0.0003 & 76.3\,$\pm$\,1.9 \\
\bottomrule
\end{tabular}
\end{adjustbox}
\end{minipage}
\par\bigskip
\begin{minipage}{\linewidth}\centering
\TableContinued{tab:mlp}{tab:mlp:6}
\begin{adjustbox}{max width=\linewidth, max totalheight=0.92\textheight}
\begin{tabular}{>{\raggedright\arraybackslash}p{\tablabw}cc}
\toprule
 & \multicolumn{2}{c}{Herwig} \\
\cmidrule(lr){2-3}
 & AUC & $R_{30}$ \\
\midrule
\mbox{pretrained on \Kfull{} classes,} \mbox{MLP probe} & 0.8077\,$\pm$\,0.0006 & 36.6\,$\pm$\,0.3 \\
\mbox{pretrained on \Kres{} classes,} \mbox{MLP probe} & 0.8060\,$\pm$\,0.0014 & 35.6\,$\pm$\,0.6 \\
\mbox{pretrained on \Kmid{} classes,} \mbox{MLP probe} & 0.8071\,$\pm$\,0.0007 & 35.7\,$\pm$\,0.3 \\
\mbox{pretrained on \Kcoarse{} classes,} \mbox{MLP probe} & 0.8062\,$\pm$\,0.0024 & 36.2\,$\pm$\,0.5 \\
\mbox{pretrained on \Kres{} classes + mass,} \mbox{MLP probe} & 0.8094, 0.8071 & 37.6, 36.8 \\
\mbox{pretrained on \Kcoarse{} classes + mass,} \mbox{MLP probe} & 0.8079, 0.8084 & 35.9, 36.5 \\
\mbox{pretrained on 17 random classes (partition 1),} \mbox{MLP probe,} \mbox{run 1} & 0.8087 & 36.6 \\
\mbox{pretrained on 17 random classes (partition 2),} \mbox{MLP probe,} \mbox{run 2} & 0.8057 & 35.2 \\
\mbox{pretrained on 17 random classes (partition 3),} \mbox{MLP probe,} \mbox{run 1} & 0.8059 & 35.2 \\
\mbox{pretrained on 17 random classes (partition 4),} \mbox{MLP probe,} \mbox{run 2} & 0.8062 & 35.5 \\
\mbox{pretrained on 17 random classes (partition 5),} \mbox{MLP probe,} \mbox{run 2} & 0.8049 & 35.4 \\
\mbox{pretrained on 17 classes (no flavor boundary),} \mbox{MLP probe} & 0.8073, 0.8071 & 35.3, 35.0 \\
\mbox{pretrained on 17 classes ($b$-quark boundary),} \mbox{MLP probe} & 0.8069, 0.8067 & 36.1, 35.6 \\
\mbox{pretrained on 17 classes (random flavor boundary),} \mbox{MLP probe} & 0.8074, 0.8052 & 35.3, 34.5 \\
\mbox{random network,} \mbox{MLP probe} & 0.8074\,$\pm$\,0.0007 & 38.2\,$\pm$\,0.4 \\
\mbox{random network, BatchNorm recomputed,} \mbox{MLP probe} & 0.8087\,$\pm$\,0.0009 & 37.8\,$\pm$\,1.5 \\
\bottomrule
\end{tabular}
\end{adjustbox}
\end{minipage}
\par\bigskip
\begin{minipage}{\linewidth}\centering
\TableContinued{tab:mlp}{tab:mlp:7}
\begin{adjustbox}{max width=\linewidth, max totalheight=0.92\textheight}
\begin{tabular}{>{\raggedright\arraybackslash}p{\tablabw}cc}
\toprule
 & \multicolumn{2}{c}{$X\to bc$ vs $bq$, $cs$} \\
\cmidrule(lr){2-3}
 & AUC & $\rho$ \\
\midrule
\mbox{pretrained on \Kfull{} classes,} \mbox{MLP probe} & 0.9434\,$\pm$\,0.0001 & 1.000 \\
\mbox{pretrained on \Kres{} classes,} \mbox{MLP probe} & 0.9432\,$\pm$\,0.0002 & 1.002\,$\pm$\,0.006 \\
\mbox{pretrained on \Kmid{} classes,} \mbox{MLP probe} & 0.9375\,$\pm$\,0.0008 & 1.104\,$\pm$\,0.017 \\
\mbox{pretrained on \Kcoarse{} classes,} \mbox{MLP probe} & 0.9080\,$\pm$\,0.0013 & 1.624\,$\pm$\,0.022 \\
\mbox{pretrained on 17 classes (no flavor boundary),} \mbox{MLP probe} & 0.9077, 0.9046 & 1.633, 1.679 \\
\mbox{pretrained on 17 classes ($b$-quark boundary),} \mbox{MLP probe} & 0.9216, 0.9227 & 1.385, 1.360 \\
\mbox{pretrained on 17 classes (random flavor boundary),} \mbox{MLP probe} & 0.9231, 0.9242 & 1.359, 1.335 \\
\mbox{random network,} \mbox{MLP probe} & 0.6304\,$\pm$\,0.0038 & -- \\
\mbox{random network, BatchNorm recomputed,} \mbox{MLP probe} & 0.6865\,$\pm$\,0.0030 & -- \\
\bottomrule
\end{tabular}
\end{adjustbox}
\end{minipage}
\par\bigskip
\begin{minipage}{\linewidth}\centering
\TableContinued{tab:mlp}{tab:mlp:8}
\begin{adjustbox}{max width=\linewidth, max totalheight=0.92\textheight}
\begin{tabular}{>{\raggedright\arraybackslash}p{\tablabw}cc}
\toprule
 & \multicolumn{2}{c}{$X\to b\bar b$ vs $c\bar c$} \\
\cmidrule(lr){2-3}
 & AUC & $\rho$ \\
\midrule
\mbox{pretrained on \Kfull{} classes,} \mbox{MLP probe} & 0.9978\,$\pm$\,$<$0.0001 & 1.000 \\
\mbox{pretrained on \Kres{} classes,} \mbox{MLP probe} & 0.9977\,$\pm$\,$<$0.0001 & 1.023\,$\pm$\,0.021 \\
\mbox{pretrained on \Kmid{} classes,} \mbox{MLP probe} & 0.9975\,$\pm$\,$<$0.0001 & 1.125\,$\pm$\,0.011 \\
\mbox{pretrained on \Kcoarse{} classes,} \mbox{MLP probe} & 0.9932\,$\pm$\,0.0002 & 3.039\,$\pm$\,0.085 \\
\mbox{pretrained on 17 classes (no flavor boundary),} \mbox{MLP probe} & 0.9926, 0.9902 & 3.329, 4.337 \\
\mbox{pretrained on 17 classes ($b$-quark boundary),} \mbox{MLP probe} & 0.9961, 0.9963 & 1.744, 1.637 \\
\mbox{pretrained on 17 classes (random flavor boundary),} \mbox{MLP probe} & 0.9951, 0.9946 & 2.207, 2.375 \\
\mbox{random network,} \mbox{MLP probe} & 0.7592\,$\pm$\,0.0091 & -- \\
\mbox{random network, BatchNorm recomputed,} \mbox{MLP probe} & 0.8655\,$\pm$\,0.0075 & -- \\
\bottomrule
\end{tabular}
\end{adjustbox}
\end{minipage}
\par\bigskip
\begin{minipage}{\linewidth}\centering
\TableContinued{tab:mlp}{tab:mlp:9}
\begin{adjustbox}{max width=\linewidth, max totalheight=0.92\textheight}
\begin{tabular}{>{\raggedright\arraybackslash}p{\tablabw}cc}
\toprule
 & \multicolumn{2}{c}{$X\to YY\to b\bar{b}qq$ vs $c\bar{c}qq$} \\
\cmidrule(lr){2-3}
 & AUC & $\rho$ \\
\midrule
\mbox{pretrained on \Kfull{} classes,} \mbox{MLP probe} & 0.9941\,$\pm$\,0.0001 & 1.000 \\
\mbox{pretrained on \Kres{} classes,} \mbox{MLP probe} & 0.9941\,$\pm$\,0.0001 & 1.011\,$\pm$\,0.030 \\
\mbox{pretrained on \Kmid{} classes,} \mbox{MLP probe} & 0.9935\,$\pm$\,$<$0.0001 & 1.112\,$\pm$\,0.006 \\
\mbox{pretrained on \Kcoarse{} classes,} \mbox{MLP probe} & 0.9828\,$\pm$\,0.0011 & 2.935\,$\pm$\,0.173 \\
\mbox{pretrained on 17 classes (no flavor boundary),} \mbox{MLP probe} & 0.9813, 0.9776 & 3.202, 3.779 \\
\mbox{pretrained on 17 classes ($b$-quark boundary),} \mbox{MLP probe} & 0.9916, 0.9919 & 1.442, 1.372 \\
\mbox{pretrained on 17 classes (random flavor boundary),} \mbox{MLP probe} & 0.9889, 0.9891 & 1.890, 1.843 \\
\mbox{random network,} \mbox{MLP probe} & 0.7472\,$\pm$\,0.0056 & -- \\
\mbox{random network, BatchNorm recomputed,} \mbox{MLP probe} & 0.8373\,$\pm$\,0.0061 & -- \\
\bottomrule
\end{tabular}
\end{adjustbox}
\end{minipage}
\par\bigskip
\begin{minipage}{\linewidth}\centering
\TableContinued{tab:mlp}{tab:mlp:10}
\begin{adjustbox}{max width=\linewidth, max totalheight=0.92\textheight}
\begin{tabular}{>{\raggedright\arraybackslash}p{\tablabw}cc}
\toprule
 & \multicolumn{2}{c}{QCD $b\bar{b}$ vs $c\bar{c}$} \\
\cmidrule(lr){2-3}
 & AUC & $\rho$ \\
\midrule
\mbox{pretrained on \Kfull{} classes,} \mbox{MLP probe} & 0.9871\,$\pm$\,0.0001 & 1.000 \\
\mbox{pretrained on \Kres{} classes,} \mbox{MLP probe} & 0.9860\,$\pm$\,0.0003 & 1.085\,$\pm$\,0.030 \\
\mbox{pretrained on \Kmid{} classes,} \mbox{MLP probe} & 0.9855\,$\pm$\,0.0001 & 1.121\,$\pm$\,0.017 \\
\mbox{pretrained on \Kcoarse{} classes,} \mbox{MLP probe} & 0.9701\,$\pm$\,0.0008 & 2.312\,$\pm$\,0.044 \\
\mbox{pretrained on 17 classes (no flavor boundary),} \mbox{MLP probe} & 0.9680, 0.9647 & 2.497, 2.715 \\
\mbox{pretrained on 17 classes ($b$-quark boundary),} \mbox{MLP probe} & 0.9817, 0.9821 & 1.432, 1.372 \\
\mbox{pretrained on 17 classes (random flavor boundary),} \mbox{MLP probe} & 0.9776, 0.9775 & 1.745, 1.725 \\
\mbox{random network,} \mbox{MLP probe} & 0.7228\,$\pm$\,0.0031 & -- \\
\mbox{random network, BatchNorm recomputed,} \mbox{MLP probe} & 0.8222\,$\pm$\,0.0049 & -- \\
\bottomrule
\end{tabular}
\end{adjustbox}
\end{minipage}
\par\bigskip
\begin{minipage}{\linewidth}\centering
\TableContinued{tab:mlp}{tab:mlp:11}
\begin{adjustbox}{max width=\linewidth, max totalheight=0.92\textheight}
\begin{tabular}{>{\raggedright\arraybackslash}p{\tablabw}cc}
\toprule
 & \multicolumn{2}{c}{QCD with $b$ vs light} \\
\cmidrule(lr){2-3}
 & AUC & $\rho$ \\
\midrule
\mbox{pretrained on \Kfull{} classes,} \mbox{MLP probe} & 0.9882\,$\pm$\,$<$0.0001 & 1.000 \\
\mbox{pretrained on \Kres{} classes,} \mbox{MLP probe} & 0.9876\,$\pm$\,0.0001 & 1.050\,$\pm$\,0.004 \\
\mbox{pretrained on \Kmid{} classes,} \mbox{MLP probe} & 0.9875\,$\pm$\,0.0001 & 1.055\,$\pm$\,0.005 \\
\mbox{pretrained on \Kcoarse{} classes,} \mbox{MLP probe} & 0.9842\,$\pm$\,0.0002 & 1.337\,$\pm$\,0.021 \\
\mbox{pretrained on 17 classes (no flavor boundary),} \mbox{MLP probe} & 0.9845, 0.9847 & 1.315, 1.288 \\
\mbox{pretrained on 17 classes ($b$-quark boundary),} \mbox{MLP probe} & 0.9863, 0.9865 & 1.158, 1.140 \\
\mbox{pretrained on 17 classes (random flavor boundary),} \mbox{MLP probe} & 0.9858, 0.9859 & 1.201, 1.187 \\
\mbox{random network,} \mbox{MLP probe} & 0.8185\,$\pm$\,0.0019 & -- \\
\mbox{random network, BatchNorm recomputed,} \mbox{MLP probe} & 0.8879\,$\pm$\,0.0033 & -- \\
\bottomrule
\end{tabular}
\end{adjustbox}
\end{minipage}
\par\bigskip
\begin{minipage}{\linewidth}\centering
\TableContinued{tab:mlp}{tab:mlp:12}
\begin{adjustbox}{max width=\linewidth, max totalheight=0.92\textheight}
\begin{tabular}{>{\raggedright\arraybackslash}p{\tablabw}cc}
\toprule
 & \multicolumn{2}{c}{QCD with $c$, no $b$, vs light} \\
\cmidrule(lr){2-3}
 & AUC & $\rho$ \\
\midrule
\mbox{pretrained on \Kfull{} classes,} \mbox{MLP probe} & 0.9329\,$\pm$\,0.0001 & 1.000 \\
\mbox{pretrained on \Kres{} classes,} \mbox{MLP probe} & 0.9303\,$\pm$\,0.0002 & 1.039\,$\pm$\,0.005 \\
\mbox{pretrained on \Kmid{} classes,} \mbox{MLP probe} & 0.9292\,$\pm$\,0.0003 & 1.055\,$\pm$\,0.006 \\
\mbox{pretrained on \Kcoarse{} classes,} \mbox{MLP probe} & 0.9145\,$\pm$\,0.0004 & 1.274\,$\pm$\,0.008 \\
\mbox{pretrained on 17 classes (no flavor boundary),} \mbox{MLP probe} & 0.9174, 0.9186 & 1.234, 1.211 \\
\mbox{pretrained on 17 classes ($b$-quark boundary),} \mbox{MLP probe} & 0.9224, 0.9229 & 1.159, 1.147 \\
\mbox{pretrained on 17 classes (random flavor boundary),} \mbox{MLP probe} & 0.9236, 0.9247 & 1.141, 1.120 \\
\mbox{random network,} \mbox{MLP probe} & 0.6823\,$\pm$\,0.0051 & -- \\
\mbox{random network, BatchNorm recomputed,} \mbox{MLP probe} & 0.7143\,$\pm$\,0.0013 & -- \\
\bottomrule
\end{tabular}
\end{adjustbox}
\end{minipage}
\par\bigskip
\begin{minipage}{\linewidth}\centering
\TableContinued{tab:mlp}{tab:mlp:13}
\begin{adjustbox}{max width=\linewidth, max totalheight=0.92\textheight}
\begin{tabular}{>{\raggedright\arraybackslash}p{\tablabw}cc}
\toprule
natural-composition & \multicolumn{2}{c}{JetClass-II} \\
\cmidrule(lr){2-3}
test set & AUC & $\rho$ \\
\midrule
\mbox{pretrained on \Kfull{} classes,} \mbox{MLP probe} & 0.9767\,$\pm$\,0.0002 & 1.000 \\
\mbox{pretrained on \Kres{} classes,} \mbox{MLP probe} & 0.9759\,$\pm$\,0.0004 & 1.033\,$\pm$\,0.016 \\
\mbox{pretrained on \Kmid{} classes,} \mbox{MLP probe} & 0.9736\,$\pm$\,0.0005 & 1.135\,$\pm$\,0.018 \\
\mbox{pretrained on \Kcoarse{} classes,} \mbox{MLP probe} & 0.9682\,$\pm$\,0.0003 & 1.364\,$\pm$\,0.008 \\
\mbox{pretrained on \Kres{} classes + mass,} \mbox{MLP probe} & 0.9752, 0.9757 & 1.061, 1.051 \\
\mbox{pretrained on \Kcoarse{} classes + mass,} \mbox{MLP probe} & 0.9674, 0.9679 & 1.393, 1.391 \\
\mbox{pretrained on 17 random classes (partition 1),} \mbox{MLP probe,} \mbox{run 1} & 0.9736 & 1.131 \\
\mbox{pretrained on 17 random classes (partition 2),} \mbox{MLP probe,} \mbox{run 2} & 0.9736 & 1.144 \\
\mbox{pretrained on 17 random classes (partition 3),} \mbox{MLP probe,} \mbox{run 1} & 0.9737 & 1.124 \\
\mbox{pretrained on 17 random classes (partition 4),} \mbox{MLP probe,} \mbox{run 2} & 0.9738 & 1.135 \\
\mbox{pretrained on 17 random classes (partition 5),} \mbox{MLP probe,} \mbox{run 2} & 0.9735 & 1.147 \\
\bottomrule
\end{tabular}
\end{adjustbox}
\end{minipage}
\par\bigskip
\begin{minipage}{\linewidth}\centering
\TableContinued{tab:mlp}{tab:mlp:14}
\begin{adjustbox}{max width=\linewidth, max totalheight=0.92\textheight}
\begin{tabular}{>{\raggedright\arraybackslash}p{\tablabw}cc}
\toprule
natural-composition & \multicolumn{2}{c}{$\rho$ on the class pairs} \\
\cmidrule(lr){2-3}
test set & $X\to bc$ vs $bq$, $cs$ & $X\to b\bar b$ vs $c\bar c$ \\
\midrule
\mbox{pretrained on \Kfull{} classes,} \mbox{MLP probe} & 1.000 & 1.00 \\
\mbox{pretrained on \Kres{} classes,} \mbox{MLP probe} & 1.003\,$\pm$\,0.004 & 1.00\,$\pm$\,0.09 \\
\mbox{pretrained on \Kmid{} classes,} \mbox{MLP probe} & 1.135\,$\pm$\,0.017 & 1.11\,$\pm$\,0.01 \\
\mbox{pretrained on \Kcoarse{} classes,} \mbox{MLP probe} & 1.762\,$\pm$\,0.043 & 3.26\,$\pm$\,0.10 \\
\mbox{pretrained on \Kres{} classes + mass,} \mbox{MLP probe} & 1.057, 1.018 & 1.21, 1.09 \\
\mbox{pretrained on \Kcoarse{} classes + mass,} \mbox{MLP probe} & 1.839, 1.797 & 3.65, 3.11 \\
\mbox{pretrained on 17 random classes (partition 1),} \mbox{MLP probe,} \mbox{run 1} & 1.201 & 1.44 \\
\mbox{pretrained on 17 random classes (partition 2),} \mbox{MLP probe,} \mbox{run 2} & 1.170 & 1.35 \\
\mbox{pretrained on 17 random classes (partition 3),} \mbox{MLP probe,} \mbox{run 1} & 1.096 & 1.44 \\
\mbox{pretrained on 17 random classes (partition 4),} \mbox{MLP probe,} \mbox{run 2} & 1.093 & 1.25 \\
\mbox{pretrained on 17 random classes (partition 5),} \mbox{MLP probe,} \mbox{run 2} & 1.212 & 1.56 \\
\bottomrule
\end{tabular}
\end{adjustbox}
\end{minipage}
\par\bigskip
\begin{minipage}{\linewidth}\centering
\TableContinued{tab:mlp}{tab:mlp:15}
\begin{adjustbox}{max width=\linewidth, max totalheight=0.92\textheight}
\begin{tabular}{>{\raggedright\arraybackslash}p{\tablabw}c}
\toprule
natural-composition & \multicolumn{1}{c}{$\rho$ on the class pairs} \\
\cmidrule(lr){2-2}
test set & $X\to YY\to b\bar{b}qq$ vs $c\bar{c}qq$ \\
\midrule
\mbox{pretrained on \Kfull{} classes,} \mbox{MLP probe} & 1.00 \\
\mbox{pretrained on \Kres{} classes,} \mbox{MLP probe} & 0.97\,$\pm$\,0.08 \\
\mbox{pretrained on \Kmid{} classes,} \mbox{MLP probe} & 1.02\,$\pm$\,0.06 \\
\mbox{pretrained on \Kcoarse{} classes,} \mbox{MLP probe} & 3.34\,$\pm$\,0.29 \\
\mbox{pretrained on \Kres{} classes + mass,} \mbox{MLP probe} & 1.17, 1.04 \\
\mbox{pretrained on \Kcoarse{} classes + mass,} \mbox{MLP probe} & 3.74, 3.34 \\
\mbox{pretrained on 17 random classes (partition 1),} \mbox{MLP probe,} \mbox{run 1} & 1.35 \\
\mbox{pretrained on 17 random classes (partition 2),} \mbox{MLP probe,} \mbox{run 2} & 1.30 \\
\mbox{pretrained on 17 random classes (partition 3),} \mbox{MLP probe,} \mbox{run 1} & 1.32 \\
\mbox{pretrained on 17 random classes (partition 4),} \mbox{MLP probe,} \mbox{run 2} & 1.25 \\
\mbox{pretrained on 17 random classes (partition 5),} \mbox{MLP probe,} \mbox{run 2} & 1.30 \\
\bottomrule
\end{tabular}
\end{adjustbox}
\end{minipage}
\par\bigskip
\begin{minipage}{\linewidth}\centering
\TableContinued{tab:mlp}{tab:mlp:16}
\begin{adjustbox}{max width=\linewidth, max totalheight=0.92\textheight}
\begin{tabular}{>{\raggedright\arraybackslash}p{\tablabw}cc}
\toprule
natural-composition & \multicolumn{2}{c}{$\rho$ on the class pairs} \\
\cmidrule(lr){2-3}
test set & QCD $b\bar{b}$ vs $c\bar{c}$ & QCD with $b$ vs light \\
\midrule
\mbox{pretrained on \Kfull{} classes,} \mbox{MLP probe} & 1.00 & 1.00 \\
\mbox{pretrained on \Kres{} classes,} \mbox{MLP probe} & 1.08\,$\pm$\,0.09 & 1.06\,$\pm$\,0.06 \\
\mbox{pretrained on \Kmid{} classes,} \mbox{MLP probe} & 1.13\,$\pm$\,0.05 & 1.05\,$\pm$\,0.09 \\
\mbox{pretrained on \Kcoarse{} classes,} \mbox{MLP probe} & 2.71\,$\pm$\,0.13 & 1.32\,$\pm$\,0.18 \\
\mbox{pretrained on \Kres{} classes + mass,} \mbox{MLP probe} & 1.30, 1.16 & 1.03, 0.94 \\
\mbox{pretrained on \Kcoarse{} classes + mass,} \mbox{MLP probe} & 3.26, 2.88 & 1.38, 1.26 \\
\mbox{pretrained on 17 random classes (partition 1),} \mbox{MLP probe,} \mbox{run 1} & 1.53 & 1.05 \\
\mbox{pretrained on 17 random classes (partition 2),} \mbox{MLP probe,} \mbox{run 2} & 1.47 & 1.00 \\
\mbox{pretrained on 17 random classes (partition 3),} \mbox{MLP probe,} \mbox{run 1} & 1.45 & 1.16 \\
\mbox{pretrained on 17 random classes (partition 4),} \mbox{MLP probe,} \mbox{run 2} & 1.37 & 1.16 \\
\mbox{pretrained on 17 random classes (partition 5),} \mbox{MLP probe,} \mbox{run 2} & 1.53 & 1.16 \\
\bottomrule
\end{tabular}
\end{adjustbox}
\end{minipage}
\par\bigskip
\begin{minipage}{\linewidth}\centering
\TableContinued{tab:mlp}{tab:mlp:17}
\begin{adjustbox}{max width=\linewidth, max totalheight=0.92\textheight}
\begin{tabular}{>{\raggedright\arraybackslash}p{\tablabw}c}
\toprule
natural-composition & \multicolumn{1}{c}{$\rho$ on the class pairs} \\
\cmidrule(lr){2-2}
test set & QCD with $c$, no $b$, vs light \\
\midrule
\mbox{pretrained on \Kfull{} classes,} \mbox{MLP probe} & 1.000 \\
\mbox{pretrained on \Kres{} classes,} \mbox{MLP probe} & 1.041\,$\pm$\,0.021 \\
\mbox{pretrained on \Kmid{} classes,} \mbox{MLP probe} & 1.057\,$\pm$\,0.016 \\
\mbox{pretrained on \Kcoarse{} classes,} \mbox{MLP probe} & 1.303\,$\pm$\,0.032 \\
\mbox{pretrained on \Kres{} classes + mass,} \mbox{MLP probe} & 1.071, 1.045 \\
\mbox{pretrained on \Kcoarse{} classes + mass,} \mbox{MLP probe} & 1.326, 1.323 \\
\mbox{pretrained on 17 random classes (partition 1),} \mbox{MLP probe,} \mbox{run 1} & 1.125 \\
\mbox{pretrained on 17 random classes (partition 2),} \mbox{MLP probe,} \mbox{run 2} & 1.101 \\
\mbox{pretrained on 17 random classes (partition 3),} \mbox{MLP probe,} \mbox{run 1} & 1.116 \\
\mbox{pretrained on 17 random classes (partition 4),} \mbox{MLP probe,} \mbox{run 2} & 1.069 \\
\mbox{pretrained on 17 random classes (partition 5),} \mbox{MLP probe,} \mbox{run 2} & 1.100 \\
\bottomrule
\end{tabular}
\end{adjustbox}
\end{minipage}
\par\bigskip
\begin{minipage}{\linewidth}\centering
\TableContinued{tab:mlp}{tab:mlp:18}
\begin{adjustbox}{max width=\linewidth, max totalheight=0.92\textheight}
\begin{tabular}{>{\raggedright\arraybackslash}p{\tablabw}c}
\toprule
 & JetClass-II AUC \\
\midrule
\mbox{pretrained on \Kfull{} classes,} \mbox{linear probe on the class token} & 0.9801\,$\pm$\,$<$0.0001 \\
\mbox{pretrained on \Kfull{} classes,} \mbox{linear probe on the 512-unit hidden layer} & 0.9802\,$\pm$\,$<$0.0001 \\
\bottomrule
\end{tabular}
\end{adjustbox}
\end{minipage}
\par\bigskip
\begin{minipage}{\linewidth}\centering
\TableContinued{tab:mlp}{tab:mlp:19}
\begin{adjustbox}{max width=\linewidth, max totalheight=0.92\textheight}
\begin{tabular}{>{\raggedright\arraybackslash}p{\tablabw}cc}
\toprule
 & quark/gluon AUC & quark/gluon $R_{30}$ \\
\midrule
\mbox{pretrained on \Kfull{} classes,} \mbox{linear probe on the class token} & 0.8780\,$\pm$\,0.0059 & 66\,$\pm$\,7 \\
\mbox{pretrained on \Kfull{} classes,} \mbox{linear probe on the 512-unit hidden layer} & 0.8802\,$\pm$\,0.0070 & 70\,$\pm$\,6 \\
\bottomrule
\end{tabular}
\end{adjustbox}
\end{minipage}
\end{center}


\subsection{BatchNorm statistics}\label{sec:app-bn}

Recomputing the BatchNorm statistics tests whether the label-set comparison depends on the statistics
each model accumulated during pretraining. Every linear probe on natural-composition JetClass-II was
repeated on the copy of the selected checkpoint whose BatchNorm statistics were recomputed on
\DesignBnJets{} training jets (table~\ref{tab:bn}); the class-balanced linear probes of the main text use
the checkpoints as trained, which is why this check runs on the natural-composition test set. The ratio of
$1-\AUC$ of a label set to the \Kfull-class model of its seed changes by at most \MlstBnMaxDev,
less than the smallest difference between label sets that the main text reports, so no result rests on
the statistics accumulated in pretraining.
% CHECK: MlstBnMaxDev well below |MlstCmpJcTwoResOverFull - 1| (2.6% today); old value 0.5%

The randomly initialized network enters every comparison in its recomputed form, because its saved
statistics are PyTorch's defaults, zero mean and unit variance, which do not normalize the inputs; a
comparison with it would otherwise set a normalized network against an unnormalized one.
Table~\ref{tab:bn} gives it both ways.

\begin{center}
\small
\begin{minipage}{\linewidth}
\captionof{table}{Linear probes on JetClass-II (natural composition, largest training set) with the BatchNorm statistics as accumulated in pretraining and recomputed on training jets: the AUC averaged over the classes and the ratio $\rho$ of $1-\mathrm{AUC}$, equation~(\ref{eq:ratio}), to the \Kfull-class run of the same seed. The class-balanced probes of the pretrained networks use their checkpoints as trained, so this check uses the natural-composition test set. Mean $\pm$ standard deviation over the \MlstNRunsWord{} pretraining runs of each level of the class hierarchy; configurations with \RandFlavRuns{} runs give both values. Every random partition has 17 classes; each numbered one (1--5) has a single run, named by the pretraining seed it shares with the \Kfull-class run of that seed. The randomly initialized network: three initializations.}
\label{tab:bn}
\end{minipage}
\par\nopagebreak[3]
\begin{adjustbox}{max width=\linewidth}
\begin{tabular}{>{\raggedright\arraybackslash}p{5.1cm}cc}
\toprule
 & \multicolumn{2}{c}{AUC} \\
\cmidrule(lr){2-3}
pretraining & as trained & recomputed \\
\midrule
\Kfull{} classes & 0.9760\,$\pm$\,$<$0.0001 & 0.9760\,$\pm$\,$<$0.0001 \\
\Kres{} classes & 0.9748\,$\pm$\,0.0003 & 0.9748\,$\pm$\,0.0003 \\
\Kmid{} classes & 0.9719\,$\pm$\,0.0006 & 0.9719\,$\pm$\,0.0006 \\
\Kcoarse{} classes & 0.9631\,$\pm$\,0.0006 & 0.9631\,$\pm$\,0.0006 \\
\Kres{} classes, mass regression & 0.9735, 0.9740 & 0.9735, 0.9740 \\
\Kcoarse{} classes, mass regression & 0.9623, 0.9632 & 0.9623, 0.9632 \\
random partition 1, seed 1 & 0.9719 & 0.9719 \\
random partition 2, seed 2 & 0.9719 & 0.9719 \\
random partition 3, seed 1 & 0.9724 & 0.9724 \\
random partition 4, seed 2 & 0.9723 & 0.9723 \\
random partition 5, seed 2 & 0.9720 & 0.9720 \\
random partition, no flavor boundary & 0.9657, 0.9645 & 0.9657, 0.9645 \\
random partition, $b$-quark boundary & 0.9698, 0.9701 & 0.9698, 0.9701 \\
random partition, random flavor boundary & 0.9698, 0.9698 & 0.9698, 0.9698 \\
randomly initialized network & 0.8058\,$\pm$\,0.0018 & 0.8525\,$\pm$\,0.0019 \\
\bottomrule
\end{tabular}
\end{adjustbox}
\par\bigskip
\begin{adjustbox}{max width=\linewidth}
\begin{tabular}{>{\raggedright\arraybackslash}p{5.1cm}cc}
\toprule
 & \multicolumn{2}{c}{$\rho$} \\
\cmidrule(lr){2-3}
pretraining & as trained & recomputed \\
\midrule
\Kfull{} classes & 1.00 & 1.00 \\
\Kres{} classes & 1.049\,$\pm$\,0.014 & 1.050\,$\pm$\,0.014 \\
\Kmid{} classes & 1.171\,$\pm$\,0.024 & 1.172\,$\pm$\,0.025 \\
\Kcoarse{} classes & 1.538\,$\pm$\,0.025 & 1.537\,$\pm$\,0.026 \\
\Kres{} classes, mass regression & 1.11, 1.08 & 1.11, 1.08 \\
\Kcoarse{} classes, mass regression & 1.57, 1.53 & 1.57, 1.53 \\
random partition 1, seed 1 & 1.17 & 1.17 \\
random partition 2, seed 2 & 1.17 & 1.17 \\
random partition 3, seed 1 & 1.15 & 1.15 \\
random partition 4, seed 2 & 1.15 & 1.15 \\
random partition 5, seed 2 & 1.16 & 1.16 \\
random partition, no flavor boundary & 1.43, 1.48 & 1.43, 1.48 \\
random partition, $b$-quark boundary & 1.26, 1.25 & 1.26, 1.25 \\
random partition, random flavor boundary & 1.26, 1.25 & 1.26, 1.25 \\
\bottomrule
\end{tabular}
\end{adjustbox}
\end{center}


\subsection{Test-sample uncertainty}\label{sec:app-boot}

A bootstrap of the test jets shows whether the finite test set or the variation between pretraining runs
is the larger uncertainty. Every model is scored on the same test jets, so the test sample moves all of
them together, and much of its effect cancels in a ratio between models. On JetClass-II the test jets are
resampled with replacement within each class \MlstBootResamples{} times, every model is scored on each
resample, and the standard deviation of the logarithm of a label-set ratio of $1-\AUC$ over the
resamples is its relative uncertainty from the test sample. On top and quark/gluon, the uncertainty of a
ratio of $R_{30}$ follows from the number of background jets that pass the cut. Table~\ref{tab:boot} sets
each beside the standard deviation over pretraining runs. The test sample contributes at most
\MlstBootTestMaxProbe\% to a ratio from the linear probes on JetClass-II and \MlstBootTestMaxFt\% after
fine-tuning there, and at most \MlstBootTestMaxRthirtyProbe\% and \MlstBootTestMaxRthirtyFt\% to a ratio
of $R_{30}$ on top and quark/gluon from the linear probe and from fine-tuning; in each case this is below
the standard deviation over pretraining runs of the same ratio.
% CHECK: pending MlstBootTestMaxProbe, MlstBootTestMaxFt (class-balanced bootstraps); each test-sample SD below the seed SD in tab:boot

\begin{table*}[t]
\centering
\caption{Linear probes on the frozen features and fine-tuning of networks pretrained on JetClass-II with 188, 162, 43 and 17 classes: uncertainty (percent) of the comparisons between their label sets at the largest training set, as the SD of the logarithm of the ratio over the pretraining runs / that from the finite test sample. JetClass-II: the ratio of $1-\mathrm{AUC}$; top and quark/gluon: the ratio of $R_{30}$.}
\label{tab:boot}
\begin{tabular*}{\linewidth}{@{\extracolsep{\fill}}lc@{}}
\toprule
 & \multicolumn{1}{c}{linear probe on frozen features} \\
\cmidrule(lr){2-2}
ratio & JetClass-II, $1-\mathrm{AUC}$ \\
\midrule
\Kres{} classes / \Kfull{} classes & 0.45 / 0.15 \\
\Kmid{} classes / \Kfull{} classes & 0.70 / 0.16 \\
\Kcoarse{} classes / \Kfull{} classes & 1.3 / 0.21 \\
\bottomrule
\end{tabular*}
\par\bigskip
\begin{tabular*}{\linewidth}{@{\extracolsep{\fill}}lcc@{}}
\toprule
 & \multicolumn{2}{c}{linear probe on frozen features} \\
\cmidrule(lr){2-3}
ratio & top, $R_{30}$ & quark/gluon, $R_{30}$ \\
\midrule
\Kres{} classes / \Kfull{} classes & 25 / 2.6 & 24 / 3.5 \\
\Kmid{} classes / \Kfull{} classes & 29 / 2.4 & 20 / 3.6 \\
\Kcoarse{} classes / \Kfull{} classes & 16 / 2.4 & 9.8 / 3.6 \\
\bottomrule
\end{tabular*}
\par\bigskip
\begin{tabular*}{\linewidth}{@{\extracolsep{\fill}}lc@{}}
\toprule
 & \multicolumn{1}{c}{Sophon's network, fine-tuned} \\
\cmidrule(lr){2-2}
ratio & JetClass-II, $1-\mathrm{AUC}$ \\
\midrule
\Kres{} classes / \Kfull{} classes & 1.1 / 0.13 \\
\Kmid{} classes / \Kfull{} classes & 0.20 / 0.13 \\
\Kcoarse{} classes / \Kfull{} classes & 0.29 / 0.15 \\
\bottomrule
\end{tabular*}
\par\bigskip
\begin{tabular*}{\linewidth}{@{\extracolsep{\fill}}lcc@{}}
\toprule
 & \multicolumn{2}{c}{Sophon's network, fine-tuned} \\
\cmidrule(lr){2-3}
ratio & top, $R_{30}$ & quark/gluon, $R_{30}$ \\
\midrule
\Kres{} classes / \Kfull{} classes & 3.2 / 13 & 2.6 / 4.6 \\
\Kmid{} classes / \Kfull{} classes & 7.3 / 13 & 2.8 / 4.6 \\
\Kcoarse{} classes / \Kfull{} classes & 9.9 / 13 & 2.6 / 4.6 \\
\bottomrule
\end{tabular*}
\end{table*}


\subsection{Class-balanced and natural-composition test sets}\label{sec:app-balance}

Comparing two test sets tests whether class balancing changes the order of the label sets on JetClass-II.
Table~\ref{tab:balance} gives the linear probes with \FtSizeESix{} training jets on the
natural-composition test set, whose class shares follow the selection, beside those on the class-balanced
set of the main text. The label sets fall in the same order on both.
% CHECK: tab:balance gives the order 188 < 162 < 43 < 17 in the ratio of 1-AUC on both test sets (v3 natural 1.171 and 1.538; v4 balanced 1.227 and 1.557 for 43 and 17)

\begin{landscape}
{\small
\setlength{\LTleft}{0pt}\setlength{\LTright}{0pt}\setlength{\LTcapwidth}{\linewidth}\makeatletter\def\LT@makecaption#1#2#3{\LT@mcol\LT@cols c{\hbox to\z@{\hss\parbox[t]\LTcapwidth{\normalsize #1{#2: }#3\par\vskip\belowcaptionskip}\hss}}}\makeatother
\begin{center}
\begin{minipage}{\linewidth}\centering
\normalsize
\captionof{table}{Linear probes on the frozen features of networks pretrained on JetClass-II with 188, 162, 43 and 17 classes and of the controls, trained on $10^{6}$ JetClass-II jets and scored on the natural-composition test set and on the class-balanced one: AUC averaged over the classes and the ratio $\rho$ of $1-\mathrm{AUC}$, equation~(\ref{eq:ratio}), to the \Kfull-class run of the same seed; (LP): linear probe on the frozen features. Mean $\pm$ SD over the \MlstNRunsWord{} pretraining runs of each level of the class hierarchy; configurations with \RandFlavRuns{} runs: over those. Every random partition has 17 classes; each numbered one (1--5) has a single run, numbered by the pretraining seed it shares with the \Kfull-class run of that seed. The random network without pretraining: three initializations.}
\label{tab:balance}
\par\medskip
\small
\begin{tabular*}{\linewidth}{@{\extracolsep{\fill}}lcccc@{}}
\toprule
 & \multicolumn{2}{c}{natural composition} & \multicolumn{2}{c}{class-balanced} \\
\cmidrule(lr){2-3}\cmidrule(lr){4-5}
 & AUC & $\rho$ & AUC & $\rho$ \\
\midrule
\Kfull{} classes (LP) & 0.9760\,$\pm$\,$<$0.0001 & 1.000 & 0.9801\,$\pm$\,$<$0.0001 & 1.000 \\
\Kres{} classes (LP) & 0.9748\,$\pm$\,0.0003 & 1.049\,$\pm$\,0.014 & 0.9796\,$\pm$\,0.0001 & 1.026\,$\pm$\,0.005 \\
\Kmid{} classes (LP) & 0.9719\,$\pm$\,0.0006 & 1.171\,$\pm$\,0.024 & 0.9756\,$\pm$\,0.0002 & 1.227\,$\pm$\,0.009 \\
\Kcoarse{} classes (LP) & 0.9631\,$\pm$\,0.0006 & 1.538\,$\pm$\,0.025 & 0.9691\,$\pm$\,0.0004 & 1.557\,$\pm$\,0.019 \\
\Kres{} classes + mass (LP) & 0.9737\,$\pm$\,0.0003 & 1.094\,$\pm$\,0.017 & 0.9792\,$\pm$\,0.0002 & 1.045\,$\pm$\,0.011 \\
\Kcoarse{} classes + mass (LP) & 0.9628\,$\pm$\,0.0006 & 1.551\,$\pm$\,0.029 & 0.9688\,$\pm$\,0.0013 & 1.573\,$\pm$\,0.063 \\
random partition 1 (LP), seed 1 & 0.9719 & 1.173 & 0.9776 & 1.129 \\
random partition 2 (LP), seed 2 & 0.9719 & 1.167 & 0.9776 & 1.130 \\
random partition 3 (LP), seed 1 & 0.9724 & 1.153 & 0.9779 & 1.115 \\
random partition 4 (LP), seed 2 & 0.9723 & 1.151 & 0.9778 & 1.117 \\
random partition 5 (LP), seed 2 & 0.9720 & 1.164 & 0.9775 & 1.134 \\
no flavor boundary (LP) & 0.9651\,$\pm$\,0.0009 & 1.455\,$\pm$\,0.032 & 0.9713\,$\pm$\,0.0005 & 1.443\,$\pm$\,0.027 \\
$b$-quark boundary (LP) & 0.9699\,$\pm$\,0.0002 & 1.253\,$\pm$\,0.011 & 0.9759\,$\pm$\,0.0001 & 1.215\,$\pm$\,0.004 \\
random flavor boundary (LP) & 0.9698\,$\pm$\,$<$0.0001 & 1.257\,$\pm$\,0.004 & 0.9760\,$\pm$\,$<$0.0001 & 1.207\,$\pm$\,0.001 \\
random network, BatchNorm recomputed (LP) & 0.8525\,$\pm$\,0.0019 & -- & 0.8664\,$\pm$\,0.0036 & -- \\
\bottomrule
\end{tabular*}
\end{minipage}
\end{center}
}
\end{landscape}


\subsection{Reproduction of ParT on top tagging}\label{sec:app-partrepro}

To locate the gap between our network trained from scratch and ParT on top tagging, ParT's top-tagging
training script~\cite{part} was run on our cluster unchanged and then with Sophon's changes added one at a
time (table~\ref{tab:partrepro}). Unchanged, it reaches a rejection at 30\% signal efficiency of
\MlstPartreproRthirtyPart, within twice the published standard deviation of ParT's
\LitPartTopRejScratch, so ParT reproduces here. Adding Sophon's output head gives
\MlstPartreproRthirtyHeadOverPartPr{} than the unchanged script, which stays within ParT's published
spread. Adding Sophon's rescaling of every jet to a common transverse momentum on top of the head gives
\MlstPartreproRthirtyRescaleOverPartPr{} than the unchanged script. Each variant is a single run. The
network trained from scratch in table~\ref{tab:benchtop} keeps the rescaling because the pretrained models use it, so the comparison of
pretrained and from-scratch networks shares every preprocessing step.
% CHECK: |MlstPartreproRthirtyPart - 1602| < 2 x 81 and |MlstPartreproRthirtyHead - 1602| < 81 (1,711 and 1,577 today)

\begin{landscape}
\begin{table}[t]
\centering
\caption{The top-tagging training script of the Particle Transformer run from scratch, compared unchanged and with one of two changes taken from Sophon, which are not combined. The changes are Sophon's output head, a 512-unit hidden layer before the output in place of a single linear layer, and its rescaling of each jet to $p_{\mathrm{T}}=500$~GeV. AUC is the area under the receiver operating characteristic curve, $R_{50}$ and $R_{30}$ the inverse background efficiencies at 50\% and 30\% signal efficiency, and the last column counts the background jets passing the $R_{30}$ selection. Table~\ref{tab:published} has the published result and our network from scratch, with both changes. Every value is from a single run, and $R_{30}$, resting on 118 to 165 background jets, has a Poisson uncertainty of about 8\%, larger than the change from Sophon's output head.}
\label{tab:partrepro}
\begin{tabular*}{\linewidth}{@{\extracolsep{\fill}}lccccc@{}}
\toprule
 & accuracy & AUC & $R_{50}$ & $R_{30}$ & background jets passing at $R_{30}$ \\
\midrule
Particle Transformer training script, unchanged & 0.93979 & 0.98585 & 424 & 1{,}710 & 118 \\
the same script with Sophon's output head & 0.94033 & 0.98586 & 416 & 1{,}580 & 128 \\
the same script with Sophon's jet-$p_{\mathrm{T}}$ rescaling & 0.93852 & 0.98504 & 357 & 1{,}220 & 165 \\
\bottomrule
\end{tabular*}
\end{table}
\end{landscape}


% ----------------------------------------------------------------------------
\section{Complete results}\label{sec:app-results}
% ----------------------------------------------------------------------------

The tables and figures of this appendix complete the results of the main text, each for the item named
with it. Their uncertainties follow the convention of the main text, and a configuration with two runs
shows both values.

\subsection{Linear probes and fine-tuning for every configuration}\label{sec:app-all}

Table~\ref{tab:probes} gives the linear probes at the largest training set of every task for every
configuration, each random partition included; on JetClass-II it gives the accuracy beside the averages
of the AUC defined in appendix~\ref{sec:app-probes}. Table~\ref{tab:probesn} gives the linear probes
against the number of training jets in each task's metric. Figure~\ref{fig:a10} draws the
four label sets against the number of training jets, to show whether their order holds at every size.
Fine-tuning is in table~\ref{tab:ft} against the number of training jets and in table~\ref{tab:ftrej} on
the full top and quark/gluon training sets, for every configuration. Table~\ref{tab:perrun} gives each
pretraining run's own values.

\providecommand{\TableContinued}[2]{\label{#2}\ifnum\getpagerefnumber{#2}>\getpagerefnumber{#1}\relax\makebox[\linewidth][l]{\textit{Table~\ref{#1} (continued)}}\par\smallskip\fi}
\begin{center}
\setlength{\tabcolsep}{4pt}
\begin{minipage}{\linewidth}\centering
\captionof{table}{Linear probes on the frozen features of the networks pretrained with 188, 162, 43 and 17 classes and of the control networks, compared on five test sets at each task's largest training set. On JetClass-II and JetClass, AUC is the area under the receiver operating characteristic curve, averaged over the classes, each against the rest. $R_{50}$ and $R_{30}$ are the inverse background efficiencies at 50\% and 30\% signal efficiency. The inputs that the top and quark/gluon data lack are replaced by \LpSubstitutionTop{} for top and by \LpSubstitutionQg{} for quark/gluon, as chosen on validation. The network with kinematic inputs is pretrained on those inputs only and lacks none, and the released Sophon checkpoint has the missing inputs set to zero. A flavor boundary is a class boundary between decays that differ only in quark flavor. Values are the mean $\pm$ standard deviation over \MlstNRunsWord{} pretraining runs per label set, \RandFlavRuns{} per flavor-boundary partition, \MlstMassNRunsWord{} per network with an added mass output, one for each of five random partitions and three initializations of the random network. Bold marks the best value of each block of a column and every value within three standard deviations of it. $^{\S}$ marks a value from a single run.}
\label{tab:probes}
\par\medskip
\small
\begin{tabular*}{\linewidth}{@{\extracolsep{\fill}}lcccc@{}}
\toprule
 & accuracy & AUC & $R_{50}$ & $R_{30}$ \\
\midrule
\multicolumn{5}{@{}l}{\textit{linear probe on the frozen features, JetClass-II}} \\
\Kfull{} classes & {\fontseries{b}\selectfont 0.4348\,$\pm$\,0.0003} & {\fontseries{b}\selectfont 0.98014\,$\pm$\,0.00001} & -- & -- \\
\Kres{} classes & {\fontseries{b}\selectfont 0.4328\,$\pm$\,0.0008} & 0.97963\,$\pm$\,0.00010 & -- & -- \\
\Kmid{} classes & 0.3609\,$\pm$\,0.0046 & 0.97562\,$\pm$\,0.00016 & -- & -- \\
\Kcoarse{} classes & 0.3341\,$\pm$\,0.0067 & 0.96908\,$\pm$\,0.00038 & -- & -- \\
\multicolumn{5}{@{}l}{17-class partitions} \\
\quad no flavor boundary & 0.3746\,$\pm$\,0.0033 & 0.97134\,$\pm$\,0.00053 & -- & -- \\
\quad $b$-quark boundary & 0.3954\,$\pm$\,0.0013 & 0.97586\,$\pm$\,0.00010 & -- & -- \\
\quad random flavor boundary & 0.4080\,$\pm$\,0.0013 & 0.97602\,$\pm$\,0.00002 & -- & -- \\
\quad random partitions & 0.4216\,$\pm$\,0.0014 & 0.97766\,$\pm$\,0.00016 & -- & -- \\
\Kres{} classes + mass & {\fontseries{b}\selectfont 0.4288\,$\pm$\,0.0024} & 0.97924\,$\pm$\,0.00023 & -- & -- \\
\Kcoarse{} classes + mass & 0.3311\,$\pm$\,0.0264 & 0.96875\,$\pm$\,0.00126 & -- & -- \\
random network & 0.0724\,$\pm$\,0.0015 & 0.81212\,$\pm$\,0.00190 & -- & -- \\
\midrule
released Sophon checkpoint & 0.4340$^{\S}$ & 0.97990$^{\S}$ & -- & -- \\
\midrule
\multicolumn{5}{@{}l}{\textit{linear probe on the frozen features, JetClass}} \\
\Kfull{} classes & {\fontseries{b}\selectfont 0.771\,$\pm$\,0.002} & {\fontseries{b}\selectfont 0.9682\,$\pm$\,0.0006} & -- & -- \\
\Kres{} classes & {\fontseries{b}\selectfont 0.771\,$\pm$\,0.001} & {\fontseries{b}\selectfont 0.9683\,$\pm$\,0.0002} & -- & -- \\
\Kmid{} classes & {\fontseries{b}\selectfont 0.770\,$\pm$\,0.001} & {\fontseries{b}\selectfont 0.9680\,$\pm$\,0.0004} & -- & -- \\
\Kcoarse{} classes & 0.761\,$\pm$\,0.003 & 0.9650\,$\pm$\,0.0008 & -- & -- \\
\multicolumn{5}{@{}l}{17-class partitions} \\
\quad no flavor boundary & 0.751\,$\pm$\,0.003 & 0.9628\,$\pm$\,0.0008 & -- & -- \\
\quad $b$-quark boundary & 0.763\,$\pm$\,0.002 & {\fontseries{b}\selectfont 0.9663\,$\pm$\,0.0003} & -- & -- \\
\quad random flavor boundary & 0.758\,$\pm$\,0.002 & 0.9649\,$\pm$\,0.0006 & -- & -- \\
\quad random partitions & 0.754\,$\pm$\,0.001 & 0.9638\,$\pm$\,0.0002 & -- & -- \\
\Kres{} classes + mass & {\fontseries{b}\selectfont 0.774\,$\pm$\,0.003} & {\fontseries{b}\selectfont 0.9690\,$\pm$\,0.0009} & -- & -- \\
\Kcoarse{} classes + mass & 0.762\,$\pm$\,0.001 & 0.9653\,$\pm$\,0.0003 & -- & -- \\
random network & 0.408\,$\pm$\,0.001 & 0.8280\,$\pm$\,0.0011 & -- & -- \\
\midrule
released Sophon checkpoint & 0.768$^{\S}$ & 0.9675$^{\S}$ & -- & -- \\
\bottomrule
\end{tabular*}
\end{minipage}
\par\bigskip
\begin{minipage}{\linewidth}\centering
\TableContinued{tab:probes}{tab:probes:2}
\small
\begin{tabular*}{\linewidth}{@{\extracolsep{\fill}}lcccc@{}}
\toprule
 & accuracy & AUC & $R_{50}$ & $R_{30}$ \\
\midrule
\multicolumn{5}{@{}l}{\textit{linear probe on the frozen features, top tagging}} \\
\Kfull{} classes & 0.9112\,$\pm$\,0.0015 & 0.9686\,$\pm$\,0.0007 & 87\,$\pm$\,4 & {\fontseries{b}\selectfont 228\,$\pm$\,25} \\
\Kres{} classes & {\fontseries{b}\selectfont 0.9098\,$\pm$\,0.0031} & 0.9671\,$\pm$\,0.0020 & 79\,$\pm$\,9 & {\fontseries{b}\selectfont 199\,$\pm$\,25} \\
\Kmid{} classes & 0.9101\,$\pm$\,0.0019 & 0.9681\,$\pm$\,0.0013 & 88\,$\pm$\,3 & {\fontseries{b}\selectfont 231\,$\pm$\,8} \\
\Kcoarse{} classes & {\fontseries{b}\selectfont 0.9152\,$\pm$\,0.0010} & {\fontseries{b}\selectfont 0.9710\,$\pm$\,0.0006} & {\fontseries{b}\selectfont 97\,$\pm$\,4} & {\fontseries{b}\selectfont 255\,$\pm$\,17} \\
\multicolumn{5}{@{}l}{17-class partitions} \\
\quad no flavor boundary & {\fontseries{b}\selectfont 0.9124\,$\pm$\,0.0029} & {\fontseries{b}\selectfont 0.9696\,$\pm$\,0.0018} & {\fontseries{b}\selectfont 93\,$\pm$\,8} & {\fontseries{b}\selectfont 253\,$\pm$\,13} \\
\quad $b$-quark boundary & 0.9075\,$\pm$\,0.0017 & 0.9665\,$\pm$\,0.0011 & 80\,$\pm$\,6 & {\fontseries{b}\selectfont 207\,$\pm$\,6} \\
\quad random flavor boundary & 0.9088\,$\pm$\,0.0004 & 0.9672\,$\pm$\,0.0002 & 84\,$\pm$\,1 & {\fontseries{b}\selectfont 224\,$\pm$\,13} \\
\quad random partitions & 0.9029\,$\pm$\,0.0035 & 0.9626\,$\pm$\,0.0019 & 64\,$\pm$\,4 & 163\,$\pm$\,11 \\
\Kres{} classes + mass & {\fontseries{b}\selectfont 0.9145\,$\pm$\,0.0030} & {\fontseries{b}\selectfont 0.9709\,$\pm$\,0.0009} & {\fontseries{b}\selectfont 100\,$\pm$\,5} & {\fontseries{b}\selectfont 283\,$\pm$\,33} \\
\Kcoarse{} classes + mass & {\fontseries{b}\selectfont 0.9196\,$\pm$\,0.0015} & {\fontseries{b}\selectfont 0.9741\,$\pm$\,0.0010} & {\fontseries{b}\selectfont 118\,$\pm$\,9} & {\fontseries{b}\selectfont 342\,$\pm$\,51} \\
\Kfull{} classes, kinematic inputs & {\fontseries{b}\selectfont 0.9163}$^{\S}$ & {\fontseries{b}\selectfont 0.9713}$^{\S}$ & {\fontseries{b}\selectfont 94}$^{\S}$ & {\fontseries{b}\selectfont 231}$^{\S}$ \\
random network & 0.8867\,$\pm$\,0.0013 & 0.9408\,$\pm$\,0.0012 & 25\,$\pm$\,1 & 57\,$\pm$\,3 \\
\midrule
released Sophon checkpoint & 0.8882$^{\S}$ & 0.9509$^{\S}$ & 39$^{\S}$ & 97$^{\S}$ \\
\midrule
\multicolumn{5}{@{}l}{\textit{linear probe on the frozen features, quark/gluon tagging}} \\
\Kfull{} classes & {\fontseries{b}\selectfont 0.8076\,$\pm$\,0.0004} & {\fontseries{b}\selectfont 0.8835\,$\pm$\,0.0002} & {\fontseries{b}\selectfont 26.17\,$\pm$\,0.35} & {\fontseries{b}\selectfont 72.7\,$\pm$\,2.4} \\
\Kres{} classes & {\fontseries{b}\selectfont 0.8051\,$\pm$\,0.0015} & {\fontseries{b}\selectfont 0.8818\,$\pm$\,0.0018} & {\fontseries{b}\selectfont 25.59\,$\pm$\,0.41} & {\fontseries{b}\selectfont 71.0\,$\pm$\,0.3} \\
\Kmid{} classes & {\fontseries{b}\selectfont 0.8055\,$\pm$\,0.0009} & {\fontseries{b}\selectfont 0.8821\,$\pm$\,0.0011} & {\fontseries{b}\selectfont 25.68\,$\pm$\,0.64} & {\fontseries{b}\selectfont 70.8\,$\pm$\,1.8} \\
\Kcoarse{} classes & {\fontseries{b}\selectfont 0.8064\,$\pm$\,0.0014} & {\fontseries{b}\selectfont 0.8821\,$\pm$\,0.0015} & {\fontseries{b}\selectfont 25.79\,$\pm$\,0.73} & {\fontseries{b}\selectfont 74.3\,$\pm$\,1.9} \\
\multicolumn{5}{@{}l}{17-class partitions} \\
\quad no flavor boundary & {\fontseries{b}\selectfont 0.8031\,$\pm$\,0.0036} & {\fontseries{b}\selectfont 0.8798\,$\pm$\,0.0035} & {\fontseries{b}\selectfont 25.09\,$\pm$\,0.65} & {\fontseries{b}\selectfont 72.8\,$\pm$\,1.5} \\
\quad $b$-quark boundary & 0.8034\,$\pm$\,0.0013 & {\fontseries{b}\selectfont 0.8805\,$\pm$\,0.0012} & 24.77\,$\pm$\,0.28 & {\fontseries{b}\selectfont 72.4\,$\pm$\,0.2} \\
\quad random flavor boundary & {\fontseries{b}\selectfont 0.8050\,$\pm$\,0.0011} & {\fontseries{b}\selectfont 0.8818\,$\pm$\,0.0016} & {\fontseries{b}\selectfont 25.34\,$\pm$\,0.22} & {\fontseries{b}\selectfont 72.0\,$\pm$\,1.2} \\
\quad random partitions & {\fontseries{b}\selectfont 0.8027\,$\pm$\,0.0031} & {\fontseries{b}\selectfont 0.8788\,$\pm$\,0.0024} & {\fontseries{b}\selectfont 24.08\,$\pm$\,0.80} & {\fontseries{b}\selectfont 68.3\,$\pm$\,2.1} \\
\Kres{} classes + mass & {\fontseries{b}\selectfont 0.8074\,$\pm$\,0.0002} & {\fontseries{b}\selectfont 0.8834\,$\pm$\,0.0005} & {\fontseries{b}\selectfont 25.61\,$\pm$\,0.11} & {\fontseries{b}\selectfont 72.3\,$\pm$\,2.3} \\
\Kcoarse{} classes + mass & {\fontseries{b}\selectfont 0.8051\,$\pm$\,0.0016} & 0.8801\,$\pm$\,0.0006 & {\fontseries{b}\selectfont 25.84\,$\pm$\,0.02} & {\fontseries{b}\selectfont 70.3\,$\pm$\,2.7} \\
random network & 0.7851\,$\pm$\,0.0018 & 0.8612\,$\pm$\,0.0016 & 19.40\,$\pm$\,0.28 & 56.4\,$\pm$\,1.0 \\
\midrule
released Sophon checkpoint & 0.7917$^{\S}$ & 0.8656$^{\S}$ & 19.30$^{\S}$ & 54.3$^{\S}$ \\
\midrule
\multicolumn{5}{@{}l}{\textit{linear probe on the frozen features, quark/gluon tagging, scored on Herwig jets}} \\
\Kfull{} classes & {\fontseries{b}\selectfont 0.7229\,$\pm$\,0.0033} & {\fontseries{b}\selectfont 0.8075\,$\pm$\,0.0006} & {\fontseries{b}\selectfont 11.81\,$\pm$\,0.10} & {\fontseries{b}\selectfont 35.0\,$\pm$\,2.6} \\
\Kres{} classes & {\fontseries{b}\selectfont 0.7234\,$\pm$\,0.0022} & {\fontseries{b}\selectfont 0.8068\,$\pm$\,0.0022} & {\fontseries{b}\selectfont 11.63\,$\pm$\,0.43} & {\fontseries{b}\selectfont 35.9\,$\pm$\,1.7} \\
\Kmid{} classes & 0.7206\,$\pm$\,0.0005 & {\fontseries{b}\selectfont 0.8058\,$\pm$\,0.0011} & {\fontseries{b}\selectfont 11.52\,$\pm$\,0.23} & {\fontseries{b}\selectfont 34.4\,$\pm$\,1.7} \\
\Kcoarse{} classes & {\fontseries{b}\selectfont 0.7201\,$\pm$\,0.0028} & {\fontseries{b}\selectfont 0.8053\,$\pm$\,0.0020} & {\fontseries{b}\selectfont 11.50\,$\pm$\,0.16} & {\fontseries{b}\selectfont 35.5\,$\pm$\,1.0} \\
\multicolumn{5}{@{}l}{17-class partitions} \\
\quad no flavor boundary & {\fontseries{b}\selectfont 0.7206\,$\pm$\,0.0009} & 0.8057\,$\pm$\,0.0002 & {\fontseries{b}\selectfont 11.47\,$\pm$\,0.16} & {\fontseries{b}\selectfont 34.6\,$\pm$\,0.4} \\
\quad $b$-quark boundary & {\fontseries{b}\selectfont 0.7214\,$\pm$\,0.0007} & 0.8058\,$\pm$\,0.0002 & {\fontseries{b}\selectfont 11.56\,$\pm$\,0.13} & {\fontseries{b}\selectfont 34.9\,$\pm$\,0.5} \\
\quad random flavor boundary & {\fontseries{b}\selectfont 0.7216\,$\pm$\,0.0022} & 0.8062\,$\pm$\,0.0003 & 11.66\,$\pm$\,0.08 & {\fontseries{b}\selectfont 35.3\,$\pm$\,0.4} \\
\quad random partitions & {\fontseries{b}\selectfont 0.7212\,$\pm$\,0.0034} & {\fontseries{b}\selectfont 0.8046\,$\pm$\,0.0025} & {\fontseries{b}\selectfont 11.34\,$\pm$\,0.35} & {\fontseries{b}\selectfont 34.0\,$\pm$\,1.4} \\
\Kres{} classes + mass & {\fontseries{b}\selectfont 0.7223\,$\pm$\,0.0006} & {\fontseries{b}\selectfont 0.8071\,$\pm$\,0.0010} & {\fontseries{b}\selectfont 11.83\,$\pm$\,0.17} & {\fontseries{b}\selectfont 35.3\,$\pm$\,0.5} \\
\Kcoarse{} classes + mass & {\fontseries{b}\selectfont 0.7248\,$\pm$\,0.0011} & {\fontseries{b}\selectfont 0.8079\,$\pm$\,0.0001} & {\fontseries{b}\selectfont 11.92\,$\pm$\,0.01} & {\fontseries{b}\selectfont 35.5\,$\pm$\,1.1} \\
random network & 0.7205\,$\pm$\,0.0021 & 0.7965\,$\pm$\,0.0009 & 10.58\,$\pm$\,0.11 & 31.6\,$\pm$\,0.8 \\
\midrule
released Sophon checkpoint & 0.7153$^{\S}$ & 0.7964$^{\S}$ & 10.32$^{\S}$ & 29.9$^{\S}$ \\
\bottomrule
\end{tabular*}
\end{minipage}
\end{center}

\begin{table*}[t]
\centering
\small
\caption{Linear probes against the number of training jets (column headers): JetClass-II AUC averaged over the classes, JetClass accuracy, and $R_{30}$ on top, quark/gluon and Herwig jets with the missing inputs substituted as chosen on validation (\LpSubstitutionTop{} for top, \LpSubstitutionQg{} for quark/gluon). Mean $\pm$ standard deviation over the \MlstNRunsWord{} pretraining runs. $^{*}$: a probe that did not converge.}
\label{tab:probesn}
\begin{adjustbox}{max width=\linewidth}
\begin{tabular}{lcccc}
\toprule
JetClass-II, AUC & 1{,}000 & 10{,}000 & 100{,}000 & 1{,}000{,}000 \\
\midrule
\Kfull{} classes & 0.9497\,$\pm$\,0.0001 & 0.9701\,$\pm$\,0.0001 & 0.9783\,$\pm$\,0.0001 & 0.9801\,$\pm$\,$<$0.0001 \\
\Kres{} classes & 0.9466\,$\pm$\,0.0006 & 0.9692\,$\pm$\,0.0002 & 0.9778\,$\pm$\,0.0001 & 0.9796\,$\pm$\,0.0001 \\
\Kmid{} classes & 0.9324\,$\pm$\,0.0029 & 0.9633\,$\pm$\,0.0011 & 0.9735\,$\pm$\,0.0003 & 0.9756\,$\pm$\,0.0002 \\
\Kcoarse{} classes & 0.9046\,$\pm$\,0.0021 & 0.9508\,$\pm$\,0.0006 & 0.9660\,$\pm$\,0.0004 & 0.9691\,$\pm$\,0.0004 \\
\Kres{} classes, mass regression & 0.9445, 0.9435 & 0.9684, 0.9684 & 0.9772, 0.9775 & 0.9791, 0.9794 \\
\Kcoarse{} classes, mass regression & 0.8999, 0.9044 & 0.9483, 0.9518 & 0.9646, 0.9665 & 0.9679, 0.9696 \\
randomly initialized network & 0.618\,$\pm$\,0.004 & 0.718\,$\pm$\,0.004 & 0.7946\,$\pm$\,0.0023 & 0.8121\,$\pm$\,0.0019 \\
randomly initialized network, BatchNorm recomputed & 0.67\,$\pm$\,0.04 & 0.800\,$\pm$\,0.013 & 0.855\,$\pm$\,0.004 & 0.866\,$\pm$\,0.004 \\
\bottomrule
\end{tabular}
\end{adjustbox}
\par\medskip
\begin{adjustbox}{max width=\linewidth}
\begin{tabular}{lcccc}
\toprule
JetClass, accuracy & 1{,}000 & 10{,}000 & 100{,}000 & 1{,}000{,}000 \\
\midrule
\Kfull{} classes & 0.711\,$\pm$\,0.006 & 0.7529\,$\pm$\,0.0008 & 0.7686\,$\pm$\,0.0022 & 0.7709\,$\pm$\,0.0023 \\
\Kres{} classes & 0.714\,$\pm$\,0.008 & 0.7537\,$\pm$\,0.0018 & 0.7687\,$\pm$\,0.0002 & 0.7709\,$\pm$\,0.0005 \\
\Kmid{} classes & 0.7109\,$\pm$\,0.0005 & 0.7521\,$\pm$\,0.0004 & 0.7679\,$\pm$\,0.0012 & 0.7695\,$\pm$\,0.0014 \\
\Kcoarse{} classes & 0.689\,$\pm$\,0.006 & 0.742\,$\pm$\,0.004 & 0.7586\,$\pm$\,0.0026 & 0.7605\,$\pm$\,0.0027 \\
\Kres{} classes, mass regression & 0.7083, 0.7093 & 0.7543, 0.7574 & 0.7694, 0.7741 & 0.7718, 0.7760 \\
\Kcoarse{} classes, mass regression & 0.6863, 0.6947 & 0.7447, 0.7426 & 0.7606, 0.7595 & 0.7624, 0.7612 \\
randomly initialized network & 0.306\,$\pm$\,0.006 & 0.3854\,$\pm$\,0.0016 & 0.4049\,$\pm$\,0.0016 & 0.4077\,$\pm$\,0.0013 \\
randomly initialized network, BatchNorm recomputed & 0.363\,$\pm$\,0.009 & 0.450\,$\pm$\,0.005 & 0.472\,$\pm$\,0.006 & 0.475\,$\pm$\,0.005 \\
\bottomrule
\end{tabular}
\end{adjustbox}
\par\medskip
\begin{adjustbox}{max width=\linewidth}
\begin{tabular}{lcccc}
\toprule
top, $R_{30}$ (missing inputs: random draws) & 1{,}000 & 10{,}000 & 100{,}000 & 1{,}200{,}000 \\
\midrule
\Kfull{} classes & 162\,$\pm$\,9 & 227\,$\pm$\,18 & 228\,$\pm$\,23 & 228\,$\pm$\,25 \\
\Kres{} classes & 130\,$\pm$\,8 & 199\,$\pm$\,21 & 203\,$\pm$\,26 & 199\,$\pm$\,25 \\
\Kmid{} classes & 146\,$\pm$\,32 & 233\,$\pm$\,18 & 236\,$\pm$\,10 & 231\,$\pm$\,8 \\
\Kcoarse{} classes & 154\,$\pm$\,5 & 252\,$\pm$\,19 & 259\,$\pm$\,20 & 255\,$\pm$\,17 \\
\Kres{} classes, mass regression & 199, 186 & 248, 306 & 266, 312 & 260, 306 \\
\Kcoarse{} classes, mass regression & 193.0, 191.0 & 286, 341 & 306, 379 & 306, 378 \\
randomly initialized network & 43.1\,$\pm$\,1.1 & 55.3\,$\pm$\,3.1 & 57.5\,$\pm$\,3.0 & 56.8\,$\pm$\,2.7 \\
randomly initialized network, BatchNorm recomputed & 36.9\,$\pm$\,2.6 & 49.0\,$\pm$\,3.3 & 49.2\,$\pm$\,3.1 & 48.6\,$\pm$\,3.0 \\
\bottomrule
\end{tabular}
\end{adjustbox}
\par\medskip
\begin{adjustbox}{max width=\linewidth}
\begin{tabular}{lcccc}
\toprule
quark/gluon, $R_{30}$ (missing inputs: the training mean) & 1{,}000 & 10{,}000 & 100{,}000 & 1{,}600{,}000 \\
\midrule
\Kfull{} classes & 66.5\,$\pm$\,1.7 & 70.5\,$\pm$\,1.8 & 73.4\,$\pm$\,2.0 & 72.7\,$\pm$\,2.4 \\
\Kres{} classes & 60.2\,$\pm$\,1.6 & 68.9\,$\pm$\,0.9 & 71.5\,$\pm$\,0.8 & 70.99\,$\pm$\,0.31 \\
\Kmid{} classes & 63.4\,$\pm$\,2.7 & 66.9\,$\pm$\,0.6 & 70.8\,$\pm$\,1.4 & 70.8\,$\pm$\,1.8 \\
\Kcoarse{} classes & 62\,$\pm$\,4 & 72.0\,$\pm$\,1.8 & 74.6\,$\pm$\,1.5 & 74.3\,$\pm$\,1.9 \\
\Kres{} classes, mass regression & 67.8, 60.6 & 70.1, 66.4 & 74.3, 71.9 & 73.9, 70.7 \\
\Kcoarse{} classes, mass regression & 66.8, 60.7 & 64.6, 69.3 & 68.9, 70.9 & 68.4, 72.2 \\
randomly initialized network & 43.0\,$\pm$\,2.5 & 55.6\,$\pm$\,1.2 & 56.7\,$\pm$\,0.5 & 56.4\,$\pm$\,1.0 \\
randomly initialized network, BatchNorm recomputed & 42.0\,$\pm$\,2.3 & 51.1\,$\pm$\,1.4 & 52.2\,$\pm$\,1.3 & 52.2\,$\pm$\,0.8 \\
\bottomrule
\end{tabular}
\end{adjustbox}
\par\medskip
\begin{adjustbox}{max width=\linewidth}
\begin{tabular}{lcccc}
\toprule
quark/gluon, Herwig, $R_{30}$ (missing inputs: the training mean) & 1{,}000 & 10{,}000 & 100{,}000 & 1{,}600{,}000 \\
\midrule
\Kfull{} classes & 32.99\,$\pm$\,0.13 & 34.5\,$\pm$\,1.2 & 34.5\,$\pm$\,2.4 & 35.0\,$\pm$\,2.6 \\
\Kres{} classes & 31.3\,$\pm$\,2.5 & 35.0\,$\pm$\,2.2 & 35.8\,$\pm$\,1.9 & 35.9\,$\pm$\,1.7 \\
\Kmid{} classes & 31.8\,$\pm$\,0.8 & 32.4\,$\pm$\,1.7 & 34.6\,$\pm$\,1.8 & 34.4\,$\pm$\,1.7 \\
\Kcoarse{} classes & 31.0\,$\pm$\,2.8 & 34.5\,$\pm$\,1.1 & 35.3\,$\pm$\,1.0 & 35.5\,$\pm$\,1.0 \\
\Kres{} classes, mass regression & 34.7, 31.7 & 33.98, 33.85 & 34.55, 35.47 & 34.93, 35.60 \\
\Kcoarse{} classes, mass regression & 33.7, 32.4 & 33.46, 33.94 & 34.5, 35.6 & 34.7, 36.3 \\
randomly initialized network & 25.2\,$\pm$\,1.0 & 30.8\,$\pm$\,0.4 & 31.9\,$\pm$\,1.2 & 31.6\,$\pm$\,0.8 \\
randomly initialized network, BatchNorm recomputed & 24.7\,$\pm$\,2.0 & 29.7\,$\pm$\,0.5 & 30.5\,$\pm$\,0.6 & 30.7\,$\pm$\,0.5 \\
\bottomrule
\end{tabular}
\end{adjustbox}
\end{table*}


\begin{figure}[tbp]
\paperfig{\linewidth}{A10}
\caption{Linear probes on frozen features against the number of training jets, for the \Kfull-,
\Kres-, \Kmid- and \Kcoarse-class models, each in its own color (solid, dashed, dash-dot and dotted lines with
circles, squares, up- and down-pointing triangles), and the randomly initialized network with BatchNorm
statistics recomputed (grey long dashes with crosses). Panels: (a) accuracy on class-balanced JetClass-II; (b) accuracy on JetClass;
$R_{30}$ on (c) top tagging and (d) quark/gluon discrimination, with the
missing inputs substituted as chosen on validation. Axes as in figure~\ref{fig:f7}. Lines join
the means; bands show the standard deviation over the \MlstNRunsWord{} pretraining runs of each label set. The
points of the configurations are offset slightly along the horizontal axis for legibility.}
\label{fig:a10}
\end{figure}
% CHECK: A10 panel (a) metric is accuracy, as F7 (paper_figures.py A10); randomly initialized network: init-s3 not pulled yet

\begin{landscape}
{\small
\setlength{\LTleft}{0pt}\setlength{\LTright}{0pt}\setlength{\LTcapwidth}{\linewidth}\makeatletter\def\LT@makecaption#1#2#3{\LT@mcol\LT@cols c{\hbox to\z@{\hss\parbox[t]\LTcapwidth{\normalsize #1{#2: }#3\par\vskip\belowcaptionskip}\hss}}}\makeatother
\begin{longtable}{@{\extracolsep{\fill}}lcccc@{}}
\caption{Fine-tuning of networks pretrained on JetClass-II with 188, 162, 43 and 17 classes and of the controls, and the network trained from scratch, on each task with the number of training jets in the column headers: on class-balanced JetClass-II the AUC averaged over the classes (all and QCD classes) and the accuracy, the accuracy on JetClass, and $R_{30}$ on top, quark/gluon and Herwig jets; (f.t.): fine-tuned. $\rho$: $1-\mathrm{AUC}$ at $10^{6}$ jets over that of the \Kfull-class run of the same seed. A dash: not run at that size. Mean $\pm$ SD over the pretraining runs of each configuration.}\label{tab:ft}\\
\toprule
\endfirsthead
\multicolumn{5}{@{}l}{\textit{Table~\ref{tab:ft} (continued)}} \\
\toprule
\endhead
\bottomrule
\endlastfoot
\multicolumn{5}{@{}l}{JetClass-II, AUC} \\*
training jets & 10{,}000 & 100{,}000 & 1{,}000{,}000 & $\rho$ \\*
\noalign{\vskip\aboverulesep\nobreak}\multispan{5}{\leaders\hrule height\lightrulewidth\hfill}\cr\noalign{\nobreak\vskip\belowrulesep\nobreak}
\mbox{\Kfull{} classes (f.t.)} & 0.9722\,$\pm$\,0.0002 & 0.9778\,$\pm$\,0.0002 & 0.9797\,$\pm$\,0.0001 & 1.000 \\*
\mbox{\Kres{} classes (f.t.)} & 0.9725\,$\pm$\,0.0001 & 0.9775\,$\pm$\,0.0001 & 0.9794\,$\pm$\,0.0002 & 1.014\,$\pm$\,0.011 \\*
\mbox{\Kmid{} classes (f.t.)} & 0.9669\,$\pm$\,0.0006 & 0.9763\,$\pm$\,0.0001 & 0.9788\,$\pm$\,0.0001 & 1.043\,$\pm$\,0.002 \\*
\mbox{\Kcoarse{} classes (f.t.)} & 0.9623\,$\pm$\,0.0003 & 0.9743\,$\pm$\,0.0002 & 0.9772\,$\pm$\,0.0001 & 1.123\,$\pm$\,0.003 \\*
\mbox{\Kres{} classes + mass (f.t.)} & 0.9723\,$\pm$\,0.0002 & 0.9770\,$\pm$\,0.0001 & 0.9791\,$\pm$\,0.0001 & 1.030\,$\pm$\,0.009 \\*
\mbox{\Kcoarse{} classes + mass (f.t.)} & 0.9626\,$\pm$\,0.0001 & 0.9737\,$\pm$\,$<$0.0001 & 0.9767\,$\pm$\,0.0001 & 1.146\,$\pm$\,0.009 \\*
\mbox{random partition 1 (f.t.),} \mbox{seed 1} & -- & -- & 0.9783 & 1.073 \\*
\mbox{random partition 2 (f.t.),} \mbox{seed 2} & -- & -- & 0.9783 & 1.066 \\*
\mbox{random partition 3 (f.t.),} \mbox{seed 1} & -- & -- & 0.9786 & 1.058 \\*
\mbox{random partition 4 (f.t.),} \mbox{seed 2} & -- & -- & 0.9786 & 1.052 \\*
\mbox{random partition 5 (f.t.),} \mbox{seed 2} & -- & -- & 0.9783 & 1.069 \\*
\mbox{no flavor boundary (f.t.)} & -- & -- & 0.9775\,$\pm$\,0.0001 & 1.106\,$\pm$\,0.002 \\*
\mbox{$b$-quark boundary (f.t.)} & -- & -- & 0.9780\,$\pm$\,0.0003 & 1.083\,$\pm$\,0.015 \\*
\mbox{random flavor boundary (f.t.)} & -- & -- & 0.9780\,$\pm$\,0.0001 & 1.083\,$\pm$\,0.009 \\*
\mbox{from scratch} & 0.9202\,$\pm$\,0.0013 & 0.9613\,$\pm$\,0.0013 & 0.9734\,$\pm$\,0.0003 & -- \\
\midrule[\heavyrulewidth]
\multicolumn{5}{@{}l}{JetClass-II, accuracy} \\*
training jets & 10{,}000 & 100{,}000 & 1{,}000{,}000 \\*
\noalign{\vskip\aboverulesep\nobreak}\multispan{5}{\leaders\hrule height\lightrulewidth\hfill}\cr\noalign{\nobreak\vskip\belowrulesep\nobreak}
\mbox{\Kfull{} classes (f.t.)} & 0.3914\,$\pm$\,0.0011 & 0.4164\,$\pm$\,0.0015 & 0.4263\,$\pm$\,0.0005 \\*
\mbox{\Kres{} classes (f.t.)} & 0.3934\,$\pm$\,0.0018 & 0.4161\,$\pm$\,0.0015 & 0.4249\,$\pm$\,0.0021 \\*
\mbox{\Kmid{} classes (f.t.)} & 0.3563\,$\pm$\,0.0025 & 0.4007\,$\pm$\,0.0013 & 0.4166\,$\pm$\,0.0011 \\*
\mbox{\Kcoarse{} classes (f.t.)} & 0.3219\,$\pm$\,0.0011 & 0.3833\,$\pm$\,0.0020 & 0.4054\,$\pm$\,0.0020 \\*
\mbox{\Kres{} classes + mass (f.t.)} & 0.3868\,$\pm$\,0.0078 & 0.4104\,$\pm$\,0.0016 & 0.4206\,$\pm$\,0.0011 \\*
\mbox{\Kcoarse{} classes + mass (f.t.)} & 0.3171\,$\pm$\,0.0022 & 0.3770\,$\pm$\,0.0016 & 0.4006\,$\pm$\,0.0008 \\*
\mbox{random partition 1 (f.t.),} \mbox{seed 1} & -- & -- & 0.4151 \\*
\mbox{random partition 2 (f.t.),} \mbox{seed 2} & -- & -- & 0.4159 \\*
\mbox{random partition 3 (f.t.),} \mbox{seed 1} & -- & -- & 0.4160 \\*
\mbox{random partition 4 (f.t.),} \mbox{seed 2} & -- & -- & 0.4189 \\*
\mbox{random partition 5 (f.t.),} \mbox{seed 2} & -- & -- & 0.4163 \\*
\mbox{no flavor boundary (f.t.)} & -- & -- & 0.4078\,$\pm$\,0.0005 \\*
\mbox{$b$-quark boundary (f.t.)} & -- & -- & 0.4101\,$\pm$\,0.0013 \\*
\mbox{random flavor boundary (f.t.)} & -- & -- & 0.4128\,$\pm$\,0.0021 \\*
\mbox{from scratch} & 0.1492\,$\pm$\,0.0011 & 0.2921\,$\pm$\,0.0040 & 0.3677\,$\pm$\,0.0012 \\
\midrule[\heavyrulewidth]
\multicolumn{5}{@{}l}{JetClass-II, QCD classes, AUC over the QCD classes} \\*
training jets & 10{,}000 & 100{,}000 & 1{,}000{,}000 \\*
\noalign{\vskip\aboverulesep\nobreak}\multispan{5}{\leaders\hrule height\lightrulewidth\hfill}\cr\noalign{\nobreak\vskip\belowrulesep\nobreak}
\mbox{\Kfull{} classes (f.t.)} & 0.9498\,$\pm$\,0.0003 & 0.9598\,$\pm$\,0.0002 & 0.9631\,$\pm$\,$<$0.0001 \\*
\mbox{\Kres{} classes (f.t.)} & 0.9479\,$\pm$\,0.0001 & 0.9574\,$\pm$\,0.0002 & 0.9618\,$\pm$\,0.0004 \\*
\mbox{\Kmid{} classes (f.t.)} & 0.9416\,$\pm$\,0.0011 & 0.9573\,$\pm$\,0.0001 & 0.9613\,$\pm$\,0.0002 \\*
\mbox{\Kcoarse{} classes (f.t.)} & 0.9368\,$\pm$\,0.0014 & 0.9543\,$\pm$\,0.0003 & 0.9586\,$\pm$\,0.0001 \\*
\mbox{\Kres{} classes + mass (f.t.)} & 0.9472\,$\pm$\,0.0012 & 0.9569\,$\pm$\,0.0004 & 0.9611\,$\pm$\,0.0002 \\*
\mbox{\Kcoarse{} classes + mass (f.t.)} & 0.9342\,$\pm$\,0.0010 & 0.9532\,$\pm$\,0.0002 & 0.9580\,$\pm$\,0.0001 \\*
\mbox{random partition 1 (f.t.),} \mbox{seed 1} & -- & -- & 0.9598 \\*
\mbox{random partition 2 (f.t.),} \mbox{seed 2} & -- & -- & 0.9598 \\*
\mbox{random partition 3 (f.t.),} \mbox{seed 1} & -- & -- & 0.9609 \\*
\mbox{random partition 4 (f.t.),} \mbox{seed 2} & -- & -- & 0.9608 \\*
\mbox{random partition 5 (f.t.),} \mbox{seed 2} & -- & -- & 0.9602 \\*
\mbox{no flavor boundary (f.t.)} & -- & -- & 0.9591\,$\pm$\,0.0002 \\*
\mbox{$b$-quark boundary (f.t.)} & -- & -- & 0.9600\,$\pm$\,0.0010 \\*
\mbox{random flavor boundary (f.t.)} & -- & -- & 0.9599\,$\pm$\,0.0006 \\*
\mbox{from scratch} & 0.8885\,$\pm$\,0.0008 & 0.9359\,$\pm$\,0.0024 & 0.9536\,$\pm$\,0.0008 \\
\midrule[\heavyrulewidth]
\multicolumn{5}{@{}l}{JetClass, accuracy} \\*
training jets & 10{,}000 & 100{,}000 & 1{,}000{,}000 \\*
\noalign{\vskip\aboverulesep\nobreak}\multispan{5}{\leaders\hrule height\lightrulewidth\hfill}\cr\noalign{\nobreak\vskip\belowrulesep\nobreak}
\mbox{\Kfull{} classes (f.t.)} & 0.8117\,$\pm$\,0.0031 & 0.8343\,$\pm$\,0.0035 & 0.8468\,$\pm$\,0.0006 \\*
\mbox{\Kres{} classes (f.t.)} & 0.8117\,$\pm$\,0.0013 & 0.8328\,$\pm$\,0.0005 & 0.8475\,$\pm$\,0.0009 \\*
\mbox{\Kmid{} classes (f.t.)} & 0.8086\,$\pm$\,0.0031 & 0.8324\,$\pm$\,0.0004 & 0.8459\,$\pm$\,0.0002 \\*
\mbox{\Kcoarse{} classes (f.t.)} & 0.7982\,$\pm$\,0.0023 & 0.8288\,$\pm$\,0.0005 & 0.8422\,$\pm$\,0.0005 \\*
\mbox{\Kres{} classes + mass (f.t.)} & 0.8103\,$\pm$\,0.0007 & 0.8337\,$\pm$\,0.0016 & 0.8465\,$\pm$\,0.0009 \\*
\mbox{\Kcoarse{} classes + mass (f.t.)} & 0.7967\,$\pm$\,0.0033 & 0.8257\,$\pm$\,0.0007 & 0.8399\,$\pm$\,0.0002 \\*
\mbox{random partition 1 (f.t.),} \mbox{seed 1} & 0.8015 & 0.8304 & 0.8435 \\*
\mbox{random partition 2 (f.t.),} \mbox{seed 2} & 0.8055 & 0.8311 & 0.8436 \\*
\mbox{random partition 3 (f.t.),} \mbox{seed 1} & 0.7991 & 0.8309 & 0.8448 \\*
\mbox{random partition 4 (f.t.),} \mbox{seed 2} & 0.8016 & 0.8291 & 0.8449 \\*
\mbox{random partition 5 (f.t.),} \mbox{seed 2} & 0.8004 & 0.8313 & 0.8448 \\*
\mbox{no flavor boundary (f.t.)} & 0.7997\,$\pm$\,0.0005 & 0.8289\,$\pm$\,0.0007 & 0.8433\,$\pm$\,0.0002 \\*
\mbox{$b$-quark boundary (f.t.)} & 0.8045\,$\pm$\,0.0018 & 0.8309\,$\pm$\,0.0011 & 0.8439\,$\pm$\,0.0007 \\*
\mbox{random flavor boundary (f.t.)} & 0.8018\,$\pm$\,0.0004 & 0.8310\,$\pm$\,0.0014 & 0.8434\,$\pm$\,0.0003 \\*
\mbox{from scratch} & 0.6773\,$\pm$\,0.0038 & 0.7773\,$\pm$\,0.0050 & 0.8267\,$\pm$\,0.0011 \\
\midrule[\heavyrulewidth]
\multicolumn{5}{@{}l}{top, $R_{30}$} \\*
training jets & 1{,}000 & 10{,}000 & 100{,}000 & 1{,}200{,}000 \\*
\noalign{\vskip\aboverulesep\nobreak}\multispan{5}{\leaders\hrule height\lightrulewidth\hfill}\cr\noalign{\nobreak\vskip\belowrulesep\nobreak}
\mbox{\Kfull{} classes (f.t.)} & 127\,$\pm$\,21 & 409\,$\pm$\,27 & 850\,$\pm$\,90 & 1{,}800\,$\pm$\,60 \\*
\mbox{\Kres{} classes (f.t.)} & 154\,$\pm$\,19 & 436\,$\pm$\,51 & 810\,$\pm$\,60 & 1{,}790\,$\pm$\,60 \\*
\mbox{\Kmid{} classes (f.t.)} & 134\,$\pm$\,25 & 448\,$\pm$\,39 & 850\,$\pm$\,80 & 1{,}780\,$\pm$\,150 \\*
\mbox{\Kcoarse{} classes (f.t.)} & 137\,$\pm$\,15 & 370\,$\pm$\,24 & 830\,$\pm$\,60 & 1{,}740\,$\pm$\,140 \\*
\mbox{\Kres{} classes + mass (f.t.)} & 140\,$\pm$\,16 & 387\,$\pm$\,19 & 850 & 1{,}770\,$\pm$\,250 \\*
\mbox{\Kcoarse{} classes + mass (f.t.)} & 88\,$\pm$\,23 & 361\,$\pm$\,10 & 810\,$\pm$\,60 & 1{,}660\,$\pm$\,40 \\*
\mbox{random partition 1 (f.t.),} \mbox{seed 1} & 101 & 367 & 710 & 1{,}660 \\*
\mbox{random partition 2 (f.t.),} \mbox{seed 2} & 104 & 401 & 930 & 1{,}900 \\*
\mbox{random partition 3 (f.t.),} \mbox{seed 1} & 99 & 377 & 790 & 1{,}640 \\*
\mbox{random partition 4 (f.t.),} \mbox{seed 2} & 186 & 375 & 740 & 1{,}820 \\*
\mbox{random partition 5 (f.t.),} \mbox{seed 2} & 109 & 322 & 950 & 1{,}740 \\*
\mbox{no flavor boundary (f.t.)} & 84\,$\pm$\,20 & 367\,$\pm$\,12 & 810\,$\pm$\,50 & 1{,}790\,$\pm$\,90 \\*
\mbox{$b$-quark boundary (f.t.)} & 136\,$\pm$\,29 & 383\,$\pm$\,70 & 900\,$\pm$\,30 & 1{,}830\,$\pm$\,330 \\*
\mbox{random flavor boundary (f.t.)} & 104\,$\pm$\,31 & 355\,$\pm$\,18 & 770\,$\pm$\,20 & 1{,}610\,$\pm$\,60 \\*
\mbox{from scratch} & 49\,$\pm$\,12 & 339\,$\pm$\,49 & 560\,$\pm$\,90 & 1{,}410\,$\pm$\,20 \\*
\mbox{\Kfull{} classes,} \mbox{kinematic inputs only (f.t.)} & 94 & 366 & 1{,}070 & 1{,}800 \\
\midrule[\heavyrulewidth]
\multicolumn{5}{@{}l}{quark/gluon, $R_{30}$} \\*
training jets & 1{,}000 & 10{,}000 & 100{,}000 & 1{,}600{,}000 \\*
\noalign{\vskip\aboverulesep\nobreak}\multispan{5}{\leaders\hrule height\lightrulewidth\hfill}\cr\noalign{\nobreak\vskip\belowrulesep\nobreak}
\mbox{\Kfull{} classes (f.t.)} & 46\,$\pm$\,10 & 75\,$\pm$\,7 & 99.3\,$\pm$\,1.4 & 107.3\,$\pm$\,2.2 \\*
\mbox{\Kres{} classes (f.t.)} & 49\,$\pm$\,5 & 76\,$\pm$\,1 & 97.9\,$\pm$\,3.3 & 108.0\,$\pm$\,1.4 \\*
\mbox{\Kmid{} classes (f.t.)} & 43\,$\pm$\,12 & 74\,$\pm$\,5 & 96.1\,$\pm$\,1.8 & 108.3\,$\pm$\,1.0 \\*
\mbox{\Kcoarse{} classes (f.t.)} & 50\,$\pm$\,3 & 80\,$\pm$\,1 & 96.4\,$\pm$\,1.0 & 106.5\,$\pm$\,0.8 \\*
\mbox{\Kres{} classes + mass (f.t.)} & 51\,$\pm$\,14 & 76\,$\pm$\,5 & 96.6\,$\pm$\,1.4 & 106.8\,$\pm$\,0.2 \\*
\mbox{\Kcoarse{} classes + mass (f.t.)} & 57\,$\pm$\,1 & 79\,$\pm$\,1 & 97.9\,$\pm$\,0.1 & 105.3\,$\pm$\,0.5 \\*
\mbox{random partition 1 (f.t.),} \mbox{seed 1} & 46 & 80 & 98.9 & 107.3 \\*
\mbox{random partition 2 (f.t.),} \mbox{seed 2} & 47 & 81 & 99.6 & 108.6 \\*
\mbox{random partition 3 (f.t.),} \mbox{seed 1} & 34 & 77 & 96.5 & 106.5 \\*
\mbox{random partition 4 (f.t.),} \mbox{seed 2} & 54 & 79 & 99.1 & 106.6 \\*
\mbox{random partition 5 (f.t.),} \mbox{seed 2} & 66 & 59 & 95.2 & 106.0 \\*
\mbox{no flavor boundary (f.t.)} & 55\,$\pm$\,5 & 83\,$\pm$\,5 & 99.4\,$\pm$\,0.7 & 109.1\,$\pm$\,0.2 \\*
\mbox{$b$-quark boundary (f.t.)} & 42\,$\pm$\,6 & 83\,$\pm$\,2 & 98.4\,$\pm$\,2.3 & 106.7\,$\pm$\,0.2 \\*
\mbox{random flavor boundary (f.t.)} & 40\,$\pm$\,3 & 74\,$\pm$\,8 & 98.0\,$\pm$\,1.0 & 106.3\,$\pm$\,2.7 \\*
\mbox{from scratch} & 50\,$\pm$\,2 & 66\,$\pm$\,9 & 88.4\,$\pm$\,0.8 & 104.0\,$\pm$\,1.0 \\
\midrule[\heavyrulewidth]
\multicolumn{5}{@{}l}{quark/gluon, Herwig, $R_{30}$} \\*
training jets & 1{,}000 & 10{,}000 & 100{,}000 & 1{,}600{,}000 \\*
\noalign{\vskip\aboverulesep\nobreak}\multispan{5}{\leaders\hrule height\lightrulewidth\hfill}\cr\noalign{\nobreak\vskip\belowrulesep\nobreak}
\mbox{\Kfull{} classes (f.t.)} & 19.7\,$\pm$\,6.4 & 27.9\,$\pm$\,2.8 & 34.3\,$\pm$\,1.0 & 35.2\,$\pm$\,2.1 \\*
\mbox{\Kres{} classes (f.t.)} & 23.8\,$\pm$\,4.7 & 27.4\,$\pm$\,0.5 & 35.5\,$\pm$\,0.5 & 34.6\,$\pm$\,2.0 \\*
\mbox{\Kmid{} classes (f.t.)} & 19.3\,$\pm$\,8.1 & 28.4\,$\pm$\,2.1 & 35.5\,$\pm$\,2.1 & 36.1\,$\pm$\,2.7 \\*
\mbox{\Kcoarse{} classes (f.t.)} & 20.3\,$\pm$\,1.9 & 32.7\,$\pm$\,1.8 & 34.9\,$\pm$\,1.1 & 36.9\,$\pm$\,0.9 \\*
\mbox{\Kres{} classes + mass (f.t.)} & 24.0\,$\pm$\,7.9 & 31.4\,$\pm$\,4.1 & 35.2\,$\pm$\,1.8 & 36.1\,$\pm$\,2.4 \\*
\mbox{\Kcoarse{} classes + mass (f.t.)} & 26.0\,$\pm$\,1.3 & 32.8\,$\pm$\,2.3 & 37.5\,$\pm$\,2.2 & 35.3\,$\pm$\,3.4 \\*
\mbox{random partition 1 (f.t.),} \mbox{seed 1} & 24.1 & 30.5 & 34.2 & 33.4 \\*
\mbox{random partition 2 (f.t.),} \mbox{seed 2} & 20.7 & 34.1 & 35.2 & 35.2 \\*
\mbox{random partition 3 (f.t.),} \mbox{seed 1} & 12.5 & 30.2 & 34.0 & 34.8 \\*
\mbox{random partition 4 (f.t.),} \mbox{seed 2} & 25.5 & 34.4 & 33.9 & 37.5 \\*
\mbox{random partition 5 (f.t.),} \mbox{seed 2} & 35.0 & 23.2 & 36.3 & 35.3 \\*
\mbox{no flavor boundary (f.t.)} & 25.8\,$\pm$\,2.3 & 32.5\,$\pm$\,1.2 & 36.1\,$\pm$\,1.3 & 36.7\,$\pm$\,3.7 \\*
\mbox{$b$-quark boundary (f.t.)} & 17.5\,$\pm$\,3.4 & 34.4\,$\pm$\,1.2 & 36.0\,$\pm$\,1.5 & 36.4\,$\pm$\,2.4 \\*
\mbox{random flavor boundary (f.t.)} & 17.1\,$\pm$\,2.0 & 29.5\,$\pm$\,6.9 & 34.0\,$\pm$\,1.6 & 36.2\,$\pm$\,1.8 \\*
\mbox{from scratch} & 27.0\,$\pm$\,1.3 & 31.3\,$\pm$\,6.8 & 36.7\,$\pm$\,0.9 & 35.5\,$\pm$\,1.2 \\
\end{longtable}
}
\end{landscape}

\begin{landscape}
{\small
\setlength{\LTleft}{0pt}\setlength{\LTright}{0pt}\setlength{\LTcapwidth}{\linewidth}\makeatletter\def\LT@makecaption#1#2#3{\LT@mcol\LT@cols c{\hbox to\z@{\hss\parbox[t]\LTcapwidth{\normalsize #1{#2: }#3\par\vskip\belowcaptionskip}\hss}}}\makeatother
\begin{longtable}{@{\extracolsep{\fill}}lcccc@{}}
\caption{Fine-tuning of networks pretrained on JetClass-II with 188, 162, 43 and 17 classes and of the controls, and the network trained from scratch, on the full training sets of top and quark/gluon (also scored on Herwig jets): accuracy, AUC and the background rejections $R_{50}$ and $R_{30}$, $1/\epsB$ at 50\% and 30\% signal efficiency; (f.t.): fine-tuned. Mean $\pm$ SD over the \MlstNRunsWord{} pretraining runs of each level of the class hierarchy; configurations with \RandFlavRuns{} runs: over those. Every random partition has 17 classes; each numbered one (1--5) has a single run, numbered by the pretraining seed it shares with the \Kfull-class run of that seed. The random network without pretraining: three initializations. From scratch: three training seeds. Every rejection rests on at least 98 background jets passing the cut.}\label{tab:ftrej}\\
\toprule
& \multicolumn{4}{c}{top} \\
\cmidrule(lr){2-5}
 & accuracy & AUC & $R_{50}$ & $R_{30}$ \\
\midrule
\endfirsthead
\multicolumn{5}{@{}l}{\textit{Table~\ref{tab:ftrej} (continued)}} \\
\toprule
& \multicolumn{4}{c}{top} \\
\cmidrule(lr){2-5}
 & accuracy & AUC & $R_{50}$ & $R_{30}$ \\
\midrule
\endhead
\bottomrule
\endlastfoot
\Kfull{} classes (f.t.) & 0.9400\,$\pm$\,0.0002 & 0.9861\,$\pm$\,0.0001 & 461\,$\pm$\,4 & 1{,}800\,$\pm$\,60 \\
\Kres{} classes (f.t.) & 0.9401\,$\pm$\,0.0002 & 0.9861\,$\pm$\,0.0001 & 443\,$\pm$\,17 & 1{,}790\,$\pm$\,60 \\
\Kmid{} classes (f.t.) & 0.9402\,$\pm$\,0.0004 & 0.9861\,$\pm$\,0.0001 & 450\,$\pm$\,13 & 1{,}780\,$\pm$\,150 \\
\Kcoarse{} classes (f.t.) & 0.9398\,$\pm$\,0.0001 & 0.9859\,$\pm$\,0.0001 & 429\,$\pm$\,10 & 1{,}740\,$\pm$\,140 \\
\Kres{} classes + mass (f.t.) & 0.9404\,$\pm$\,0.0003 & 0.9860\,$\pm$\,$<$0.0001 & 428\,$\pm$\,12 & 1{,}770\,$\pm$\,250 \\
\Kcoarse{} classes + mass (f.t.) & 0.9398\,$\pm$\,0.0001 & 0.9858\,$\pm$\,0.0001 & 428\,$\pm$\,10 & 1{,}660\,$\pm$\,40 \\
random partition 1 (f.t.), seed 1 & 0.9402 & 0.9859 & 432 & 1{,}660 \\
random partition 2 (f.t.), seed 2 & 0.9404 & 0.9861 & 430 & 1{,}900 \\
random partition 3 (f.t.), seed 1 & 0.9402 & 0.9860 & 421 & 1{,}640 \\
random partition 4 (f.t.), seed 2 & 0.9397 & 0.9859 & 458 & 1{,}820 \\
random partition 5 (f.t.), seed 2 & 0.9402 & 0.9860 & 431 & 1{,}740 \\
no flavor boundary (f.t.) & 0.9401\,$\pm$\,0.0001 & 0.9860\,$\pm$\,0.0001 & 445\,$\pm$\,26 & 1{,}790\,$\pm$\,90 \\
$b$-quark boundary (f.t.) & 0.9403\,$\pm$\,0.0003 & 0.9860\,$\pm$\,0.0001 & 438\,$\pm$\,8 & 1{,}830\,$\pm$\,330 \\
random flavor boundary (f.t.) & 0.9402\,$\pm$\,0.0007 & 0.9860\,$\pm$\,0.0002 & 420\,$\pm$\,38 & 1{,}610\,$\pm$\,60 \\
from scratch & 0.9387\,$\pm$\,0.0002 & 0.9852\,$\pm$\,0.0001 & 358\,$\pm$\,11 & 1{,}410\,$\pm$\,20 \\
\Kfull{} classes, kinematic inputs only (f.t.) & 0.9407 & 0.9862 & 443 & 1{,}800 \\
\end{longtable}
\begin{center}
\begin{minipage}{\linewidth}\centering
\makebox[\linewidth][l]{\small\textit{Table~\ref{tab:ftrej} (continued)}}
\par\smallskip
\small
\begin{tabular*}{\linewidth}{@{\extracolsep{\fill}}lcccc@{}}
\toprule
 & \multicolumn{4}{c}{quark/gluon} \\
\cmidrule(lr){2-5}
 & accuracy & AUC & $R_{50}$ & $R_{30}$ \\
\midrule
\Kfull{} classes (f.t.) & 0.8433\,$\pm$\,0.0002 & 0.9148\,$\pm$\,0.0001 & 42.56\,$\pm$\,0.10 & 107.3\,$\pm$\,2.2 \\
\Kres{} classes (f.t.) & 0.8432\,$\pm$\,$<$0.0001 & 0.9148\,$\pm$\,0.0001 & 42.39\,$\pm$\,0.22 & 108.0\,$\pm$\,1.4 \\
\Kmid{} classes (f.t.) & 0.8430\,$\pm$\,0.0005 & 0.9146\,$\pm$\,0.0003 & 42.25\,$\pm$\,0.77 & 108.3\,$\pm$\,1.0 \\
\Kcoarse{} classes (f.t.) & 0.8427\,$\pm$\,0.0003 & 0.9144\,$\pm$\,$<$0.0001 & 41.97\,$\pm$\,0.25 & 106.5\,$\pm$\,0.8 \\
\Kres{} classes + mass (f.t.) & 0.8426\,$\pm$\,0.0001 & 0.9146\,$\pm$\,0.0001 & 42.34\,$\pm$\,0.18 & 106.8\,$\pm$\,0.2 \\
\Kcoarse{} classes + mass (f.t.) & 0.8425\,$\pm$\,$<$0.0001 & 0.9141\,$\pm$\,0.0001 & 41.58\,$\pm$\,0.27 & 105.3\,$\pm$\,0.5 \\
random partition 1 (f.t.), seed 1 & 0.8434 & 0.9146 & 42.14 & 107.3 \\
random partition 2 (f.t.), seed 2 & 0.8435 & 0.9147 & 42.46 & 108.6 \\
random partition 3 (f.t.), seed 1 & 0.8428 & 0.9145 & 42.14 & 106.5 \\
random partition 4 (f.t.), seed 2 & 0.8431 & 0.9146 & 41.84 & 106.6 \\
random partition 5 (f.t.), seed 2 & 0.8431 & 0.9146 & 42.14 & 106.0 \\
no flavor boundary (f.t.) & 0.8428\,$\pm$\,0.0008 & 0.9145\,$\pm$\,0.0001 & 41.81\,$\pm$\,0.20 & 109.1\,$\pm$\,0.2 \\
$b$-quark boundary (f.t.) & 0.8431\,$\pm$\,0.0001 & 0.9147\,$\pm$\,$<$0.0001 & 42.60\,$\pm$\,0.24 & 106.7\,$\pm$\,0.2 \\
random flavor boundary (f.t.) & 0.8427\,$\pm$\,0.0001 & 0.9146\,$\pm$\,0.0002 & 42.11\,$\pm$\,0.03 & 106.3\,$\pm$\,2.7 \\
from scratch & 0.8397\,$\pm$\,0.0007 & 0.9121\,$\pm$\,0.0002 & 41.47\,$\pm$\,0.29 & 104.0\,$\pm$\,1.0 \\
\bottomrule
\end{tabular*}
\end{minipage}
\end{center}
\begin{center}
\begin{minipage}{\linewidth}\centering
\makebox[\linewidth][l]{\small\textit{Table~\ref{tab:ftrej} (continued)}}
\par\smallskip
\small
\begin{tabular*}{\linewidth}{@{\extracolsep{\fill}}lcccc@{}}
\toprule
 & \multicolumn{4}{c}{quark/gluon, Herwig} \\
\cmidrule(lr){2-5}
 & accuracy & AUC & $R_{50}$ & $R_{30}$ \\
\midrule
\Kfull{} classes (f.t.) & 0.7123\,$\pm$\,0.0050 & 0.8078\,$\pm$\,0.0023 & 12.3\,$\pm$\,0.4 & 35.2\,$\pm$\,2.1 \\
\Kres{} classes (f.t.) & 0.7100\,$\pm$\,0.0017 & 0.8074\,$\pm$\,0.0010 & 12.3\,$\pm$\,0.3 & 34.6\,$\pm$\,2.0 \\
\Kmid{} classes (f.t.) & 0.7134\,$\pm$\,0.0058 & 0.8087\,$\pm$\,0.0046 & 12.5\,$\pm$\,0.6 & 36.1\,$\pm$\,2.7 \\
\Kcoarse{} classes (f.t.) & 0.7175\,$\pm$\,0.0033 & 0.8103\,$\pm$\,0.0018 & 12.8\,$\pm$\,0.2 & 36.9\,$\pm$\,0.9 \\
\Kres{} classes + mass (f.t.) & 0.7149\,$\pm$\,0.0040 & 0.8091\,$\pm$\,0.0025 & 12.6\,$\pm$\,0.6 & 36.1\,$\pm$\,2.4 \\
\Kcoarse{} classes + mass (f.t.) & 0.7133\,$\pm$\,0.0037 & 0.8076\,$\pm$\,0.0034 & 12.3\,$\pm$\,0.4 & 35.3\,$\pm$\,3.4 \\
random partition 1 (f.t.), seed 1 & 0.7120 & 0.8072 & 12.2 & 33.4 \\
random partition 2 (f.t.), seed 2 & 0.7121 & 0.8078 & 12.3 & 35.2 \\
random partition 3 (f.t.), seed 1 & 0.7106 & 0.8072 & 12.2 & 34.8 \\
random partition 4 (f.t.), seed 2 & 0.7141 & 0.8090 & 12.7 & 37.5 \\
random partition 5 (f.t.), seed 2 & 0.7132 & 0.8076 & 12.2 & 35.3 \\
no flavor boundary (f.t.) & 0.7106\,$\pm$\,0.0021 & 0.8070\,$\pm$\,0.0034 & 12.5\,$\pm$\,0.6 & 36.7\,$\pm$\,3.7 \\
$b$-quark boundary (f.t.) & 0.7132\,$\pm$\,0.0030 & 0.8087\,$\pm$\,0.0030 & 12.6\,$\pm$\,0.5 & 36.4\,$\pm$\,2.4 \\
random flavor boundary (f.t.) & 0.7143\,$\pm$\,0.0040 & 0.8077\,$\pm$\,0.0042 & 12.5\,$\pm$\,0.4 & 36.2\,$\pm$\,1.8 \\
from scratch & 0.7153\,$\pm$\,0.0024 & 0.8061\,$\pm$\,0.0007 & 12.2\,$\pm$\,0.1 & 35.5\,$\pm$\,1.2 \\
\bottomrule
\end{tabular*}
\end{minipage}
\end{center}
}
\end{landscape}

\begin{landscape}   % wider than the 15 cm text block
{\scriptsize
\setlength{\tabcolsep}{2pt}
\begin{longtable}{lccccccccc}
\caption{Every pretraining run's values at the largest training set: linear probe and fine-tuning. Linear probes on top and quark/gluon with the missing inputs substituted as chosen on validation (\LpSubstitutionTop{} for top, \LpSubstitutionQg{} for quark/gluon); fine-tuning with them set to zero.}\label{tab:perrun}\\
\toprule
run & \multicolumn{4}{c}{linear probe} & \multicolumn{5}{c}{fine-tuning} \\
 & JetClass-II AUC & JetClass AUC & top $R_{30}$ & q/g $R_{30}$ & JetClass-II AUC & JetClass accuracy & top $R_{30}$ & q/g $R_{30}$ & Herwig $R_{30}$ \\
\midrule
\endfirsthead
\toprule
run & \multicolumn{4}{c}{linear probe} & \multicolumn{5}{c}{fine-tuning} \\
 & JetClass-II AUC & JetClass AUC & top $R_{30}$ & q/g $R_{30}$ & JetClass-II AUC & JetClass accuracy & top $R_{30}$ & q/g $R_{30}$ & Herwig $R_{30}$ \\
\midrule
\endhead
\bottomrule
\endlastfoot
\Kfull{} classes, seed 1 & 0.9801 & 0.9683 & 47.3 & 57.2 & -- & 0.8489 & 1{,}740 & 108 & 33.7 \\
\Kfull{} classes, seed 2 & 0.9801 & 0.9687 & 61.1 & 70.5 & -- & 0.8497 & 1{,}850 & 109 & 37.5 \\
\Kfull{} classes, seed 3 & 0.9801 & 0.9675 & 70.3 & 69.1 & -- & 0.8487 & 1{,}800 & 105 & 34.2 \\
\Kres{} classes, seed 1 & 0.9795 & 0.9682 & 74.9 & 60.5 & -- & 0.8492 & 1{,}740 & 110 & 33.3 \\
\Kres{} classes, seed 2 & 0.9797 & 0.9683 & 59.0 & 47.6 & -- & 0.8496 & 1{,}790 & 107 & 36.9 \\
\Kres{} classes, seed 3 & 0.9796 & 0.9685 & 92.7 & 65.6 & -- & 0.8496 & 1{,}850 & 108 & 33.5 \\
\Kmid{} classes, seed 1 & 0.9758 & 0.9678 & 59.7 & 61.8 & -- & 0.8475 & 1{,}760 & 107 & 33.1 \\
\Kmid{} classes, seed 2 & 0.9755 & 0.9678 & 43.1 & 52.0 & -- & 0.8479 & 1{,}940 & 108 & 38.2 \\
\Kmid{} classes, seed 3 & 0.9755 & 0.9685 & 67.8 & 69.0 & -- & 0.8486 & 1{,}640 & 109 & 37.1 \\
\Kcoarse{} classes, seed 1 & 0.9695 & 0.9658 & 52.6 & 57.6 & -- & 0.8434 & 1{,}870 & 106 & 37.5 \\
\Kcoarse{} classes, seed 2 & 0.9689 & 0.9643 & 51.3 & 70.0 & -- & 0.8435 & 1{,}770 & 106 & 37.4 \\
\Kcoarse{} classes, seed 3 & 0.9688 & 0.9651 & 59.8 & 58.3 & -- & 0.8433 & 1{,}590 & 107 & 35.8 \\
\Kres{} classes, mass regression, seed 1 & 0.9791 & 0.9684 & 76.6 & 65.4 & -- & 0.8478 & 1{,}940 & 107 & 37.8 \\
\Kres{} classes, mass regression, seed 2 & 0.9794 & 0.9697 & 94.8 & 61.1 & -- & 0.8490 & 1{,}590 & 107 & 34.4 \\
\Kcoarse{} classes, mass regression, seed 1 & 0.9679 & 0.9655 & 103 & 66.0 & -- & 0.8420 & 1{,}680 & 105 & 37.7 \\
\Kcoarse{} classes, mass regression, seed 2 & 0.9696 & 0.9651 & 68.5 & 65.5 & -- & 0.8419 & 1{,}630 & 106 & 32.9 \\
random partition 1, seed 1 & 0.9776 & 0.9636 & 44.2 & 66.0 & -- & 0.8450 & 1{,}660 & 107 & 33.4 \\
random partition 2, seed 2 & 0.9776 & 0.9641 & 60.5 & 59.0 & -- & 0.8458 & 1{,}900 & 109 & 35.2 \\
random partition 3, seed 1 & 0.9779 & 0.9638 & 50.0 & 65.2 & -- & 0.8475 & 1{,}640 & 106 & 34.8 \\
random partition 4, seed 2 & 0.9778 & 0.9636 & 30.2 & 51.0 & -- & 0.8464 & 1{,}820 & 107 & 37.5 \\
random partition 5, seed 2 & 0.9775 & 0.9639 & 73.9 & 54.9 & -- & 0.8459 & 1{,}740 & 106 & 35.3 \\
random partition, no flavor boundary, seed 1 & 0.9717 & 0.9633 & 69.7 & 63.1 & -- & 0.8446 & 1{,}730 & 109 & 34.1 \\
random partition, no flavor boundary, seed 2 & 0.9710 & 0.9622 & 69.3 & 66.4 & -- & 0.8442 & 1{,}850 & 109 & 39.3 \\
random partition, $b$-quark boundary, seed 1 & 0.9758 & 0.9661 & 73.7 & 54.6 & -- & 0.8451 & 1{,}590 & 106 & 38.1 \\
random partition, $b$-quark boundary, seed 2 & 0.9759 & 0.9665 & 80.7 & 59.4 & -- & 0.8469 & 2{,}060 & 107 & 34.7 \\
random partition, random flavor boundary, seed 1 & 0.9760 & 0.9654 & 64.5 & 65.6 & -- & 0.8456 & 1{,}570 & 108 & 34.9 \\
random partition, random flavor boundary, seed 2 & 0.9760 & 0.9645 & 40.9 & 54.3 & -- & 0.8454 & 1{,}660 & 104 & 37.4 \\
\Kfull{} classes, kinematic inputs only & -- & -- & -- & -- & -- & -- & -- & -- & -- \\
from scratch, seed 1 & -- & -- & -- & -- & -- & 0.8273 & 1{,}410 & 103 & 34.9 \\
from scratch, seed 2 & -- & -- & -- & -- & -- & 0.8276 & 1{,}430 & 105 & 36.9 \\
from scratch, seed 3 & -- & -- & -- & -- & -- & 0.8258 & 1{,}390 & 104 & 34.7 \\
\end{longtable}
}

\end{landscape}

\subsection{JetClass}\label{sec:app-jetclass}

Table~\ref{tab:jetclassall} extends table~\ref{tab:jetclass} to every configuration, and
table~\ref{tab:jetclassprobe} gives the linear probes in the same columns; figure~\ref{fig:a6j} draws the
ROC curves of these linear probes in the layout of figure~\ref{fig:f6j}. The rejection of every
signal class rests on the QCD jets that pass its working point, and table~\ref{tab:jetclassnbkg} gives
these counts, from which the daggers of both tables follow.

To show which boundaries the random partitions lose on JetClass, table~\ref{tab:randprong} lists the
prong-count boundaries each partition keeps beside its fine-tuned rejection of each signal class,
relative to the \Kfull-class model of its seed.

\begin{landscape}
{\small
\setlength{\LTleft}{0pt}\setlength{\LTright}{0pt}\setlength{\LTcapwidth}{\linewidth}\makeatletter\def\LT@makecaption#1#2#3{\LT@mcol\LT@cols c{\hbox to\z@{\hss\parbox[t]\LTcapwidth{\normalsize #1{#2: }#3\par\vskip\belowcaptionskip}\hss}}}\makeatother
\begin{center}
\begin{minipage}{\linewidth}\centering
\normalsize
\captionof{table}{JetClass fine-tuning as table~\ref{tab:jetclass} of networks pretrained on JetClass-II with 188, 162, 43 and 17 classes and of the controls; the row of all five random partitions gives the mean $\pm$ SD over them; (f.t.): fine-tuned. Mean $\pm$ SD over the pretraining runs of each configuration. $^\dagger$ and $^{\ddagger}$ as in table~\ref{tab:jetclassrej}.}
\label{tab:jetclassall}
\par\medskip
\small
\begin{tabular*}{\linewidth}{@{\extracolsep{\fill}}llcccccc@{}}
\toprule
 &  &  &  &  & \multicolumn{3}{c}{rejection of QCD} \\
\cmidrule(lr){6-8}
 & jets & \# params & accuracy & AUC & $H\to b\bar b$ & $H\to c\bar c$ & $H\to gg$ \\
 &  &  &  &  & $R_{50}$ & $R_{50}$ & $R_{50}$ \\
\midrule
from scratch & 1 M & 2.21\,M & 0.8267\,$\pm$\,0.0011 & 0.9818\,$\pm$\,0.0003 & 4{,}800\,$\pm$\,100 & 1{,}570\,$\pm$\,50 & 86.4\,$\pm$\,3.2 \\
\Kfull{} classes (f.t.) & 1 M & 2.21\,M & 0.8468\,$\pm$\,0.0006 & 0.9854\,$\pm$\,0.0001 & 6{,}600\,$\pm$\,300 & 2{,}650\,$\pm$\,260 & 103.6\,$\pm$\,3.8 \\
\Kres{} classes (f.t.) & 1 M & 2.21\,M & 0.8475\,$\pm$\,0.0009 & 0.9856\,$\pm$\,$<$0.0001 & 6{,}100$^{\ddagger}$ & 2{,}820\,$\pm$\,40 & 101.8\,$\pm$\,1.5 \\
\Kmid{} classes (f.t.) & 1 M & 2.21\,M & 0.8459\,$\pm$\,0.0002 & 0.9853\,$\pm$\,0.0001 & 6{,}500\,$\pm$\,500 & 2{,}610\,$\pm$\,160 & 103.5\,$\pm$\,1.3 \\
\Kcoarse{} classes (f.t.) & 1 M & 2.21\,M & 0.8422\,$\pm$\,0.0005 & 0.9847\,$\pm$\,0.0001 & 6{,}700\,$\pm$\,200 & 2{,}390\,$\pm$\,40 & 100.8\,$\pm$\,6.2 \\
\Kres{} classes + mass (f.t.) & 1 M & 2.21\,M & 0.8465\,$\pm$\,0.0009 & 0.9853\,$\pm$\,0.0001 & 6{,}600\,$\pm$\,200 & 2{,}420\,$\pm$\,160 & 103.6\,$\pm$\,2.4 \\
\Kcoarse{} classes + mass (f.t.) & 1 M & 2.21\,M & 0.8399\,$\pm$\,0.0002 & 0.9843\,$\pm$\,0.0001 & 6{,}000\,$\pm$\,600 & 2{,}160\,$\pm$\,200 & 104.1\,$\pm$\,2.0 \\
random partitions (f.t.) & 1 M & 2.21\,M & 0.8443\,$\pm$\,0.0007 & 0.9851\,$\pm$\,0.0001 & 6{,}100\,$\pm$\,900 & 2{,}580\,$\pm$\,160 & 102.6\,$\pm$\,1.6 \\
no flavor boundary (f.t.) & 1 M & 2.21\,M & 0.8433\,$\pm$\,0.0002 & 0.9848\,$\pm$\,$<$0.0001 & 6{,}100$^{\ddagger}$ & 2{,}160\,$\pm$\,20 & 105.5\,$\pm$\,1.1 \\
$b$-quark boundary (f.t.) & 1 M & 2.21\,M & 0.8439\,$\pm$\,0.0007 & 0.9849\,$\pm$\,0.0001 & 6{,}600\,$\pm$\,1{,}200 & 2{,}380\,$\pm$\,40 & 106.5\,$\pm$\,2.2 \\
random flavor boundary (f.t.) & 1 M & 2.21\,M & 0.8434\,$\pm$\,0.0003 & 0.9849\,$\pm$\,$<$0.0001 & 6{,}200\,$\pm$\,400 & 2{,}740\,$\pm$\,160 & 101.4\,$\pm$\,6.5 \\
\bottomrule
\end{tabular*}
\end{minipage}
\end{center}
\begin{center}
\begin{minipage}{\linewidth}\centering
\makebox[\linewidth][l]{\small\textit{Table~\ref{tab:jetclassall} (continued)}}
\par\smallskip
\small
\begin{tabular*}{\linewidth}{@{\extracolsep{\fill}}llcccccc@{}}
\toprule
 &  & \multicolumn{6}{c}{rejection of QCD} \\
\cmidrule(lr){3-8}
 & jets & $H\to 4q$ & $H\to \ell\nu qq'$ & $t\to bqq'$ & $t\to b\ell\nu$ & $W\to qq'$ & $Z\to q\bar q$ \\
 &  & $R_{50}$ & $R_{99}$ & $R_{50}$ & $R_{99.5}$ & $R_{50}$ & $R_{50}$ \\
\midrule
from scratch & 1 M & 636\,$\pm$\,45 & 1{,}290\,$\pm$\,100 & 6{,}400\,$\pm$\,200 & 3{,}600\,$\pm$\,1{,}000 & 262\,$\pm$\,15 & 210\,$\pm$\,12 \\
\Kfull{} classes (f.t.) & 1 M & 1{,}091\,$\pm$\,15 & 3{,}060\,$\pm$\,390 & 12{,}900\,$\pm$\,1{,}600$^\dagger$ & 9{,}400\,$\pm$\,1{,}500$^\dagger$ & 374\,$\pm$\,8 & 312\,$\pm$\,6 \\
\Kres{} classes (f.t.) & 1 M & 1{,}150\,$\pm$\,67 & 3{,}070\,$\pm$\,330 & 14{,}100\,$\pm$\,3{,}500$^\dagger$ & 7{,}800\,$\pm$\,300 & 389\,$\pm$\,9 & 308\,$\pm$\,1 \\
\Kmid{} classes (f.t.) & 1 M & 1{,}034\,$\pm$\,40 & 3{,}160\,$\pm$\,280 & 14{,}200\,$\pm$\,2{,}100$^\dagger$ & 8{,}500\,$\pm$\,200 & 376\,$\pm$\,6 & 307\,$\pm$\,7 \\
\Kcoarse{} classes (f.t.) & 1 M & 1{,}049\,$\pm$\,17 & 2{,}800\,$\pm$\,200 & 10{,}600\,$\pm$\,2{,}700$^\dagger$ & 8{,}600\,$\pm$\,1{,}300 & 378\,$\pm$\,22 & 285\,$\pm$\,11 \\
\Kres{} classes + mass (f.t.) & 1 M & 1{,}114\,$\pm$\,22 & 2{,}860\,$\pm$\,120 & 12{,}700\,$\pm$\,2{,}200$^\dagger$ & 7{,}600\,$\pm$\,1{,}000 & 386\,$\pm$\,9 & 302\,$\pm$\,1 \\
\Kcoarse{} classes + mass (f.t.) & 1 M & 950\,$\pm$\,22 & 2{,}920\,$\pm$\,30 & 12{,}400\,$\pm$\,2{,}700$^\dagger$ & 9{,}600\,$\pm$\,2{,}200$^\dagger$ & 358\,$\pm$\,7 & 283\,$\pm$\,10 \\
random partitions (f.t.) & 1 M & 1{,}006\,$\pm$\,31 & 2{,}690\,$\pm$\,300 & 10{,}900\,$\pm$\,2{,}000$^\dagger$ & 7{,}700\,$\pm$\,1{,}300 & 379\,$\pm$\,18 & 299\,$\pm$\,7 \\
no flavor boundary (f.t.) & 1 M & 1{,}106\,$\pm$\,99 & 2{,}820$^{\ddagger}$ & 15{,}000\,$\pm$\,2{,}400$^\dagger$ & 6{,}900\,$\pm$\,300 & 351$^{\ddagger}$ & 292\,$\pm$\,$<$1 \\
$b$-quark boundary (f.t.) & 1 M & 1{,}036\,$\pm$\,8 & 2{,}580\,$\pm$\,20 & 13{,}400\,$\pm$\,1{,}300$^\dagger$ & 7{,}100$^{\ddagger}$ & 356\,$\pm$\,14 & 295\,$\pm$\,4 \\
random flavor boundary (f.t.) & 1 M & 1{,}024\,$\pm$\,33 & 2{,}740\,$\pm$\,110 & 12{,}500\,$\pm$\,1{,}100$^\dagger$ & 7{,}300\,$\pm$\,200 & 371\,$\pm$\,$<$1 & 299\,$\pm$\,4 \\
\bottomrule
\end{tabular*}
\end{minipage}
\end{center}
}
\end{landscape}

\begin{landscape}
{
\setlength{\LTleft}{0pt}\setlength{\LTright}{0pt}\setlength{\LTcapwidth}{\linewidth}\makeatletter\def\LT@makecaption#1#2#3{\LT@mcol\LT@cols c{\hbox to\z@{\hss\parbox[t]\LTcapwidth{\normalsize #1{#2. }#3\par\vskip\belowcaptionskip}\hss}}}\makeatother
\addtolength{\belowcaptionskip}{\medskipamount}
\begin{longtable}{@{\extracolsep{\fill}}lccccc@{}}
\caption{Linear probes on the frozen features of the networks pretrained with 188, 162, 43 and 17 classes, with an added mass output and with every 17-class partition, of the random networks and of the released Sophon checkpoint, trained and scored on JetClass at the largest training set, in the columns of table~\ref{tab:jetclassall}. The released Sophon checkpoint, set apart by a rule, is outside the controlled comparison. AUC is the area under the receiver operating characteristic curve, averaged over the classes, each against the rest. With JetClass's equal class sizes this equals the average over class pairs that ParT reports. The random network with BatchNorm recomputed re-estimates its BatchNorm statistics on \DesignBnJets{} training jets. A flavor boundary keeps apart decays that differ only in quark flavor. Values are the mean $\pm$ standard deviation over \MlstNRunsWord{} runs, \MlstMassNRunsWord{} per mass network and flavor-boundary partition. $^{\S}$~marks a value from a single run. $^{\P}$~marks a value whose runs agree to the digits shown.}\label{tab:jetclassprobe}\\
\toprule
&  &  & \multicolumn{3}{c}{rejection of QCD} \\
\cmidrule(lr){4-6}
 & accuracy & AUC & $H\to b\bar b$ & $H\to c\bar c$ & $H\to gg$ \\
 &  &  & $R_{50}$ & $R_{50}$ & $R_{50}$ \\
\midrule
\endfirsthead
\multicolumn{6}{@{}l}{\textit{Table~\ref{tab:jetclassprobe} (continued)}} \\[\smallskipamount]
\toprule
&  &  & \multicolumn{3}{c}{rejection of QCD} \\
\cmidrule(lr){4-6}
 & accuracy & AUC & $H\to b\bar b$ & $H\to c\bar c$ & $H\to gg$ \\
 &  &  & $R_{50}$ & $R_{50}$ & $R_{50}$ \\
\midrule
\endhead
\bottomrule
\endlastfoot
\multicolumn{6}{@{}l}{\textit{linear probe on the frozen features}} \\
\Kfull{} classes & 0.771\,$\pm$\,0.002 & 0.9682\,$\pm$\,0.0006 & 836\,$\pm$\,93 & 382\,$\pm$\,34 & 19.98\,$\pm$\,0.97 \\
\Kres{} classes & 0.771\,$\pm$\,0.001 & 0.9683\,$\pm$\,0.0002 & 797\,$\pm$\,131 & 415\,$\pm$\,52 & 20.70\,$\pm$\,0.81 \\
\Kres{} classes + mass & 0.774\,$\pm$\,0.003 & 0.9690\,$\pm$\,0.0009 & 815\,$\pm$\,12 & 395\,$\pm$\,25 & 21.18\,$\pm$\,1.84 \\
\Kmid{} classes & 0.770\,$\pm$\,0.001 & 0.9680\,$\pm$\,0.0004 & 875\,$\pm$\,184 & 440\,$\pm$\,31 & 20.78\,$\pm$\,0.71 \\
\Kcoarse{} classes & 0.761\,$\pm$\,0.003 & 0.9650\,$\pm$\,0.0008 & 831\,$\pm$\,16 & 294\,$\pm$\,20 & 21.60\,$\pm$\,0.70 \\
\Kcoarse{} classes + mass & 0.762\,$\pm$\,0.001 & 0.9653\,$\pm$\,0.0003 & 781\,$\pm$\,204 & 274\,$\pm$\,19 & 22.24\,$\pm$\,0.48 \\
\multicolumn{6}{@{}l}{17-class partitions} \\
\quad no flavor boundary & 0.751\,$\pm$\,0.003 & 0.9628\,$\pm$\,0.0008 & 641\,$\pm$\,189 & 239\,$\pm$\,8 & 20.33\,$\pm$\,1.65 \\
\quad $b$-quark boundary & 0.763\,$\pm$\,0.002 & 0.9663\,$\pm$\,0.0003 & 1{,}045\,$\pm$\,77 & 349\,$\pm$\,68 & 19.28\,$\pm$\,1.87 \\
\quad random flavor boundary & 0.758\,$\pm$\,0.002 & 0.9649\,$\pm$\,0.0006 & 710\,$\pm$\,67 & 325\,$\pm$\,10 & 18.83\,$\pm$\,0.05 \\
\quad random partition 1, seed 1 & 0.753$^{\S}$ & 0.9636$^{\S}$ & 712$^{\S}$ & 323$^{\S}$ & 18.25$^{\S}$ \\
\quad random partition 2, seed 2 & 0.755$^{\S}$ & 0.9641$^{\S}$ & 658$^{\S}$ & 357$^{\S}$ & 18.70$^{\S}$ \\
\quad random partition 3, seed 1 & 0.755$^{\S}$ & 0.9638$^{\S}$ & 877$^{\S}$ & 382$^{\S}$ & 18.46$^{\S}$ \\
\quad random partition 4, seed 2 & 0.754$^{\S}$ & 0.9636$^{\S}$ & 797$^{\S}$ & 348$^{\S}$ & 17.71$^{\S}$ \\
\quad random partition 5, seed 2 & 0.754$^{\S}$ & 0.9639$^{\S}$ & 749$^{\S}$ & 353$^{\S}$ & 16.49$^{\S}$ \\
random network & 0.408\,$\pm$\,0.001 & 0.8280\,$\pm$\,0.0011 & 25\,$\pm$\,1 & 20\,$\pm$\,2 & 17.91\,$\pm$\,0.83 \\
random network, BatchNorm recomputed & 0.475\,$\pm$\,0.005 & 0.8655\,$\pm$\,0.0040 & 46\,$\pm$\,2 & 14\,$\pm$\,1 & 15.37\,$\pm$\,0.53 \\
\midrule
released Sophon checkpoint & 0.768$^{\S}$ & 0.9675$^{\S}$ & 752$^{\S}$ & 374$^{\S}$ & 19.70$^{\S}$ \\
\end{longtable}
\newpage
\addtocounter{table}{-1}
\begin{longtable}{@{\extracolsep{\fill}}lcccccc@{}}
\multicolumn{7}{@{}l}{\textit{Table~\ref{tab:jetclassprobe} (continued)}} \\[\smallskipamount]
\toprule
& \multicolumn{6}{c}{rejection of QCD} \\
\cmidrule(lr){2-7}
 & $H\to 4q$ & $H\to \ell\nu qq'$ & $t\to bqq'$ & $t\to b\ell\nu$ & $W\to qq'$ & $Z\to q\bar q$ \\
 & $R_{50}$ & $R_{99}$ & $R_{50}$ & $R_{99.5}$ & $R_{50}$ & $R_{50}$ \\
\midrule
\endfirsthead
\multicolumn{7}{@{}l}{\textit{Table~\ref{tab:jetclassprobe} (continued)}} \\[\smallskipamount]
\toprule
& \multicolumn{6}{c}{rejection of QCD} \\
\cmidrule(lr){2-7}
 & $H\to 4q$ & $H\to \ell\nu qq'$ & $t\to bqq'$ & $t\to b\ell\nu$ & $W\to qq'$ & $Z\to q\bar q$ \\
 & $R_{50}$ & $R_{99}$ & $R_{50}$ & $R_{99.5}$ & $R_{50}$ & $R_{50}$ \\
\midrule
\endhead
\bottomrule
\endlastfoot
\multicolumn{7}{@{}l}{\textit{linear probe on the frozen features}} \\
\Kfull{} classes & 200\,$\pm$\,15 & 918\,$\pm$\,36 & 1{,}097\,$\pm$\,82 & 2{,}689\,$\pm$\,306 & 76.1\,$\pm$\,1.7 & 63.1\,$\pm$\,2.5 \\
\Kres{} classes & 207\,$\pm$\,15 & 933\,$\pm$\,46 & 1{,}221\,$\pm$\,150 & 3{,}361\,$\pm$\,583 & 81.6\,$\pm$\,2.1 & 65.8\,$\pm$\,3.2 \\
\Kres{} classes + mass & 227\,$\pm$\,23 & 929\,$\pm$\,40 & 1{,}082\,$\pm$\,144 & 4{,}939\,$\pm$\,86 & 75.4\,$\pm$\,3.5 & 68.3\,$\pm$\,1.1 \\
\Kmid{} classes & 207\,$\pm$\,13 & 919\,$\pm$\,75 & 1{,}298\,$\pm$\,203 & 3{,}682\,$\pm$\,484 & 79.1\,$\pm$\,1.6 & 69.5\,$\pm$\,0.9 \\
\Kcoarse{} classes & 183\,$\pm$\,12 & 814\,$\pm$\,113 & 1{,}184\,$\pm$\,308 & 3{,}094\,$\pm$\,275 & 70.8\,$\pm$\,5.0 & 68.9\,$\pm$\,2.8 \\
\Kcoarse{} classes + mass & 170\,$\pm$\,1 & 791\,$\pm$\,13 & 843\,$\pm$\,43 & 2{,}152\,$\pm$\,324 & 67.8\,$\pm$\,2.8 & 69.5\,$\pm$\,1.9 \\
\multicolumn{7}{@{}l}{17-class partitions} \\
\quad no flavor boundary & 156\,$\pm$\,2 & 697\,$\pm$\,21 & 899\,$\pm$\,14 & 2{,}058\,$\pm$\,861 & 67.3\,$\pm$\,2.6 & 68.7\,$\pm$\,0.7 \\
\quad $b$-quark boundary & 200\,$\pm$\,7 & 706\,$\pm$\,23 & 948\,$\pm$\,157 & 1{,}973\,$\pm$\,352 & 72.8\,$\pm$\,3.1 & 67.0\,$\pm$\,1.2 \\
\quad random flavor boundary & 180\,$\pm$\,7 & 672\,$\pm$\,18 & 937\,$\pm$\,254 & 1{,}951\,$\pm$\,187 & 73.7\,$\pm$\,1.9 & 66.0\,$\pm$\,0.5 \\
\quad random partition 1, seed 1 & 161$^{\S}$ & 557$^{\S}$ & 847$^{\S}$ & 2{,}247$^{\S}$ & 67.2$^{\S}$ & 56.1$^{\S}$ \\
\quad random partition 2, seed 2 & 186$^{\S}$ & 667$^{\S}$ & 664$^{\S}$ & 1{,}550$^{\S}$ & 71.3$^{\S}$ & 57.7$^{\S}$ \\
\quad random partition 3, seed 1 & 180$^{\S}$ & 617$^{\S}$ & 1{,}117$^{\S}$ & 2{,}020$^{\S}$ & 72.4$^{\S}$ & 58.7$^{\S}$ \\
\quad random partition 4, seed 2 & 154$^{\S}$ & 662$^{\S}$ & 858$^{\S}$ & 2{,}500$^{\S}$ & 70.1$^{\S}$ & 56.5$^{\S}$ \\
\quad random partition 5, seed 2 & 164$^{\S}$ & 610$^{\S}$ & 840$^{\S}$ & 2{,}326$^{\S}$ & 69.0$^{\S}$ & 56.0$^{\S}$ \\
random network & 16\,$\pm$\,1 & 1$^{\P}$ & 20\,$\pm$\,1 & 1$^{\P}$ & 17.3\,$\pm$\,1.3 & 14.3\,$\pm$\,1.0 \\
random network, BatchNorm recomputed & 14\,$\pm$\,1 & 2\,$\pm$\,1 & 23\,$\pm$\,2 & 2$^{\P}$ & 13.5\,$\pm$\,0.3 & 11.0\,$\pm$\,0.6 \\
\midrule
released Sophon checkpoint & 199$^{\S}$ & 791$^{\S}$ & 1{,}005$^{\S}$ & 2{,}857$^{\S}$ & 78.1$^{\S}$ & 62.3$^{\S}$ \\
\end{longtable}
}
\end{landscape}


\begin{figure}[tbp]
\paperfig{\linewidth}{A6J}
\caption{Linear probes on frozen features on JetClass, trained on \FtSizeESix{} jets: background rejection
$1/\epsB$ against signal efficiency $\epsS$ for each signal class against the QCD background, with the
two-class discriminant of equation~(\ref{eq:discr}). Colored lines (solid, dashed, dash-dot and dotted): the \Kfull-, \Kres-, \Kmid- and
\Kcoarse-class models; grey long dashes: the randomly initialized network with its BatchNorm statistics recomputed,
the reference for what the architecture alone gives. Lines are the mean of the runs at fixed $\epsS$, with bands
of one standard deviation over the \MlstNRunsWord{} pretraining runs and over the \MlstNRunsWord{} initial
weights of the randomly initialized network. Lower panels: ratio to the \Kfull-class model, run by run against
the \Kfull-class run of the same seed. Each curve stops where fewer than $\nmin$ QCD jets pass, equation~(\ref{eq:cut}), marked by
the dashed black line at $N_B/\nmin$. The values at ParT's working points are in table~\ref{tab:jetclassprobe}.}
\label{fig:a6j}
\end{figure}
% CHECK: pending init-s3 (randomly initialized network: two of three initial weights pulled today).

\begin{table*}[t]
\centering
\small
\caption{QCD jets passing each signal class's working point after fine-tuning, the fewest over the runs of each configuration: the counts behind the rejections of table~\ref{tab:jetclass}.}
\label{tab:jetclassnbkg}
\begin{adjustbox}{max width=\linewidth}
\begin{tabular}{lccccccccc}
\toprule
fine-tuned from & $H\to b\bar b$ & $H\to c\bar c$ & $H\to gg$ & $H\to 4q$ & $H\to \ell\nu qq'$ & $t\to bqq'$ & $t\to b\ell\nu$ & $W\to qq'$ & $Z\to q\bar q$ \\
\midrule
from scratch & 39 & 123 & 2216 & 286 & 145 & 30 & 48 & 724 & 897 \\
\Kfull{} classes & 29 & 68 & 1884 & 181 & 57 & 14 & 19 & 524 & 628 \\
\Kres{} classes & 33 & 70 & 1933 & 166 & 58 & 11 & 25 & 508 & 648 \\
\Kmid{} classes & 29 & 72 & 1914 & 187 & 58 & 13 & 23 & 525 & 636 \\
\Kcoarse{} classes & 29 & 82 & 1873 & 188 & 66 & 15 & 21 & 503 & 679 \\
\Kres{} classes, mass regression & 30 & 79 & 1899 & 177 & 68 & 14 & 24 & 510 & 661 \\
\Kcoarse{} classes, mass regression & 31 & 87 & 1896 & 207 & 68 & 14 & 18 & 551 & 691 \\
random partitions (all five) & 26 & 71 & 1921 & 194 & 64 & 15 & 22 & 492 & 649 \\
random partition, no flavor boundary & 33 & 92 & 1882 & 170 & 71 & 12 & 28 & 569 & 685 \\
random partition, $b$-quark boundary & 27 & 83 & 1851 & 192 & 77 & 14 & 28 & 546 & 672 \\
random partition, random flavor boundary & 31 & 70 & 1886 & 191 & 71 & 15 & 27 & 538 & 664 \\
\bottomrule
\end{tabular}
\end{adjustbox}
\end{table*}

\begin{landscape}
\begin{table*}[t]
\centering
\small
\caption{Networks pretrained on JetClass-II with random partitions into 17 classes, fine-tuned on JetClass: none keeps a boundary between two, three and four prongs (for each pair of these prong counts, every partition has a group holding classes of both); the groups that mix prong counts, and each partition's fine-tuned rejection of QCD per signal class over that of the \Kfull-class run of the same seed; (f.t.): fine-tuned. $^\dagger$: fewer than $\nmin$ QCD jets pass in the partition's run or in that \Kfull-class run (equation~(\ref{eq:cut})).}
\label{tab:randprong}
\begin{tabular*}{\linewidth}{@{\extracolsep{\fill}}lcccccccccc@{}}
\toprule
 & & \multicolumn{9}{c}{rejection over \Kfull{} classes} \\
\cmidrule(lr){3-11}
partition & mixed groups & $H\to b\bar b$ & $H\to c\bar c$ & $H\to gg$ & $H\to 4q$ & $H\to \ell\nu qq'$ & $t\to bqq'$ & $t\to b\ell\nu$ & $W\to qq'$ & $Z\to q\bar q$ \\
\midrule
random partition 1 (f.t.), seed 1 & 12 & 0.842 & 1.04 & 0.984 & 0.862 & 1.12 & 0.652$^\dagger$ & 1.13 & 0.991 & 0.929 \\
random partition 2 (f.t.), seed 2 & 12 & 0.853 & 1.10 & 0.944 & 0.935 & 0.958 & 0.875$^\dagger$ & 0.606 & 1.03 & 0.929 \\
random partition 3 (f.t.), seed 1 & 13 & 0.889 & 1.09 & 0.966 & 0.928 & 0.935 & 1.00$^\dagger$ & 1.18 & 0.934 & 0.939 \\
random partition 4 (f.t.), seed 2 & 12 & 0.906 & 0.940 & 0.966 & 0.959 & 0.908 & 0.636$^\dagger$ & 0.667 & 1.11 & 0.971 \\
random partition 5 (f.t.), seed 2 & 11 & 1.12 & 0.963 & 0.981 & 0.944 & 0.793 & 0.778$^\dagger$ & 0.800 & 1.03 & 0.991 \\
\bottomrule
\end{tabular*}
\end{table*}
\end{landscape}


\subsection{JetClass-II class pairs and classes}\label{sec:app-pairs}

Table~\ref{tab:pairs} completes figure~\ref{fig:f4} with every configuration and seven class pairs. Beside
the four pairs of the main text it gives QCD $b\bar b$ against QCD $c\bar c$, QCD jets with a $b$ quark against
QCD jets without $b$, $c$ or $s$ quarks, and QCD jets with a $c$ but no $b$ quark against the same light QCD
jets. Each pair has its AUC and, in a separate column, the ratio $\rho$ of equation~(\ref{eq:ratio}) to the
\Kfull-class model. Table~\ref{tab:perclass} gives the numbers behind figure~\ref{fig:f3}, one row per
JetClass-II class.

Figure~\ref{fig:a3r} draws the ROC curves of selected JetClass-II classes, each against every other class, to
show at which signal efficiencies the per-class losses of figure~\ref{fig:f3} lie. The classes are chosen from
the per-class AUCs by a fixed rule: the eight with the largest loss of AUC at \Kcoarse{} classes, and the four
whose $1-\AUC$ stays closest to that of the \Kfull-class model at every level of the hierarchy, among the
classes not hatched in figure~\ref{fig:f3}. The ratio decides the second set because a class whose AUC is
already near one can lose much of its rejection with almost no change in AUC, as $X\to ee$ and $X\to\mu\mu$ do.

\providecommand{\TableContinued}[2]{\label{#2}\ifnum\getpagerefnumber{#2}>\getpagerefnumber{#1}\relax{\footnotesize\textit{Table~\ref{#1} (continued)}}\par\smallskip\fi}
\begin{center}
\small
\begin{minipage}{\linewidth}\centering
\captionof{table}{Pairs of JetClass-II classes separated by networks pretrained on JetClass-II with 188, 162, 43 and 17 classes and by the controls: (LP), a binary linear probe on the frozen features, trained on up to \JcBalPairTrainPerSide{} jets of each side from held-out JetClass-II files and scored on up to \JcBalPairTestPerSide{} of each side from the JetClass-II test files, each side in its natural class mix; (f.t.), the network fine-tuned on $10^{6}$ class-balanced JetClass-II jets with the \Kfull{} classes, scored on the pair's jets among \LpNTestJcii{} JetClass-II test jets at their natural composition by the summed output scores of the first side over those of both sides. The two modes are scored on different test jets, so they are compared within a mode only. The AUC and the ratio $\rho$ of $1-\mathrm{AUC}$, equation~(\ref{eq:ratio}), to the \Kfull-class run of the same seed in the same mode. No pretraining: fine-tuned from scratch, three seeds, no probe. $^{\dagger}$: the label set merges the two sides into one class; $^{\ddagger}$: it merges the first side with part of the second. Mean $\pm$ standard deviation over the \MlstNRunsWord{} pretraining runs of each level of the class hierarchy; configurations with \RandFlavRuns{} runs give both values. Every random partition has 17 classes; each numbered one (1--5) has a single run, numbered by the pretraining seed it shares with the \Kfull-class run of that seed. The random network without pretraining: three initializations. The $b$-quark boundary and the random flavor boundary lie only in the four-quark decays; every label set but the \Kfull-class one merges all QCD classes. The random network without pretraining, as initialized and with its BatchNorm statistics recomputed, has no label set, no $\rho$ and no fine-tuning. The $X\to ee$ vs $\mu\mu$ probe is trained on the jets of the two classes in the class-balanced training set (\MlstPairEemmNTrain{} jets) and scored on every JetClass-II test jet of the two classes (\MlstPairEemmNTestPerSide{} per side). A mark after two values applies to the configuration. A dash in a ratio column: the \Kfull-class run of that seed separates the test jets almost completely ($1-\mathrm{AUC} < 10^{-4}$), so a ratio to it would mean nothing.}
\label{tab:pairs}
\par\medskip
\begin{adjustbox}{max width=\linewidth, max totalheight=0.55\textheight}
\begin{tabular}{>{\raggedright\arraybackslash}p{4.2cm}cc}
\toprule
 & \multicolumn{2}{c}{$X\to bc$ vs $bq$, $cs$} \\
\cmidrule(lr){2-3}
 & AUC & $\rho$ \\
\midrule
\multicolumn{3}{l}{\textit{(LP): linear probe on the frozen features}} \\
\mbox{\Kfull{} classes} & 0.9409\,$\pm$\,0.0003 & 1.000 \\
\mbox{\Kres{} classes} & 0.9410\,$\pm$\,0.0002 & 0.998\,$\pm$\,0.007 \\
\mbox{\Kmid{} classes} & 0.9319\,$\pm$\,0.0014$^{\ddagger}$ & 1.151\,$\pm$\,0.029 \\
\mbox{\Kcoarse{} classes} & 0.8845\,$\pm$\,0.0039$^{\dagger}$ & 1.953\,$\pm$\,0.063 \\
\mbox{\Kres{} classes + mass} & 0.9375, 0.9400 & 1.057, 1.010 \\
\mbox{\Kcoarse{} classes + mass} & 0.8808, 0.8833$^{\dagger}$ & 2.018, 1.964 \\
\mbox{random partition 1,} \mbox{seed 1} & 0.9311$^{\dagger}$ & 1.166 \\
\mbox{random partition 2,} \mbox{seed 2} & 0.9298$^{\ddagger}$ & 1.181 \\
\mbox{random partition 3,} \mbox{seed 1} & 0.9369 & 1.069 \\
\mbox{random partition 4,} \mbox{seed 2} & 0.9361 & 1.075 \\
\mbox{random partition 5,} \mbox{seed 2} & 0.9266$^{\dagger}$ & 1.236 \\
\mbox{no flavor boundary} & 0.8819, 0.8697$^{\dagger}$ & 1.999, 2.193 \\
\mbox{$b$-quark boundary} & 0.9103, 0.9091$^{\dagger}$ & 1.518, 1.530 \\
\mbox{random flavor boundary} & 0.9096, 0.9100$^{\dagger}$ & 1.530, 1.515 \\
\mbox{random network} & 0.6011\,$\pm$\,0.0028 & -- \\
\mbox{random network,} \mbox{BatchNorm recomputed} & 0.6528\,$\pm$\,0.0090 & -- \\
\midrule
\multicolumn{3}{l}{\textit{(f.t.): fine-tuned}} \\
\mbox{\Kfull{} classes} & 0.9380\,$\pm$\,0.0011 & 1.000 \\
\mbox{\Kres{} classes} & 0.9377\,$\pm$\,0.0007 & 1.005\,$\pm$\,0.029 \\
\mbox{\Kmid{} classes} & 0.9339\,$\pm$\,0.0004 & 1.065\,$\pm$\,0.018 \\
\mbox{\Kcoarse{} classes} & 0.9238\,$\pm$\,0.0014 & 1.228\,$\pm$\,0.022 \\
\mbox{\Kres{} classes + mass} & 0.9363, 0.9376 & 1.044, 0.988 \\
\mbox{\Kcoarse{} classes + mass} & 0.9224, 0.9227 & 1.271, 1.222 \\
\mbox{random partition 1,} \mbox{seed 1} & 0.9293 & 1.159 \\
\mbox{random partition 2,} \mbox{seed 2} & 0.9314 & 1.085 \\
\mbox{random partition 3,} \mbox{seed 1} & 0.9334 & 1.092 \\
\mbox{random partition 4,} \mbox{seed 2} & 0.9341 & 1.043 \\
\mbox{random partition 5,} \mbox{seed 2} & 0.9293 & 1.118 \\
\mbox{no flavor boundary} & 0.9246, 0.9238 & 1.235, 1.205 \\
\mbox{$b$-quark boundary} & 0.9236, 0.9275 & 1.251, 1.147 \\
\mbox{random flavor boundary} & 0.9255, 0.9271 & 1.220, 1.153 \\
\mbox{no pretraining} & 0.9163\,$\pm$\,0.0006 & -- \\
\bottomrule
\end{tabular}
\end{adjustbox}
\end{minipage}
\par\bigskip
\begin{minipage}{\linewidth}\centering
\TableContinued{tab:pairs}{tab:pairs:2}
\begin{adjustbox}{max width=\linewidth, max totalheight=0.92\textheight}
\begin{tabular}{>{\raggedright\arraybackslash}p{4.2cm}cccc}
\toprule
 & \multicolumn{2}{c}{$X\to b\bar b$ vs $c\bar c$} & \multicolumn{2}{c}{$X\to YY\to b\bar{b}qq$ vs $c\bar{c}qq$} \\
\cmidrule(lr){2-3}\cmidrule(lr){4-5}
 & AUC & $\rho$ & AUC & $\rho$ \\
\midrule
\multicolumn{5}{l}{\textit{(LP): linear probe on the frozen features}} \\
\mbox{\Kfull{} classes} & 0.9977\,$\pm$\,$<$0.0001 & 1.00 & 0.9939\,$\pm$\,$<$0.0001 & 1.000 \\
\mbox{\Kres{} classes} & 0.9976\,$\pm$\,0.0001 & 1.01\,$\pm$\,0.03 & 0.9939\,$\pm$\,0.0002 & 1.005\,$\pm$\,0.032 \\
\mbox{\Kmid{} classes} & 0.9974\,$\pm$\,$<$0.0001 & 1.10\,$\pm$\,0.01 & 0.9932\,$\pm$\,$<$0.0001 & 1.117\,$\pm$\,0.008 \\
\mbox{\Kcoarse{} classes} & 0.9902\,$\pm$\,0.0007$^{\dagger}$ & 4.22\,$\pm$\,0.29 & 0.9716\,$\pm$\,0.0033$^{\dagger}$ & 4.676\,$\pm$\,0.513 \\
\mbox{\Kres{} classes + mass} & 0.9973, 0.9975 & 1.18, 1.09 & 0.9930, 0.9935 & 1.164, 1.070 \\
\mbox{\Kcoarse{} classes + mass} & 0.9891, 0.9904$^{\dagger}$ & 4.70, 4.13 & 0.9676, 0.9706$^{\dagger}$ & 5.358, 4.815 \\
\mbox{random partition 1,} \mbox{seed 1} & 0.9965$^{\dagger}$ & 1.51 & 0.9914$^{\dagger}$ & 1.427 \\
\mbox{random partition 2,} \mbox{seed 2} & 0.9964 & 1.54 & 0.9911 & 1.457 \\
\mbox{random partition 3,} \mbox{seed 1} & 0.9965$^{\dagger}$ & 1.53 & 0.9914$^{\dagger}$ & 1.430 \\
\mbox{random partition 4,} \mbox{seed 2} & 0.9968 & 1.37 & 0.9919 & 1.320 \\
\mbox{random partition 5,} \mbox{seed 2} & 0.9961$^{\dagger}$ & 1.67 & 0.9910 & 1.475 \\
\mbox{no flavor boundary} & 0.9864, 0.9687$^{\dagger}$ & 5.87, 13.40 & 0.9614, 0.9426$^{\dagger}$ & 6.394, 9.384 \\
\mbox{$b$-quark boundary} & 0.9957, 0.9959$^{\dagger}$ & 1.84, 1.75 & 0.9909, 0.9912 & 1.510, 1.447 \\
\mbox{random flavor boundary} & 0.9940, 0.9924$^{\dagger}$ & 2.58, 3.25 & 0.9863, 0.9860$^{\dagger}$ & 2.272, 2.281 \\
\mbox{random network} & 0.7012\,$\pm$\,0.0122 & -- & 0.6803\,$\pm$\,0.0068 & -- \\
\mbox{random network,} \mbox{BatchNorm recomputed} & 0.8169\,$\pm$\,0.0204 & -- & 0.7822\,$\pm$\,0.0189 & -- \\
\midrule
\multicolumn{5}{l}{\textit{(f.t.): fine-tuned}} \\
\mbox{\Kfull{} classes} & 0.9970\,$\pm$\,0.0002 & 1.00 & 0.9925\,$\pm$\,0.0003 & 1.000 \\
\mbox{\Kres{} classes} & 0.9968\,$\pm$\,0.0001 & 1.05\,$\pm$\,0.03 & 0.9925\,$\pm$\,0.0003 & 1.010\,$\pm$\,0.052 \\
\mbox{\Kmid{} classes} & 0.9966\,$\pm$\,0.0002 & 1.13\,$\pm$\,0.04 & 0.9918\,$\pm$\,0.0001 & 1.100\,$\pm$\,0.035 \\
\mbox{\Kcoarse{} classes} & 0.9943\,$\pm$\,0.0003 & 1.89\,$\pm$\,0.05 & 0.9884\,$\pm$\,0.0009 & 1.545\,$\pm$\,0.066 \\
\mbox{\Kres{} classes + mass} & 0.9966, 0.9968 & 1.15, 1.11 & 0.9915, 0.9919 & 1.159, 1.115 \\
\mbox{\Kcoarse{} classes + mass} & 0.9941, 0.9943 & 2.02, 1.95 & 0.9879, 0.9878 & 1.652, 1.680 \\
\mbox{random partition 1,} \mbox{seed 1} & 0.9957 & 1.48 & 0.9903 & 1.320 \\
\mbox{random partition 2,} \mbox{seed 2} & 0.9960 & 1.38 & 0.9903 & 1.328 \\
\mbox{random partition 3,} \mbox{seed 1} & 0.9959 & 1.41 & 0.9907 & 1.267 \\
\mbox{random partition 4,} \mbox{seed 2} & 0.9960 & 1.37 & 0.9911 & 1.228 \\
\mbox{random partition 5,} \mbox{seed 2} & 0.9956 & 1.53 & 0.9904 & 1.324 \\
\mbox{no flavor boundary} & 0.9949, 0.9948 & 1.74, 1.78 & 0.9893, 0.9886 & 1.466, 1.566 \\
\mbox{$b$-quark boundary} & 0.9954, 0.9954 & 1.59, 1.57 & 0.9906, 0.9910 & 1.284, 1.240 \\
\mbox{random flavor boundary} & 0.9950, 0.9952 & 1.71, 1.64 & 0.9890, 0.9901 & 1.499, 1.365 \\
\mbox{no pretraining} & 0.9910\,$\pm$\,0.0003 & -- & 0.9826\,$\pm$\,0.0005 & -- \\
\bottomrule
\end{tabular}
\end{adjustbox}
\end{minipage}
\par\bigskip
\begin{minipage}{\linewidth}\centering
\TableContinued{tab:pairs}{tab:pairs:3}
\begin{adjustbox}{max width=\linewidth, max totalheight=0.92\textheight}
\begin{tabular}{>{\raggedright\arraybackslash}p{4.2cm}ccc}
\toprule
 & \multicolumn{1}{c}{$X\to ee$ vs $\mu\mu$} & \multicolumn{2}{c}{QCD $b\bar{b}$ vs $c\bar{c}$} \\
\cmidrule(lr){2-2}\cmidrule(lr){3-4}
 & AUC & AUC & $\rho$ \\
\midrule
\multicolumn{4}{l}{\textit{(LP): linear probe on the frozen features}} \\
\mbox{\Kfull{} classes} & 1.0000\,$\pm$\,$<$0.0001 & 0.9867\,$\pm$\,0.0001 & 1.000 \\
\mbox{\Kres{} classes} & 1.0000\,$\pm$\,$<$0.0001 & 0.9848\,$\pm$\,0.0004$^{\dagger}$ & 1.148\,$\pm$\,0.037 \\
\mbox{\Kmid{} classes} & 0.9610\,$\pm$\,0.0030$^{\dagger}$ & 0.9846\,$\pm$\,0.0001$^{\dagger}$ & 1.163\,$\pm$\,0.016 \\
\mbox{\Kcoarse{} classes} & 0.9602\,$\pm$\,0.0202$^{\dagger}$ & 0.9563\,$\pm$\,0.0003$^{\dagger}$ & 3.294\,$\pm$\,0.033 \\
\mbox{\Kres{} classes + mass} & 1.0000, 1.0000 & 0.9829, 0.9842$^{\dagger}$ & 1.304, 1.185 \\
\mbox{\Kcoarse{} classes + mass} & 0.9693, 0.9983$^{\dagger}$ & 0.9481, 0.9531$^{\dagger}$ & 3.951, 3.508 \\
\mbox{random partition 1,} \mbox{seed 1} & 1.0000$^{\dagger}$ & 0.9795$^{\dagger}$ & 1.557 \\
\mbox{random partition 2,} \mbox{seed 2} & 1.0000$^{\dagger}$ & 0.9789$^{\dagger}$ & 1.580 \\
\mbox{random partition 3,} \mbox{seed 1} & 1.0000 & 0.9802$^{\dagger}$ & 1.509 \\
\mbox{random partition 4,} \mbox{seed 2} & 1.0000$^{\dagger}$ & 0.9809$^{\dagger}$ & 1.432 \\
\mbox{random partition 5,} \mbox{seed 2} & 1.0000 & 0.9786$^{\dagger}$ & 1.602 \\
\mbox{no flavor boundary} & 1.0000, 1.0000 & 0.9443, 0.9303$^{\dagger}$ & 4.244, 5.213 \\
\mbox{$b$-quark boundary} & 1.0000, 1.0000 & 0.9795, 0.9795$^{\dagger}$ & 1.561, 1.532 \\
\mbox{random flavor boundary} & 1.0000, 1.0000 & 0.9721, 0.9707$^{\dagger}$ & 2.124, 2.189 \\
\mbox{random network} & 0.9995\,$\pm$\,0.0003 & 0.6550\,$\pm$\,0.0048 & -- \\
\mbox{random network,} \mbox{BatchNorm recomputed} & 1.0000\,$\pm$\,$<$0.0001 & 0.7653\,$\pm$\,0.0162 & -- \\
\midrule
\multicolumn{4}{l}{\textit{(f.t.): fine-tuned}} \\
\mbox{\Kfull{} classes} & 1.0000\,$\pm$\,$<$0.0001 & 0.9841\,$\pm$\,0.0001 & 1.000 \\
\mbox{\Kres{} classes} & 1.0000\,$\pm$\,$<$0.0001 & 0.9826\,$\pm$\,0.0007 & 1.096\,$\pm$\,0.037 \\
\mbox{\Kmid{} classes} & 1.0000\,$\pm$\,$<$0.0001 & 0.9824\,$\pm$\,0.0003 & 1.105\,$\pm$\,0.027 \\
\mbox{\Kcoarse{} classes} & 1.0000\,$\pm$\,0 & 0.9765\,$\pm$\,0.0006 & 1.481\,$\pm$\,0.030 \\
\mbox{\Kres{} classes + mass} & 1.0000, 1.0000 & 0.9813, 0.9818 & 1.173, 1.153 \\
\mbox{\Kcoarse{} classes + mass} & 1.0000, 1.0000 & 0.9753, 0.9754 & 1.549, 1.559 \\
\mbox{random partition 1,} \mbox{seed 1} & 1.0000 & 0.9800 & 1.253 \\
\mbox{random partition 2,} \mbox{seed 2} & 1.0000 & 0.9793 & 1.316 \\
\mbox{random partition 3,} \mbox{seed 1} & 1.0000 & 0.9813 & 1.175 \\
\mbox{random partition 4,} \mbox{seed 2} & 1.0000 & 0.9805 & 1.236 \\
\mbox{random partition 5,} \mbox{seed 2} & 1.0000 & 0.9795 & 1.301 \\
\mbox{no flavor boundary} & 1.0000, 1.0000 & 0.9774, 0.9765 & 1.422, 1.490 \\
\mbox{$b$-quark boundary} & 1.0000, 1.0000 & 0.9796, 0.9809 & 1.279, 1.214 \\
\mbox{random flavor boundary} & 1.0000, 1.0000 & 0.9768, 0.9787 & 1.459, 1.351 \\
\mbox{no pretraining} & 1.0000\,$\pm$\,$<$0.0001 & 0.9707\,$\pm$\,0.0014 & -- \\
\bottomrule
\end{tabular}
\end{adjustbox}
\end{minipage}
\par\bigskip
\begin{minipage}{\linewidth}\centering
\TableContinued{tab:pairs}{tab:pairs:4}
\begin{adjustbox}{max width=\linewidth, max totalheight=0.92\textheight}
\begin{tabular}{>{\raggedright\arraybackslash}p{4.2cm}cccc}
\toprule
 & \multicolumn{2}{c}{QCD with $b$ vs light} & \multicolumn{2}{c}{QCD with $c$, no $b$, vs light} \\
\cmidrule(lr){2-3}\cmidrule(lr){4-5}
 & AUC & $\rho$ & AUC & $\rho$ \\
\midrule
\multicolumn{5}{l}{\textit{(LP): linear probe on the frozen features}} \\
\mbox{\Kfull{} classes} & 0.9880\,$\pm$\,0.0001 & 1.000 & 0.9324\,$\pm$\,0.0003 & 1.000 \\
\mbox{\Kres{} classes} & 0.9870\,$\pm$\,0.0001$^{\dagger}$ & 1.085\,$\pm$\,0.004 & 0.9282\,$\pm$\,0.0001$^{\dagger}$ & 1.061\,$\pm$\,0.005 \\
\mbox{\Kmid{} classes} & 0.9870\,$\pm$\,0.0001$^{\dagger}$ & 1.081\,$\pm$\,0.008 & 0.9275\,$\pm$\,0.0003$^{\dagger}$ & 1.072\,$\pm$\,0.008 \\
\mbox{\Kcoarse{} classes} & 0.9815\,$\pm$\,0.0007$^{\dagger}$ & 1.544\,$\pm$\,0.069 & 0.9048\,$\pm$\,0.0010$^{\dagger}$ & 1.407\,$\pm$\,0.021 \\
\mbox{\Kres{} classes + mass} & 0.9868, 0.9870$^{\dagger}$ & 1.107, 1.078 & 0.9262, 0.9273$^{\dagger}$ & 1.093, 1.070 \\
\mbox{\Kcoarse{} classes + mass} & 0.9804, 0.9800$^{\dagger}$ & 1.638, 1.662 & 0.9046, 0.9002$^{\dagger}$ & 1.414, 1.470 \\
\mbox{random partition 1,} \mbox{seed 1} & 0.9857$^{\dagger}$ & 1.201 & 0.9219$^{\dagger}$ & 1.157 \\
\mbox{random partition 2,} \mbox{seed 2} & 0.9854$^{\dagger}$ & 1.212 & 0.9210$^{\dagger}$ & 1.163 \\
\mbox{random partition 3,} \mbox{seed 1} & 0.9856$^{\dagger}$ & 1.209 & 0.9229$^{\dagger}$ & 1.143 \\
\mbox{random partition 4,} \mbox{seed 2} & 0.9860$^{\dagger}$ & 1.163 & 0.9240$^{\dagger}$ & 1.119 \\
\mbox{random partition 5,} \mbox{seed 2} & 0.9857$^{\dagger}$ & 1.186 & 0.9223$^{\dagger}$ & 1.144 \\
\mbox{no flavor boundary} & 0.9814, 0.9825$^{\dagger}$ & 1.557, 1.453 & 0.9088, 0.9122$^{\dagger}$ & 1.352, 1.293 \\
\mbox{$b$-quark boundary} & 0.9855, 0.9854$^{\dagger}$ & 1.212, 1.213 & 0.9166, 0.9171$^{\dagger}$ & 1.237, 1.220 \\
\mbox{random flavor boundary} & 0.9848, 0.9845$^{\dagger}$ & 1.277, 1.286 & 0.9204, 0.9211$^{\dagger}$ & 1.179, 1.162 \\
\mbox{random network} & 0.7419\,$\pm$\,0.0077 & -- & 0.6561\,$\pm$\,0.0049 & -- \\
\mbox{random network,} \mbox{BatchNorm recomputed} & 0.8294\,$\pm$\,0.0116 & -- & 0.6821\,$\pm$\,0.0003 & -- \\
\midrule
\multicolumn{5}{l}{\textit{(f.t.): fine-tuned}} \\
\mbox{\Kfull{} classes} & 0.9882\,$\pm$\,0.0003 & 1.000 & 0.9285\,$\pm$\,0.0006 & 1.000 \\
\mbox{\Kres{} classes} & 0.9879\,$\pm$\,0.0004 & 1.023\,$\pm$\,0.017 & 0.9268\,$\pm$\,0.0001 & 1.024\,$\pm$\,0.009 \\
\mbox{\Kmid{} classes} & 0.9878\,$\pm$\,0.0002 & 1.034\,$\pm$\,0.030 & 0.9253\,$\pm$\,0.0006 & 1.045\,$\pm$\,0.001 \\
\mbox{\Kcoarse{} classes} & 0.9870\,$\pm$\,0.0003 & 1.103\,$\pm$\,0.025 & 0.9219\,$\pm$\,0.0011 & 1.093\,$\pm$\,0.013 \\
\mbox{\Kres{} classes + mass} & 0.9878, 0.9882 & 1.062, 0.994 & 0.9254, 0.9265 & 1.047, 1.035 \\
\mbox{\Kcoarse{} classes + mass} & 0.9869, 0.9874 & 1.142, 1.058 & 0.9200, 0.9206 & 1.122, 1.118 \\
\mbox{random partition 1,} \mbox{seed 1} & 0.9875 & 1.088 & 0.9230 & 1.080 \\
\mbox{random partition 2,} \mbox{seed 2} & 0.9874 & 1.058 & 0.9229 & 1.086 \\
\mbox{random partition 3,} \mbox{seed 1} & 0.9875 & 1.088 & 0.9244 & 1.060 \\
\mbox{random partition 4,} \mbox{seed 2} & 0.9873 & 1.069 & 0.9244 & 1.063 \\
\mbox{random partition 5,} \mbox{seed 2} & 0.9870 & 1.089 & 0.9230 & 1.084 \\
\mbox{no flavor boundary} & 0.9875, 0.9874 & 1.087, 1.055 & 0.9230, 0.9231 & 1.079, 1.082 \\
\mbox{$b$-quark boundary} & 0.9872, 0.9874 & 1.109, 1.060 & 0.9225, 0.9232 & 1.086, 1.081 \\
\mbox{random flavor boundary} & 0.9873, 0.9874 & 1.099, 1.056 & 0.9223, 0.9237 & 1.089, 1.074 \\
\mbox{no pretraining} & 0.9864\,$\pm$\,0.0001 & -- & 0.9166\,$\pm$\,0.0006 & -- \\
\bottomrule
\end{tabular}
\end{adjustbox}
\end{minipage}
\end{center}

\begin{landscape}   % wider than the 15 cm text block
{\scriptsize
\setlength{\tabcolsep}{2pt}
\begin{longtable}{lcccccccccccc}
\caption{Linear probes on the frozen features of networks pretrained on JetClass-II with 188, 162, 43 and 17 classes and of the controls, trained on JetClass-II (class-balanced test set, largest training set), per class: the AUC of the class against the rest for the \Kfull-class model (mean $\pm$ SD over its runs) and the ratio of $1-\mathrm{AUC}$ to it of every other pretraining, run by run and averaged; p1--p5 are the random partitions, then the partitions with no, a $b$-quark and a random flavor boundary. Classes grouped by their \Kcoarse-class group. $^{\S}$ marks a value from a single run. $^{\P}$ marks a value whose standard deviation over the runs rounds to zero at the digits shown.}\label{tab:perclass}\\
\toprule
class & AUC, \Kfull{} & \Kres{} & \Kmid{} & \Kcoarse{} & p1 & p2 & p3 & p4 & p5 & no flavor & $b$ quark & random flavor \\
\midrule
\endfirsthead
\toprule
class & AUC, \Kfull{} & \Kres{} & \Kmid{} & \Kcoarse{} & p1 & p2 & p3 & p4 & p5 & no flavor & $b$ quark & random flavor \\
\midrule
\endhead
\bottomrule
\endlastfoot
$X\to bb$ & 0.98932\,$\pm$\,0.00018 & 1.0112\,$\pm$\,0.0294 & 1.066\,$\pm$\,0.018 & 1.445\,$\pm$\,0.007 & 1.14$^{\S}$ & 1.11$^{\S}$ & 1.15$^{\S}$ & 1.13$^{\S}$ & 1.18$^{\S}$ & 1.562\,$\pm$\,0.107 & 1.1628\,$\pm$\,0.0377 & 1.2485\,$\pm$\,0.0685 \\
$X\to cc$ & 0.97782\,$\pm$\,0.00013 & 0.9977\,$\pm$\,0.0109 & 1.048\,$\pm$\,0.012 & 1.563\,$\pm$\,0.020 & 1.16$^{\S}$ & 1.17$^{\S}$ & 1.15$^{\S}$ & 1.15$^{\S}$ & 1.15$^{\S}$ & 1.638\,$\pm$\,0.114 & 1.2365\,$\pm$\,0.0215 & 1.2543\,$\pm$\,0.0319 \\
$X\to ss$ & 0.97603\,$\pm$\,0.00005 & 1.0023\,$\pm$\,0.0048 & 1.078\,$\pm$\,0.004 & 1.401\,$\pm$\,0.030 & 1.08$^{\S}$ & 1.08$^{\S}$ & 1.07$^{\S}$ & 1.04$^{\S}$ & 1.10$^{\S}$ & 1.390\,$\pm$\,0.001 & 1.2164\,$\pm$\,0.0078 & 1.1581\,$\pm$\,0.0037 \\
$X\to qq$ & 0.97505\,$\pm$\,0.00024 & 1.0088\,$\pm$\,0.0054 & 1.062\,$\pm$\,0.012 & 1.284\,$\pm$\,0.061 & 1.06$^{\S}$ & 1.09$^{\S}$ & 1.06$^{\S}$ & 1.08$^{\S}$ & 1.08$^{\S}$ & 1.227\,$\pm$\,0.009 & 1.1753\,$\pm$\,0.0097 & 1.1196\,$\pm$\,0.0059 \\
$X\to bc$ & 0.98102\,$\pm$\,0.00013 & 1.0097\,$\pm$\,0.0129 & 1.084\,$\pm$\,0.017 & 1.458\,$\pm$\,0.026 & 1.19$^{\S}$ & 1.18$^{\S}$ & 1.14$^{\S}$ & 1.19$^{\S}$ & 1.20$^{\S}$ & 1.484\,$\pm$\,0.021 & 1.2463\,$\pm$\,0.0104 & 1.3249\,$\pm$\,0.0004 \\
$X\to cs$ & 0.96902\,$\pm$\,0.00007 & 1.0022\,$\pm$\,0.0089 & 1.069\,$\pm$\,0.005 & 1.343\,$\pm$\,0.015 & 1.11$^{\S}$ & 1.13$^{\S}$ & 1.09$^{\S}$ & 1.10$^{\S}$ & 1.13$^{\S}$ & 1.366\,$\pm$\,0.015 & 1.2189\,$\pm$\,0.0266 & 1.2230\,$\pm$\,0.0047 \\
$X\to bq$ & 0.98120\,$\pm$\,0.00011 & 1.0249\,$\pm$\,0.0148 & 1.141\,$\pm$\,0.019 & 1.607\,$\pm$\,0.027 & 1.22$^{\S}$ & 1.22$^{\S}$ & 1.17$^{\S}$ & 1.18$^{\S}$ & 1.26$^{\S}$ & 1.590\,$\pm$\,0.030 & 1.3460\,$\pm$\,0.0288 & 1.3781\,$\pm$\,0.0094 \\
$X\to cq$ & 0.97004\,$\pm$\,0.00033 & 1.0046\,$\pm$\,0.0203 & 1.102\,$\pm$\,0.017 & 1.358\,$\pm$\,0.014 & 1.11$^{\S}$ & 1.12$^{\S}$ & 1.10$^{\S}$ & 1.10$^{\S}$ & 1.15$^{\S}$ & 1.396\,$\pm$\,0.030 & 1.2218\,$\pm$\,0.0258 & 1.2421\,$\pm$\,0.0151 \\
$X\to sq$ & 0.97454\,$\pm$\,0.00011 & 1.0015\,$\pm$\,0.0051 & 1.044\,$\pm$\,0.004 & 1.299\,$\pm$\,0.027 & 1.08$^{\S}$ & 1.08$^{\S}$ & 1.08$^{\S}$ & 1.05$^{\S}$ & 1.10$^{\S}$ & 1.264\,$\pm$\,0.018 & 1.1745\,$\pm$\,0.0048 & 1.1613\,$\pm$\,0.0058 \\
$X\to gg$ & 0.96142\,$\pm$\,0.00010 & 1.0082\,$\pm$\,0.0065 & 1.037\,$\pm$\,0.017 & 1.125\,$\pm$\,0.003 & 1.04$^{\S}$ & 1.04$^{\S}$ & 1.05$^{\S}$ & 1.06$^{\S}$ & 1.05$^{\S}$ & 1.086\,$\pm$\,0.004 & 1.0474\,$\pm$\,0.0043 & 1.0395\,$\pm$\,0.0045 \\
\midrule
$X\to \tau_h\tau_h$ & 0.99939\,$\pm$\,0.00003 & 1.0412\,$\pm$\,0.0443 & 1.028\,$\pm$\,0.029 & 1.321\,$\pm$\,0.224 & 1.29$^{\S}$ & 1.11$^{\S}$ & 1.38$^{\S}$ & 1.09$^{\S}$ & 1.26$^{\S}$ & 1.128\,$\pm$\,0.089 & 1.0897\,$\pm$\,0.0311 & 1.2171\,$\pm$\,0.1112 \\
\midrule
$X\to \tau_h\tau_e$ & 0.99935\,$\pm$\,0.00002 & 0.9885\,$\pm$\,0.0645 & 4.020\,$\pm$\,0.632 & 3.359\,$\pm$\,0.610 & 1.09$^{\S}$ & 1.07$^{\S}$ & 1.21$^{\S}$ & 1.04$^{\S}$ & 1.26$^{\S}$ & 1.078\,$\pm$\,0.003 & 1.1029\,$\pm$\,0.0847 & 1.0430\,$\pm$\,0.0580 \\
$X\to \tau_h\tau_\mu$ & 0.99930\,$\pm$\,0.00003 & 1.0459\,$\pm$\,0.0987 & 3.804\,$\pm$\,0.293 & 3.320\,$\pm$\,0.296 & 1.20$^{\S}$ & 1.02$^{\S}$ & 1.20$^{\S}$ & 1.03$^{\S}$ & 1.13$^{\S}$ & 1.202\,$\pm$\,0.061 & 1.0803\,$\pm$\,0.0227 & 1.2541\,$\pm$\,0.1839 \\
\midrule
$X\to ee$ & 0.99963\,$\pm$\,0.00002 & 1.0057\,$\pm$\,0.0472 & 3.638\,$\pm$\,0.162 & 3.600\,$\pm$\,0.821 & 1.00$^{\S}$ & 1.03$^{\S}$ & 1.00$^{\S}$ & 1.09$^{\S}$ & 1.06$^{\S}$ & 1.069\,$\pm$\,0.027 & 1.0941\,$\pm$\,0.0801 & 1.0304\,$\pm$\,0.0659 \\
$X\to \mu\mu$ & 0.99950\,$\pm$\,0.00001 & 1.0007\,$\pm$\,0.0501 & 3.455\,$\pm$\,0.322 & 3.533\,$\pm$\,0.481 & 0.99$^{\S}$ & 1.04$^{\S}$ & 1.07$^{\S}$ & 1.16$^{\S}$ & 1.06$^{\S}$ & 1.214\,$\pm$\,0.175 & 1.0807\,$\pm$\,0.0333 & 1.0924\,$\pm$\,0.0551 \\
\midrule
$X\to YY\to bbb$ & 0.99121\,$\pm$\,0.00010 & 1.0043\,$\pm$\,0.0042 & 1.105\,$\pm$\,0.031 & 1.516\,$\pm$\,0.022 & 1.14$^{\S}$ & 1.15$^{\S}$ & 1.15$^{\S}$ & 1.15$^{\S}$ & 1.15$^{\S}$ & 1.720\,$\pm$\,0.181 & 1.2434\,$\pm$\,0.0336 & 1.3189\,$\pm$\,0.0767 \\
$X\to YY\to bbc$ & 0.97741\,$\pm$\,0.00013 & 1.0013\,$\pm$\,0.0093 & 1.102\,$\pm$\,0.003 & 1.462\,$\pm$\,0.018 & 1.12$^{\S}$ & 1.12$^{\S}$ & 1.12$^{\S}$ & 1.11$^{\S}$ & 1.09$^{\S}$ & 1.578\,$\pm$\,0.069 & 1.2363\,$\pm$\,0.0525 & 1.2034\,$\pm$\,0.0097 \\
$X\to YY\to bbs$ & 0.97558\,$\pm$\,0.00005 & 0.9978\,$\pm$\,0.0028 & 1.123\,$\pm$\,0.010 & 1.533\,$\pm$\,0.028 & 1.13$^{\S}$ & 1.11$^{\S}$ & 1.11$^{\S}$ & 1.10$^{\S}$ & 1.12$^{\S}$ & 1.572\,$\pm$\,0.066 & 1.2267\,$\pm$\,0.0194 & 1.2180\,$\pm$\,0.0207 \\
$X\to YY\to bbq$ & 0.97401\,$\pm$\,0.00008 & 1.0104\,$\pm$\,0.0049 & 1.109\,$\pm$\,0.003 & 1.484\,$\pm$\,0.014 & 1.13$^{\S}$ & 1.13$^{\S}$ & 1.10$^{\S}$ & 1.11$^{\S}$ & 1.13$^{\S}$ & 1.607\,$\pm$\,0.078 & 1.2098\,$\pm$\,0.0134 & 1.2127\,$\pm$\,0.0059 \\
$X\to YY\to ccb$ & 0.97048\,$\pm$\,0.00036 & 0.9962\,$\pm$\,0.0129 & 1.079\,$\pm$\,0.020 & 1.575\,$\pm$\,0.029 & 1.14$^{\S}$ & 1.19$^{\S}$ & 1.11$^{\S}$ & 1.17$^{\S}$ & 1.14$^{\S}$ & 1.598\,$\pm$\,0.026 & 1.2594\,$\pm$\,0.0461 & 1.2653\,$\pm$\,0.0491 \\
$X\to YY\to ccc$ & 0.97565\,$\pm$\,0.00009 & 0.9964\,$\pm$\,0.0078 & 1.039\,$\pm$\,0.016 & 1.753\,$\pm$\,0.057 & 1.15$^{\S}$ & 1.18$^{\S}$ & 1.11$^{\S}$ & 1.15$^{\S}$ & 1.19$^{\S}$ & 1.889\,$\pm$\,0.160 & 1.2387\,$\pm$\,0.0121 & 1.3000\,$\pm$\,0.0487 \\
$X\to YY\to ccs$ & 0.95974\,$\pm$\,0.00012 & 1.0031\,$\pm$\,0.0090 & 1.058\,$\pm$\,0.009 & 1.488\,$\pm$\,0.024 & 1.11$^{\S}$ & 1.10$^{\S}$ & 1.09$^{\S}$ & 1.08$^{\S}$ & 1.09$^{\S}$ & 1.583\,$\pm$\,0.067 & 1.1868\,$\pm$\,0.0171 & 1.1928\,$\pm$\,0.0128 \\
$X\to YY\to ccq$ & 0.95693\,$\pm$\,0.00014 & 0.9974\,$\pm$\,0.0032 & 1.058\,$\pm$\,0.005 & 1.491\,$\pm$\,0.039 & 1.09$^{\S}$ & 1.11$^{\S}$ & 1.09$^{\S}$ & 1.07$^{\S}$ & 1.10$^{\S}$ & 1.587\,$\pm$\,0.118 & 1.1642\,$\pm$\,0.0216 & 1.1857\,$\pm$\,0.0057 \\
$X\to YY\to ssb$ & 0.97147\,$\pm$\,0.00016 & 0.9972\,$\pm$\,0.0080 & 1.149\,$\pm$\,0.002 & 1.922\,$\pm$\,0.093 & 1.13$^{\S}$ & 1.13$^{\S}$ & 1.10$^{\S}$ & 1.16$^{\S}$ & 1.18$^{\S}$ & 1.775\,$\pm$\,0.007 & 1.3905\,$\pm$\,0.0329 & 1.2830\,$\pm$\,0.0150 \\
$X\to YY\to ssc$ & 0.95610\,$\pm$\,0.00046 & 0.9981\,$\pm$\,0.0143 & 1.080\,$\pm$\,0.009 & 1.434\,$\pm$\,0.035 & 1.05$^{\S}$ & 1.06$^{\S}$ & 1.05$^{\S}$ & 1.07$^{\S}$ & 1.07$^{\S}$ & 1.447\,$\pm$\,0.025 & 1.1896\,$\pm$\,0.0117 & 1.1174\,$\pm$\,0.0020 \\
$X\to YY\to sss$ & 0.97265\,$\pm$\,0.00014 & 1.0060\,$\pm$\,0.0050 & 1.093\,$\pm$\,0.007 & 1.540\,$\pm$\,0.050 & 1.07$^{\S}$ & 1.07$^{\S}$ & 1.05$^{\S}$ & 1.05$^{\S}$ & 1.06$^{\S}$ & 1.490\,$\pm$\,0.048 & 1.2461\,$\pm$\,0.0366 & 1.1213\,$\pm$\,0.0033 \\
$X\to YY\to ssq$ & 0.96308\,$\pm$\,0.00007 & 1.0004\,$\pm$\,0.0029 & 1.051\,$\pm$\,0.009 & 1.317\,$\pm$\,0.046 & 1.05$^{\S}$ & 1.06$^{\S}$ & 1.03$^{\S}$ & 1.06$^{\S}$ & 1.05$^{\S}$ & 1.280\,$\pm$\,0.018 & 1.1470\,$\pm$\,0.0063 & 1.0984\,$\pm$\,0.0151 \\
$X\to YY\to qqb$ & 0.96662\,$\pm$\,0.00028 & 0.9998\,$\pm$\,0.0070 & 1.091\,$\pm$\,0.020 & 1.597\,$\pm$\,0.038 & 1.11$^{\S}$ & 1.11$^{\S}$ & 1.08$^{\S}$ & 1.13$^{\S}$ & 1.11$^{\S}$ & 1.543\,$\pm$\,0.040 & 1.2894\,$\pm$\,0.0136 & 1.2271\,$\pm$\,0.0134 \\
$X\to YY\to qqc$ & 0.95188\,$\pm$\,0.00012 & 0.9998\,$\pm$\,0.0058 & 1.083\,$\pm$\,0.006 & 1.337\,$\pm$\,0.007 & 1.08$^{\S}$ & 1.08$^{\S}$ & 1.10$^{\S}$ & 1.06$^{\S}$ & 1.07$^{\S}$ & 1.372\,$\pm$\,0.054 & 1.1847\,$\pm$\,0.0013 & 1.1691\,$\pm$\,0.0030 \\
$X\to YY\to qqs$ & 0.96263\,$\pm$\,0.00020 & 1.0059\,$\pm$\,0.0058 & 1.034\,$\pm$\,0.005 & 1.260\,$\pm$\,0.044 & 1.04$^{\S}$ & 1.05$^{\S}$ & 1.04$^{\S}$ & 1.05$^{\S}$ & 1.04$^{\S}$ & 1.209\,$\pm$\,0.012 & 1.1271\,$\pm$\,0.0027 & 1.0698\,$\pm$\,0.0092 \\
$X\to YY\to qqq$ & 0.96909\,$\pm$\,0.00015 & 1.0098\,$\pm$\,0.0066 & 1.085\,$\pm$\,0.006 & 1.402\,$\pm$\,0.075 & 1.05$^{\S}$ & 1.07$^{\S}$ & 1.03$^{\S}$ & 1.07$^{\S}$ & 1.05$^{\S}$ & 1.292\,$\pm$\,0.019 & 1.2256\,$\pm$\,0.0030 & 1.1002\,$\pm$\,0.0073 \\
$X\to YY\to bcs$ & 0.95828\,$\pm$\,0.00020 & 1.0013\,$\pm$\,0.0052 & 1.112\,$\pm$\,0.011 & 1.429\,$\pm$\,0.015 & 1.10$^{\S}$ & 1.13$^{\S}$ & 1.10$^{\S}$ & 1.13$^{\S}$ & 1.11$^{\S}$ & 1.409\,$\pm$\,0.014 & 1.2627\,$\pm$\,0.0211 & 1.2225\,$\pm$\,0.0121 \\
$X\to YY\to bcq$ & 0.95816\,$\pm$\,0.00030 & 1.0095\,$\pm$\,0.0102 & 1.119\,$\pm$\,0.010 & 1.388\,$\pm$\,0.012 & 1.15$^{\S}$ & 1.15$^{\S}$ & 1.13$^{\S}$ & 1.13$^{\S}$ & 1.12$^{\S}$ & 1.418\,$\pm$\,0.001 & 1.2084\,$\pm$\,0.0074 & 1.2117\,$\pm$\,0.0096 \\
$X\to YY\to bsq$ & 0.96477\,$\pm$\,0.00031 & 0.9982\,$\pm$\,0.0154 & 1.081\,$\pm$\,0.007 & 1.491\,$\pm$\,0.023 & 1.13$^{\S}$ & 1.10$^{\S}$ & 1.11$^{\S}$ & 1.14$^{\S}$ & 1.15$^{\S}$ & 1.470\,$\pm$\,0.013 & 1.2292\,$\pm$\,0.0444 & 1.2414\,$\pm$\,0.0166 \\
$X\to YY\to csq$ & 0.94770\,$\pm$\,0.00039 & 1.0013\,$\pm$\,0.0053 & 1.053\,$\pm$\,0.019 & 1.257\,$\pm$\,0.015 & 1.07$^{\S}$ & 1.07$^{\S}$ & 1.07$^{\S}$ & 1.07$^{\S}$ & 1.06$^{\S}$ & 1.280\,$\pm$\,0.040 & 1.1104\,$\pm$\,0.0005 & 1.1267\,$\pm$\,0.0080 \\
$X\to YY\to bbg$ & 0.97705\,$\pm$\,0.00022 & 1.0145\,$\pm$\,0.0191 & 1.095\,$\pm$\,0.025 & 1.345\,$\pm$\,0.029 & 1.10$^{\S}$ & 1.11$^{\S}$ & 1.10$^{\S}$ & 1.12$^{\S}$ & 1.13$^{\S}$ & 1.372\,$\pm$\,0.089 & 1.0985\,$\pm$\,0.0020 & 1.1548\,$\pm$\,0.0269 \\
$X\to YY\to ccg$ & 0.96355\,$\pm$\,0.00012 & 1.0053\,$\pm$\,0.0034 & 1.063\,$\pm$\,0.002 & 1.525\,$\pm$\,0.053 & 1.11$^{\S}$ & 1.11$^{\S}$ & 1.10$^{\S}$ & 1.10$^{\S}$ & 1.10$^{\S}$ & 1.504\,$\pm$\,0.105 & 1.1571\,$\pm$\,0.0014 & 1.1875\,$\pm$\,0.0210 \\
$X\to YY\to ssg$ & 0.96078\,$\pm$\,0.00003 & 1.0032\,$\pm$\,0.0013 & 1.104\,$\pm$\,0.008 & 1.386\,$\pm$\,0.016 & 1.05$^{\S}$ & 1.07$^{\S}$ & 1.06$^{\S}$ & 1.07$^{\S}$ & 1.06$^{\S}$ & 1.306\,$\pm$\,0.006 & 1.1546\,$\pm$\,0.0178 & 1.1014\,$\pm$\,0.0136 \\
$X\to YY\to qqg$ & 0.95995\,$\pm$\,0.00007 & 0.9977\,$\pm$\,0.0023 & 1.089\,$\pm$\,0.007 & 1.277\,$\pm$\,0.036 & 1.04$^{\S}$ & 1.03$^{\S}$ & 1.04$^{\S}$ & 1.06$^{\S}$ & 1.05$^{\S}$ & 1.177\,$\pm$\,0.011 & 1.1417\,$\pm$\,0.0109 & 1.0686\,$\pm$\,0.0085 \\
$X\to YY\to ggb$ & 0.96771\,$\pm$\,0.00012 & 0.9999\,$\pm$\,0.0081 & 1.100\,$\pm$\,0.024 & 1.384\,$\pm$\,0.007 & 1.08$^{\S}$ & 1.08$^{\S}$ & 1.06$^{\S}$ & 1.10$^{\S}$ & 1.10$^{\S}$ & 1.272\,$\pm$\,0.070 & 1.1049\,$\pm$\,0.0110 & 1.1523\,$\pm$\,0.0213 \\
$X\to YY\to ggc$ & 0.95547\,$\pm$\,0.00030 & 0.9989\,$\pm$\,0.0046 & 1.104\,$\pm$\,0.012 & 1.337\,$\pm$\,0.012 & 1.10$^{\S}$ & 1.08$^{\S}$ & 1.07$^{\S}$ & 1.07$^{\S}$ & 1.11$^{\S}$ & 1.251\,$\pm$\,0.030 & 1.1133\,$\pm$\,0.0110 & 1.1438\,$\pm$\,0.0034 \\
$X\to YY\to ggs$ & 0.95473\,$\pm$\,0.00009 & 0.9989\,$\pm$\,0.0010 & 1.069\,$\pm$\,0.010 & 1.233\,$\pm$\,0.020 & 1.04$^{\S}$ & 1.04$^{\S}$ & 1.04$^{\S}$ & 1.04$^{\S}$ & 1.06$^{\S}$ & 1.148\,$\pm$\,0.009 & 1.0946\,$\pm$\,0.0020 & 1.0695\,$\pm$\,0.0058 \\
$X\to YY\to ggq$ & 0.95626\,$\pm$\,0.00012 & 1.0035\,$\pm$\,0.0095 & 1.059\,$\pm$\,0.007 & 1.190\,$\pm$\,0.037 & 1.05$^{\S}$ & 1.04$^{\S}$ & 1.05$^{\S}$ & 1.06$^{\S}$ & 1.04$^{\S}$ & 1.119\,$\pm$\,0.007 & 1.0891\,$\pm$\,0.0011 & 1.0725\,$\pm$\,0.0112 \\
$X\to YY\to ggg$ & 0.96812\,$\pm$\,0.00010 & 1.0020\,$\pm$\,0.0004 & 1.036\,$\pm$\,0.020 & 1.120\,$\pm$\,0.021 & 1.02$^{\S}$ & 1.04$^{\S}$ & 1.04$^{\S}$ & 1.05$^{\S}$ & 1.05$^{\S}$ & 1.057\,$\pm$\,0.004 & 1.0321\,$\pm$\,0.0023 & 1.0287\,$\pm$\,0.0110 \\
\midrule
$X\to YY\to b\tau_h\tau_h$ & 0.99649\,$\pm$\,0.00010 & 0.9750\,$\pm$\,0.0481 & 0.987\,$\pm$\,0.028 & 1.366\,$\pm$\,0.024 & 1.07$^{\S}$ & 1.08$^{\S}$ & 1.13$^{\S}$ & 1.07$^{\S}$ & 1.11$^{\S}$ & 1.552\,$\pm$\,0.074 & 1.2109\,$\pm$\,0.0239 & 1.2974\,$\pm$\,0.0110 \\
$X\to YY\to c\tau_h\tau_h$ & 0.99291\,$\pm$\,0.00013 & 1.0025\,$\pm$\,0.0041 & 1.018\,$\pm$\,0.031 & 1.345\,$\pm$\,0.039 & 1.06$^{\S}$ & 1.09$^{\S}$ & 1.07$^{\S}$ & 1.13$^{\S}$ & 1.11$^{\S}$ & 1.383\,$\pm$\,0.053 & 1.2183\,$\pm$\,0.0220 & 1.2282\,$\pm$\,0.0173 \\
$X\to YY\to s\tau_h\tau_h$ & 0.99199\,$\pm$\,0.00003 & 1.0042\,$\pm$\,0.0057 & 1.068\,$\pm$\,0.015 & 1.303\,$\pm$\,0.029 & 1.06$^{\S}$ & 1.05$^{\S}$ & 1.05$^{\S}$ & 1.05$^{\S}$ & 1.07$^{\S}$ & 1.267\,$\pm$\,0.026 & 1.2070\,$\pm$\,0.0255 & 1.1199\,$\pm$\,0.0207 \\
$X\to YY\to q\tau_h\tau_h$ & 0.99185\,$\pm$\,0.00008 & 0.9968\,$\pm$\,0.0183 & 1.047\,$\pm$\,0.024 & 1.212\,$\pm$\,0.011 & 1.06$^{\S}$ & 1.05$^{\S}$ & 1.05$^{\S}$ & 1.03$^{\S}$ & 1.04$^{\S}$ & 1.159\,$\pm$\,0.012 & 1.1507\,$\pm$\,0.0019 & 1.1160\,$\pm$\,0.0080 \\
$X\to YY\to g\tau_h\tau_h$ & 0.99076\,$\pm$\,0.00007 & 1.0071\,$\pm$\,0.0079 & 1.030\,$\pm$\,0.013 & 1.147\,$\pm$\,0.016 & 1.05$^{\S}$ & 1.07$^{\S}$ & 1.04$^{\S}$ & 1.08$^{\S}$ & 1.07$^{\S}$ & 1.071\,$\pm$\,0.005 & 1.0236\,$\pm$\,0.0232 & 1.0244\,$\pm$\,0.0279 \\
\midrule
$X\to YY\to bbe$ & 0.99767\,$\pm$\,0.00007 & 1.0453\,$\pm$\,0.0286 & 1.962\,$\pm$\,0.256 & 2.102\,$\pm$\,0.299 & 1.20$^{\S}$ & 1.15$^{\S}$ & 1.18$^{\S}$ & 1.17$^{\S}$ & 1.28$^{\S}$ & 1.713\,$\pm$\,0.234 & 1.2550\,$\pm$\,0.0153 & 1.3919\,$\pm$\,0.0265 \\
$X\to YY\to cce$ & 0.99389\,$\pm$\,0.00006 & 1.0087\,$\pm$\,0.0153 & 1.704\,$\pm$\,0.145 & 2.018\,$\pm$\,0.194 & 1.19$^{\S}$ & 1.13$^{\S}$ & 1.16$^{\S}$ & 1.13$^{\S}$ & 1.20$^{\S}$ & 1.689\,$\pm$\,0.077 & 1.2899\,$\pm$\,0.0176 & 1.3378\,$\pm$\,0.0001 \\
$X\to YY\to sse$ & 0.99218\,$\pm$\,0.00006 & 1.0207\,$\pm$\,0.0108 & 1.553\,$\pm$\,0.082 & 1.751\,$\pm$\,0.128 & 1.09$^{\S}$ & 1.10$^{\S}$ & 1.10$^{\S}$ & 1.07$^{\S}$ & 1.10$^{\S}$ & 1.466\,$\pm$\,0.002 & 1.2821\,$\pm$\,0.0007 & 1.1842\,$\pm$\,0.0476 \\
$X\to YY\to qqe$ & 0.99112\,$\pm$\,0.00009 & 1.0074\,$\pm$\,0.0101 & 1.475\,$\pm$\,0.091 & 1.539\,$\pm$\,0.138 & 1.06$^{\S}$ & 1.08$^{\S}$ & 1.07$^{\S}$ & 1.05$^{\S}$ & 1.05$^{\S}$ & 1.242\,$\pm$\,0.002 & 1.1976\,$\pm$\,0.0276 & 1.1245\,$\pm$\,0.0060 \\
$X\to YY\to bb\mu$ & 0.99763\,$\pm$\,0.00006 & 1.0193\,$\pm$\,0.0251 & 2.043\,$\pm$\,0.076 & 1.985\,$\pm$\,0.195 & 1.14$^{\S}$ & 1.12$^{\S}$ & 1.23$^{\S}$ & 1.15$^{\S}$ & 1.19$^{\S}$ & 1.779\,$\pm$\,0.135 & 1.2127\,$\pm$\,0.0634 & 1.4177\,$\pm$\,0.0025 \\
$X\to YY\to cc\mu$ & 0.99420\,$\pm$\,0.00012 & 0.9964\,$\pm$\,0.0139 & 1.874\,$\pm$\,0.118 & 2.270\,$\pm$\,0.249 & 1.09$^{\S}$ & 1.17$^{\S}$ & 1.12$^{\S}$ & 1.24$^{\S}$ & 1.24$^{\S}$ & 1.729\,$\pm$\,0.169 & 1.3472\,$\pm$\,0.0333 & 1.3977\,$\pm$\,0.0166 \\
$X\to YY\to ss\mu$ & 0.99230\,$\pm$\,0.00015 & 1.0176\,$\pm$\,0.0368 & 1.636\,$\pm$\,0.018 & 1.796\,$\pm$\,0.095 & 1.07$^{\S}$ & 1.10$^{\S}$ & 1.10$^{\S}$ & 1.07$^{\S}$ & 1.10$^{\S}$ & 1.480\,$\pm$\,0.021 & 1.2995\,$\pm$\,0.0043 & 1.1823\,$\pm$\,0.0324 \\
$X\to YY\to qq\mu$ & 0.99026\,$\pm$\,0.00006 & 1.0014\,$\pm$\,0.0031 & 1.552\,$\pm$\,0.075 & 1.561\,$\pm$\,0.148 & 1.04$^{\S}$ & 1.05$^{\S}$ & 1.04$^{\S}$ & 1.06$^{\S}$ & 1.06$^{\S}$ & 1.239\,$\pm$\,0.005 & 1.1930\,$\pm$\,0.0281 & 1.1001\,$\pm$\,0.0170 \\
$X\to YY\to gge$ & 0.99346\,$\pm$\,0.00017 & 0.9950\,$\pm$\,0.0286 & 1.553\,$\pm$\,0.067 & 1.487\,$\pm$\,0.133 & 1.08$^{\S}$ & 1.09$^{\S}$ & 1.12$^{\S}$ & 1.03$^{\S}$ & 1.13$^{\S}$ & 1.136\,$\pm$\,0.021 & 1.0758\,$\pm$\,0.0451 & 1.0955\,$\pm$\,0.0183 \\
$X\to YY\to gg\mu$ & 0.99356\,$\pm$\,0.00009 & 1.0027\,$\pm$\,0.0156 & 1.673\,$\pm$\,0.026 & 1.455\,$\pm$\,0.175 & 1.07$^{\S}$ & 1.02$^{\S}$ & 1.10$^{\S}$ & 1.06$^{\S}$ & 1.09$^{\S}$ & 1.082\,$\pm$\,0.001 & 1.0370\,$\pm$\,0.0086 & 1.0366\,$\pm$\,0.0188 \\
\midrule
$X\to YY\to bee$ & 0.99813\,$\pm$\,0.00004 & 0.9923\,$\pm$\,0.0203 & 3.073\,$\pm$\,0.111 & 3.479\,$\pm$\,0.175 & 1.14$^{\S}$ & 1.12$^{\S}$ & 1.10$^{\S}$ & 1.09$^{\S}$ & 1.10$^{\S}$ & 1.490\,$\pm$\,0.016 & 1.2415\,$\pm$\,0.0114 & 1.3149\,$\pm$\,0.0291 \\
$X\to YY\to cee$ & 0.99542\,$\pm$\,0.00001 & 0.9990\,$\pm$\,0.0150 & 2.434\,$\pm$\,0.111 & 3.030\,$\pm$\,0.125 & 1.09$^{\S}$ & 1.09$^{\S}$ & 1.06$^{\S}$ & 1.07$^{\S}$ & 1.09$^{\S}$ & 1.454\,$\pm$\,0.002 & 1.2287\,$\pm$\,0.0224 & 1.2407\,$\pm$\,0.0042 \\
$X\to YY\to see$ & 0.99432\,$\pm$\,0.00002 & 0.9983\,$\pm$\,0.0089 & 2.290\,$\pm$\,0.079 & 2.553\,$\pm$\,0.086 & 1.04$^{\S}$ & 1.02$^{\S}$ & 1.02$^{\S}$ & 1.05$^{\S}$ & 1.04$^{\S}$ & 1.218\,$\pm$\,0.008 & 1.1570\,$\pm$\,0.0109 & 1.0861\,$\pm$\,0.0116 \\
$X\to YY\to qee$ & 0.99418\,$\pm$\,0.00002 & 0.9963\,$\pm$\,0.0036 & 2.300\,$\pm$\,0.088 & 2.411\,$\pm$\,0.206 & 1.03$^{\S}$ & 1.01$^{\S}$ & 1.04$^{\S}$ & 1.03$^{\S}$ & 1.04$^{\S}$ & 1.135\,$\pm$\,0.007 & 1.1374\,$\pm$\,0.0109 & 1.0815\,$\pm$\,0.0053 \\
$X\to YY\to b\mu\mu$ & 0.99825\,$\pm$\,0.00002 & 0.9968\,$\pm$\,0.0114 & 3.130\,$\pm$\,0.100 & 3.736\,$\pm$\,0.355 & 1.10$^{\S}$ & 1.07$^{\S}$ & 1.07$^{\S}$ & 1.10$^{\S}$ & 1.10$^{\S}$ & 1.555\,$\pm$\,0.010 & 1.2454\,$\pm$\,0.0304 & 1.3192\,$\pm$\,0.0323 \\
$X\to YY\to c\mu\mu$ & 0.99518\,$\pm$\,0.00003 & 1.0039\,$\pm$\,0.0150 & 2.326\,$\pm$\,0.050 & 2.874\,$\pm$\,0.199 & 1.05$^{\S}$ & 1.07$^{\S}$ & 1.04$^{\S}$ & 1.07$^{\S}$ & 1.09$^{\S}$ & 1.388\,$\pm$\,0.030 & 1.2031\,$\pm$\,0.0239 & 1.2319\,$\pm$\,0.0173 \\
$X\to YY\to s\mu\mu$ & 0.99430\,$\pm$\,0.00003 & 1.0030\,$\pm$\,0.0044 & 2.276\,$\pm$\,0.100 & 2.546\,$\pm$\,0.131 & 1.04$^{\S}$ & 1.03$^{\S}$ & 1.03$^{\S}$ & 1.05$^{\S}$ & 1.04$^{\S}$ & 1.233\,$\pm$\,0.005 & 1.2065\,$\pm$\,0.0109 & 1.0999\,$\pm$\,0.0142 \\
$X\to YY\to q\mu\mu$ & 0.99398\,$\pm$\,0.00003 & 0.9952\,$\pm$\,0.0047 & 2.253\,$\pm$\,0.071 & 2.343\,$\pm$\,0.093 & 1.04$^{\S}$ & 1.01$^{\S}$ & 1.02$^{\S}$ & 1.04$^{\S}$ & 1.03$^{\S}$ & 1.146\,$\pm$\,0.011 & 1.1395\,$\pm$\,0.0085 & 1.0795\,$\pm$\,0.0048 \\
$X\to YY\to gee$ & 0.99418\,$\pm$\,0.00002 & 1.0047\,$\pm$\,0.0063 & 2.282\,$\pm$\,0.095 & 2.286\,$\pm$\,0.126 & 1.03$^{\S}$ & 1.01$^{\S}$ & 1.01$^{\S}$ & 1.05$^{\S}$ & 1.02$^{\S}$ & 1.056\,$\pm$\,0.009 & 1.0233\,$\pm$\,0.0068 & 1.0342\,$\pm$\,0.0151 \\
$X\to YY\to g\mu\mu$ & 0.99410\,$\pm$\,0.00007 & 1.0108\,$\pm$\,0.0293 & 2.229\,$\pm$\,0.107 & 2.268\,$\pm$\,0.116 & 1.03$^{\S}$ & 1.04$^{\S}$ & 1.02$^{\S}$ & 1.02$^{\S}$ & 1.03$^{\S}$ & 1.055\,$\pm$\,0.006 & 1.0288\,$\pm$\,0.0211 & 1.0288\,$\pm$\,0.0076 \\
\midrule
$X\to YY\to b\tau_h\tau_e$ & 0.99596\,$\pm$\,0.00005 & 0.9879\,$\pm$\,0.0158 & 2.128\,$\pm$\,0.069 & 2.460\,$\pm$\,0.137 & 1.11$^{\S}$ & 1.18$^{\S}$ & 1.15$^{\S}$ & 1.16$^{\S}$ & 1.18$^{\S}$ & 1.600\,$\pm$\,0.055 & 1.2159\,$\pm$\,0.0292 & 1.2948\,$\pm$\,0.0393 \\
$X\to YY\to c\tau_h\tau_e$ & 0.99158\,$\pm$\,0.00004 & 0.9997\,$\pm$\,0.0203 & 1.939\,$\pm$\,0.050 & 2.214\,$\pm$\,0.125 & 1.12$^{\S}$ & 1.13$^{\S}$ & 1.09$^{\S}$ & 1.11$^{\S}$ & 1.12$^{\S}$ & 1.355\,$\pm$\,0.009 & 1.1953\,$\pm$\,0.0329 & 1.2045\,$\pm$\,0.0050 \\
$X\to YY\to s\tau_h\tau_e$ & 0.99157\,$\pm$\,0.00007 & 1.0022\,$\pm$\,0.0137 & 2.045\,$\pm$\,0.082 & 2.253\,$\pm$\,0.180 & 1.09$^{\S}$ & 1.09$^{\S}$ & 1.08$^{\S}$ & 1.08$^{\S}$ & 1.09$^{\S}$ & 1.347\,$\pm$\,0.012 & 1.2514\,$\pm$\,0.0319 & 1.1513\,$\pm$\,0.0066 \\
$X\to YY\to q\tau_h\tau_e$ & 0.99073\,$\pm$\,0.00007 & 1.0240\,$\pm$\,0.0169 & 2.060\,$\pm$\,0.060 & 2.042\,$\pm$\,0.198 & 1.07$^{\S}$ & 1.09$^{\S}$ & 1.06$^{\S}$ & 1.07$^{\S}$ & 1.11$^{\S}$ & 1.200\,$\pm$\,0.023 & 1.1832\,$\pm$\,0.0178 & 1.1285\,$\pm$\,0.0028 \\
$X\to YY\to b\tau_h\tau_\mu$ & 0.99649\,$\pm$\,0.00005 & 1.0171\,$\pm$\,0.0345 & 2.286\,$\pm$\,0.032 & 2.596\,$\pm$\,0.090 & 1.15$^{\S}$ & 1.18$^{\S}$ & 1.20$^{\S}$ & 1.13$^{\S}$ & 1.12$^{\S}$ & 1.535\,$\pm$\,0.051 & 1.1911\,$\pm$\,0.0041 & 1.3216\,$\pm$\,0.0018 \\
$X\to YY\to c\tau_h\tau_\mu$ & 0.99211\,$\pm$\,0.00008 & 1.0097\,$\pm$\,0.0016 & 1.985\,$\pm$\,0.036 & 2.292\,$\pm$\,0.185 & 1.07$^{\S}$ & 1.12$^{\S}$ & 1.06$^{\S}$ & 1.10$^{\S}$ & 1.11$^{\S}$ & 1.360\,$\pm$\,0.049 & 1.2056\,$\pm$\,0.0362 & 1.2418\,$\pm$\,0.0109 \\
$X\to YY\to s\tau_h\tau_\mu$ & 0.99177\,$\pm$\,0.00011 & 0.9953\,$\pm$\,0.0212 & 2.089\,$\pm$\,0.117 & 2.230\,$\pm$\,0.175 & 1.08$^{\S}$ & 1.09$^{\S}$ & 1.06$^{\S}$ & 1.08$^{\S}$ & 1.09$^{\S}$ & 1.268\,$\pm$\,0.029 & 1.2049\,$\pm$\,0.0302 & 1.1465\,$\pm$\,0.0285 \\
$X\to YY\to q\tau_h\tau_\mu$ & 0.99145\,$\pm$\,0.00006 & 1.0010\,$\pm$\,0.0105 & 2.050\,$\pm$\,0.078 & 2.055\,$\pm$\,0.228 & 1.06$^{\S}$ & 1.08$^{\S}$ & 1.06$^{\S}$ & 1.06$^{\S}$ & 1.06$^{\S}$ & 1.174\,$\pm$\,0.003 & 1.1519\,$\pm$\,0.0265 & 1.1271\,$\pm$\,0.0030 \\
$X\to YY\to g\tau_h\tau_e$ & 0.99048\,$\pm$\,0.00005 & 1.0068\,$\pm$\,0.0045 & 1.937\,$\pm$\,0.072 & 1.885\,$\pm$\,0.128 & 1.07$^{\S}$ & 1.08$^{\S}$ & 1.07$^{\S}$ & 1.06$^{\S}$ & 1.11$^{\S}$ & 1.079\,$\pm$\,0.035 & 1.0400\,$\pm$\,0.0034 & 1.0397\,$\pm$\,0.0277 \\
$X\to YY\to g\tau_h\tau_\mu$ & 0.99068\,$\pm$\,0.00010 & 1.0166\,$\pm$\,0.0223 & 1.945\,$\pm$\,0.055 & 1.879\,$\pm$\,0.143 & 1.06$^{\S}$ & 1.09$^{\S}$ & 1.06$^{\S}$ & 1.06$^{\S}$ & 1.11$^{\S}$ & 1.077\,$\pm$\,0.016 & 1.0230\,$\pm$\,0.0177 & 1.0562\,$\pm$\,0.0215 \\
\midrule
$X\to YY\to bbbb$ & 0.99642\,$\pm$\,0.00016 & 1.0201\,$\pm$\,0.0832 & 1.210\,$\pm$\,0.022 & 2.119\,$\pm$\,0.147 & 1.44$^{\S}$ & 1.35$^{\S}$ & 1.34$^{\S}$ & 1.31$^{\S}$ & 1.25$^{\S}$ & 2.433\,$\pm$\,0.016 & 1.5672\,$\pm$\,0.0344 & 1.7151\,$\pm$\,0.0245 \\
$X\to YY\to bbcc$ & 0.98460\,$\pm$\,0.00004 & 0.9969\,$\pm$\,0.0166 & 1.154\,$\pm$\,0.001 & 1.768\,$\pm$\,0.037 & 1.20$^{\S}$ & 1.20$^{\S}$ & 1.14$^{\S}$ & 1.20$^{\S}$ & 1.18$^{\S}$ & 1.836\,$\pm$\,0.083 & 1.3074\,$\pm$\,0.0041 & 1.2630\,$\pm$\,0.0097 \\
$X\to YY\to bbss$ & 0.98147\,$\pm$\,0.00004 & 1.0080\,$\pm$\,0.0114 & 1.256\,$\pm$\,0.011 & 1.975\,$\pm$\,0.056 & 1.18$^{\S}$ & 1.19$^{\S}$ & 1.18$^{\S}$ & 1.16$^{\S}$ & 1.18$^{\S}$ & 2.012\,$\pm$\,0.122 & 1.3989\,$\pm$\,0.0365 & 1.2709\,$\pm$\,0.0173 \\
$X\to YY\to bbqq$ & 0.97696\,$\pm$\,0.00018 & 0.9973\,$\pm$\,0.0169 & 1.157\,$\pm$\,0.024 & 1.713\,$\pm$\,0.020 & 1.17$^{\S}$ & 1.14$^{\S}$ & 1.14$^{\S}$ & 1.13$^{\S}$ & 1.12$^{\S}$ & 1.814\,$\pm$\,0.034 & 1.3223\,$\pm$\,0.0251 & 1.2322\,$\pm$\,0.0149 \\
$X\to YY\to cccc$ & 0.98510\,$\pm$\,0.00017 & 1.0057\,$\pm$\,0.0323 & 1.098\,$\pm$\,0.008 & 2.315\,$\pm$\,0.154 & 1.30$^{\S}$ & 1.29$^{\S}$ & 1.22$^{\S}$ & 1.20$^{\S}$ & 1.26$^{\S}$ & 2.409\,$\pm$\,0.133 & 1.3700\,$\pm$\,0.0232 & 1.4831\,$\pm$\,0.0260 \\
$X\to YY\to ccss$ & 0.96566\,$\pm$\,0.00017 & 0.9990\,$\pm$\,0.0070 & 1.119\,$\pm$\,0.011 & 1.622\,$\pm$\,0.057 & 1.11$^{\S}$ & 1.11$^{\S}$ & 1.12$^{\S}$ & 1.12$^{\S}$ & 1.12$^{\S}$ & 1.708\,$\pm$\,0.035 & 1.2491\,$\pm$\,0.0085 & 1.1751\,$\pm$\,0.0163 \\
$X\to YY\to ccqq$ & 0.96128\,$\pm$\,0.00005 & 0.9971\,$\pm$\,0.0050 & 1.087\,$\pm$\,0.008 & 1.541\,$\pm$\,0.024 & 1.10$^{\S}$ & 1.12$^{\S}$ & 1.08$^{\S}$ & 1.09$^{\S}$ & 1.10$^{\S}$ & 1.643\,$\pm$\,0.110 & 1.2082\,$\pm$\,0.0069 & 1.1848\,$\pm$\,0.0083 \\
$X\to YY\to ssss$ & 0.98005\,$\pm$\,0.00006 & 0.9984\,$\pm$\,0.0070 & 1.140\,$\pm$\,0.020 & 1.720\,$\pm$\,0.042 & 1.09$^{\S}$ & 1.10$^{\S}$ & 1.05$^{\S}$ & 1.06$^{\S}$ & 1.07$^{\S}$ & 1.629\,$\pm$\,0.044 & 1.3180\,$\pm$\,0.0279 & 1.1482\,$\pm$\,0.0122 \\
$X\to YY\to ssqq$ & 0.96823\,$\pm$\,0.00008 & 1.0022\,$\pm$\,0.0094 & 1.051\,$\pm$\,0.006 & 1.373\,$\pm$\,0.030 & 1.06$^{\S}$ & 1.07$^{\S}$ & 1.04$^{\S}$ & 1.04$^{\S}$ & 1.06$^{\S}$ & 1.321\,$\pm$\,0.001 & 1.1544\,$\pm$\,0.0024 & 1.0796\,$\pm$\,0.0078 \\
$X\to YY\to qqqq$ & 0.97547\,$\pm$\,0.00017 & 0.9996\,$\pm$\,0.0142 & 1.096\,$\pm$\,0.011 & 1.468\,$\pm$\,0.037 & 1.08$^{\S}$ & 1.06$^{\S}$ & 1.05$^{\S}$ & 1.05$^{\S}$ & 1.04$^{\S}$ & 1.343\,$\pm$\,0.042 & 1.2518\,$\pm$\,0.0059 & 1.0858\,$\pm$\,0.0126 \\
$X\to YY\to qqqb$ & 0.97637\,$\pm$\,0.00029 & 1.0083\,$\pm$\,0.0034 & 1.143\,$\pm$\,0.011 & 1.748\,$\pm$\,0.059 & 1.18$^{\S}$ & 1.17$^{\S}$ & 1.14$^{\S}$ & 1.16$^{\S}$ & 1.21$^{\S}$ & 1.712\,$\pm$\,0.004 & 1.3748\,$\pm$\,0.0086 & 1.3271\,$\pm$\,0.0120 \\
$X\to YY\to qqqc$ & 0.96096\,$\pm$\,0.00027 & 1.0053\,$\pm$\,0.0181 & 1.130\,$\pm$\,0.005 & 1.458\,$\pm$\,0.007 & 1.12$^{\S}$ & 1.11$^{\S}$ & 1.10$^{\S}$ & 1.09$^{\S}$ & 1.11$^{\S}$ & 1.476\,$\pm$\,0.028 & 1.2355\,$\pm$\,0.0056 & 1.2236\,$\pm$\,0.0193 \\
$X\to YY\to qqqs$ & 0.97027\,$\pm$\,0.00017 & 1.0007\,$\pm$\,0.0115 & 1.046\,$\pm$\,0.010 & 1.362\,$\pm$\,0.003 & 1.07$^{\S}$ & 1.06$^{\S}$ & 1.05$^{\S}$ & 1.06$^{\S}$ & 1.05$^{\S}$ & 1.296\,$\pm$\,0.020 & 1.1618\,$\pm$\,0.0180 & 1.0696\,$\pm$\,0.0007 \\
$X\to YY\to bbcq$ & 0.97785\,$\pm$\,0.00026 & 0.9990\,$\pm$\,0.0092 & 1.176\,$\pm$\,0.014 & 1.645\,$\pm$\,0.050 & 1.18$^{\S}$ & 1.22$^{\S}$ & 1.17$^{\S}$ & 1.12$^{\S}$ & 1.16$^{\S}$ & 1.785\,$\pm$\,0.120 & 1.3099\,$\pm$\,0.0406 & 1.2342\,$\pm$\,0.0148 \\
$X\to YY\to ccbs$ & 0.96956\,$\pm$\,0.00020 & 0.9950\,$\pm$\,0.0156 & 1.157\,$\pm$\,0.013 & 1.636\,$\pm$\,0.019 & 1.16$^{\S}$ & 1.17$^{\S}$ & 1.12$^{\S}$ & 1.15$^{\S}$ & 1.17$^{\S}$ & 1.618\,$\pm$\,0.024 & 1.3261\,$\pm$\,0.0020 & 1.2421\,$\pm$\,0.0060 \\
$X\to YY\to ccbq$ & 0.96912\,$\pm$\,0.00044 & 1.0155\,$\pm$\,0.0308 & 1.151\,$\pm$\,0.018 & 1.628\,$\pm$\,0.009 & 1.22$^{\S}$ & 1.15$^{\S}$ & 1.12$^{\S}$ & 1.13$^{\S}$ & 1.15$^{\S}$ & 1.617\,$\pm$\,0.001 & 1.2801\,$\pm$\,0.0411 & 1.2268\,$\pm$\,0.0278 \\
$X\to YY\to ccsq$ & 0.95974\,$\pm$\,0.00024 & 1.0044\,$\pm$\,0.0103 & 1.079\,$\pm$\,0.015 & 1.514\,$\pm$\,0.013 & 1.14$^{\S}$ & 1.12$^{\S}$ & 1.12$^{\S}$ & 1.13$^{\S}$ & 1.12$^{\S}$ & 1.562\,$\pm$\,0.070 & 1.1659\,$\pm$\,0.0149 & 1.1837\,$\pm$\,0.0017 \\
$X\to YY\to sscq$ & 0.95891\,$\pm$\,0.00018 & 1.0052\,$\pm$\,0.0081 & 1.117\,$\pm$\,0.004 & 1.442\,$\pm$\,0.038 & 1.12$^{\S}$ & 1.12$^{\S}$ & 1.10$^{\S}$ & 1.11$^{\S}$ & 1.11$^{\S}$ & 1.482\,$\pm$\,0.009 & 1.1985\,$\pm$\,0.0003 & 1.1731\,$\pm$\,0.0106 \\
$X\to YY\to qqbc$ & 0.96422\,$\pm$\,0.00024 & 1.0062\,$\pm$\,0.0102 & 1.154\,$\pm$\,0.002 & 1.518\,$\pm$\,0.004 & 1.17$^{\S}$ & 1.18$^{\S}$ & 1.13$^{\S}$ & 1.15$^{\S}$ & 1.16$^{\S}$ & 1.495\,$\pm$\,0.013 & 1.2365\,$\pm$\,0.0244 & 1.2140\,$\pm$\,0.0002 \\
$X\to YY\to qqbs$ & 0.97253\,$\pm$\,0.00031 & 1.0160\,$\pm$\,0.0186 & 1.120\,$\pm$\,0.011 & 1.695\,$\pm$\,0.053 & 1.21$^{\S}$ & 1.18$^{\S}$ & 1.17$^{\S}$ & 1.16$^{\S}$ & 1.20$^{\S}$ & 1.646\,$\pm$\,0.011 & 1.2957\,$\pm$\,0.0260 & 1.3190\,$\pm$\,0.0003 \\
$X\to YY\to qqcs$ & 0.95420\,$\pm$\,0.00005 & 1.0047\,$\pm$\,0.0131 & 1.058\,$\pm$\,0.012 & 1.309\,$\pm$\,0.012 & 1.10$^{\S}$ & 1.11$^{\S}$ & 1.08$^{\S}$ & 1.08$^{\S}$ & 1.09$^{\S}$ & 1.323\,$\pm$\,0.033 & 1.1096\,$\pm$\,0.0099 & 1.1649\,$\pm$\,0.0067 \\
$X\to YY\to bcsq$ & 0.96323\,$\pm$\,0.00027 & 1.0084\,$\pm$\,0.0164 & 1.132\,$\pm$\,0.013 & 1.462\,$\pm$\,0.007 & 1.19$^{\S}$ & 1.16$^{\S}$ & 1.15$^{\S}$ & 1.14$^{\S}$ & 1.17$^{\S}$ & 1.449\,$\pm$\,0.008 & 1.2348\,$\pm$\,0.0066 & 1.2208\,$\pm$\,0.0180 \\
$X\to YY\to bbgg$ & 0.97819\,$\pm$\,0.00019 & 1.0025\,$\pm$\,0.0077 & 1.125\,$\pm$\,0.016 & 1.512\,$\pm$\,0.087 & 1.08$^{\S}$ & 1.12$^{\S}$ & 1.09$^{\S}$ & 1.08$^{\S}$ & 1.11$^{\S}$ & 1.445\,$\pm$\,0.118 & 1.1041\,$\pm$\,0.0105 & 1.1857\,$\pm$\,0.0078 \\
$X\to YY\to ccgg$ & 0.96410\,$\pm$\,0.00028 & 1.0052\,$\pm$\,0.0195 & 1.091\,$\pm$\,0.009 & 1.555\,$\pm$\,0.042 & 1.08$^{\S}$ & 1.08$^{\S}$ & 1.07$^{\S}$ & 1.06$^{\S}$ & 1.07$^{\S}$ & 1.460\,$\pm$\,0.055 & 1.1502\,$\pm$\,0.0233 & 1.1773\,$\pm$\,0.0004 \\
$X\to YY\to ssgg$ & 0.96380\,$\pm$\,0.00012 & 1.0010\,$\pm$\,0.0056 & 1.088\,$\pm$\,0.007 & 1.337\,$\pm$\,0.004 & 1.05$^{\S}$ & 1.06$^{\S}$ & 1.03$^{\S}$ & 1.03$^{\S}$ & 1.05$^{\S}$ & 1.264\,$\pm$\,0.014 & 1.1284\,$\pm$\,0.0038 & 1.0825\,$\pm$\,0.0102 \\
$X\to YY\to qqgg$ & 0.96228\,$\pm$\,0.00016 & 0.9981\,$\pm$\,0.0028 & 1.069\,$\pm$\,0.002 & 1.290\,$\pm$\,0.034 & 1.04$^{\S}$ & 1.05$^{\S}$ & 1.02$^{\S}$ & 1.04$^{\S}$ & 1.03$^{\S}$ & 1.153\,$\pm$\,0.001 & 1.1203\,$\pm$\,0.0086 & 1.0722\,$\pm$\,0.0032 \\
$X\to YY\to gggg$ & 0.97374\,$\pm$\,0.00009 & 0.9984\,$\pm$\,0.0130 & 1.058\,$\pm$\,0.011 & 1.188\,$\pm$\,0.004 & 1.05$^{\S}$ & 1.04$^{\S}$ & 1.03$^{\S}$ & 1.03$^{\S}$ & 1.05$^{\S}$ & 1.088\,$\pm$\,0.011 & 1.0255\,$\pm$\,0.0080 & 1.0356\,$\pm$\,0.0010 \\
\midrule
$X\to YY\to bb\tau_h\tau_h$ & 0.99705\,$\pm$\,0.00018 & 1.0089\,$\pm$\,0.0818 & 1.039\,$\pm$\,0.097 & 1.861\,$\pm$\,0.127 & 1.36$^{\S}$ & 1.15$^{\S}$ & 1.46$^{\S}$ & 1.09$^{\S}$ & 1.24$^{\S}$ & 2.183\,$\pm$\,0.088 & 1.3495\,$\pm$\,0.0144 & 1.5533\,$\pm$\,0.1304 \\
$X\to YY\to cc\tau_h\tau_h$ & 0.99340\,$\pm$\,0.00010 & 0.9894\,$\pm$\,0.0123 & 1.022\,$\pm$\,0.024 & 1.666\,$\pm$\,0.046 & 1.11$^{\S}$ & 1.17$^{\S}$ & 1.12$^{\S}$ & 1.14$^{\S}$ & 1.17$^{\S}$ & 1.696\,$\pm$\,0.078 & 1.3333\,$\pm$\,0.0439 & 1.3810\,$\pm$\,0.0424 \\
$X\to YY\to ss\tau_h\tau_h$ & 0.99140\,$\pm$\,0.00003 & 1.0179\,$\pm$\,0.0118 & 1.099\,$\pm$\,0.015 & 1.504\,$\pm$\,0.028 & 1.08$^{\S}$ & 1.08$^{\S}$ & 1.10$^{\S}$ & 1.02$^{\S}$ & 1.09$^{\S}$ & 1.433$^{\P}$ & 1.3022\,$\pm$\,0.0033 & 1.1620\,$\pm$\,0.0115 \\
$X\to YY\to qq\tau_h\tau_h$ & 0.99133\,$\pm$\,0.00007 & 1.0084\,$\pm$\,0.0027 & 1.077\,$\pm$\,0.017 & 1.341\,$\pm$\,0.039 & 1.07$^{\S}$ & 1.08$^{\S}$ & 1.06$^{\S}$ & 1.08$^{\S}$ & 1.10$^{\S}$ & 1.297\,$\pm$\,0.018 & 1.2392\,$\pm$\,0.0327 & 1.1570\,$\pm$\,0.0090 \\
$X\to YY\to gg\tau_h\tau_h$ & 0.98959\,$\pm$\,0.00021 & 0.9911\,$\pm$\,0.0309 & 1.029\,$\pm$\,0.011 & 1.203\,$\pm$\,0.020 & 1.05$^{\S}$ & 1.05$^{\S}$ & 1.08$^{\S}$ & 1.03$^{\S}$ & 1.04$^{\S}$ & 1.096\,$\pm$\,0.002 & 1.0533\,$\pm$\,0.0043 & 1.0494\,$\pm$\,0.0167 \\
\midrule
$X\to YY\to bbee$ & 0.99913\,$\pm$\,0.00003 & 0.9722\,$\pm$\,0.0184 & 4.010\,$\pm$\,0.327 & 5.188\,$\pm$\,0.386 & 1.37$^{\S}$ & 1.14$^{\S}$ & 1.34$^{\S}$ & 1.18$^{\S}$ & 1.71$^{\S}$ & 2.305\,$\pm$\,0.122 & 1.4508\,$\pm$\,0.1845 & 1.6000\,$\pm$\,0.0982 \\
$X\to YY\to ccee$ & 0.99763$^{\P}$ & 0.9922\,$\pm$\,0.0192 & 2.843\,$\pm$\,0.055 & 4.164\,$\pm$\,0.304 & 1.12$^{\S}$ & 1.12$^{\S}$ & 1.08$^{\S}$ & 1.14$^{\S}$ & 1.14$^{\S}$ & 1.845\,$\pm$\,0.101 & 1.3887\,$\pm$\,0.0066 & 1.3607\,$\pm$\,0.0359 \\
$X\to YY\to ssee$ & 0.99613\,$\pm$\,0.00001 & 0.9940\,$\pm$\,0.0092 & 2.581\,$\pm$\,0.071 & 3.107\,$\pm$\,0.199 & 1.04$^{\S}$ & 1.03$^{\S}$ & 1.03$^{\S}$ & 1.05$^{\S}$ & 1.04$^{\S}$ & 1.484\,$\pm$\,0.031 & 1.3260\,$\pm$\,0.0146 & 1.1506\,$\pm$\,0.0488 \\
$X\to YY\to qqee$ & 0.99596\,$\pm$\,0.00002 & 1.0028\,$\pm$\,0.0083 & 2.522\,$\pm$\,0.100 & 2.856\,$\pm$\,0.103 & 1.07$^{\S}$ & 1.04$^{\S}$ & 1.03$^{\S}$ & 1.08$^{\S}$ & 1.04$^{\S}$ & 1.265\,$\pm$\,0.043 & 1.2701\,$\pm$\,0.0796 & 1.1096\,$\pm$\,0.0173 \\
$X\to YY\to bb\mu\mu$ & 0.99906$^{\P}$ & 0.9802\,$\pm$\,0.0108 & 4.065\,$\pm$\,0.142 & 4.788\,$\pm$\,0.384 & 1.15$^{\S}$ & 1.15$^{\S}$ & 1.19$^{\S}$ & 1.07$^{\S}$ & 1.26$^{\S}$ & 2.168\,$\pm$\,0.231 & 1.3648\,$\pm$\,0.0057 & 1.6064\,$\pm$\,0.1631 \\
$X\to YY\to cc\mu\mu$ & 0.99722\,$\pm$\,0.00006 & 1.0018\,$\pm$\,0.0383 & 2.746\,$\pm$\,0.057 & 3.962\,$\pm$\,0.394 & 1.08$^{\S}$ & 1.13$^{\S}$ & 1.07$^{\S}$ & 1.12$^{\S}$ & 1.13$^{\S}$ & 1.805\,$\pm$\,0.001 & 1.3311\,$\pm$\,0.0268 & 1.3818\,$\pm$\,0.0318 \\
$X\to YY\to ss\mu\mu$ & 0.99565\,$\pm$\,0.00004 & 1.0047\,$\pm$\,0.0181 & 2.642\,$\pm$\,0.084 & 3.051\,$\pm$\,0.230 & 1.05$^{\S}$ & 1.05$^{\S}$ & 1.04$^{\S}$ & 1.04$^{\S}$ & 1.06$^{\S}$ & 1.448\,$\pm$\,0.011 & 1.2925\,$\pm$\,0.0026 & 1.1253\,$\pm$\,0.0078 \\
$X\to YY\to qq\mu\mu$ & 0.99560\,$\pm$\,0.00003 & 1.0004\,$\pm$\,0.0012 & 2.611\,$\pm$\,0.042 & 2.740\,$\pm$\,0.201 & 1.07$^{\S}$ & 1.05$^{\S}$ & 1.05$^{\S}$ & 1.05$^{\S}$ & 1.06$^{\S}$ & 1.286\,$\pm$\,0.001 & 1.2799\,$\pm$\,0.0369 & 1.1561\,$\pm$\,0.0203 \\
$X\to YY\to ggee$ & 0.99610\,$\pm$\,0.00002 & 1.0029\,$\pm$\,0.0101 & 2.473\,$\pm$\,0.081 & 2.474\,$\pm$\,0.229 & 1.03$^{\S}$ & 1.03$^{\S}$ & 1.02$^{\S}$ & 1.03$^{\S}$ & 1.03$^{\S}$ & 1.060\,$\pm$\,0.010 & 1.0357\,$\pm$\,0.0218 & 1.0331\,$\pm$\,0.0149 \\
$X\to YY\to gg\mu\mu$ & 0.99550\,$\pm$\,0.00004 & 1.0060\,$\pm$\,0.0041 & 2.609\,$\pm$\,0.104 & 2.572\,$\pm$\,0.273 & 1.06$^{\S}$ & 1.07$^{\S}$ & 1.05$^{\S}$ & 1.05$^{\S}$ & 1.04$^{\S}$ & 1.076\,$\pm$\,0.009 & 1.0664\,$\pm$\,0.0133 & 1.0536\,$\pm$\,0.0204 \\
\midrule
$X\to YY\to bb\tau_h\tau_e$ & 0.99734\,$\pm$\,0.00003 & 0.9957\,$\pm$\,0.0381 & 2.394\,$\pm$\,0.162 & 3.132\,$\pm$\,0.095 & 1.22$^{\S}$ & 1.22$^{\S}$ & 1.29$^{\S}$ & 1.16$^{\S}$ & 1.27$^{\S}$ & 2.170\,$\pm$\,0.228 & 1.3434\,$\pm$\,0.0496 & 1.6055\,$\pm$\,0.0851 \\
$X\to YY\to cc\tau_h\tau_e$ & 0.99316\,$\pm$\,0.00008 & 1.0042\,$\pm$\,0.0071 & 1.968\,$\pm$\,0.073 & 2.895\,$\pm$\,0.050 & 1.20$^{\S}$ & 1.17$^{\S}$ & 1.13$^{\S}$ & 1.11$^{\S}$ & 1.21$^{\S}$ & 1.782\,$\pm$\,0.070 & 1.4106\,$\pm$\,0.0580 & 1.4221\,$\pm$\,0.0054 \\
$X\to YY\to ss\tau_h\tau_e$ & 0.99090\,$\pm$\,0.00007 & 1.0047\,$\pm$\,0.0110 & 2.086\,$\pm$\,0.092 & 2.586\,$\pm$\,0.121 & 1.07$^{\S}$ & 1.08$^{\S}$ & 1.11$^{\S}$ & 1.08$^{\S}$ & 1.08$^{\S}$ & 1.479\,$\pm$\,0.027 & 1.3327\,$\pm$\,0.0397 & 1.1992\,$\pm$\,0.0331 \\
$X\to YY\to qq\tau_h\tau_e$ & 0.98992\,$\pm$\,0.00012 & 0.9910\,$\pm$\,0.0120 & 2.034\,$\pm$\,0.054 & 2.268\,$\pm$\,0.171 & 1.08$^{\S}$ & 1.10$^{\S}$ & 1.10$^{\S}$ & 1.07$^{\S}$ & 1.06$^{\S}$ & 1.335\,$\pm$\,0.006 & 1.2966\,$\pm$\,0.0182 & 1.1731\,$\pm$\,0.0301 \\
$X\to YY\to bb\tau_h\tau_\mu$ & 0.99737\,$\pm$\,0.00008 & 0.9912\,$\pm$\,0.0451 & 2.439\,$\pm$\,0.013 & 3.165\,$\pm$\,0.158 & 1.14$^{\S}$ & 1.25$^{\S}$ & 1.16$^{\S}$ & 1.17$^{\S}$ & 1.15$^{\S}$ & 2.027\,$\pm$\,0.375 & 1.3037\,$\pm$\,0.0984 & 1.4952\,$\pm$\,0.0918 \\
$X\to YY\to cc\tau_h\tau_\mu$ & 0.99428\,$\pm$\,0.00009 & 1.0123\,$\pm$\,0.0270 & 2.160\,$\pm$\,0.080 & 3.190\,$\pm$\,0.223 & 1.17$^{\S}$ & 1.18$^{\S}$ & 1.10$^{\S}$ & 1.15$^{\S}$ & 1.23$^{\S}$ & 1.815\,$\pm$\,0.120 & 1.4233\,$\pm$\,0.0260 & 1.4661\,$\pm$\,0.0337 \\
$X\to YY\to ss\tau_h\tau_\mu$ & 0.99190\,$\pm$\,0.00004 & 1.0102\,$\pm$\,0.0059 & 2.288\,$\pm$\,0.058 & 2.687\,$\pm$\,0.183 & 1.11$^{\S}$ & 1.10$^{\S}$ & 1.07$^{\S}$ & 1.09$^{\S}$ & 1.06$^{\S}$ & 1.449\,$\pm$\,0.019 & 1.3118\,$\pm$\,0.0170 & 1.1694\,$\pm$\,0.0158 \\
$X\to YY\to qq\tau_h\tau_\mu$ & 0.99106\,$\pm$\,0.00006 & 1.0052\,$\pm$\,0.0162 & 2.162\,$\pm$\,0.070 & 2.310\,$\pm$\,0.192 & 1.08$^{\S}$ & 1.10$^{\S}$ & 1.11$^{\S}$ & 1.09$^{\S}$ & 1.09$^{\S}$ & 1.289\,$\pm$\,0.024 & 1.2432\,$\pm$\,0.0243 & 1.1336\,$\pm$\,0.0353 \\
$X\to YY\to gg\tau_h\tau_e$ & 0.98906\,$\pm$\,0.00012 & 0.9960\,$\pm$\,0.0102 & 1.926\,$\pm$\,0.077 & 1.950\,$\pm$\,0.151 & 1.05$^{\S}$ & 1.12$^{\S}$ & 1.04$^{\S}$ & 1.07$^{\S}$ & 1.09$^{\S}$ & 1.099\,$\pm$\,0.032 & 1.0457\,$\pm$\,0.0172 & 1.0559\,$\pm$\,0.0177 \\
$X\to YY\to gg\tau_h\tau_\mu$ & 0.99079\,$\pm$\,0.00004 & 1.0018\,$\pm$\,0.0129 & 2.137\,$\pm$\,0.083 & 2.119\,$\pm$\,0.199 & 1.03$^{\S}$ & 1.12$^{\S}$ & 1.09$^{\S}$ & 1.09$^{\S}$ & 1.08$^{\S}$ & 1.136\,$\pm$\,0.020 & 1.0740\,$\pm$\,0.0100 & 1.0766\,$\pm$\,0.0241 \\
\midrule
$X\to YY\to bce\nu$ & 0.99457\,$\pm$\,0.00007 & 1.0075\,$\pm$\,0.0136 & 1.773\,$\pm$\,0.122 & 1.840\,$\pm$\,0.165 & 1.11$^{\S}$ & 1.11$^{\S}$ & 1.06$^{\S}$ & 1.06$^{\S}$ & 1.09$^{\S}$ & 1.618\,$\pm$\,0.001 & 1.2983\,$\pm$\,0.0284 & 1.3723\,$\pm$\,0.0157 \\
$X\to YY\to cse\nu$ & 0.98970\,$\pm$\,0.00012 & 1.0033\,$\pm$\,0.0129 & 1.608\,$\pm$\,0.076 & 1.694\,$\pm$\,0.127 & 1.08$^{\S}$ & 1.06$^{\S}$ & 1.06$^{\S}$ & 1.05$^{\S}$ & 1.05$^{\S}$ & 1.447\,$\pm$\,0.013 & 1.3253\,$\pm$\,0.0476 & 1.2434\,$\pm$\,0.0210 \\
$X\to YY\to bqe\nu$ & 0.99512\,$\pm$\,0.00007 & 0.9975\,$\pm$\,0.0266 & 1.869\,$\pm$\,0.134 & 2.151\,$\pm$\,0.176 & 1.14$^{\S}$ & 1.21$^{\S}$ & 1.08$^{\S}$ & 1.10$^{\S}$ & 1.16$^{\S}$ & 1.839\,$\pm$\,0.073 & 1.4269\,$\pm$\,0.0320 & 1.4605\,$\pm$\,0.0324 \\
$X\to YY\to cqe\nu$ & 0.98991\,$\pm$\,0.00006 & 1.0103\,$\pm$\,0.0126 & 1.638\,$\pm$\,0.095 & 1.815\,$\pm$\,0.127 & 1.09$^{\S}$ & 1.09$^{\S}$ & 1.06$^{\S}$ & 1.08$^{\S}$ & 1.11$^{\S}$ & 1.589\,$\pm$\,0.048 & 1.3735\,$\pm$\,0.0033 & 1.3646\,$\pm$\,0.0129 \\
$X\to YY\to sqe\nu$ & 0.98982\,$\pm$\,0.00009 & 1.0013\,$\pm$\,0.0163 & 1.421\,$\pm$\,0.073 & 1.552\,$\pm$\,0.086 & 1.03$^{\S}$ & 1.04$^{\S}$ & 1.04$^{\S}$ & 1.04$^{\S}$ & 1.06$^{\S}$ & 1.328\,$\pm$\,0.023 & 1.2504\,$\pm$\,0.0231 & 1.1948\,$\pm$\,0.0170 \\
$X\to YY\to qqe\nu$ & 0.99013\,$\pm$\,0.00001 & 1.0097\,$\pm$\,0.0127 & 1.442\,$\pm$\,0.079 & 1.465\,$\pm$\,0.112 & 1.08$^{\S}$ & 1.04$^{\S}$ & 1.04$^{\S}$ & 1.03$^{\S}$ & 1.07$^{\S}$ & 1.223\,$\pm$\,0.019 & 1.1899\,$\pm$\,0.0144 & 1.1396\,$\pm$\,0.0100 \\
$X\to YY\to bc\mu\nu$ & 0.99487\,$\pm$\,0.00006 & 1.0076\,$\pm$\,0.0192 & 2.013\,$\pm$\,0.071 & 2.052\,$\pm$\,0.060 & 1.14$^{\S}$ & 1.12$^{\S}$ & 1.12$^{\S}$ & 1.08$^{\S}$ & 1.13$^{\S}$ & 1.716\,$\pm$\,0.018 & 1.3691\,$\pm$\,0.0079 & 1.4601\,$\pm$\,0.0262 \\
$X\to YY\to cs\mu\nu$ & 0.98895\,$\pm$\,0.00010 & 0.9957\,$\pm$\,0.0023 & 1.779\,$\pm$\,0.078 & 1.768\,$\pm$\,0.150 & 1.05$^{\S}$ & 1.07$^{\S}$ & 1.06$^{\S}$ & 1.09$^{\S}$ & 1.12$^{\S}$ & 1.461\,$\pm$\,0.039 & 1.3272\,$\pm$\,0.0018 & 1.2705\,$\pm$\,0.0318 \\
$X\to YY\to bq\mu\nu$ & 0.99461\,$\pm$\,0.00003 & 0.9909\,$\pm$\,0.0153 & 1.931\,$\pm$\,0.101 & 2.178\,$\pm$\,0.149 & 1.08$^{\S}$ & 1.08$^{\S}$ & 1.11$^{\S}$ & 1.08$^{\S}$ & 1.10$^{\S}$ & 1.793\,$\pm$\,0.008 & 1.3908\,$\pm$\,0.0284 & 1.4737\,$\pm$\,0.0161 \\
$X\to YY\to cq\mu\nu$ & 0.98910\,$\pm$\,0.00006 & 1.0031\,$\pm$\,0.0103 & 1.774\,$\pm$\,0.099 & 1.843\,$\pm$\,0.114 & 1.09$^{\S}$ & 1.08$^{\S}$ & 1.07$^{\S}$ & 1.08$^{\S}$ & 1.09$^{\S}$ & 1.567\,$\pm$\,0.022 & 1.3862\,$\pm$\,0.0049 & 1.3604\,$\pm$\,0.0122 \\
$X\to YY\to sq\mu\nu$ & 0.98957\,$\pm$\,0.00010 & 1.0055\,$\pm$\,0.0094 & 1.496\,$\pm$\,0.047 & 1.562\,$\pm$\,0.098 & 1.04$^{\S}$ & 1.08$^{\S}$ & 1.03$^{\S}$ & 1.06$^{\S}$ & 1.06$^{\S}$ & 1.332\,$\pm$\,0.031 & 1.2628\,$\pm$\,0.0057 & 1.2037\,$\pm$\,0.0073 \\
$X\to YY\to qq\mu\nu$ & 0.98962\,$\pm$\,0.00011 & 1.0151\,$\pm$\,0.0081 & 1.558\,$\pm$\,0.043 & 1.537\,$\pm$\,0.072 & 1.06$^{\S}$ & 1.06$^{\S}$ & 1.04$^{\S}$ & 1.06$^{\S}$ & 1.07$^{\S}$ & 1.254\,$\pm$\,0.005 & 1.2558\,$\pm$\,0.0470 & 1.1640\,$\pm$\,0.0166 \\
\midrule
$X\to YY\to bc\tau_e\nu$ & 0.99078\,$\pm$\,0.00014 & 1.0069\,$\pm$\,0.0186 & 1.843\,$\pm$\,0.076 & 2.235\,$\pm$\,0.092 & 1.18$^{\S}$ & 1.18$^{\S}$ & 1.12$^{\S}$ & 1.11$^{\S}$ & 1.18$^{\S}$ & 1.611\,$\pm$\,0.015 & 1.2766\,$\pm$\,0.0042 & 1.4122\,$\pm$\,0.0004 \\
$X\to YY\to cs\tau_e\nu$ & 0.98478\,$\pm$\,0.00016 & 1.0158\,$\pm$\,0.0208 & 1.870\,$\pm$\,0.079 & 2.091\,$\pm$\,0.108 & 1.07$^{\S}$ & 1.08$^{\S}$ & 1.08$^{\S}$ & 1.08$^{\S}$ & 1.13$^{\S}$ & 1.435\,$\pm$\,0.040 & 1.2629\,$\pm$\,0.0013 & 1.2692\,$\pm$\,0.0143 \\
$X\to YY\to bq\tau_e\nu$ & 0.99068\,$\pm$\,0.00004 & 1.0228\,$\pm$\,0.0087 & 1.890\,$\pm$\,0.037 & 2.299\,$\pm$\,0.048 & 1.12$^{\S}$ & 1.22$^{\S}$ & 1.11$^{\S}$ & 1.16$^{\S}$ & 1.19$^{\S}$ & 1.695\,$\pm$\,0.125 & 1.3627\,$\pm$\,0.0119 & 1.4743\,$\pm$\,0.0199 \\
$X\to YY\to cq\tau_e\nu$ & 0.98351\,$\pm$\,0.00014 & 1.0175\,$\pm$\,0.0201 & 1.857\,$\pm$\,0.046 & 2.114\,$\pm$\,0.080 & 1.12$^{\S}$ & 1.10$^{\S}$ & 1.11$^{\S}$ & 1.12$^{\S}$ & 1.15$^{\S}$ & 1.490\,$\pm$\,0.057 & 1.2876\,$\pm$\,0.0037 & 1.3334\,$\pm$\,0.0099 \\
$X\to YY\to sq\tau_e\nu$ & 0.98690\,$\pm$\,0.00005 & 0.9983\,$\pm$\,0.0060 & 1.794\,$\pm$\,0.063 & 2.038\,$\pm$\,0.112 & 1.08$^{\S}$ & 1.09$^{\S}$ & 1.09$^{\S}$ & 1.07$^{\S}$ & 1.11$^{\S}$ & 1.348\,$\pm$\,0.005 & 1.2400\,$\pm$\,0.0231 & 1.2379\,$\pm$\,0.0326 \\
$X\to YY\to qq\tau_e\nu$ & 0.98684\,$\pm$\,0.00004 & 1.0065\,$\pm$\,0.0048 & 1.825\,$\pm$\,0.074 & 1.968\,$\pm$\,0.129 & 1.07$^{\S}$ & 1.07$^{\S}$ & 1.05$^{\S}$ & 1.08$^{\S}$ & 1.08$^{\S}$ & 1.283\,$\pm$\,0.034 & 1.2114\,$\pm$\,0.0133 & 1.1771\,$\pm$\,0.0076 \\
$X\to YY\to bc\tau_\mu\nu$ & 0.99170\,$\pm$\,0.00011 & 1.0032\,$\pm$\,0.0266 & 1.863\,$\pm$\,0.062 & 2.302\,$\pm$\,0.010 & 1.20$^{\S}$ & 1.18$^{\S}$ & 1.15$^{\S}$ & 1.14$^{\S}$ & 1.15$^{\S}$ & 1.585\,$\pm$\,0.073 & 1.2690\,$\pm$\,0.0244 & 1.3913\,$\pm$\,0.0256 \\
$X\to YY\to cs\tau_\mu\nu$ & 0.98654\,$\pm$\,0.00009 & 1.0059\,$\pm$\,0.0098 & 1.945\,$\pm$\,0.044 & 2.194\,$\pm$\,0.154 & 1.09$^{\S}$ & 1.09$^{\S}$ & 1.11$^{\S}$ & 1.14$^{\S}$ & 1.13$^{\S}$ & 1.479\,$\pm$\,0.029 & 1.2856\,$\pm$\,0.0018 & 1.3102\,$\pm$\,0.0253 \\
$X\to YY\to bq\tau_\mu\nu$ & 0.99183\,$\pm$\,0.00017 & 1.0136\,$\pm$\,0.0248 & 1.989\,$\pm$\,0.072 & 2.654\,$\pm$\,0.113 & 1.11$^{\S}$ & 1.22$^{\S}$ & 1.11$^{\S}$ & 1.18$^{\S}$ & 1.23$^{\S}$ & 1.837\,$\pm$\,0.060 & 1.4393\,$\pm$\,0.0538 & 1.5009\,$\pm$\,0.0116 \\
$X\to YY\to cq\tau_\mu\nu$ & 0.98631\,$\pm$\,0.00006 & 1.0051\,$\pm$\,0.0117 & 1.961\,$\pm$\,0.049 & 2.281\,$\pm$\,0.115 & 1.10$^{\S}$ & 1.12$^{\S}$ & 1.08$^{\S}$ & 1.11$^{\S}$ & 1.12$^{\S}$ & 1.543\,$\pm$\,0.044 & 1.3478\,$\pm$\,0.0071 & 1.3666\,$\pm$\,0.0255 \\
$X\to YY\to sq\tau_\mu\nu$ & 0.98866\,$\pm$\,0.00008 & 1.0041\,$\pm$\,0.0101 & 1.904\,$\pm$\,0.056 & 2.095\,$\pm$\,0.161 & 1.08$^{\S}$ & 1.11$^{\S}$ & 1.08$^{\S}$ & 1.09$^{\S}$ & 1.10$^{\S}$ & 1.364\,$\pm$\,0.007 & 1.2777\,$\pm$\,0.0118 & 1.2632\,$\pm$\,0.0049 \\
$X\to YY\to qq\tau_\mu\nu$ & 0.98849\,$\pm$\,0.00006 & 1.0033\,$\pm$\,0.0105 & 1.925\,$\pm$\,0.035 & 2.102\,$\pm$\,0.116 & 1.10$^{\S}$ & 1.09$^{\S}$ & 1.11$^{\S}$ & 1.10$^{\S}$ & 1.10$^{\S}$ & 1.321\,$\pm$\,0.033 & 1.2773\,$\pm$\,0.0164 & 1.2121\,$\pm$\,0.0164 \\
\midrule
$X\to YY\to bc\tau_h\nu$ & 0.99016\,$\pm$\,0.00011 & 0.9936\,$\pm$\,0.0197 & 1.049\,$\pm$\,0.039 & 1.499\,$\pm$\,0.084 & 1.26$^{\S}$ & 1.25$^{\S}$ & 1.17$^{\S}$ & 1.21$^{\S}$ & 1.24$^{\S}$ & 1.561\,$\pm$\,0.071 & 1.3378\,$\pm$\,0.0355 & 1.3601\,$\pm$\,0.0598 \\
$X\to YY\to cs\tau_h\nu$ & 0.98361\,$\pm$\,0.00012 & 1.0049\,$\pm$\,0.0074 & 1.058\,$\pm$\,0.016 & 1.391\,$\pm$\,0.013 & 1.13$^{\S}$ & 1.10$^{\S}$ & 1.10$^{\S}$ & 1.12$^{\S}$ & 1.14$^{\S}$ & 1.393\,$\pm$\,0.016 & 1.2568\,$\pm$\,0.0137 & 1.2226\,$\pm$\,0.0310 \\
$X\to YY\to bq\tau_h\nu$ & 0.98990\,$\pm$\,0.00007 & 0.9948\,$\pm$\,0.0240 & 1.072\,$\pm$\,0.023 & 1.530\,$\pm$\,0.044 & 1.16$^{\S}$ & 1.16$^{\S}$ & 1.08$^{\S}$ & 1.10$^{\S}$ & 1.22$^{\S}$ & 1.563\,$\pm$\,0.077 & 1.2812\,$\pm$\,0.0142 & 1.3847\,$\pm$\,0.0163 \\
$X\to YY\to cq\tau_h\nu$ & 0.98301\,$\pm$\,0.00014 & 1.0058\,$\pm$\,0.0115 & 1.063\,$\pm$\,0.011 & 1.358\,$\pm$\,0.001 & 1.12$^{\S}$ & 1.09$^{\S}$ & 1.09$^{\S}$ & 1.13$^{\S}$ & 1.13$^{\S}$ & 1.393\,$\pm$\,0.033 & 1.2473\,$\pm$\,0.0210 & 1.2653\,$\pm$\,0.0033 \\
$X\to YY\to sq\tau_h\nu$ & 0.98758\,$\pm$\,0.00005 & 1.0150\,$\pm$\,0.0025 & 1.056\,$\pm$\,0.006 & 1.409\,$\pm$\,0.028 & 1.11$^{\S}$ & 1.11$^{\S}$ & 1.11$^{\S}$ & 1.07$^{\S}$ & 1.09$^{\S}$ & 1.379\,$\pm$\,0.038 & 1.3098\,$\pm$\,0.0514 & 1.2350\,$\pm$\,0.0040 \\
$X\to YY\to qq\tau_h\nu$ & 0.98711\,$\pm$\,0.00001 & 1.0017\,$\pm$\,0.0139 & 1.072\,$\pm$\,0.006 & 1.349\,$\pm$\,0.046 & 1.08$^{\S}$ & 1.09$^{\S}$ & 1.08$^{\S}$ & 1.08$^{\S}$ & 1.10$^{\S}$ & 1.295\,$\pm$\,0.046 & 1.2518\,$\pm$\,0.0325 & 1.2051\,$\pm$\,0.0103 \\
\midrule
QCD, $b$ & 0.96990\,$\pm$\,0.00026 & 1.0835\,$\pm$\,0.0188 & 1.107\,$\pm$\,0.015 & 1.446\,$\pm$\,0.025 & 1.20$^{\S}$ & 1.20$^{\S}$ & 1.18$^{\S}$ & 1.17$^{\S}$ & 1.20$^{\S}$ & 1.453\,$\pm$\,0.023 & 1.2338\,$\pm$\,0.0108 & 1.2764\,$\pm$\,0.0191 \\
QCD, $bb$ & 0.97223\,$\pm$\,0.00015 & 1.0855\,$\pm$\,0.0062 & 1.132\,$\pm$\,0.021 & 1.512\,$\pm$\,0.015 & 1.21$^{\S}$ & 1.21$^{\S}$ & 1.18$^{\S}$ & 1.16$^{\S}$ & 1.21$^{\S}$ & 1.618\,$\pm$\,0.042 & 1.2158\,$\pm$\,0.0090 & 1.2595\,$\pm$\,0.0154 \\
QCD, $bbc$ & 0.96939\,$\pm$\,0.00043 & 1.0908\,$\pm$\,0.0083 & 1.144\,$\pm$\,0.019 & 1.578\,$\pm$\,0.030 & 1.21$^{\S}$ & 1.23$^{\S}$ & 1.19$^{\S}$ & 1.21$^{\S}$ & 1.26$^{\S}$ & 1.731\,$\pm$\,0.077 & 1.2537\,$\pm$\,0.0103 & 1.3081\,$\pm$\,0.0224 \\
QCD, $bbcc$ & 0.97609\,$\pm$\,0.00002 & 1.0770\,$\pm$\,0.0070 & 1.128\,$\pm$\,0.022 & 1.660\,$\pm$\,0.060 & 1.24$^{\S}$ & 1.27$^{\S}$ & 1.24$^{\S}$ & 1.22$^{\S}$ & 1.25$^{\S}$ & 1.735\,$\pm$\,0.009 & 1.2784\,$\pm$\,0.0183 & 1.3038\,$\pm$\,0.0115 \\
QCD, $bbccs$ & 0.98029\,$\pm$\,0.00040 & 1.0885\,$\pm$\,0.0054 & 1.198\,$\pm$\,0.030 & 1.632\,$\pm$\,0.118 & 1.29$^{\S}$ & 1.24$^{\S}$ & 1.30$^{\S}$ & 1.23$^{\S}$ & 1.27$^{\S}$ & 1.704\,$\pm$\,0.018 & 1.2564\,$\pm$\,0.0440 & 1.3470\,$\pm$\,0.0127 \\
QCD, $bbccss$ & 0.98422\,$\pm$\,0.00020 & 1.1153\,$\pm$\,0.0308 & 1.188\,$\pm$\,0.014 & 1.633\,$\pm$\,0.095 & 1.28$^{\S}$ & 1.27$^{\S}$ & 1.32$^{\S}$ & 1.23$^{\S}$ & 1.33$^{\S}$ & 1.779\,$\pm$\,0.013 & 1.2624\,$\pm$\,0.0079 & 1.3991\,$\pm$\,0.0198 \\
QCD, $bbcs$ & 0.96609\,$\pm$\,0.00055 & 1.0531\,$\pm$\,0.0164 & 1.129\,$\pm$\,0.005 & 1.581\,$\pm$\,0.055 & 1.20$^{\S}$ & 1.25$^{\S}$ & 1.23$^{\S}$ & 1.16$^{\S}$ & 1.13$^{\S}$ & 1.624\,$\pm$\,0.005 & 1.2389\,$\pm$\,0.0300 & 1.3128\,$\pm$\,0.0204 \\
QCD, $bbcss$ & 0.97162\,$\pm$\,0.00077 & 1.0865\,$\pm$\,0.0161 & 1.158\,$\pm$\,0.054 & 1.531\,$\pm$\,0.071 & 1.14$^{\S}$ & 1.23$^{\S}$ & 1.23$^{\S}$ & 1.23$^{\S}$ & 1.25$^{\S}$ & 1.623\,$\pm$\,0.014 & 1.2178\,$\pm$\,0.0087 & 1.2450\,$\pm$\,0.0122 \\
QCD, $bbs$ & 0.96944\,$\pm$\,0.00013 & 1.0824\,$\pm$\,0.0021 & 1.134\,$\pm$\,0.006 & 1.491\,$\pm$\,0.037 & 1.19$^{\S}$ & 1.21$^{\S}$ & 1.18$^{\S}$ & 1.17$^{\S}$ & 1.19$^{\S}$ & 1.579\,$\pm$\,0.017 & 1.1991\,$\pm$\,0.0051 & 1.2659\,$\pm$\,0.0155 \\
QCD, $bbss$ & 0.97367\,$\pm$\,0.00013 & 1.1292\,$\pm$\,0.0158 & 1.191\,$\pm$\,0.001 & 1.684\,$\pm$\,0.029 & 1.25$^{\S}$ & 1.25$^{\S}$ & 1.27$^{\S}$ & 1.22$^{\S}$ & 1.28$^{\S}$ & 1.757\,$\pm$\,0.004 & 1.2838\,$\pm$\,0.0101 & 1.3481\,$\pm$\,0.0120 \\
QCD, $bc$ & 0.95198\,$\pm$\,0.00026 & 1.0622\,$\pm$\,0.0041 & 1.096\,$\pm$\,0.003 & 1.337\,$\pm$\,0.029 & 1.15$^{\S}$ & 1.14$^{\S}$ & 1.10$^{\S}$ & 1.14$^{\S}$ & 1.17$^{\S}$ & 1.367\,$\pm$\,0.041 & 1.1703\,$\pm$\,0.0141 & 1.2105\,$\pm$\,0.0035 \\
QCD, $bcc$ & 0.95342\,$\pm$\,0.00056 & 1.0674\,$\pm$\,0.0065 & 1.162\,$\pm$\,0.009 & 1.598\,$\pm$\,0.020 & 1.25$^{\S}$ & 1.22$^{\S}$ & 1.22$^{\S}$ & 1.20$^{\S}$ & 1.20$^{\S}$ & 1.605\,$\pm$\,0.001 & 1.3192\,$\pm$\,0.0516 & 1.3436\,$\pm$\,0.0032 \\
QCD, $bccs$ & 0.94777\,$\pm$\,0.00076 & 1.0879\,$\pm$\,0.0222 & 1.159\,$\pm$\,0.061 & 1.522\,$\pm$\,0.071 & 1.23$^{\S}$ & 1.26$^{\S}$ & 1.20$^{\S}$ & 1.22$^{\S}$ & 1.26$^{\S}$ & 1.443\,$\pm$\,0.067 & 1.2817\,$\pm$\,0.0419 & 1.3222\,$\pm$\,0.0053 \\
QCD, $bccss$ & 0.95705\,$\pm$\,0.00080 & 1.0910\,$\pm$\,0.0397 & 1.186\,$\pm$\,0.039 & 1.503\,$\pm$\,0.004 & 1.25$^{\S}$ & 1.27$^{\S}$ & 1.20$^{\S}$ & 1.26$^{\S}$ & 1.29$^{\S}$ & 1.476\,$\pm$\,0.107 & 1.3202\,$\pm$\,0.0034 & 1.2882\,$\pm$\,0.0369 \\
QCD, $bcs$ & 0.94563\,$\pm$\,0.00196 & 1.1039\,$\pm$\,0.0605 & 1.139\,$\pm$\,0.047 & 1.383\,$\pm$\,0.148 & 1.37$^{\S}$ & 1.10$^{\S}$ & 1.22$^{\S}$ & 1.11$^{\S}$ & 1.25$^{\S}$ & 1.416\,$\pm$\,0.034 & 1.2332\,$\pm$\,0.0156 & 1.2617\,$\pm$\,0.0669 \\
QCD, $bcss$ & 0.93997\,$\pm$\,0.00057 & 1.1218\,$\pm$\,0.0320 & 1.217\,$\pm$\,0.045 & 1.524\,$\pm$\,0.048 & 1.24$^{\S}$ & 1.27$^{\S}$ & 1.30$^{\S}$ & 1.32$^{\S}$ & 1.30$^{\S}$ & 1.465\,$\pm$\,0.031 & 1.3272\,$\pm$\,0.0019 & 1.3051\,$\pm$\,0.0573 \\
QCD, $bs$ & 0.96104\,$\pm$\,0.00023 & 1.0777\,$\pm$\,0.0117 & 1.105\,$\pm$\,0.009 & 1.408\,$\pm$\,0.015 & 1.18$^{\S}$ & 1.19$^{\S}$ & 1.14$^{\S}$ & 1.16$^{\S}$ & 1.18$^{\S}$ & 1.408\,$\pm$\,0.040 & 1.2283\,$\pm$\,0.0083 & 1.2452\,$\pm$\,0.0284 \\
QCD, $bss$ & 0.96152\,$\pm$\,0.00003 & 1.1128\,$\pm$\,0.0065 & 1.167\,$\pm$\,0.005 & 1.615\,$\pm$\,0.025 & 1.26$^{\S}$ & 1.24$^{\S}$ & 1.22$^{\S}$ & 1.24$^{\S}$ & 1.27$^{\S}$ & 1.600\,$\pm$\,0.029 & 1.3349\,$\pm$\,0.0033 & 1.3334\,$\pm$\,0.0057 \\
QCD, $c$ & 0.95932\,$\pm$\,0.00019 & 1.0590\,$\pm$\,0.0053 & 1.074\,$\pm$\,0.008 & 1.288\,$\pm$\,0.006 & 1.12$^{\S}$ & 1.12$^{\S}$ & 1.11$^{\S}$ & 1.12$^{\S}$ & 1.13$^{\S}$ & 1.277\,$\pm$\,0.021 & 1.1450\,$\pm$\,0.0111 & 1.1563\,$\pm$\,0.0007 \\
QCD, $cc$ & 0.95582\,$\pm$\,0.00023 & 1.0697\,$\pm$\,0.0104 & 1.075\,$\pm$\,0.007 & 1.475\,$\pm$\,0.013 & 1.18$^{\S}$ & 1.17$^{\S}$ & 1.14$^{\S}$ & 1.16$^{\S}$ & 1.18$^{\S}$ & 1.506\,$\pm$\,0.050 & 1.2063\,$\pm$\,0.0244 & 1.2379\,$\pm$\,0.0080 \\
QCD, $ccs$ & 0.95432\,$\pm$\,0.00023 & 1.1129\,$\pm$\,0.0191 & 1.105\,$\pm$\,0.002 & 1.563\,$\pm$\,0.011 & 1.21$^{\S}$ & 1.24$^{\S}$ & 1.22$^{\S}$ & 1.21$^{\S}$ & 1.22$^{\S}$ & 1.557\,$\pm$\,0.004 & 1.2822\,$\pm$\,0.0128 & 1.3038\,$\pm$\,0.0024 \\
QCD, $ccss$ & 0.95953\,$\pm$\,0.00029 & 1.1295\,$\pm$\,0.0098 & 1.137\,$\pm$\,0.005 & 1.540\,$\pm$\,0.007 & 1.22$^{\S}$ & 1.25$^{\S}$ & 1.23$^{\S}$ & 1.20$^{\S}$ & 1.22$^{\S}$ & 1.540\,$\pm$\,0.019 & 1.2816\,$\pm$\,0.0114 & 1.2931\,$\pm$\,0.0045 \\
QCD, $cs$ & 0.95169\,$\pm$\,0.00020 & 1.0835\,$\pm$\,0.0086 & 1.094\,$\pm$\,0.004 & 1.291\,$\pm$\,0.022 & 1.14$^{\S}$ & 1.16$^{\S}$ & 1.14$^{\S}$ & 1.15$^{\S}$ & 1.18$^{\S}$ & 1.280\,$\pm$\,0.013 & 1.1651\,$\pm$\,0.0051 & 1.1926\,$\pm$\,0.0067 \\
QCD, $css$ & 0.95217\,$\pm$\,0.00011 & 1.1011\,$\pm$\,0.0097 & 1.126\,$\pm$\,0.005 & 1.376\,$\pm$\,0.009 & 1.16$^{\S}$ & 1.19$^{\S}$ & 1.17$^{\S}$ & 1.16$^{\S}$ & 1.20$^{\S}$ & 1.386\,$\pm$\,0.005 & 1.2190\,$\pm$\,0.0052 & 1.2148\,$\pm$\,0.0061 \\
QCD, light & 0.97913\,$\pm$\,0.00008 & 1.0465\,$\pm$\,0.0042 & 1.050\,$\pm$\,0.006 & 1.148\,$\pm$\,0.011 & 1.10$^{\S}$ & 1.08$^{\S}$ & 1.10$^{\S}$ & 1.07$^{\S}$ & 1.07$^{\S}$ & 1.137\,$\pm$\,0.032 & 1.1016\,$\pm$\,0.0020 & 1.0763\,$\pm$\,0.0093 \\
QCD, $s$ & 0.97268\,$\pm$\,0.00002 & 1.0477\,$\pm$\,0.0001 & 1.044\,$\pm$\,0.007 & 1.171\,$\pm$\,0.009 & 1.09$^{\S}$ & 1.08$^{\S}$ & 1.09$^{\S}$ & 1.07$^{\S}$ & 1.09$^{\S}$ & 1.174\,$\pm$\,0.003 & 1.1079\,$\pm$\,0.0068 & 1.0962\,$\pm$\,0.0037 \\
QCD, $ss$ & 0.96946\,$\pm$\,0.00020 & 1.0551\,$\pm$\,0.0072 & 1.071\,$\pm$\,0.014 & 1.237\,$\pm$\,0.019 & 1.09$^{\S}$ & 1.10$^{\S}$ & 1.09$^{\S}$ & 1.09$^{\S}$ & 1.10$^{\S}$ & 1.197\,$\pm$\,0.009 & 1.1243\,$\pm$\,0.0062 & 1.1033\,$\pm$\,0.0204 \\
\end{longtable}
}

\end{landscape}

\begin{figure}[tbp]
\paperfig{\linewidth}{A3R}
\caption{Linear probes on JetClass-II, one class against the rest: background rejection $1/\epsB$ against
signal efficiency $\epsS$ for selected classes, with the discriminant $D_S = g_S$ of equation~(\ref{eq:discr})
and every other class of the class-balanced test set as the background, after training on \FtSizeESix{} jets.
Upper two rows: the eight classes with the largest loss of AUC at \Kcoarse{} classes, \MlstRocJcTwoHurtList{},
which lose between \MlstRocJcTwoHurtMinDrop{} and \MlstRocJcTwoHurtMaxDrop{} percentage points; lower row: the
four classes whose $1-\AUC$ stays closest to that of the \Kfull-class model at every level of the hierarchy,
\MlstRocJcTwoStableList{}, within \MlstRocJcTwoStableMaxChange\%. Both sets are chosen from the AUCs of
table~\ref{tab:perclass} by the figure code, among the classes not hatched in figure~\ref{fig:f3}, and each
is ordered as in figure~\ref{fig:f3}. Lines (solid, dashed, dash-dot and dotted): the \Kfull-, \Kres-, \Kmid- and \Kcoarse-class models, the mean
of the runs at fixed $\epsS$ with bands of one standard deviation over \MlstNRunsWord{} pretraining runs.
Lower panels: ratio to the \Kfull-class run of the same seed, drawn flat at one for that model. Each curve
stops where fewer than $\nmin$ background jets pass, equation~(\ref{eq:cut}), marked by the dashed black line
at $N_B/\nmin$.}
\label{fig:a3r}
\end{figure}

\subsection{Jet mass}\label{sec:app-mass}

Table~\ref{tab:massdetail} completes table~\ref{tab:mass} and figure~\ref{fig:f8a}. It gives the
resolution and the median response of the regressed mass and of ridge and MLP regressions on frozen
features, over all test jets and over resonant and QCD jets separately, and against jet $p_{\mathrm{T}}$
and generator soft-drop mass. The median response of the regressed mass lies close to one because its
target $\theta$, equation~(\ref{eq:masstarget}), is the logarithm of the mass ratio itself. Figure~\ref{fig:a17} shows the distributions of
the response for resonant and QCD jets.
% CHECK: pending MlstMassHeadResMassResponse, MlstMassHeadCoarseMassResponse within a few per cent of 1 (0.995 and 1.003 in the older head files)

\providecommand{\TableContinued}[2]{\label{#2}\ifnum\getpagerefnumber{#2}>\getpagerefnumber{#1}\relax{\footnotesize\textit{Table~\ref{#1} (continued)}}\par\smallskip\fi}
\begin{center}
\small
\begin{minipage}{\linewidth}\centering
\captionof{table}{Jet-mass regression on JetClass-II from networks pretrained on JetClass-II with 188, 162, 43 and 17 classes, in detail, over all test jets and over resonant and QCD jets. Upper part: resolution and median response of the regression output of the models pretrained with the mass regression, and of ridge and multilayer perceptron regressions on frozen features. Lower parts: resolution and median response of the soft-drop mass and of the regression output against jet $p_{\mathrm{T}}$ and against the generator-level soft-drop mass; every bin holds at least \MlstMassMinJets{} jets. The $p_{\mathrm{T}}$ bins hold every jet, also those with a generator-level mass below 20~GeV, whose response is very broad: hence a soft-drop resolution of QCD jets above 1 at high $p_{\mathrm{T}}$; the generator-mass bins span 20--500~GeV. Mean $\pm$ standard deviation over the \MlstNRunsWord{} pretraining runs; the models with the mass regression have \MlstMassNRunsWord{} runs, and both values are given. The random network without pretraining: three initializations.}
\label{tab:massdetail}
\par\medskip
\begin{adjustbox}{max width=\linewidth, max totalheight=0.55\textheight}
\begin{tabular}{>{\raggedright\arraybackslash}p{3.2cm}ccc}
\toprule
 & \multicolumn{3}{c}{resolution} \\
\cmidrule(lr){2-4}
jets & all & resonant & QCD \\
\midrule
soft-drop mass & 0.2448 & 0.2128 & 1.525 \\
\multicolumn{4}{l}{\textit{regression output of the pretrained network}} \\
\mbox{\Kres{} classes + mass} & 0.1471, 0.1463 & 0.1332, 0.1320 & 0.226, 0.226 \\
\mbox{\Kcoarse{} classes + mass} & 0.1450, 0.1461 & 0.1312, 0.1320 & 0.224, 0.224 \\
\multicolumn{4}{l}{\textit{ridge regression on the frozen features}} \\
\mbox{random network} & 0.3124\,$\pm$\,0.0072 & 0.2863\,$\pm$\,0.0063 & 0.913\,$\pm$\,0.031 \\
\mbox{\Kfull{} classes} & 0.3211\,$\pm$\,0.0034 & 0.2921\,$\pm$\,0.0026 & 0.559\,$\pm$\,0.013 \\
\mbox{\Kres{} classes} & 0.3156\,$\pm$\,0.0090 & 0.2876\,$\pm$\,0.0080 & 0.548\,$\pm$\,0.021 \\
\mbox{\Kmid{} classes} & 0.3041\,$\pm$\,0.0158 & 0.2797\,$\pm$\,0.0125 & 0.487\,$\pm$\,0.062 \\
\mbox{\Kcoarse{} classes} & 0.2929\,$\pm$\,0.0118 & 0.2708\,$\pm$\,0.0116 & 0.448\,$\pm$\,0.030 \\
\mbox{\Kres{} classes + mass} & 0.1885, 0.1922 & 0.1751, 0.1785 & 0.267, 0.270 \\
\mbox{\Kcoarse{} classes + mass} & 0.1842, 0.1824 & 0.1696, 0.1687 & 0.263, 0.260 \\
\multicolumn{4}{l}{\textit{multilayer perceptron on the frozen features}} \\
\mbox{random network} & 0.2421\,$\pm$\,0.0018 & 0.2214\,$\pm$\,0.0017 & 0.502\,$\pm$\,0.014 \\
\mbox{\Kfull{} classes} & 0.2200\,$\pm$\,0.0014 & 0.1990\,$\pm$\,0.0016 & 0.349\,$\pm$\,0.003 \\
\mbox{\Kres{} classes} & 0.2195\,$\pm$\,0.0017 & 0.1989\,$\pm$\,0.0020 & 0.343\,$\pm$\,0.003 \\
\mbox{\Kmid{} classes} & 0.2158\,$\pm$\,0.0029 & 0.1949\,$\pm$\,0.0023 & 0.338\,$\pm$\,0.007 \\
\mbox{\Kcoarse{} classes} & 0.2098\,$\pm$\,0.0031 & 0.1894\,$\pm$\,0.0035 & 0.327\,$\pm$\,0.006 \\
\mbox{\Kres{} classes + mass} & 0.1591, 0.1586 & 0.1436, 0.1430 & 0.247, 0.245 \\
\mbox{\Kcoarse{} classes + mass} & 0.1583, 0.1588 & 0.1420, 0.1431 & 0.247, 0.246 \\
\bottomrule
\end{tabular}
\end{adjustbox}
\end{minipage}
\par\bigskip
\begin{minipage}{\linewidth}\centering
\TableContinued{tab:massdetail}{tab:massdetail:2}
\begin{adjustbox}{max width=\linewidth, max totalheight=0.92\textheight}
\begin{tabular}{>{\raggedright\arraybackslash}p{3.2cm}ccc}
\toprule
 & \multicolumn{3}{c}{median response} \\
\cmidrule(lr){2-4}
jets & all & resonant & QCD \\
\midrule
soft-drop mass & 0.929 & 0.929 & 0.931 \\
\multicolumn{4}{l}{\textit{regression output of the pretrained network}} \\
\mbox{\Kres{} classes + mass} & 0.995, 1.003 & 0.996, 1.004 & 0.990, 1.000 \\
\mbox{\Kcoarse{} classes + mass} & 1.005, 1.005 & 1.006, 1.006 & 1.001, 1.000 \\
\multicolumn{4}{l}{\textit{ridge regression on the frozen features}} \\
\mbox{random network} & 0.941\,$\pm$\,0.003 & 0.936\,$\pm$\,0.004 & 0.965\,$\pm$\,0.003 \\
\mbox{\Kfull{} classes} & 0.964\,$\pm$\,0.001 & 0.975\,$\pm$\,0.001 & 0.905\,$\pm$\,0.006 \\
\mbox{\Kres{} classes} & 0.965\,$\pm$\,$<$0.001 & 0.977\,$\pm$\,0.001 & 0.909\,$\pm$\,0.007 \\
\mbox{\Kmid{} classes} & 0.968\,$\pm$\,0.005 & 0.977\,$\pm$\,0.004 & 0.922\,$\pm$\,0.012 \\
\mbox{\Kcoarse{} classes} & 0.973\,$\pm$\,0.005 & 0.980\,$\pm$\,0.004 & 0.937\,$\pm$\,0.010 \\
\mbox{\Kres{} classes + mass} & 1.000, 1.000 & 1.003, 1.003 & 0.983, 0.983 \\
\mbox{\Kcoarse{} classes + mass} & 1.002, 1.002 & 1.005, 1.003 & 0.989, 0.994 \\
\multicolumn{4}{l}{\textit{multilayer perceptron on the frozen features}} \\
\mbox{random network} & 0.983\,$\pm$\,0.005 & 0.984\,$\pm$\,0.005 & 0.980\,$\pm$\,0.007 \\
\mbox{\Kfull{} classes} & 0.991\,$\pm$\,0.005 & 0.994\,$\pm$\,0.004 & 0.975\,$\pm$\,0.007 \\
\mbox{\Kres{} classes} & 0.988\,$\pm$\,0.001 & 0.990\,$\pm$\,0.002 & 0.974\,$\pm$\,0.006 \\
\mbox{\Kmid{} classes} & 0.991\,$\pm$\,0.005 & 0.993\,$\pm$\,0.004 & 0.976\,$\pm$\,0.010 \\
\mbox{\Kcoarse{} classes} & 0.991\,$\pm$\,0.004 & 0.993\,$\pm$\,0.004 & 0.977\,$\pm$\,0.008 \\
\mbox{\Kres{} classes + mass} & 1.004, 1.006 & 1.006, 1.008 & 0.988, 0.993 \\
\mbox{\Kcoarse{} classes + mass} & 1.009, 1.006 & 1.012, 1.008 & 0.994, 0.993 \\
\bottomrule
\end{tabular}
\end{adjustbox}
\end{minipage}
\par\bigskip
\begin{minipage}{\linewidth}\centering
\TableContinued{tab:massdetail}{tab:massdetail:3}
\begin{adjustbox}{max width=\linewidth, max totalheight=0.92\textheight}
\begin{tabular}{>{\raggedright\arraybackslash}p{3.2cm}cccc}
\toprule
\textit{resolution} & \multicolumn{4}{c}{jet $p_{\mathrm{T}}$ [GeV]} \\
\cmidrule(lr){2-5}
estimator & 200--400 & 400--700 & 700--1000 & 1000--2500 \\
\midrule
\multicolumn{5}{l}{\textit{all jets}} \\
soft-drop mass & 0.2421 & 0.2252 & 0.23 & 0.33 \\
\mbox{\Kres{} classes + mass} & 0.1720, 0.1703 & 0.1483, 0.1468 & 0.14, 0.14 & 0.14, 0.14 \\
\mbox{\Kcoarse{} classes + mass} & 0.1691, 0.1702 & 0.1455, 0.1470 & 0.14, 0.14 & 0.14, 0.14 \\
\multicolumn{5}{l}{\textit{resonant jets}} \\
soft-drop mass & 0.2287 & 0.2146 & 0.22 & 0.21 \\
\mbox{\Kres{} classes + mass} & 0.1642, 0.1621 & 0.1412, 0.1393 & 0.13, 0.13 & 0.12, 0.12 \\
\mbox{\Kcoarse{} classes + mass} & 0.1610, 0.1620 & 0.1382, 0.1396 & 0.13, 0.13 & 0.12, 0.12 \\
\multicolumn{5}{l}{\textit{QCD jets}} \\
soft-drop mass & 0.2786 & 0.3329 & 1.10 & 3.61 \\
\mbox{\Kres{} classes + mass} & 0.1945, 0.1928 & 0.1854, 0.1846 & 0.21, 0.21 & 0.28, 0.28 \\
\mbox{\Kcoarse{} classes + mass} & 0.1918, 0.1947 & 0.1816, 0.1830 & 0.21, 0.21 & 0.28, 0.28 \\
\bottomrule
\end{tabular}
\end{adjustbox}
\end{minipage}
\par\bigskip
\begin{minipage}{\linewidth}\centering
\TableContinued{tab:massdetail}{tab:massdetail:4}
\begin{adjustbox}{max width=\linewidth, max totalheight=0.92\textheight}
\begin{tabular}{>{\raggedright\arraybackslash}p{3.2cm}cccc}
\toprule
\textit{median response} & \multicolumn{4}{c}{jet $p_{\mathrm{T}}$ [GeV]} \\
\cmidrule(lr){2-5}
estimator & 200--400 & 400--700 & 700--1000 & 1000--2500 \\
\midrule
\multicolumn{5}{l}{\textit{all jets}} \\
soft-drop mass & 0.8926 & 0.91 & 0.92 & 0.95 \\
\mbox{\Kres{} classes + mass} & 0.9862, 0.9970 & 0.99, 1.00 & 1.00, 1.00 & 1.00, 1.00 \\
\mbox{\Kcoarse{} classes + mass} & 0.9930, 0.9974 & 1.00, 1.00 & 1.01, 1.01 & 1.01, 1.01 \\
\multicolumn{5}{l}{\textit{resonant jets}} \\
soft-drop mass & 0.9036 & 0.92 & 0.92 & 0.95 \\
\mbox{\Kres{} classes + mass} & 0.9868, 0.9978 & 1.00, 1.00 & 1.00, 1.00 & 1.00, 1.00 \\
\mbox{\Kcoarse{} classes + mass} & 0.9933, 0.9980 & 1.00, 1.00 & 1.01, 1.01 & 1.01, 1.01 \\
\multicolumn{5}{l}{\textit{QCD jets}} \\
soft-drop mass & 0.8634 & 0.89 & 0.93 & 1.02 \\
\mbox{\Kres{} classes + mass} & 0.9845, 0.9949 & 0.99, 1.00 & 0.99, 1.00 & 0.99, 1.00 \\
\mbox{\Kcoarse{} classes + mass} & 0.9921, 0.9955 & 1.00, 1.00 & 1.00, 1.00 & 1.01, 1.00 \\
\bottomrule
\end{tabular}
\end{adjustbox}
\end{minipage}
\par\bigskip
\begin{minipage}{\linewidth}\centering
\TableContinued{tab:massdetail}{tab:massdetail:5}
\begin{adjustbox}{max width=\linewidth, max totalheight=0.92\textheight}
\begin{tabular}{>{\raggedright\arraybackslash}p{3.2cm}cc}
\toprule
\textit{resolution} & \multicolumn{2}{c}{generator-level soft-drop mass [GeV]} \\
\cmidrule(lr){2-3}
estimator & 20--50 & 50--100 \\
\midrule
\multicolumn{3}{l}{\textit{all jets}} \\
soft-drop mass & 0.2808 & 0.2064 \\
\mbox{\Kres{} classes + mass} & 0.1751, 0.1756 & 0.1425, 0.1416 \\
\mbox{\Kcoarse{} classes + mass} & 0.1755, 0.1740 & 0.1418, 0.1411 \\
\multicolumn{3}{l}{\textit{resonant jets}} \\
soft-drop mass & 0.2746 & 0.1939 \\
\mbox{\Kres{} classes + mass} & 0.1632, 0.1634 & 0.1366, 0.1359 \\
\mbox{\Kcoarse{} classes + mass} & 0.1631, 0.1621 & 0.1357, 0.1351 \\
\multicolumn{3}{l}{\textit{QCD jets}} \\
soft-drop mass & 0.2920 & 0.2393 \\
\mbox{\Kres{} classes + mass} & 0.2033, 0.2030 & 0.1653, 0.1639 \\
\mbox{\Kcoarse{} classes + mass} & 0.2032, 0.2006 & 0.1630, 0.1631 \\
\bottomrule
\end{tabular}
\end{adjustbox}
\end{minipage}
\par\bigskip
\begin{minipage}{\linewidth}\centering
\TableContinued{tab:massdetail}{tab:massdetail:6}
\begin{adjustbox}{max width=\linewidth, max totalheight=0.92\textheight}
\begin{tabular}{>{\raggedright\arraybackslash}p{3.2cm}ccc}
\toprule
\textit{resolution} & \multicolumn{3}{c}{generator-level soft-drop mass [GeV]} \\
\cmidrule(lr){2-4}
estimator & 100--150 & 150--250 & 250--500 \\
\midrule
\multicolumn{4}{l}{\textit{all jets}} \\
soft-drop mass & 0.2045 & 0.2102 & 0.2245 \\
\mbox{\Kres{} classes + mass} & 0.1302, 0.1291 & 0.1233, 0.1213 & 0.1163, 0.1135 \\
\mbox{\Kcoarse{} classes + mass} & 0.1284, 0.1288 & 0.1210, 0.1220 & 0.1134, 0.1145 \\
\multicolumn{4}{l}{\textit{resonant jets}} \\
soft-drop mass & 0.1924 & 0.2010 & 0.2165 \\
\mbox{\Kres{} classes + mass} & 0.1261, 0.1252 & 0.1198, 0.1175 & 0.1133, 0.1107 \\
\mbox{\Kcoarse{} classes + mass} & 0.1246, 0.1250 & 0.1174, 0.1183 & 0.1107, 0.1118 \\
\multicolumn{4}{l}{\textit{QCD jets}} \\
soft-drop mass & 0.2457 & 0.2646 & 0.2782 \\
\mbox{\Kres{} classes + mass} & 0.1508, 0.1515 & 0.1515, 0.1509 & 0.1511, 0.1470 \\
\mbox{\Kcoarse{} classes + mass} & 0.1505, 0.1514 & 0.1483, 0.1515 & 0.1434, 0.1482 \\
\bottomrule
\end{tabular}
\end{adjustbox}
\end{minipage}
\par\bigskip
\begin{minipage}{\linewidth}\centering
\TableContinued{tab:massdetail}{tab:massdetail:7}
\begin{adjustbox}{max width=\linewidth, max totalheight=0.92\textheight}
\begin{tabular}{>{\raggedright\arraybackslash}p{3.2cm}ccccc}
\toprule
\textit{median response} & \multicolumn{5}{c}{generator-level soft-drop mass [GeV]} \\
\cmidrule(lr){2-6}
estimator & 20--50 & 50--100 & 100--150 & 150--250 & 250--500 \\
\midrule
\multicolumn{6}{l}{\textit{all jets}} \\
soft-drop mass & 1.00 & 0.94 & 0.91 & 0.89 & 0.8519 \\
\mbox{\Kres{} classes + mass} & 1.03, 1.04 & 1.01, 1.01 & 0.99, 1.00 & 0.99, 1.00 & 0.9710, 0.9772 \\
\mbox{\Kcoarse{} classes + mass} & 1.03, 1.04 & 1.01, 1.02 & 1.00, 1.00 & 1.00, 1.00 & 0.9837, 0.9792 \\
\multicolumn{6}{l}{\textit{resonant jets}} \\
soft-drop mass & 1.01 & 0.95 & 0.92 & 0.90 & 0.8621 \\
\mbox{\Kres{} classes + mass} & 1.03, 1.04 & 1.01, 1.02 & 1.00, 1.00 & 0.99, 1.00 & 0.9722, 0.9783 \\
\mbox{\Kcoarse{} classes + mass} & 1.04, 1.05 & 1.02, 1.02 & 1.01, 1.01 & 1.00, 1.00 & 0.9847, 0.9804 \\
\multicolumn{6}{l}{\textit{QCD jets}} \\
soft-drop mass & 0.97 & 0.86 & 0.82 & 0.78 & 0.7382 \\
\mbox{\Kres{} classes + mass} & 1.01, 1.02 & 0.98, 0.99 & 0.97, 0.98 & 0.96, 0.97 & 0.9565, 0.9631 \\
\mbox{\Kcoarse{} classes + mass} & 1.02, 1.02 & 0.99, 0.99 & 0.98, 0.98 & 0.98, 0.97 & 0.9704, 0.9641 \\
\bottomrule
\end{tabular}
\end{adjustbox}
\end{minipage}
\end{center}


\begin{figure}[tbp]
\paperfig{\linewidth}{\FigIdAseventeen}
\caption{Distribution of the response $r = m_{\mathrm{est}}/\mgen$ of equation~(\ref{eq:resolution}), the
estimated over the generator-level soft-drop mass, for (a) resonant jets and (b) QCD jets of the JetClass-II test set with a
generator-level match: the reconstructed soft-drop mass (filled), and the regressed mass of the \Kres- and
\Kcoarse-class models with the auxiliary mass regression (orange long dashes and red long dash-dot); the lower panels give the ratio to the
reconstructed soft-drop mass, drawn only in bins that hold at least \MlstMassMinJets{} jets of the reconstructed
soft-drop mass. Bands span the \MlstMassNRunsWord{}
pretraining runs of each model with the auxiliary mass regression. The resolution and the median response in each group, and against jet transverse
momentum and generator mass, are in table~\ref{tab:massdetail}.}
\label{fig:a17}
\end{figure}

Table~\ref{tab:sculpt} gives the values behind figure~\ref{fig:f8b} at QCD efficiencies of
\SculptEffHigh{} and \SculptEffLow, beside the distance that a random selection of as many QCD test jets
gives.

\begin{landscape}
{
\setlength{\LTleft}{0pt}\setlength{\LTright}{0pt}\setlength{\LTcapwidth}{\linewidth}\makeatletter\def\LT@makecaption#1#2#3{\LT@mcol\LT@cols c{\hbox to\z@{\hss\parbox[t]\LTcapwidth{\normalsize #1{#2. }#3\par\vskip\belowcaptionskip}\hss}}}\makeatother
\begin{center}
\begin{minipage}{\linewidth}\centering
\setlength{\abovecaptionskip}{0pt}
\captionof{table}{Mass sculpting of QCD jets by the output scores of the networks pretrained with 188, 162, 43 and 17 classes and of the control networks. The distance is the Jensen--Shannon distance (base 2) $\times10^{2}$ between the soft-drop mass distributions of all QCD test jets and of those passing a cut on $1-g_{\mathrm{QCD}}$, one minus the summed output of the QCD classes. The cut keeps the fraction $\epsB$ of the QCD jets or half of the resonant test jets, the \LitJetClassTwoNResonantClasses{} resonant classes taken together ($\epsS = 50\%$), where $R_{50} = 1/\epsB$, and $\epsS$ beside a distance is the fraction of the resonant test jets that pass. The last row is the distance expected for a random selection of as many QCD test jets. A flavor boundary is a class boundary between decays that differ only in quark flavor. Values are the mean $\pm$ standard deviation over \MlstNRunsWord{} pretraining runs per label set, \RandFlavRuns{} per flavor-boundary partition and \MlstMassNRunsWord{} per network with an added mass output. Bold marks the best value of a column and every value within three standard deviations of it.}
\label{tab:sculpt}
\par\medskip
\begin{tabular*}{\linewidth}{@{\extracolsep{\fill}}lcccccc@{}}
\toprule
 & \multicolumn{2}{c}{$\epsB = 10\%$} & \multicolumn{2}{c}{$\epsB = 1\%$} & \multicolumn{2}{c}{$\varepsilon_S = 50\%$} \\
\cmidrule(lr){2-3}\cmidrule(lr){4-5}\cmidrule(lr){6-7}
 & distance & $\varepsilon_S$ & distance & $\varepsilon_S$ & $R_{50}$ & distance \\
\midrule
\Kfull{} classes & {\fontseries{b}\selectfont 5.24\,$\pm$\,0.34} & {\fontseries{b}\selectfont 0.8178\,$\pm$\,0.0011} & {\fontseries{b}\selectfont 4.39\,$\pm$\,0.24} & {\fontseries{b}\selectfont 0.5483\,$\pm$\,0.0014} & {\fontseries{b}\selectfont 169\,$\pm$\,2} & {\fontseries{b}\selectfont 5.5\,$\pm$\,0.5} \\
\Kres{} classes & {\fontseries{b}\selectfont 5.32\,$\pm$\,0.09} & {\fontseries{b}\selectfont 0.8181\,$\pm$\,0.0009} & {\fontseries{b}\selectfont 4.29\,$\pm$\,0.29} & {\fontseries{b}\selectfont 0.5490\,$\pm$\,0.0009} & {\fontseries{b}\selectfont 170\,$\pm$\,3} & {\fontseries{b}\selectfont 5.5\,$\pm$\,0.1} \\
\Kmid{} classes & {\fontseries{b}\selectfont 5.60\,$\pm$\,0.27} & {\fontseries{b}\selectfont 0.8179\,$\pm$\,0.0004} & {\fontseries{b}\selectfont 5.05\,$\pm$\,0.74} & {\fontseries{b}\selectfont 0.5486\,$\pm$\,0.0007} & {\fontseries{b}\selectfont 169\,$\pm$\,2} & {\fontseries{b}\selectfont 6.4\,$\pm$\,0.2} \\
\Kcoarse{} classes & 5.43\,$\pm$\,0.10 & 0.8147\,$\pm$\,0.0005 & 5.61\,$\pm$\,0.17 & 0.5437\,$\pm$\,0.0009 & {\fontseries{b}\selectfont 159\,$\pm$\,3} & 6.7\,$\pm$\,0.2 \\
\multicolumn{7}{@{}l}{17-class partitions} \\
\quad no flavor boundary & 5.48\,$\pm$\,0.02 & {\fontseries{b}\selectfont 0.8162\,$\pm$\,0.0004} & {\fontseries{b}\selectfont 5.09\,$\pm$\,0.01} & {\fontseries{b}\selectfont 0.5466\,$\pm$\,0.0013} & {\fontseries{b}\selectfont 164\,$\pm$\,2} & {\fontseries{b}\selectfont 6.0\,$\pm$\,0.4} \\
\quad $b$-quark boundary & {\fontseries{b}\selectfont 5.41\,$\pm$\,0.44} & {\fontseries{b}\selectfont 0.8174\,$\pm$\,0.0004} & {\fontseries{b}\selectfont 5.21\,$\pm$\,0.07} & {\fontseries{b}\selectfont 0.5472\,$\pm$\,0.0011} & {\fontseries{b}\selectfont 166\,$\pm$\,4} & {\fontseries{b}\selectfont 5.9\,$\pm$\,0.5} \\
\quad random flavor boundary & {\fontseries{b}\selectfont 5.33\,$\pm$\,0.08} & {\fontseries{b}\selectfont 0.8167\,$\pm$\,0.0010} & {\fontseries{b}\selectfont 4.86\,$\pm$\,0.06} & {\fontseries{b}\selectfont 0.5462\,$\pm$\,0.0022} & {\fontseries{b}\selectfont 166\,$\pm$\,5} & {\fontseries{b}\selectfont 6.0\,$\pm$\,0.4} \\
\Kres{} classes + mass & {\fontseries{b}\selectfont 4.43\,$\pm$\,0.30} & {\fontseries{b}\selectfont 0.8156\,$\pm$\,0.0016} & {\fontseries{b}\selectfont 4.33\,$\pm$\,0.12} & 0.5447\,$\pm$\,0.0011 & {\fontseries{b}\selectfont 160\,$\pm$\,3} & {\fontseries{b}\selectfont 5.5\,$\pm$\,0.3} \\
\Kcoarse{} classes + mass & {\fontseries{b}\selectfont 4.72\,$\pm$\,0.70} & 0.8124\,$\pm$\,0.0012 & {\fontseries{b}\selectfont 4.28\,$\pm$\,0.33} & 0.5387\,$\pm$\,0.0013 & 151\,$\pm$\,2 & {\fontseries{b}\selectfont 5.1\,$\pm$\,0.4} \\
\midrule
random selection of as many QCD jets & 0.94 & -- & 3.13 & -- & -- & -- \\
\bottomrule
\end{tabular*}
\end{minipage}
\end{center}
}
\end{landscape}


\subsection{Resonance tagging of $X\to b\bar b$ and $X\to bs$}\label{sec:app-xbxres}
% CHECK: pending X->bb/X->bs samples; drop this subsection with F9 if they are not staged by 2026-11-12

Table~\ref{tab:xbx} gives the signal efficiencies at fixed QCD efficiency behind figure~\ref{fig:f9}, and
table~\ref{tab:xbxrej} the rejection of QCD jets at 50\% and 30\% signal efficiency for every
configuration, with the transfer head, the direct discriminant and the closest-output-node control.

\begin{table}[t]
\centering
\small
\caption{Resonance tagging in the soft-drop mass window, \XbxWindowLo--\XbxWindowHi~GeV: signal efficiency (\%) at the QCD efficiency given, mean $\pm$ standard deviation over the pretraining runs. Transfer: a classifier on the frozen features, its seeds averaged per run (the randomly initialized network: the linear probe). Direct: Sophon's discriminant from the output scores, where $X\to b\bar b$ is a class of its own.}
\label{tab:xbx}
\begin{adjustbox}{max width=\linewidth}
\begin{tabular}{>{\raggedright\arraybackslash}p{5.0cm}c}
\toprule
Pretraining & $X\to b\bar b$, transfer ($\epsilon_B=0.01$) \\
\midrule
\mbox{\Kfull{} classes} & 71.79\,$\pm$\,0.14 \\
\mbox{\Kres{} classes} & 71.69\,$\pm$\,0.27 \\
\mbox{\Kmid{} classes} & 70.81\,$\pm$\,0.33 \\
\mbox{\Kcoarse{} classes} & 65.9\,$\pm$\,1.0 \\
\mbox{\Kres{} classes + mass} & 71.4, 70.7 \\
\mbox{\Kcoarse{} classes + mass} & 66.6, 67.0 \\
random partitions (all five) & 69.1\,$\pm$\,0.8 \\
\mbox{no flavor boundary} & 63.1, 61.9 \\
\mbox{$b$-quark boundary} & 66.9, 67.6 \\
\mbox{random flavor boundary} & 67.1, 66.7 \\
\mbox{random network} & 17.9\,$\pm$\,1.6 \\
\bottomrule
\end{tabular}
\end{adjustbox}
\par\bigskip
\begin{adjustbox}{max width=\linewidth}
\begin{tabular}{>{\raggedright\arraybackslash}p{5.0cm}cc}
\toprule
Pretraining & $X\to b\bar b$, direct ($\epsilon_B=0.01$) & $X\to bs$, transfer ($\epsilon_B=0.01$) \\
\midrule
\mbox{\Kfull{} classes} & 69.84\,$\pm$\,0.05 & 69.22\,$\pm$\,0.19 \\
\mbox{\Kres{} classes} & 69.95\,$\pm$\,0.19 & 68.90\,$\pm$\,0.06 \\
\mbox{\Kmid{} classes} & --- & 67.14\,$\pm$\,0.07 \\
\mbox{\Kcoarse{} classes} & --- & 59.4\,$\pm$\,1.6 \\
\mbox{\Kres{} classes + mass} & 69.6, 69.9 & 68.1, 68.0 \\
\mbox{\Kcoarse{} classes + mass} & --- & 59.1, 60.7 \\
random partitions (all five) & --- & 65.1\,$\pm$\,1.2 \\
\mbox{no flavor boundary} & --- & 59.5, 59.1 \\
\mbox{$b$-quark boundary} & --- & 62.0, 61.2 \\
\mbox{random flavor boundary} & --- & 63.0, 63.4 \\
\mbox{random network} & --- & 13.4\,$\pm$\,0.6 \\
\bottomrule
\end{tabular}
\end{adjustbox}
\end{table}

\begin{landscape}
\begin{table*}[t]
\centering
\small
\caption{Resonance tagging in the mass window: rejection $1/\epsB$ of QCD jets at 50\% and 30\% signal efficiency. Transfer: the classifier on frozen features; direct: Sophon's discriminant from the output scores; closest node: the pretrained $X\to bq$ output against QCD; from scratch: the network trained on each signal, three seeds. A dash: the read-out cannot be formed. Mean $\pm$ SD over the pretraining runs; two values: the two runs of a configuration. Random partitions: mean $\pm$ SD over the five partitions into 17 classes, one run each.}
\label{tab:xbxrej}
\begin{tabular*}{\linewidth}{@{\extracolsep{\fill}}lcccccccc@{}}
\toprule
 & \multicolumn{2}{c}{$X\to b\bar b$, transfer} & \multicolumn{2}{c}{$X\to b\bar b$, direct} & \multicolumn{2}{c}{$X\to bs$, transfer} & \multicolumn{2}{c}{$X\to bs$, closest node} \\
\cmidrule(lr){2-3}\cmidrule(lr){4-5}\cmidrule(lr){6-7}\cmidrule(lr){8-9}
 & $R_{50}$ & $R_{30}$ & $R_{50}$ & $R_{30}$ & $R_{50}$ & $R_{30}$ & $R_{50}$ & $R_{30}$ \\
\midrule
\Kfull{} classes & 350\,$\pm$\,4 & 1{,}010\,$\pm$\,40 & 288.8\,$\pm$\,3.2 & 792\,$\pm$\,27 & 435\,$\pm$\,27 & 2{,}070\,$\pm$\,150 & 144.1\,$\pm$\,3.5 & 377\,$\pm$\,19 \\
\Kres{} classes & 332\,$\pm$\,10 & 920\,$\pm$\,40 & 289.1\,$\pm$\,1.7 & 776\,$\pm$\,7 & 425\,$\pm$\,4 & 2{,}020\,$\pm$\,170 & 142.0\,$\pm$\,2.8 & 366\,$\pm$\,10 \\
\Kmid{} classes & 321\,$\pm$\,10 & 890\,$\pm$\,50 & -- & -- & 357\,$\pm$\,16 & 1{,}610\,$\pm$\,60 & 91.9\,$\pm$\,3.5 & 238\,$\pm$\,10 \\
\Kcoarse{} classes & 258\,$\pm$\,19 & 770\,$\pm$\,80 & -- & -- & 191\,$\pm$\,26 & 690\,$\pm$\,90 & 38.0\,$\pm$\,0.8 & 112.2\,$\pm$\,3.1 \\
\Kres{} classes + mass & 333, 318 & 1{,}005, 865 & 291, 289 & 807, 786 & 370, 375 & 1{,}713, 1{,}788 & 141, 138 & 369, 370 \\
\Kcoarse{} classes + mass & 266, 261 & 803, 740 & -- & -- & 179, 211 & 713, 820 & 38.1, 39.5 & 116, 117 \\
no flavor boundary & 223, 212 & 689, 697 & -- & -- & 199, 188 & 760, 774 & 38.3, 40.0 & 112, 118 \\
$b$-quark boundary & 262, 267 & 753, 742 & -- & -- & 224, 197 & 776, 665 & 39.8, 37.3 & 115, 106 \\
random flavor boundary & 279, 275 & 788, 830 & -- & -- & 256, 258 & 983, 1{,}071 & 40.3, 38.7 & 119, 111 \\
random partitions & 290\,$\pm$\,18 & 790\,$\pm$\,40 & -- & -- & 310\,$\pm$\,40 & 1{,}360\,$\pm$\,310 & 34\,$\pm$\,8 & 94\,$\pm$\,19 \\
random network & 15.9\,$\pm$\,0.9 & 42.8\,$\pm$\,3.5 & -- & -- & 12.9\,$\pm$\,0.6 & 31.5\,$\pm$\,1.8 & -- & -- \\
from scratch & -- & -- & -- & -- & 146\,$\pm$\,11 & 610\,$\pm$\,40 & -- & -- \\
\bottomrule
\end{tabular*}
\end{table*}
\end{landscape}

